\chapter*{Overview}

\subsection*{The result}

Let $f \in \ZZ[x]$ be monic and squarefree of odd degree $2g+1$ with $g \ge 1$, and let
$C \colon y^2 = f(x)$ be the associated hyperelliptic curve of genus $g$, with its single
point $\infty$ at infinity. Let $J$ denote its Jacobian, realized below as a class group,
and let $p$ be an odd prime such that the reduction $\bar f \in \Fp[x]$ is squarefree
(equivalently, $p \nmid \operatorname{disc} f$). This project formalizes the
method of Chabauty and Coleman for such curves: if
\[
  \operatorname{rk} J(\QQ) \;:=\; \dim_{\QQ} \bigl(J(\QQ) \otimes_{\ZZ} \QQ\bigr) \;\le\; r
\]
for a natural number $r < g$, and $p > 2r + 2$, then
\[
  \#C(\QQ) \;\le\; \#\bar C(\Fp) + 2r
\]
(Theorem~\ref{thm:main}); in particular $C(\QQ)$ is finite
(Corollary~\ref{cor:finiteness}). The choice $r = g-1$, admissible as soon as
$\operatorname{rk} J(\QQ) < g$, recovers Coleman's bound
$\#C(\QQ) \le \#\bar C(\Fp) + 2g-2$ for $p > 2g$; the improvement to $2r$ is Theorem~6.4
of M.~Stoll, \emph{Independence of rational points on twists of a given curve}, Compositio
Math.\ 142 (2006), 1201--1214. No finite generation of $J(\QQ)$ is assumed: the rank is
a cardinal, required to be $\le r$. The main theorems are proved over an abstract complete
discrete valuation ring of characteristic $0$ with algebraically closed residue field of
characteristic $\neq 2$, with an arithmetic base field $K_0$ in the role of $\QQ$; the
classical statement over $\ZZ$ is an instantiation at the Witt vectors $W(\Fpbar)$. That
ring is not assumed unramified: it carries a uniformizer $\pi$ and a natural number $e$
with $v(p) \le e$ bounding the absolute ramification index, and the condition on the prime
is then $p > 2r + e + 1$ --- Corollary~6.7 of \emph{loc.\ cit.}, which is what admits
ramified primes over a general base field. Over $\QQ$ one has $\pi = p$ and $e = 1$.

The blueprint also covers $2$-descent on the Jacobian (Chapter~\ref{ch:descent}): the $x - T$ map
$J(K) \to L^\times/(L^\times)^2$, $L = K[x]/(f)$, whose kernel is $2J(K)$, certificates for halving
points, the $2$-Selmer group over a number field with the bound on the rank of~$J(K)$ that it
gives when~$J(K)$ is finitely generated, and two examples of genus~$2$ over~$\QQ$.

\subsection*{The strategy}

The proof follows the classical outline. The $p$-adic points of $C$ decompose into
\emph{residue disks}, the fibres of the reduction map $C(\Qp) \to \bar C(\Fp)$
(Lemma~\ref{lem:disk-decomp}); each disk is parametrized by a power series chart
(Lemma~\ref{lem:disk-param}). A $p$-adic logarithm on the kernel of reduction of the
Jacobian pairs regular differentials $\omega$ against divisor classes; the pairing
$\lambda_\omega$ is additive in the class, linear in $\omega$, and on a fixed residue disk
it is computed by the antiderivative of the power-series expansion of $\omega$
(Theorem~\ref{thm:glue}). If the rank is at most $r$, a dimension count produces a whole
\emph{space of annihilating differentials}, of codimension at most $r$ in the regular
differentials: every $\omega$ in it has $\lambda_\omega = 0$ on the relevant subgroup
(Theorem~\ref{thm:annihilator}). The rational points in one residue disk then all map to
zeros of a single $p$-adic power series, whose number is bounded through Strassmann's
theorem by $1 + \ord_{\bar P}(\bar\omega)$ (Theorem~\ref{thm:disk-bound}). Nothing forces
the same $\omega$ in every disk: each may use the differential of that space vanishing to
the smallest order at its own centre, and those minimal orders sum to at most $2r$
(Corollary~\ref{cor:sum-ord-space}), where one fixed differential would only give the
degree $2g-2$ of a canonical divisor (Corollary~\ref{cor:sum-ord}).

\subsection*{The Jacobian as a class group}

Mathlib has no algebraic geometry of abelian varieties: no Jacobian variety, no smooth
projective models, no N\'eron models, and no Riemann--Roch theorem. The project therefore
works with an entirely ring-theoretic model of the Jacobian. For a field $K$ of
characteristic $\neq 2$, the affine coordinate ring $\A{K} = K[x][y]/(y^2 - f)$ is a
Dedekind domain (Theorem~\ref{thm:hypring-dedekind}), and the Jacobian is \emph{defined} as
its ideal class group,
\[
  \J{K} \;:=\; \Cl(\A{K})
\]
(Definition~\ref{def:hypjac}). Because $\deg f$ is odd there is exactly one point at
infinity, and the degree-zero divisor classes of the smooth projective curve correspond
exactly to the ideal classes of the affine ring, so nothing is lost. The Abel--Jacobi map
sends a point to the class of its maximal ideal, and every class is represented by a unique
\emph{reduced Mumford ideal} $(u,\, y - v)$ with $u$ monic, $\deg v < \deg u \le g$ and
$u \mid v^2 - f$ (Lemmas~\ref{lem:mumford-reduce} and~\ref{lem:mumford-exists}); the group
law is computed by Cantor's composition-and-reduction algorithm on these representatives.
All of Chapter~1 works over an arbitrary field, replacing the geometry of the curve by
commutative algebra of $\A{K}$: norms and traces in the quadratic extension, ideal
factorizations, and elementary degree and parity arguments.

Three consequences of this model shape the whole development.

\emph{Reduction of classes is constructed by hand.} With no N\'eron model, the reduction
map $\redJ \colon \J{K} \to \Cl(\A{\bar\kappa})$ on classes is defined by reducing points:
a class is represented over a finite extension by a divisor supported on rational points,
each point is reduced through its residue disk, and well-definedness comes from a
specialization theorem for principal divisors, proved with Weierstrass preparation
(Theorem~\ref{thm:spec}, Definition~\ref{def:redJ}, Theorem~\ref{thm:redJ-hom}).

\emph{The formal group is built from Cantor's algorithm.} The identity of the Jacobian
lies on the theta locus, where reduced Mumford representatives degenerate below degree
$g$, so there is no coordinate chart at the origin. Instead, an auxiliary class $\tau$ in
general position is constructed (Lemmas~\ref{lem:tau-generic} and~\ref{lem:tau-exists}),
and the kernel of reduction is coordinatized by the chart $c \mapsto c + \tau$: the
translated group law
$F_\tau(X,Y)$, expressing the Mumford data of $\mathrm{ch}(X) + \mathrm{ch}(Y) - \tau$ in
chart coordinates, is a $g$-dimensional commutative formal group law over the base ring
(Theorem~\ref{thm:Ftau-assoc}), and the kernel of reduction is isomorphic to the group of
its points on a polydisk (Theorem~\ref{thm:chart-mem}). Cantor's algorithm enters as a
family of exact polynomial identities: each composition step carries a B\'ezout identity
between the participating $u$-polynomials as a hypothesis, and each reduction step an
inverse of a leading coefficient, so that all constructions are polynomial in the data and
these certificates, and specialize along arbitrary ring homomorphisms. Identities between
power series (commutativity, associativity, and the various identifications of charts) are
proved by evaluating at all points of the polydisk and interpolating
(Theorem~\ref{thm:interpolation}).

\emph{Coleman integration is replaced by the formal logarithm.} The logarithm of the
formal group law $F_\tau$ (Chapter~2, Theorem~\ref{thm:log-iso}) provides the homomorphism
$\Lam$ on the kernel of reduction (Definition~\ref{def:Lam}). Its differential pulls back
along the curve-to-chart maps to explicit rational differentials $\omega'_1, \dots,
\omega'_g$ on the curve, which are shown --- by a closed trace-form computation, with no
Riemann--Roch --- to be a basis of the regular differentials, and whose disk expansions
integrate $\Lam$ on each residue disk (Theorem~\ref{thm:omega-prime-global}). This is the
substitute for Coleman's integration theory, and Theorem~\ref{thm:glue} is its summary
interface: the pairing $\lambda$, its linearity, and the glue identity expressing
$\lambda_\omega$ on a same-disk difference as a difference of antiderivative values.

\subsection*{Guide to the chapters}

\emph{Chapter 1} develops the class-group model over an arbitrary field: the Dedekind
property (Theorem~\ref{thm:hypring-dedekind}), Mumford representation with existence and
uniqueness of reduced representatives (Lemmas~\ref{lem:mumford-reduce},
\ref{lem:mumford-exists}), and base change (Definition~\ref{def:basechange}).

\emph{Chapter 2} contains the $p$-adic analysis, independent of the curve: evaluation of
power series on the maximal ideal (Definition~\ref{def:eval}), Weierstrass preparation and
zero counts (Theorem~\ref{thm:wprep}, Corollary~\ref{cor:wprep-count}), Strassmann's
theorem (Theorem~\ref{thm:strassmann}), the zero bound for antiderivatives
(Theorem~\ref{thm:antideriv-zeros}), the interpolation theorem
(Theorem~\ref{thm:interpolation}), formal group laws with their invariant differentials
and logarithms (Definitions~\ref{def:fgl}--\ref{def:fgl-log},
Theorem~\ref{thm:log-iso}), and Hensel's lemma in the smooth-point form
(Theorem~\ref{thm:hensel}).

\emph{Chapter 3} constructs the residue disks and the reduction of classes: the disk
parametrizations (Lemma~\ref{lem:disk-param}), integral expansions of differentials with
the order matching $\ord_{\bar P}(\bar\omega) = \ord_z(\bar w)$
(Lemmas~\ref{lem:omega-expansion}, \ref{lem:ord-match}), the vanishing budget
$\sum m_{\bar P} \le 2g-2$ (Corollary~\ref{cor:sum-ord}) and its refinement to a space of
differentials (Corollary~\ref{cor:sum-ord-space}), the specialization theorem for
principal divisors (Theorem~\ref{thm:spec}), and the reduction homomorphism $\redJ$
(Theorem~\ref{thm:redJ-hom}).

\emph{Chapter 4} is the heart: the auxiliary class $\tau$ and its genericity --- proved
for every genus $g \ge 1$ through the point-by-point fold and its cuspidal degeneration
(Definition~\ref{def:centre-fold} through Theorem~\ref{thm:fold-witness}) --- the
translated formal group law $F_\tau$ (Theorem~\ref{thm:Ftau-assoc}), the chart
isomorphism for the kernel of reduction (Theorem~\ref{thm:chart-mem}), the logarithm
$\Lam$ (Definition~\ref{def:Lam}), the differentials $\omega'_i$ with the trace-form
computation of their regularity (Theorem~\ref{thm:omega-prime-global}), the genericity of
the global configuration (Lemma~\ref{lem:global-generic}), and the interface theorem for
the pairing $\lambda$ (Theorem~\ref{thm:glue}).

\emph{Chapter 5} concludes: the space of annihilating differentials from the rank
hypothesis (Theorem~\ref{thm:annihilator}), its integral numerators
(Lemma~\ref{lem:normalize}),
the per-disk bound (Theorem~\ref{thm:disk-bound}), and the main theorem with its
arithmetic instantiations (Theorem~\ref{thm:main}, Corollary~\ref{cor:finiteness}).

\emph{Chapter 6} instantiates the arithmetic base over an arbitrary number field $K$: for
a prime $\mathfrak p$ of good reduction with ramification index $e$, the ring playing the
role of the completed maximal unramified extension of $K_{\mathfrak p}$ is constructed as
a direct factor of $W(\Fpbar) \otimes_{\ZZ} O_K$ (Definition~\ref{def:local-factor},
Lemma~\ref{lem:factor-dvr}), giving the bound
$\#C(K) \le \#\bar C(\Fq) + 2r$ for $p > 2r + e + 1$ (Theorem~\ref{thm:main-nf}).

\emph{Chapter 7} develops $2$-descent over an arbitrary field~$K$ in which $2 \ne 0$, for $f$
squarefree of odd degree and not necessarily monic: the classes of Mumford ideals of arbitrary
pairs of polynomials (Definition~\ref{def:hv-class}), the $x - T$ map~$\mu$
(Definition~\ref{def:ds-mu}, Theorem~\ref{thm:ds-formula}), certificates for halving points
(Theorem~\ref{thm:hv-general}, Corollary~\ref{cor:hv-record}), the kernel $\ker\mu = 2J(K)$
(Theorem~\ref{thm:ds-kernel}) and the rank formula (Theorem~\ref{thm:ds-rank}), the image of~$\mu$
over the fraction field of a Dedekind domain (Theorem~\ref{thm:ds-image-selmer}), the $2$-Selmer
group and the rank bound over a number field (Theorems~\ref{thm:ds-sel2}
and~\ref{thm:ds-rank-bound}), with the Mordell--Weil theorem and Mattuck's theorem as cited inputs
(Definition~\ref{def:ds-cited}), and two examples over~$\QQ$ (Theorems~\ref{thm:ds-e1}
and~\ref{thm:ds-e2}).

\subsection*{Notation and standing hypotheses}

Throughout: $f \in \ZZ[x]$ is monic and squarefree with $\deg f = 2g+1$, $g \ge 1$ (the
rank hypothesis of the main theorem forces $g \ge 1$, and no lower genus bound is used
anywhere); $C \colon y^2 = f(x)$ with its point $\infty$; $p$ is a prime with
$p \nmid \operatorname{disc}(f)$ and $p > 2r + 2$, where $r < g$ is the bound on the rank
in the main theorem (so $p$ is odd and $\bar f \in \Fp[x]$ is squarefree of degree
$2g+1$; the classical condition $p > 2g$ is the case $r = g-1$). Over an arbitrary arithmetic base
(Theorem~\ref{thm:main-arith}) the condition is $p > 2r + e + 1$, with $e$ bounding the
absolute ramification index; over $\QQ$, $e = 1$. Fields named $K$ are of characteristic $0$ or of
characteristic $p$ as indicated; the results of Chapter~1 hold over an arbitrary field of
characteristic $\ne 2$ unless stated otherwise. $W(\Fq)$ denotes the Witt vectors,
$\Oq = W(\Fq)$ and $\Qq = \Oq[1/p]$ for $q = p^N$; in the formalization, $\Oq$ is an
abstract complete discrete valuation ring of characteristic $0$ with algebraically closed
residue field of characteristic $\neq 2$.

Statements are given in the form and generality in which they are formalized. Remarks
labelled \emph{Formalization} record, concisely, how a statement or proof is realized in
Lean; the typewriter-font names in them are the corresponding Lean declarations, and the
lists attached to each node are machine-checked against the code. Some statements are
formalized although the main theorems do not depend on them; their Lean declarations are in
the modules \texttt{HyperellipticCurves.Archive.*}, and a remark says so at each of them.

\chapter{The Jacobian as a class group}
This chapter develops the class-group model of the Jacobian over an arbitrary field $K$
with $2 \in K^\times$; we write $f_K$ for the image of $f$ in $K[x]$.

\begin{definition}[Hyperelliptic curve]
\label{def:hypcurve}
\lean{HyperellipticCurve, HyperellipticCurve.genus, HyperellipticCurve.IsOdd,
  HyperellipticCurve.map, HyperellipticCurve.ofSquarefree, HyperellipticCurve.ofMonicOdd,
  Polynomial.discrOfDegree, Polynomial.discrOfDegree_add_two}
\leanok
Let $R$ be a commutative ring in which $2$ is invertible. A hyperelliptic curve of genus $g$
over $R$ is given by a polynomial $f \in R[x]$ of degree at most $2g + 2$ such that the
discriminant of the binary form $Z^{2g+2} f(X/Z)$ of degree $2g+2$ is a unit of $R$; we write
$C \colon y^2 = f(x)$. When $R \ne 0$, the degree of $f$ is $2g+1$ or $2g+2$, and in the first
case the leading coefficient of $f$ is a unit. The model is of odd degree if $f$ is monic of
degree $2g+1$; the curve then has exactly one point at infinity. Every ring homomorphism
$\varphi \colon R \to S$ (with $2$ invertible in $S$) maps $C$ to the curve
$y^2 = f^\varphi(x)$ of the same genus over $S$, and a nonconstant squarefree polynomial over a
perfect field defines a curve.
\end{definition}

\begin{remark}[Formalization]
The discriminant of the binary form is \texttt{Polynomial.discrOfDegree}: Mathlib's
discriminant matrix with the degree replaced by $2g+2$, so that it is a polynomial in the
coefficients of $f$ and commutes with every ring homomorphism. It equals the discriminant of $f$
when $\deg f = 2g+2$ and $\operatorname{lc}(f)^2$ times it when $\deg f = 2g+1$; the general
form of the latter is the drop-by-one formula $D_{n+2}(f) = a_{n+1}^2\, D_{n+1}(f)$ for
$\deg f \le n+1$ (\texttt{Polynomial.discrOfDegree\_add\_two}), proved by expanding Mathlib's
modified Sylvester matrix along its normalized last row. The invertibility of $2$ is a parameter
of the type (\texttt{[Invertible (2 : R)]}); the genus is part of the data, so base change needs
no hypothesis. The coordinate ring, function field and Jacobian below are still indexed by the
polynomial $f$; a curve supplies their hypotheses (odd degree, squarefreeness and the
non-square property over a field) through instances.
\end{remark}

\begin{definition}[Coordinate ring and function field]
\label{def:hypring}
\lean{ChabautyColeman.HypRing, ChabautyColeman.HypField,
  ChabautyColeman.HypRing.iota, ChabautyColeman.HypField.iota}
\leanok
$\A{K} := K[x][Y]/(Y^2 - f_K)$, written $K[x,y]$, and
$\F{K} := K(x)[Y]/(Y^2 - f_K)$. The hyperelliptic involution $\iota$ is the $K[x]$-
(resp.\ $K(x)$-) algebra automorphism with $\iota(y) = -y$.
\end{definition}

\begin{lemma}[Irreducibility]
\label{lem:hyp-irred}
\lean{ChabautyColeman.irreducible_Ysq_sub_f, ChabautyColeman.not_isSquare_algebraMap_ratFunc,
  ChabautyColeman.instFactNotIsSquareOfOdd, IsIntegrallyClosed.isSquare_algebraMap_iff,
  HyperellipticCurve.instFactNotIsSquareF, Squarefree.isUnit_of_isSquare,
  ChabautyColeman.instFactHypField, QuadraticAlgebra.finrank_eq_two,
  ChabautyColeman.HypRing.toHypField, ChabautyColeman.HypRing.toHypField_injective,
  ChabautyColeman.HypRing.instIsDomain}
\leanok
\uses{def:hypring}
For every field $K$, the polynomial $Y^2 - f_K$ is irreducible over $K(x)$; hence $\F{K}$
is a field, quadratic over $K(x)$, and $\A{K}$ is a domain, as it embeds into $\F{K}$.
(The extension is moreover Galois with group $\{1,\iota\}$; this refinement is not used
below.)
\end{lemma}
\begin{proof}
\leanok
A root would give $f_K = r^2$ with $0 \ne r \in K(x)$; comparing degrees,
$\deg f = 2\deg r$ is even, contradicting $\deg f = 2g+1$. \emph{Formalization.} The
statements assume only that~$f_K$ is not a square in~$K[x]$: as $K[x]$ is integrally closed,
$f_K$ is then not a square in~$K(x)$ either. One instance derives this hypothesis from the odd
degree, by the parity argument on \texttt{RatFunc.intDegree}, and another one from a curve over
a field, whose polynomial is squarefree of positive degree.
\end{proof}

\begin{lemma}[Norm form]
\label{lem:normform}
\lean{ChabautyColeman.HypField.norm_eq, ChabautyColeman.HypField.trace_eq,
  ChabautyColeman.HypField.mul_star_eq_norm, ChabautyColeman.HypRing.norm_eq,
  ChabautyColeman.HypRing.trace_eq, ChabautyColeman.HypRing.mul_star_eq_norm}
\leanok
\uses{lem:hyp-irred}
For $\alpha = a + by \in \F{K}$ one has
$\Nm_{\F{K}/K(x)}(\alpha) = \alpha\,\iota(\alpha) = a^2 - b^2 f_K$
and $\operatorname{Tr}(\alpha) = 2a$. The same formulas hold for the norm and trace of
$\A{K}$ over $K[x]$ (and over any base ring).
\end{lemma}
\begin{proof}
\leanok
Computation in the basis $(1,y)$.
\end{proof}

\begin{lemma}[Integral elements]
\label{lem:integral-char}
\lean{ChabautyColeman.HypField.isIntegral_iff, ChabautyColeman.HypField.isIntegral_y,
  ChabautyColeman.HypRing.isIntegral_toHypField,
  RatFunc.exists_algebraMap_eq_of_sq_mul_squarefree}
\leanok
\uses{lem:hyp-irred, lem:normform}
If $f_K$ is squarefree (and $2 \in K^\times$), then $a + by \in \F{K}$ is integral over
$K[x]$ if and only if $a, b \in K[x]$, i.e.\ if and only if $a + by$ lies in the image of
$\A{K}$.
\end{lemma}
\begin{proof}
\leanok
($\Leftarrow$) $y$ satisfies $Y^2 = f_K$. ($\Rightarrow$) $\alpha$ and $\iota\alpha$ are
integral, so $2a \in K[x]$ and $a^2 - b^2 f_K \in K[x]$ ($K[x]$ being integrally closed);
hence $b^2 f_K \in K[x]$ forces $b \in K[x]$: writing $b = p/q$ in lowest terms with $q$
monic, $p^2 f_K = c\,q^2$ and $\gcd(p,q) = 1$ give $q^2 \mid f_K$, so $q = 1$ by
squarefreeness. (This lowest-terms argument replaces a per-prime valuation computation.)
\end{proof}

\begin{theorem}[The coordinate ring is Dedekind]
\label{thm:hypring-dedekind}
\lean{ChabautyColeman.HypRing.instIsIntegralClosure,
  ChabautyColeman.HypRing.instIsDedekindDomain, ChabautyColeman.HypField.instIsSeparable,
  Algebra.IsSeparable.of_finrank_le_two}
\leanok
\uses{lem:integral-char}
If $f_K$ is squarefree, then $\A{K}$ is the integral closure of $K[x]$ in $\F{K}$ and is a
Dedekind domain, free over $K[x]$ with basis $(1,y)$.
\end{theorem}
\begin{proof}
\leanok
\uses{lem:integral-char, lem:hyp-irred}
By Lemma~\ref{lem:integral-char} the integral closure is $K[x] \oplus K[x]\,y = \A{K}$;
the integral closure of the principal ideal domain $K[x]$ in the finite separable
extension $\F{K}/K(x)$ is Dedekind. Separability holds because the extension is quadratic
and $2 \in K^\times$: any extension of degree $\le 2$ in characteristic $\ne 2$ is
separable. Freeness with basis $(1,y)$ is immediate from the construction.
\emph{Formalization.} The Dedekind property is Mathlib's
\texttt{IsIntegralClosure.isDedekindDomain}.
\end{proof}

\begin{lemma}[Units]
\label{lem:hypring-units}
\lean{ChabautyColeman.HypRing.isUnit_iff}
\leanok
\uses{lem:normform}
$(\A{K})^\times = K^\times$ for every field $K$. Only the odd degree of $f$ is used; no
assumption on the characteristic or on squarefreeness is needed.
\end{lemma}
\begin{proof}
\leanok
$\Nm(u) = a^2 - b^2 f_K \in K[x]^\times = K^\times$; if $b \ne 0$, the right-hand side of
$a^2 = c + b^2 f_K$ has odd degree $2\deg b + 2g + 1$, while the left-hand side has even
degree. So $b = 0$ and $a$ is constant.
\end{proof}

\begin{definition}[Jacobian and Abel--Jacobi map]
\label{def:hypjac}
\lean{ChabautyColeman.HypJac, ChabautyColeman.aj, ChabautyColeman.pointIdeal}
\leanok
\uses{thm:hypring-dedekind}
$\J{K} := \Cl(\A{K})$. For an affine point $P = (x_0,y_0) \in C(K)$ set
$\aj(P) := [\mathfrak m_P] \in \J{K}$, where $\mathfrak m_P = (x - x_0,\, y - y_0)$ is the
maximal ideal of $P$, and $\aj(\infty) := 0$. More generally, the class of a $K$-rational
divisor supported on affine points is the class of the corresponding fractional ideal.
\end{definition}

\begin{remark}[Formalization]
In Lean, $\Cl$ is Mathlib's multiplicative \texttt{ClassGroup}, so the identity class is
written $1$; classes of general ideals are obtained through \texttt{ClassGroup.mk0} and
need no separate definition.
\end{remark}

\begin{lemma}[The involution acts as inversion on point classes]
\label{lem:aj-inverse}
\lean{ChabautyColeman.aj_inv, ChabautyColeman.pointIdeal_mul_pointIdeal_neg}
\leanok
\uses{def:hypjac, lem:mumford-norm}
$\divisor(x - x_0)_{\mathrm{aff}} = \mathfrak m_P \cdot \mathfrak m_{\iota P}$, hence
$\aj(\iota P) = -\aj(P)$. (That $\iota$ inverts \emph{all} of $\J{K}$ is
Lemma~\ref{lem:iota-norm}.)
\end{lemma}
\begin{proof}
\leanok
This is the special case $u = x - x_0$, $v = y_0$ of the Mumford norm identity
(Lemma~\ref{lem:mumford-norm}), since classes of principal ideals vanish.
\end{proof}

\begin{lemma}[Conjugate ideal and norm]
\label{lem:iota-norm}
\lean{ChabautyColeman.mul_map_star_eq_map_relNorm, ChabautyColeman.classGroup_mk0_map_star,
  ChabautyColeman.classGroup_mk0_mumfordIdeal_map_star, Squarefree.dvd_of_sq_dvd_mul_sq}
\leanok
\uses{thm:hypring-dedekind, lem:normform, lem:integral-char}
For every ideal $\mathfrak a$ of $\A{K}$:
$\mathfrak a \cdot \iota(\mathfrak a) = (\Nm \mathfrak a)\,\A{K}$, where
$\Nm \mathfrak a \subseteq K[x]$ is the relative ideal norm. Hence $\iota$ induces
inversion on $\J{K}$.
\end{lemma}
\begin{proof}
\leanok
Since $K[x]$ is a principal ideal domain, $\Nm \mathfrak a = (n_0)$ is principal, and no
factorization into primes is needed: the inclusion $\supseteq$ holds because
$N(a) = a\,\iota(a) \in \mathfrak a\,\iota(\mathfrak a)$ for $a \in \mathfrak a$. For
$\subseteq$, take $t = a\,\iota(b)$ with $a, b \in \mathfrak a$; then
$t + \iota t = N(a+b) - N(a) - N(b) \in (n_0)$ gives $n_0 \mid 2\,t_{\mathrm{re}}$, and
$t\,\iota t = N(a)N(b) \in (n_0^2)$ gives
$n_0^2 \mid t_{\mathrm{re}}^2 - f\,t_{\mathrm{im}}^2$; since $2 \in K^\times$ and $f$ is
squarefree, $n_0$ divides both components of $t$ (using that $n^2 \mid f b^2$ with $f$
squarefree forces $n \mid b$ in a unique factorization domain).
\end{proof}

\begin{remark}[Formalization]
Lemma~\ref{lem:iota-norm} is formalized in
\texttt{HyperellipticCurves.Archive.ConjugateIdealNorm}; the main theorems do not depend on it,
since both places that invert a class invert a \emph{Mumford} ideal, where the Mumford norm
identity of Lemma~\ref{lem:mumford-norm} gives the inversion directly.
\end{remark}

\begin{definition}[Mumford ideals]
\label{def:mumford-ideal}
\lean{ChabautyColeman.mumfordIdeal, ChabautyColeman.mumfordHom,
  ChabautyColeman.mumfordQuotEquiv, ChabautyColeman.ker_mumfordHom}
\leanok
\uses{def:hypjac}
For $u, v \in K[x]$, the \emph{Mumford ideal} is
$\mathfrak a_{u,v} := (u,\, y - v) \subseteq \A{K}$; point ideals are the case
$(x - x_0,\, y - y_0)$. If $u \mid v^2 - f_K$, then the evaluation
$p + qy \mapsto p + qv \bmod u$ is a surjective ring homomorphism
$\A{K} \to K[x]/(u)$ with kernel $\mathfrak a_{u,v}$, so
$\A{K}/\mathfrak a_{u,v} \cong K[x]/(u)$; in particular $\mathfrak a_{u,v}$ is a
proper nonzero ideal for nonconstant $u$.
\end{definition}

\begin{lemma}[Mumford norm identity]
\label{lem:mumford-norm}
\lean{ChabautyColeman.mumfordIdeal_mul_map_star, ChabautyColeman.map_star_mumfordIdeal}
\leanok
\uses{def:mumford-ideal}
If $u \mid v^2 - f_K$ with $f_K$ squarefree (and $2 \in K^\times$), then
$\iota(\mathfrak a_{u,v}) = \mathfrak a_{u,-v}$ and
$\mathfrak a_{u,v} \cdot \iota(\mathfrak a_{u,v}) = (u)$.
\end{lemma}
\begin{proof}
\leanok
$\subseteq$: all four generator products are divisible by $u$, using
$(y-v)(y+v) = f_K - v^2 = -uw$ with $w := (v^2-f_K)/u$. $\supseteq$: a common prime $\pi$
of $(u, 2v, w)$ would give $\pi \mid v$ and $\pi^2 \mid uw = v^2 - f_K$, so
$\pi^2 \mid f_K$, contradicting squarefreeness; a B\'ezout identity then exhibits $u$ as
a combination of $u^2$, $u\cdot 2v$ and $uw$, all of which lie in the product.
\end{proof}

\begin{lemma}[The divisor of $y - v$]
\label{lem:mumford-divisor-y}
\lean{ChabautyColeman.mumfordIdeal_mul_cofactor,
  ChabautyColeman.classGroup_mk0_mumfordIdeal_cofactor}
\leanok
\uses{def:mumford-ideal, def:hypjac}
If $v^2 - f_K = u\,w$ (with $f_K$ squarefree, $2 \in K^\times$), then
$\mathfrak a_{u,v} \cdot \mathfrak a_{w,v} = (y - v)$ is principal; hence
$[\mathfrak a_{w,v}] = [\mathfrak a_{u,v}]^{-1}$ in $\J{K}$. Combined with
Lemma~\ref{lem:mumford-norm} this yields Cantor's reduction step:
$[\mathfrak a_{u,v}] = [\mathfrak a_{w,-v}]$ with
$\deg w \le \max(2\deg v, 2g+1) - \deg u$.
\end{lemma}
\begin{proof}
\leanok
$\subseteq$: the generator products are divisible by $y - v$, using
$uw = v^2 - f_K = -(y-v)(y+v)$. $\supseteq$: the same B\'ezout identity for
$(u, 2v, w)$ as in Lemma~\ref{lem:mumford-norm} exhibits $y - v$ in the product,
using $(y-v)\cdot 2v = -uw - (y-v)^2$.
\end{proof}

\begin{lemma}[Semi-reduction]
\label{lem:mumford-semireduced}
\lean{ChabautyColeman.exists_span_singleton_mul_mumfordIdeal,
  ChabautyColeman.exists_mk0_eq_mk0_mumfordIdeal}
\leanok
\uses{def:mumford-ideal, lem:normform, def:hypjac}
Every nonzero ideal $\mathfrak a$ of $\A{K}$ has the form
$\mathfrak a = (c)\,\mathfrak a_{u,v}$ with $c, u, v \in K[x]$, $c, u \ne 0$ and
$u \mid v^2 - f_K$; hence every class in $\J{K}$ is the class of a Mumford ideal.
\end{lemma}
\begin{proof}
\leanok
Hermite normal form over the principal ideal domain $K[x]$, with no factorization into
primes: let $c$ generate the image of $\mathfrak a$ under the coefficient projection
$a + by \mapsto b$ and let $a_0$ generate $\mathfrak a \cap K[x]$; multiplying elements
by $y$ shows $c \mid a_0$ and that $c$ divides the $K[x]$-part of a lift $z_0$ of $c$,
giving $\mathfrak a = (a_0, z_0) = (c)(u, y - v)$ with $a_0 = cu$, $z_0 = c(y-v)$.
Stability of $\mathfrak a$ under multiplication by $y$, applied to $z_0$, yields
$u \mid v^2 - f_K$. (Nonvanishing of $a_0$ uses the norm form:
$0 \ne \Nm(z) \in \mathfrak a \cap K[x]$.)
\end{proof}

\begin{lemma}[Existence of reduced representatives]
\label{lem:mumford-reduce}
\lean{ChabautyColeman.exists_mk0_mumfordIdeal_reduced}
\leanok
\uses{lem:mumford-semireduced, lem:mumford-divisor-y, lem:mumford-norm}
Every class in $\J{K}$ is the class of a \emph{reduced} Mumford ideal
$\mathfrak a_{u,v}$: $u$ monic, $\deg v < \deg u \le g$, $u \mid v^2 - f_K$.
\end{lemma}
\begin{proof}
\leanok
By Lemma~\ref{lem:mumford-semireduced} the class contains a Mumford ideal. The Mumford
ideal is unchanged by unit rescaling of $u$ and by replacing $v$ with $v \bmod u$;
while $\deg u > g$, the class-preserving replacement
$[\mathfrak a_{u,v}] = [\mathfrak a_{w,-v}]$ with $v^2 - f_K = uw$
(Lemmas~\ref{lem:mumford-divisor-y} and~\ref{lem:mumford-norm}) strictly decreases
$\deg u$: after reducing $v$ modulo $u$,
$\deg w \le \max(2\deg v, 2g+1) - \deg u < \deg u$. Conclude by induction on $\deg u$.
\end{proof}

\begin{lemma}[Uniqueness of reduced representatives]
\label{lem:mumford-exists}
\lean{ChabautyColeman.mumfordIdeal_reduced_unique, ChabautyColeman.mem_mumfordIdeal_iff,
  ChabautyColeman.comap_mumfordIdeal}
\leanok
\uses{thm:hypring-dedekind, lem:mumford-norm, def:mumford-ideal,
  lem:mumford-norm, lem:mumford-divisor-y, lem:mumford-reduce}
The reduced Mumford ideal representing a class (Lemma~\ref{lem:mumford-reduce}) is
unique. (Cantor's composition-and-reduction algorithm computes the reduced
representative of a product of two reduced classes; existence is
Lemma~\ref{lem:mumford-reduce}.)
\end{lemma}
\begin{proof}
\leanok
Let $\mathfrak a_1, \mathfrak a_2$ be reduced Mumford ideals in the same class. Then
$\mathfrak a_1\,\iota(\mathfrak a_2)$ is an integral principal ideal $(h)$, and
$(\Nm h) = (u_1 u_2)$ by the norm identity, of degree $\le 2g$. Since $\Nm(a + by) =
a^2 - b^2 f_K$ has degree $\ge 2g+1$ whenever $b \ne 0$ (parity), $h = a \in K[x]$.
Multiplying $(h) = \mathfrak a_1 \iota(\mathfrak a_2)$ and its conjugate by
$\mathfrak a_2$ resp.\ $\mathfrak a_1$ and taking traces to $K[x]$ (the trace of
$\mathfrak a_{u,v}$ is $(u)$, immediate from the quotient description) gives
$u_1 \sim a \sim u_2$, hence $u_1 = u_2$ by monicity; cancelling the principal factor
yields $\mathfrak a_1 = \mathfrak a_2$, and then $u \mid v_2 - v_1$ with
$\deg(v_2 - v_1) < \deg u$ forces $v_1 = v_2$. (The norm-degree parity argument
replaces the usual pole count at the place at infinity.)
\end{proof}

\begin{definition}[Base change]
\label{def:basechange}
\lean{ChabautyColeman.HypJac.map, ChabautyColeman.HypRing.baseChange,
  ChabautyColeman.HypJac.map_aj, ChabautyColeman.pointIdeal_map}
\leanok
\uses{def:hypjac}
For a field extension $K \subseteq K'$, the homomorphism $\J{K} \to \J{K'}$ induced by
ideal extension along $\A{K} \to \A{K'}$. It is compatible with the Abel--Jacobi maps.
\emph{Formalization:} this is Mathlib's \texttt{ClassGroup.extendedHom} applied to the
base change of coordinate rings.
\end{definition}

\begin{lemma}[Injectivity criterion]
\label{lem:basechange-crit}
\lean{ChabautyColeman.hypJac_map_injective_iff}
\leanok
\uses{def:basechange, def:hypjac}
$\J{K} \to \J{K'}$ is injective if and only if every nonzero ideal of $\A{K}$ whose
extension to $\A{K'}$ is principal is itself principal.
\end{lemma}
\begin{proof}
\leanok
Kernel triviality for the group homomorphism, using the surjectivity of
$\mathfrak a \mapsto [\mathfrak a]$ and the compatibility of the base-change map with
ideal extension.
\end{proof}

\begin{lemma}[Contraction]
\label{lem:basechange-contract}
\lean{ChabautyColeman.comap_map_baseChange}
\leanok
\uses{def:basechange}
For every ideal $\mathfrak a$ of $\A{K}$ and every field extension $K \subseteq K'$:
$(\mathfrak a \A{K'}) \cap \A{K} = \mathfrak a$.
\end{lemma}
\begin{proof}
\leanok
Choose a $K$-linear splitting $p \colon K' \to K$ of the inclusion and apply it to all
polynomial coefficients; the resulting additive retraction
$\pi \colon \A{K'} \to \A{K}$ satisfies $\pi(\iota(a)w) = a\,\pi(w)$, and applying
it to an element of $\mathfrak a \A{K'} \cap \A{K}$, written as a combination
$\sum \iota(a_j) w_j$, lands in $\mathfrak a$. No flatness theory is needed.
\end{proof}

\begin{lemma}[Galois action and the cocycle]
\label{lem:basechange-galact}
\lean{ChabautyColeman.HypRing.galAct, ChabautyColeman.HypRing.galAct_algebraMap,
  ChabautyColeman.HypRing.exists_unit_galAct_eq}
\leanok
\uses{def:basechange, lem:hypring-units}
A $K$-automorphism $\sigma$ of $K'$ acts on $\A{K'}$ coefficientwise, fixing the
image of $\A{K}$ and stabilizing extended ideals. If
$\mathfrak a\,\A{K'} = (h)$ is principal, then $\sigma(h) = c_\sigma\, h$ for a
constant $c_\sigma \in (K')^\times$.
\end{lemma}
\begin{proof}
\leanok
$(\sigma(h)) = \sigma(\mathfrak a \A{K'}) = \mathfrak a \A{K'} = (h)$, so
$\sigma(h)$ and $h$ are associated; the unit is a constant by
Lemma~\ref{lem:hypring-units}.
\end{proof}

\begin{lemma}[Fixed-point descent]
\label{lem:basechange-fixed}
\lean{ChabautyColeman.HypRing.exists_algebraMap_eq_of_forall_galAct}
\leanok
\uses{lem:basechange-galact}
For a finite Galois extension $K \subseteq K'$, an element of $\A{K'}$ fixed by the
full Galois action lies in the image of $\A{K}$.
\end{lemma}
\begin{proof}
\leanok
Coefficientwise from $(K')^{\mathrm{Gal}(K'/K)} = K$.
\end{proof}

\begin{lemma}[Descent of principality along Galois extensions]
\label{lem:basechange-galois}
\lean{ChabautyColeman.HypRing.isPrincipal_of_map_isPrincipal,
  ChabautyColeman.hypJac_map_injective_of_isGalois}
\leanok
\uses{lem:basechange-crit, lem:basechange-contract, lem:basechange-galact,
  lem:basechange-fixed}
Let $K \subseteq K'$ be a finite Galois extension. If an ideal $\mathfrak a$ of $\A{K}$
extends to a principal ideal of $\A{K'}$, then $\mathfrak a$ is principal. Consequently
(by Lemma~\ref{lem:basechange-crit}) the map $\J{K} \to \J{K'}$ is injective.
\end{lemma}
\begin{proof}
\leanok
Let $\mathfrak a\,\A{K'} = (h)$; then $h \ne 0$ unless $\mathfrak a = 0$ (which is
trivially principal). Lemma~\ref{lem:basechange-galact} gives
$c \colon \mathrm{Gal}(K'/K) \to (K')^\times$ with $\sigma(h) = c_\sigma h$; cancelling
$h$ from the two evaluations of $(\sigma\tau)(h)$ shows $c$ is a $1$-cocycle. By
Hilbert's Theorem~90 it is a coboundary, $c_\sigma = \sigma(\beta)/\beta$ for some
$\beta \in (K')^\times$. Then $h' = \beta^{-1} h$ is a Galois-fixed generator of
$\mathfrak a\,\A{K'}$, so $h'$ comes from some $w \in \A{K}$ by
Lemma~\ref{lem:basechange-fixed}. Contracting via Lemma~\ref{lem:basechange-contract}
gives $\mathfrak a = \mathfrak a\,\A{K'} \cap \A{K} = w\A{K'} \cap \A{K} = w\A{K}$.
\end{proof}

\begin{theorem}[Base-change injectivity]
\label{thm:basechange-inj}
\lean{ChabautyColeman.hypJac_map_injective, ChabautyColeman.mumfordIdeal_map}
\leanok
\uses{def:basechange, lem:mumford-reduce, lem:mumford-exists}
For \emph{every} field extension $K \subseteq K'$ (with the standing hypotheses on $f$
over both fields), the map $\J{K} \to \J{K'}$ is injective. In particular this applies
to $K' \in \{\QQ_p, \Qq, \Cp\}$ over $K = \QQ$; no finiteness, separability, or
characteristic assumption on the extension is needed.
\end{theorem}
\begin{proof}
\leanok
\uses{lem:mumford-reduce, lem:mumford-exists}
Represent two classes with equal images by their reduced Mumford representatives
$(u_1, v_1)$ and $(u_2, v_2)$ (Lemma~\ref{lem:mumford-reduce}). Base change sends
Mumford ideals to Mumford ideals and preserves the reducedness conditions (monicity,
the degree bounds, $u \mid v^2 - f$), so the base-changed pairs are reduced
representatives of the common image class, and by uniqueness
(Lemma~\ref{lem:mumford-exists}) they coincide over $K'$. Since all coefficients lie in
$K$ and $K \to K'$ is injective, $(u_1, v_1) = (u_2, v_2)$, so the classes agree.
\end{proof}

\begin{remark}[Formalization]
Base-change injectivity and its Galois-descent chain
(Lemmas~\ref{lem:basechange-crit}--\ref{lem:basechange-galois},
Theorem~\ref{thm:basechange-inj}) are formalized in
\texttt{HyperellipticCurves.Archive.BaseChangeInjectivity}; the main theorems do not depend on
them, since the main development only uses the base-change \emph{homomorphism} of
Definition~\ref{def:basechange}. The comparison of reduced representatives in the proof of
Theorem~\ref{thm:basechange-inj} needs no descent and works for any field extension, whereas
descent through finite Galois subextensions (Hilbert~90) covers only algebraic extensions: a
transcendental extension such as $\QQ_p/\QQ$ would require an additional specialization
argument.
\end{remark}

\chapter{$p$-adic analysis and formal group laws}
Throughout this chapter, $O$ is a complete discrete valuation ring with fraction field of
characteristic $0$, residue characteristic $p$, maximal ideal $\mathfrak m = (\pi)$ and
normalized valuation $v$. A natural number $e$ with $v(p) \le e$ bounds the absolute
ramification index; the unramified case is $e = 1$, where $\pi = p$. Nothing in this chapter
mentions the curve.

\begin{definition}[Evaluation of power series]
\label{def:eval}
\lean{ChabautyColeman.PSeries.eval, ChabautyColeman.PSeries.hasSum_eval,
  ChabautyColeman.PSeries.hasEval_of_mem, PadicInt.isAdic}
\leanok
For $h = \sum a_k z^k \in O[[z]]$ and $z_0 \in \mathfrak m$, set
$h(z_0) := \sum a_k z_0^k$; the series converges since $v(a_k z_0^k) \to \infty$.
\emph{Formalization:} a thin layer over Mathlib's power-series evaluation; the standing
setting is a complete Hausdorff uniform ring $O$, local with the $\mathfrak m$-adic
topology, so that elements of $\mathfrak m$ are topologically nilpotent. The
identification of the metric topology of $\ZZ_p$ with the $p\ZZ_p$-adic one makes
$O = \ZZ_p$ an instance.
\end{definition}

\begin{lemma}[Evaluation calculus]
\label{lem:eval-hom}
\lean{ChabautyColeman.PSeries.evalRingHom, ChabautyColeman.PSeries.comp,
  ChabautyColeman.PSeries.eval_comp, ChabautyColeman.PSeries.constantCoeff_comp,
  ChabautyColeman.PSeries.eval_deriv, ChabautyColeman.PSeries.summable_of_mem_pow}
\leanok
\uses{def:eval}
The map $h \mapsto h(z_0)$ is a ring homomorphism; evaluation commutes with composition
$h \circ k$ for $k(0) \in \mathfrak m$; and the formal derivative controls first
differences: $h(z_0 + w) = h(z_0) + h'(z_0)\,w + r\,w^2$ for some $r \in O$ and all
$z_0, w \in \mathfrak m$. (That a power series vanishing identically on $\mathfrak m$ is
$0$ when the residue field is infinite is the one-variable case of
Theorem~\ref{thm:interpolation}.)
\end{lemma}
\begin{proof}
\leanok
The ring-homomorphism and composition statements follow from continuity of evaluation
(the composition is defined by evaluating $h$ at the point $k$ of $O[[z]]$ for the
product topology). For the Taylor statement, peel the first two terms off the binomial
expansion of $(z_0 + w)^n$ termwise; the remainder series has terms in
$\mathfrak m^{n-2}$, hence converges.
\end{proof}

\begin{theorem}[Weierstrass preparation]
\label{thm:wprep}
\lean{PowerSeries.exists_isWeierstrassFactorization,
  PowerSeries.IsWeierstrassFactorization.elim,
  ChabautyColeman.PSeries.instIsAdicComplete}
\leanok
\uses{def:eval}
If $h \in O[[z]]$ has $\bar h \ne 0$ with $m := \ord(\bar h) < \infty$, then $h = u W$
with $u \in O[[z]]^\times$ and $W \in O[z]$ monic of degree $m$ with
$W \equiv z^m \pmod{\mathfrak m}$, and this factorization is unique.
\end{theorem}
\begin{proof}
\leanok
This is in Mathlib for local rings that are adically complete with respect to the
maximal ideal (distinguished polynomials are by definition monic with lower coefficients
in $\mathfrak m$, and their degree is $\ord(\bar h)$). What remains is that the standing
topological hypotheses imply $\mathfrak m$-adic completeness.
\end{proof}

\begin{corollary}[Zeros in the open unit disk]
\label{cor:wprep-count}
\lean{ChabautyColeman.PSeries.zeros_le_ord, ChabautyColeman.PSeries.finite_setOf_eval_eq_zero,
  ChabautyColeman.PSeries.setOf_eval_eq_zero_eq, ChabautyColeman.PSeries.isUnit_eval_of_isUnit}
\leanok
\uses{thm:wprep, lem:eval-hom}
For $O$ a domain, such an $h$ has at most $m$ distinct zeros in $\mathfrak m$; the zeros
are exactly the roots of $W$, all of which lie in $\mathfrak m$. Since the statement
applies to every $O$ satisfying the standing hypotheses, it applies in particular to
every complete extension $O'/O$. (A refinement counting zeros with multiplicity is
possible, but only the count of distinct zeros is used: rational points inject into
zeros.)
\end{corollary}
\begin{proof}
\leanok
Values of a power series on $\mathfrak m$ are congruent to the constant coefficient
modulo $\mathfrak m$ (write $h - h(0) = z\,g$ and evaluate), so units of $O[[z]]$ do
not vanish on $\mathfrak m$: $h$ and $W$ have the same zeros there. Roots of a
distinguished polynomial lie in $\mathfrak m$ ($z^m \in \mathfrak m$ at a root, and
$\mathfrak m$ is prime), and a monic polynomial of degree $m$ over a domain has at most
$m$ roots; $\deg W = m$ by comparison with the reduction.
\end{proof}

\begin{theorem}[Strassmann]
\label{thm:strassmann}
\lean{ChabautyColeman.PSeries.strassmann, ChabautyColeman.PSeries.strassmann_finite,
  ChabautyColeman.PSeries.evalT, ChabautyColeman.PSeries.recenter,
  ChabautyColeman.PSeries.evalT_add_eq, ChabautyColeman.PSeries.map_recenter}
\leanok
\uses{thm:wprep, cor:wprep-count, lem:eval-hom}
Let $O$ be a domain and $h = \sum a_k z^k \in O[[z]]$ with $a_k \to 0$ (so that $h$
converges on all of $O$), $a_N$ a unit, and $a_k \in \mathfrak m$ for all $k > N$. Then
$h$ has at most $N$ zeros in $O$. (The decay hypothesis $a_k \to 0$ is necessary for
convergence outside $\mathfrak m$; finiteness of the number of unit coefficients alone
does not suffice.)
\end{theorem}
\begin{proof}
\leanok
Instead of preparation at radius $1$, which would require the Tate algebra, decompose
the zero set by residue classes: the reduction $\bar h$ is a polynomial $q$ of degree
$N$; zeros of $h$ lie over roots of $q$; over a root $\bar c$ with lift $c$, the zeros
of $h$ are the zeros on $\mathfrak m$ of the recentred series $h(c + z)$ (defined
coefficientwise by convergent series; the identity between $h(c+w)$ and the recentring
evaluated at $w$ is an exchange of summation). The reduction of $h(c + z)$ is the
Taylor shift $q(\bar c + z)$, whose vanishing order is the multiplicity of $q$ at
$\bar c$, so Corollary~\ref{cor:wprep-count} bounds the zeros in the class by that
multiplicity, and $\sum_{\bar c} \mathrm{mult}_{\bar c}(q) \le \deg q = N$.
\end{proof}

\begin{lemma}[An index-valuation inequality]
\label{lem:padic-index}
\lean{Nat.add_mul_multiplicity_lt, Nat.mul_multiplicity_le_of_lt}
\leanok
Let $p > m + e + 1$ and $k + 1 > m + 1$. Then $(k+1) - e\,v_p(k+1) > m + 1$. (No primality
of $p$ is needed; in Lean the inequality is stated without natural subtraction, as
$m + 1 + e\,v_p(k+1) < k+1$.)
\end{lemma}
\begin{proof}
\leanok
Let $j := v_p(k+1)$. For $j = 0$ this is the hypothesis $k + 1 > m + 1$. For $j \ge 1$,
$k + 1 \ge p^j \ge (m+e+2)^j \ge (m+e+2)\,j \ge m + 2 + e\,j$, the middle step because
$c\,j \le c^j$ for $c \ge 2$ and the last because $(m+2)(j-1) \ge 0$.
\end{proof}

\begin{theorem}[Zeros of antiderivatives]
\label{thm:antideriv-zeros}
\lean{ChabautyColeman.PSeries.antideriv_zeros_le,
  ChabautyColeman.PSeries.antideriv_zeros_finite,
  ChabautyColeman.PSeries.pow_dvd_coeff_antideriv,
  IsLocalRing.IsUniformizer.exists_pow_eq,
  Nat.tendsto_sub_mul_multiplicity_atTop}
\leanok
\uses{thm:strassmann, lem:padic-index}
Let $\mathfrak m = (\pi)$ with $v(p) \le e$, let
$w = \sum_{k\ge0} b_k z^k \in O[[z]]$ with $b_m$ a unit, and let
$G(z) = c + \sum_k \tfrac{b_k}{k+1} z^{k+1}$ be an antiderivative, $c \in O$. If
$p > m + e + 1$, then $G$ has at most $m + 1$ zeros in $\mathfrak m$. ($m$ need not be
the exact reduction order of $w$: any index with unit coefficient works, and
$m = \ord(\bar w)$ gives the sharpest bound. For $v(c) < 0$ the statement is trivial, so
requiring $c \in O$ loses nothing.)
\end{theorem}
\begin{proof}
\leanok
\uses{lem:padic-index, thm:strassmann}
Work with the substituted series $H(u) := G(\pi u)$, which has coefficients in $O$ and is
characterized division-free by $(k+1)\,\mathrm{coeff}_{k+1}(H) = \pi^{k+1} b_k$; the claim
becomes a bound on its zeros in all of $O$. Integrality is the divisibility
$(k+1) \mid \pi^{k+1}$, which holds because $v(k+1) = v(p)\,v_p(k+1) \le e\,v_p(k+1)$ and
$(e+2)\,v_p(n) \le (e+2)^{v_p(n)} \le p^{v_p(n)} \le n$ for $p \ge e+2$. The coefficient of
$H$ at index $k+1$ has valuation $v(b_k) + (k+1) - v(k+1)$, hence at least
$v(b_k) + (k+1) - e\,v_p(k+1)$: it equals $m+1$ at index $m+1$ (using $m + 1 < p$, so
integers prime to $p$ are units in $O$), exceeds $m+1$ for every index $> m+1$ by
Lemma~\ref{lem:padic-index}, and the coefficients tend to $0$ (by the same estimate, since
$(e+2)\,v_p(k+1) \le k+1$). Dividing by $\pi^{s}$ with $s$ the minimal valuation ($\le m+1$)
gives a series with unit content and top unit index $\le m+1$; Strassmann's theorem
bounds its zeros in $O$ by $m+1$.
\end{proof}

\begin{definition}[Multivariate evaluation]
\label{def:mv-eval}
\lean{ChabautyColeman.MvPSeries.eval, ChabautyColeman.MvPSeries.hasEval_of_mem,
  ChabautyColeman.MvPSeries.evalRingHom, ChabautyColeman.MvPSeries.hasSum_eval,
  MvPowerSeries.continuous_subst', ChabautyColeman.MvPSeries.eval_subst,
  ChabautyColeman.MvPSeries.hasEval_eval}
\leanok
\uses{def:eval}
For $h \in O[[z_1,\dots,z_n]]$ and $\mathbf z_0 \in \mathfrak m^n$, the value
$h(\mathbf z_0)$ is defined by the same estimates. Evaluation is a ring homomorphism and
is compatible with formal substitution: substitution is continuous for the product
topologies (each coefficient of a substitution is a fixed finite linear combination of
the input coefficients), so evaluating a substitution equals evaluating at the evaluated
family, which is itself a legitimate evaluation point.
\end{definition}

\begin{definition}[Evaluation of series with decaying coefficients]
\label{def:mv-evalT}
\lean{ChabautyColeman.MvPSeries.evalT, ChabautyColeman.MvPSeries.hasSum_evalT,
  ChabautyColeman.MvPSeries.evalT_eq_eval, ChabautyColeman.MvPSeries.evalT_add,
  ChabautyColeman.MvPSeries.tendsto_coeff_zero_of_mem_pow,
  Finsupp.tendsto_degree_cofinite}
\leanok
\uses{def:mv-eval}
For $h \in O[[z_1,\dots,z_n]]$ whose coefficients tend to $0$ (along the cofinite filter
on multi-indices) and \emph{any} family $\mathbf z \in O^n$ --- not restricted to
$\mathfrak m^n$ --- the sum $h(\mathbf z) = \sum_d (\mathrm{coeff}_d h)\,\mathbf z^d$
converges, since each term lies in a power of $\mathfrak m$ determined by
$\mathrm{coeff}_d h$, the point family being bounded. This total evaluation is additive
in $h$ and agrees with the evaluation of Definition~\ref{def:mv-eval} wherever the
latter converges. A convenient sufficient condition for the decay hypothesis:
$\mathrm{coeff}_d h \in \mathfrak m^{e(d)}$ with $e(d) \to \infty$ cofinitely, which
holds once $e$ grows with the total degree, since only finitely many multi-indices have
bounded degree. This is the $n$-variable analogue of the coefficient-decay evaluation
behind Strassmann's theorem; it is what evaluates the logarithm on~$\mathfrak m^n$ after
scaling by~$\pi$.
\end{definition}

\begin{theorem}[Interpolation]
\label{thm:interpolation}
\lean{ChabautyColeman.MvPSeries.eq_zero_of_forall_eval_zero,
  ChabautyColeman.PSeries.eq_zero_of_forall_eval_zero}
\leanok
\uses{def:mv-eval, cor:wprep-count}
Let $O$ be a domain of the standing setting with $\mathfrak m = (\pi)$ principal,
$\pi \ne 0$, and $O$ infinite. If $h \in O[[z_1,\dots,z_n]]$ vanishes at every point of
$\mathfrak m^n$, then $h = 0$. No condition on the residue field is needed; in
particular the statement applies directly over $\ZZ_p$.
\end{theorem}
\begin{proof}
\leanok
One variable: if the reduction of $h$ were nonzero, the zero set of $h$ in $\mathfrak m$
would be finite (Corollary~\ref{cor:wprep-count}), but it is all of the infinite set
$\mathfrak m$; hence $\pi$ divides all coefficients, and iterating, all coefficients lie
in $\bigcap_s \mathfrak m^s = 0$. Several variables, by induction on $n$: fixing all but
one variable, the coefficient slices collapse into a one-variable series computing the
evaluation (an exchange of summation), and the one-variable case applies
coefficientwise.
\end{proof}

\begin{definition}[Formal group law]
\label{def:fgl}
\lean{ChabautyColeman.FormalGroupLaw, ChabautyColeman.FormalGroupLaw.Points}
\leanok
\uses{def:mv-eval}
An $n$-dimensional commutative formal group law over $O$ is a tuple
$F \in O[[X;Y]]^n$ with $F(X,0)=X$, $F(0,Y)=Y$, $F(X,Y)=F(Y,X)$ and
$F(F(X,Y),Z)=F(X,F(Y,Z))$, as formal substitution identities. It makes $\mathfrak m^n$ a
commutative \emph{monoid} via Definition~\ref{def:mv-eval}: each monoid axiom is the
corresponding formal axiom evaluated, using that evaluation of a substitution is
evaluation at the evaluated family, and that values of series without constant term stay
in $\mathfrak m$. (Inverses, which would make this monoid a group, are never needed: the
group laws below parametrize residue polydisks of a group and inherit inverses from it.)
\end{definition}

\begin{remark}[Formalization]
The monoid of points (\texttt{ChabautyColeman.FormalGroupLaw.Points}) is formalized in
\texttt{HyperellipticCurves.Archive.FormalGroupPoints}; the main theorems do not depend on it,
since the main development works with the evaluated identities directly, in the $\pi$-scaled
form of Theorem~\ref{thm:log-iso}.
\end{remark}

\begin{lemma}[Formal partial derivatives]
\label{lem:pderiv}
\lean{MvPowerSeries.pderiv, MvPowerSeries.pderiv_subst, MvPowerSeries.continuous_subst'}
\leanok
Formal partial derivatives of multivariate power series satisfy the Leibniz rule and the
chain rule for formal substitution.
\end{lemma}
\begin{proof}
\leanok
The Leibniz rule holds by construction. For the chain rule, both sides are continuous for the
product topologies (with discrete coefficients) and agree on the dense subring of polynomials,
where the identity follows by induction from the polynomial chain rule.
\end{proof}

\begin{remark}[Formalization]
The partial derivative itself is Mathlib's \texttt{MvPowerSeries.pderiv}, which is a derivation,
so the Leibniz rule is \texttt{Derivation.leibniz} and needs no separate statement; the project
adds the chain rule for substitution and the continuity of substitution over an arbitrary
topological coefficient ring.
\end{remark}

\begin{definition}[Differential matrix]
\label{def:fgl-diff}
\lean{ChabautyColeman.FormalGroupLaw.diffMatrix,
  ChabautyColeman.FormalGroupLaw.diffMatrix_map_constantCoeff,
  ChabautyColeman.FormalGroupLaw.isUnit_diffMatrix}
\leanok
\uses{def:fgl, lem:pderiv}
$M(X) := (\partial F/\partial Y)(X,0) \in \mathrm{GL}_n(O[[X]])$: applying the chain
rule to $F(0,Y) = Y$ gives $M(0) = I$, so $\det M$ has unit constant term and $M$ is
invertible over the local ring $O[[X]]$.
\end{definition}

\begin{definition}[Base change of a formal group law]
\label{def:fgl-basechange}
\lean{ChabautyColeman.FormalGroupLaw.map,
  ChabautyColeman.FormalGroupLaw.diffMatrix_map,
  ChabautyColeman.FormalGroupLaw.diffMatrixInv_map}
\leanok
\uses{def:fgl, def:fgl-diff}
For a ring homomorphism $f \colon O \to O'$, applying $f$ to the coefficients of each
$F_i$ gives a formal group law $f_*\Phi$ over $O'$. The differential matrix and its
inverse commute with base change: $M(f_*\Phi) = f_*(M(\Phi))$ and
$M(f_*\Phi)^{-1} = f_*(M(\Phi)^{-1})$. In particular, since $M(\Phi)^{-1}$ has entries
in $O[[X]]$ already over $O$, the entries of the invariant differential $M^{-1}$ of the
base change to $\mathrm{Frac}\,O$ are integral; only the logarithm below acquires
denominators.
\end{definition}

\begin{definition}[Invariant differential and logarithm]
\label{def:fgl-log}
\lean{ChabautyColeman.FormalGroupLaw.log, ChabautyColeman.eulerAntideriv,
  ChabautyColeman.FormalGroupLaw.coeff_single_log,
  ChabautyColeman.FormalGroupLaw.diffMatrixInv_map_constantCoeff}
\leanok
\uses{def:fgl, def:fgl-diff}
The invariant differential is $\delta := M(X)^{-1}dX$. Over a field of coefficients, the
logarithm $\Lam$ is \emph{defined} as the radial (Euler) antiderivative of
$\sum_j (M^{-1})_{ij} X_j$: divide the degree-$k$ homogeneous part by $k$. Then
$\Lam \in X + (\deg \ge 2)$ and $\Lam(0) = 0$, the normalization holding because
$M^{-1}(0) = I$. That $d\Lam = \delta$ --- for which $\Lam$ is then the unique solution
with $\Lam(0) = 0$, and which requires the closedness of $\delta$ --- is
Lemma~\ref{lem:fgl-log-diff}.
\end{definition}

\begin{lemma}[Matrix invariance identity]
\label{lem:fgl-invariance-matrix}
\lean{ChabautyColeman.FormalGroupLaw.diffMatrix_subst_F}
\leanok
\uses{def:fgl-diff, lem:pderiv}
$M(F(X,Y)) = (D_Y F)(X,Y)\cdot M(Y)$ as matrices over $O[[X,Y]]$, where
$(D_Y F)_{ij} = \partial F_i/\partial Y_j$ and $M(Y)$ is the differential matrix in the
variables of the second block.
\end{lemma}
\begin{proof}
\leanok
Differentiate the associativity law $F(F(X,Y),C) = F(X,F(Y,C))$ with respect to $C_k$
and set $C = 0$. By the chain rule (Lemma~\ref{lem:pderiv}), the left side becomes
$M(F(X,Y))_{ik}$ and the right side $\sum_j \partial_{Y_j}F_i(X,Y)\cdot M(Y)_{jk}$,
using $F(Y,0) = Y$ to collapse the inner substitutions.
\end{proof}

\begin{lemma}[Closedness of the invariant differential]
\label{lem:fgl-invariance-closed}
\lean{ChabautyColeman.FormalGroupLaw.pderiv_diffMatrixInv_symm}
\leanok
\uses{def:fgl-diff, lem:fgl-invariance-matrix, lem:pderiv}
The invariant differential $\delta_i = \sum_a (M^{-1})_{ia}\,dX_a$ is closed:
$\partial_l (M^{-1})_{ik} = \partial_k (M^{-1})_{il}$. Commutativity of $F$ is
essential.
\end{lemma}
\begin{proof}
\leanok
Inverting the matrix identity of Lemma~\ref{lem:fgl-invariance-matrix} gives
$N(F(X,Y))\,(D_Y F)(X,Y) = N(Y)$ for $N = M^{-1}$. Differentiate with respect to $Y_l$
and antisymmetrize in $(k,l)$: the second-derivative terms cancel since partial
derivatives commute, leaving, with $A_{iab} = \partial_b N_{ia} - \partial_a N_{ib}$,
\[ A_{ikl}(Y) = \sum_{a,b} A_{iab}(F(X,Y))\,\partial_{Y_k}F_a\,\partial_{Y_l}F_b. \]
Setting $Y = 0$ (so $F(X,0) = X$ and $\partial_{Y_k}F_a(X,0) = M_{ak}$) shows
$M^{\mathsf T} A_i M = A_i(0)$ is the \emph{constant} matrix of linear coefficients.
From $M^{-1}M = 1$ and $M(0) = M^{-1}(0) = 1$, the linear coefficients of $M^{-1}$ and
$M$ are opposite, and those of $M$ are the mixed second derivatives
$\partial_{X_l}\partial_{Y_k} F_i(0)$, symmetric in $(k,l)$ by commutativity of $F$
(applied as a renaming of variables). Hence $A_i(0) = 0$, and multiplying by the
inverses of $M^{\mathsf T}$ and $M$ gives $A_i = 0$.
\end{proof}

\begin{lemma}[Differential of the logarithm]
\label{lem:fgl-log-diff}
\lean{ChabautyColeman.FormalGroupLaw.pderiv_log}
\leanok
\uses{def:fgl-log, lem:fgl-invariance-closed, lem:pderiv}
$d\Lam = \delta$, i.e.\ $\partial_k \Lam_i = (M^{-1})_{ik}$, in characteristic zero.
\end{lemma}
\begin{proof}
\leanok
For a closed form, the radial antiderivative is a primitive: with the Euler operator
$E = \sum_j X_j\partial_j$, which acts on the degree-$n$ part as multiplication by $n$,
one has $E\Lam_i = \sum_j (M^{-1})_{ij}X_j$ by construction, and applying $\partial_k$
gives $(E+1)(\partial_k\Lam_i) = \partial_k\sum_j (M^{-1})_{ij}X_j
= (E+1)\bigl((M^{-1})_{ik}\bigr)$, the last equality by closedness
(Lemma~\ref{lem:fgl-invariance-closed}). Since $E + 1$ multiplies the degree-$n$ part by
$n + 1 \ne 0$, it is injective on coefficients in characteristic zero.
\end{proof}

\begin{lemma}[Invariance and additivity]
\label{lem:fgl-invariance}
\lean{ChabautyColeman.FormalGroupLaw.diffMatrixInv_subst_F,
  ChabautyColeman.FormalGroupLaw.log_subst_F}
\leanok
\uses{def:fgl-log, lem:fgl-invariance-matrix, lem:fgl-log-diff}
$\delta(F(X,Y))\cdot d_Y F(X,Y) = \delta(Y)$ as formal identities; consequently
$\Lam(F(X,Y)) = \Lam(X) + \Lam(Y)$.
\end{lemma}
\begin{proof}
\leanok
Inverting the matrix identity of Lemma~\ref{lem:fgl-invariance-matrix} gives
$M(F(X,Y))^{-1}\,(D_Y F)(X,Y) = M(Y)^{-1}$, which is the first identity. For
additivity, consider $G := \Lam(F(X,Y)) - \Lam(X) - \Lam(Y)$. Its $Y$-partials vanish:
by the chain rule and $d\Lam = \delta$ (Lemma~\ref{lem:fgl-log-diff}), both
$\partial_{Y_k}\Lam(F(X,Y))$ and $\partial_{Y_k}\Lam(Y)$ equal the $Y$-substituted
entry $(M^{-1})_{ik}(Y)$ by the first identity, while $\Lam(X)$ has no $Y$-dependence.
The $X$-partials vanish by the same computation reflected through commutativity of $F$
(under the block swap, the first identity turns into its $X$-block counterpart
$M(F(X,Y))^{-1}\,(D_X F)(X,Y) = M(X)^{-1}$). Finally $G$ has zero constant term, and in
characteristic zero a series with vanishing constant term and vanishing first partials
is zero.
\end{proof}

\begin{lemma}[Integral coefficients of the logarithm]
\label{lem:log-coeff-integral}
\lean{ChabautyColeman.FormalGroupLaw.coeff_log_map}
\leanok
\uses{def:fgl-log, def:fgl-basechange}
For $\Phi$ over $O$ and any homomorphism $f \colon O \to K$ into a field, the
coefficient of the logarithm of the base change $f_*\Phi$ at multi-index $d$ is
$(\deg d)^{-1}\, f(c_d)$, where $c_d \in O$ is the corresponding coefficient of the
integral row $\sum_j (M^{-1})_{ij} X_j$ of the invariant differential. Thus the
numerator is integral and the only denominator is the total degree $\deg d$.
\end{lemma}
\begin{proof}
\leanok
The differential matrix and its inverse commute with base change
(Definition~\ref{def:fgl-basechange}), so the row $\sum_j (M^{-1})_{ij} X_j$ of
$f_*\Phi$ is $f$ applied coefficientwise to that of $\Phi$; the Euler antiderivative
then divides the degree-$d$ part by $\deg d$.
\end{proof}

\begin{lemma}[$p$-power denominators of the logarithm]
\label{lem:log-coeff-padic}
\lean{ChabautyColeman.FormalGroupLaw.coeff_log_map_padic}
\leanok
\uses{lem:log-coeff-integral}
Over a local ring $O$ of residue characteristic $p$ (so that every integer prime to $p$
is a unit), the degree-$d$ coefficient of the logarithm of $f_*\Phi$ becomes integral
after multiplication by $f(p)^{v_p(\deg d)}$: the denominator is a pure power of $p$,
with exponent $v_p(\deg d)$.
\end{lemma}
\begin{proof}
\leanok
By Lemma~\ref{lem:log-coeff-integral} the denominator is $\deg d$; writing
$\deg d = p^{v_p(\deg d)}\cdot m$ with $p \nmid m$, the prime-to-$p$ factor $m$ is a
unit in $O$, so only the $p$-power part survives as a denominator.
\end{proof}

\begin{lemma}[Integrality of the scaled logarithm]
\label{lem:log-scale-integral}
\lean{ChabautyColeman.FormalGroupLaw.coeff_log_map_scale}
\leanok
\uses{lem:log-coeff-padic}
Assume $p > e + 1$. Scaling the variables by the uniformizer $\pi$ makes the logarithm
integral: the degree-$d$ coefficient of $\Lam(\pi U)$ is $f(c)$ for some
$c \in \mathfrak m^{\deg d - e\,v_p(\deg d)}$. So $\Lam(\pi U) \in O[[U]]$, and since
$\deg d - e\,v_p(\deg d) \to \infty$ the coefficients tend to $0$, which drives the
convergence of $\Lam$ on $\mathfrak m^n$.
\end{lemma}
\begin{proof}
\leanok
Multiply the $p$-power denominator bound of Lemma~\ref{lem:log-coeff-padic} by
$\pi^{\deg d}$ and split the denominator off the numerator: writing $v := v_p(\deg d)$, we
have $e\,v \le \deg d$ because $(e+2)\,v \le (e+2)^v \le p^v \le \deg d$, so
$\pi^{\deg d} = p^{v}\,c'$ with $c' \in \mathfrak m^{\deg d - e v}$ --- the ramification
bound is spent exactly here.
\end{proof}

\begin{lemma}[Convergence of the logarithm on $\mathfrak m^n$]
\label{lem:log-converge}
\lean{ChabautyColeman.FormalGroupLaw.exists_integral_scaled_log}
\leanok
\uses{lem:log-scale-integral, def:mv-evalT}
Over a local ring $O$ of residue characteristic $p$ with $p > e + 1$: after scaling the
variables by the uniformizer $\pi$, the logarithm of $f_*\Phi$ descends to an integral
power series $H \in O[[U]]$ --- its base change is the $\pi$-scaled logarithm,
$f(\mathrm{coeff}_d H) = f(\pi)^{\deg d}\,\mathrm{coeff}_d \Lam$ --- whose coefficients
tend to $0$. Hence $\Lam$ converges on $\mathfrak m^n$, its value at
$\mathbf x = \pi\mathbf y$ being $f$ of the total evaluation of $H$ at $\mathbf y$, with
$v(\Lam(\mathbf x)) \ge 1$.
\end{lemma}
\begin{proof}
\leanok
By Lemma~\ref{lem:log-scale-integral} the degree-$d$ coefficient of $\Lam(\pi U)$ is
$f(c_d)$ for some $c_d \in \mathfrak m^{\deg d - e\,v_p(\deg d)}$; take $H$ with these
integral coefficients. Since $\deg d - e\,v_p(\deg d) \to \infty$ (the $p$-adic valuation
is sublinear and only finitely many multi-indices have bounded degree), the coefficients
tend to~$0$ and $H$ converges by Definition~\ref{def:mv-evalT}.
\end{proof}

\begin{lemma}[The formal exponential]
\label{lem:fgl-exp}
\lean{ChabautyColeman.FormalGroupLaw.exp, ChabautyColeman.FormalGroupLaw.log_comp_exp,
  ChabautyColeman.FormalGroupLaw.exp_comp_log}
\leanok
\uses{def:fgl-log, def:fgl}
The formal exponential $\mathrm{Exp}$ is the compositional inverse of $\Lam$:
$\Lam \circ \mathrm{Exp} = \mathrm{Exp} \circ \Lam = X$ as formal power series. Since
$\Lam \in X + (\deg \ge 2)$ has identity linear part, this inverse exists and is unique.
\end{lemma}
\begin{proof}
\leanok
The compositional inverse of a formal self-map with zero constant term and identity
linear part exists (by the fixed-point iteration $E^{k+1} = X - (\Lam - X)(E^k)$, which
stabilizes one coefficient degree per step because $\Lam - X$ has order $\ge 2$), is
two-sided, and is unique. Apply this to $\Lam$, whose linear part is the identity
because $M^{-1}(0) = I$.
\end{proof}

\begin{lemma}[Convergence of the exponential for odd $p$]
\label{lem:exp-converge}
\lean{ChabautyColeman.FormalGroupLaw.exp_converges, ChabautyColeman.FormalGroupLaw.exp_bounded,
  ChabautyColeman.Dexp}
\leanok
\uses{lem:fgl-exp, lem:log-converge}
The inverse series $\mathrm{Exp}$ of $\Lam$ converges on~$(pO)^n$, by the classical
estimate $v_p(k!) = (k - s(k))/(p-1)$, where $s(k)$ is the sum of the base-$p$ digits
of~$k$, and $p \ge 3$: the degree-$d$ coefficient of $\mathrm{Exp}$ becomes integral after
multiplication by $p^{\lfloor (\deg d - 1)/(p-1) \rfloor}$, and $\mathrm{Exp}(pU)$
descends to an integral power series with coefficients tending to~$0$, as for the
logarithm in Lemma~\ref{lem:log-converge}.
\end{lemma}
\begin{proof}
\leanok
The coefficient bound is proved by strong induction on $\deg d$ through the recursion
$\mathrm{Exp} = X - (\Lam - X)\circ\mathrm{Exp}$: the product slack from
$\prod_i \mathrm{Exp}_i^{m_i}$ cancels the logarithm's $1/\deg m$ denominator, and
$(p-1)\,v_p(m) \le m - 1$ closes the estimate. Convergence needs $p \ge 3$ via
$\deg d - (\deg d - 1)/(p-1) \to \infty$.
\end{proof}

\begin{remark}[Formalization]
The exponential (Lemmas~\ref{lem:fgl-exp} and~\ref{lem:exp-converge}) serves only the
injectivity of the logarithm. All three are formalized, in
\texttt{HyperellipticCurves.Archive.FormalGroupExp},
\texttt{HyperellipticCurves.Archive.ExpConverge} and
\texttt{HyperellipticCurves.Archive.LogInjective}; the main theorems do not depend on them.
\end{remark}

\begin{theorem}[The logarithm as a homomorphism]
\label{thm:log-iso}
\lean{ChabautyColeman.FormalGroupLaw.log_addHom, ChabautyColeman.FormalGroupLaw.scaledF,
  ChabautyColeman.FormalGroupLaw.scaledLog,
  ChabautyColeman.FormalGroupLaw.subst_scaledF_descent,
  ChabautyColeman.MvPSeries.evalT_subst}
\leanok
\uses{lem:fgl-invariance, lem:log-converge, def:mv-eval}
For $p > e + 1$, $\Lam \colon \mathcal G(F)(\mathfrak m) \to (\mathrm{Frac}\,O)^n$ is a
group homomorphism. ($\Lam$ is moreover injective --- indeed, in the unramified case, a
bijection onto $(pO)^n$ via $\mathrm{Exp}$ --- but this is never used.)
\end{theorem}
\begin{proof}
\leanok
Descend the formal additivity $\Lam(F(\mathbf X,\mathbf Y)) = \Lam(\mathbf X) +
\Lam(\mathbf Y)$ (Lemma~\ref{lem:fgl-invariance}) to the $\pi$-scaled integral forms and
evaluate. \emph{Formalization.} The value of $\Lam$ is realized as the total evaluation
(Definition~\ref{def:mv-evalT}) of the $\pi$-scaled logarithm $H$ (\texttt{scaledLog}), and
the $\pi$-scaled group sum $F(\pi\mathbf X, \pi\mathbf Y)/\pi$ (\texttt{scaledF}, integral
with decaying coefficients for any scaling factor in $\mathfrak m$) is turned into addition
(\texttt{log\_addHom}): the additivity identity descends along the $\pi$-rescaling
substitution --- injectivity of the coefficient map and associativity of substitution --- to
$H\circ F_\pi = H(\mathbf X) + H(\mathbf Y)$ for the scaled group law $F_\pi$
(\texttt{subst\_scaledF\_descent}), and total evaluation commutes with substitution
(\texttt{evalT\_subst}).
\end{proof}

\begin{theorem}[Hensel's lemma, smooth-point form]
\label{thm:hensel}
\lean{ChabautyColeman.exists_smooth_root, ChabautyColeman.hensel_unique,
  ChabautyColeman.exists_weierstrass_root, IsAdicComplete.henselianLocalRing}
\leanok
Let $g \in O[X]$ (not necessarily monic) with $g(a_0) \in \mathfrak m$ and
$g'(a_0) \in O^\times$. Then $g$ has a unique root $x \equiv a_0 \pmod{\mathfrak m}$. In
particular, for $f$ with $f(\omega) = 0$, $f'(\omega) \in O^\times$ and any
$y \in \mathfrak m$, there is a unique $x \equiv \omega$ with $f(x) = y^2$ --- the
implicit-function step of the disk charts. This is the full Hensel content the proof
consumes: every chart root is a single one-variable Hensel root, over $O$ or over the
$z$-adically complete base $O[[z]]$.
\end{theorem}
\begin{proof}
\leanok
Newton iteration $x \mapsto x - c\, g(x)$ with $c\,g'(a_0) = 1$, contracting in the
$\mathfrak m$-adic metric: the residual after $k$ steps lies in $\mathfrak m^{k+1}$, the
increments telescope to a limit by nonarchimedean summability, and the limit is a root
by Hausdorffness. Uniqueness: the difference of two roots is killed by a unit times
itself modulo $\mathfrak m$. \emph{Formalization.} For monic polynomials this is
Mathlib's Henselian-ring property: a local ring that is adically complete at its maximal
ideal is Henselian (\texttt{IsAdicComplete.henselianLocalRing}); a multivariate smooth-point
form is formalized in \texttt{HyperellipticCurves.Archive.HenselSmooth}, and the main theorems
do not depend on it.
\end{proof}

\chapter{Residue disks and the reduction of classes}
In this chapter $O$ is $\Zp$ or $\Oq = W(\Fq)$ unless stated otherwise; $K' :=
\mathrm{Frac}(O)$, with residue field $\kappa$. Recall the standing hypotheses:
$p \nmid \operatorname{disc} f$ and $p$ odd.

\begin{lemma}[Two kinds of points]
\label{lem:two-charts}
\lean{ChabautyColeman.point_int_or_inf}
\leanok
Every $P \in C(K')$ (affine or $\infty$) is exactly one of the following: of
\emph{finite type}, i.e.\ $P = (x_0, y_0)$ with $x_0, y_0 \in O$; or of \emph{infinite
type}, i.e.\ $v(x_0) = -2s < 0$, $v(y_0) = -(2g+1)s$, and
$t(P) := x_0^g/y_0 \in \mathfrak m$ (with $t(\infty) := 0$).
\end{lemma}
\begin{proof}
\leanok
If $v(x_0) < 0$ then $v(f(x_0)) = (2g+1)\,v(x_0)$ ($f$ is monic, so the leading term
dominates), hence $2v(y_0) = (2g+1)v(x_0)$ and
$2\,v(t(P)) = -v(x_0) > 0$, giving $t(P) \in \mathfrak m$ directly. If $v(x_0)\ge 0$
then $v(y_0^2) = v(f(x_0)) \ge 0$, so $y_0 \in O$.
\end{proof}

\begin{definition}[Reduction and residue disks]
\label{def:disk}
\lean{ChabautyColeman.redPoint, ChabautyColeman.ResidueDisk}
\leanok
\uses{lem:two-charts}
$\red(P) := (\bar x_0, \bar y_0) \in \bar C(\kappa)$ for points of finite type and
$\red(P) := \bar\infty$ for points of infinite type; the \emph{residue disk} of
$\bar P \in \bar C(\kappa)$ is $D_{\bar P} := \red^{-1}(\bar P) \subseteq C(K')$.
\end{definition}

\begin{lemma}[Square-root series]
\label{lem:sqrt-series}
\lean{ChabautyColeman.sqrtSeries}
\leanok
Let $c \in O[[z]]$ with $c(0)$ a square of a unit, $p$ odd. Then there is
$\eta \in O[[z]]$ with $\eta^2 = c$ and prescribed $\eta(0)$; explicitly, the binomial
series.
\end{lemma}
\begin{proof}
\leanok
The coefficient recursion
$\eta_n = \bigl(c_n - \sum_{0<i,j,\ i+j=n}\eta_i\eta_j\bigr)/(2\eta_0)$, solvable since
$2\eta_0 \in O^\times$, defines $\eta$; no completeness is needed. \emph{Formalization:}
\texttt{sqrtSeries} is constructed over an arbitrary commutative ring under the sole
hypothesis $2\eta_0 \in O^\times$.
\end{proof}

\begin{lemma}[Disk parametrizations]
\label{lem:disk-param}
\lean{ChabautyColeman.diskParamGeneric, ChabautyColeman.diskParamWeierstrass,
  ChabautyColeman.diskParamInf}
\leanok
\uses{def:disk, lem:sqrt-series, thm:hensel, lem:two-charts}
For each $\bar P \in \bar C(\kappa)$ there is a bijection
$z \colon D_{\bar P} \to \mathfrak m$ with inverse given by power series over $O$:
\begin{itemize}
\item[(i)] \emph{generic} $\bar P = (\bar a, \bar b)$ with $\bar b \ne 0$: fix the Hensel
lift $(a,b)$; then $z = x - a$ and $y = b\,\eta(z)$, with $\eta$ the square-root series
of $f(a+z)/f(a)$ (Lemma~\ref{lem:sqrt-series});
\item[(ii)] \emph{Weierstrass} $\bar P = (\bar\omega, 0)$: $z = y$ and
$x = \omega + \theta(z)$ with $\theta \in z^2 O[[z^2]]$ solving $z^2 = f(\omega +
\theta)$ by Theorem~\ref{thm:hensel} in series form; here $f'(\omega) \in O^\times$
because $\bar f$ is squarefree;
\item[(iii)] \emph{infinity}: $z = t = x^g/y$, and $x = z^{-2}\xi(z)$,
$y = z^{-(2g+1)}\psi(z)$ with $\xi, \psi \in 1 + zO[[z]]$ obtained by reverting
$t^2 = x^{2g}/f(x)$ in the variable $1/x$.
\end{itemize}
Weierstrass and generic disks over $\Fq$-points of the residue curve are treated over
$W(\Fq')$ for the field of definition $\Fq'$ of $\bar P$.
\end{lemma}
\begin{proof}
\leanok
In each chart, Hensel's lemma gives existence and uniqueness of the lift with prescribed
parameter value, and bijectivity onto $\mathfrak m$ follows from the uniqueness. For
(iii), revert by the standard inverse-series lemma in $O[[z]]$ (unit linear coefficient
after normalization). \emph{Formalization.} The Weierstrass and infinity charts need no
explicit power series in the parameter: for each parameter value the coordinate is the
unique Hensel root of a (in the $\infty$ case non-monic) polynomial equation.
\end{proof}

\begin{lemma}[Integral expansions of differentials]
\label{lem:omega-expansion}
\lean{ChabautyColeman.omegaExpansionGeneric, ChabautyColeman.omegaExpansionWeierstrass,
  ChabautyColeman.omegaExpansionInf}
\leanok
\uses{lem:disk-param}
Let $\omega = h(x)\,dx/y$ with $h \in O[x]$, $\deg h \le g-1$, and let $D$ be one of the
disks of Lemma~\ref{lem:disk-param}. Then $\omega|_{D} = w_D(z)\,dz$ with
$w_D \in O'[[z]]$ integral over the disk's base ring $O'$:
(i) $w = h(a+z)\,b^{-1}\eta(z)^{-1}$;
(ii) $w = 2h(x(z))/f'(x(z))$ times a unit series;
(iii) $w$ is the Laurent computation of
$h(z^{-2}\xi)\,d(z^{-2}\xi)\big/\bigl(z^{-(2g+1)}\psi\bigr)$, which is pole-free, of
order $2(g-1) - 2\deg h$ at lowest order.
\end{lemma}
\begin{proof}
\leanok
(i), (ii): chain rule (Lemma~\ref{lem:eval-hom}); the denominators are units because
$\bar f$ is squarefree and $p$ is odd. (iii): a valuation computation transcribing
$v_\infty(x^i\, dx/y) = 2g-2-2i \ge 0$; the leading coefficient is a unit multiple of
the leading coefficient of $h$. \emph{Formalization.} Each $w_D$ is a power series over
$O$ by construction, so integrality is definitional; the identity
``$\omega|_D = w_D\,dz$'' is a cleared-denominator differential identity in each chart,
and the order computation in (iii) is deferred to Lemma~\ref{lem:ord-match}.
\end{proof}

\begin{definition}[Disk antiderivatives]
\label{def:disk-antider}
\lean{ChabautyColeman.antideriv, ChabautyColeman.diskAntiderivGeneric,
  ChabautyColeman.diskAntiderivWeierstrass, ChabautyColeman.diskAntiderivInf}
\leanok
\uses{lem:omega-expansion}
$G^{\omega}_{D}(z) := \sum_k \tfrac{(w_D)_k}{k+1}\,z^{k+1}$ (so $G(0) = 0$), a
convergent function on $\mathfrak m$ by the estimates of
Theorem~\ref{thm:antideriv-zeros}.
\end{definition}

\begin{lemma}[Order matching]
\label{lem:ord-match}
\lean{ChabautyColeman.ordMatch_generic, ChabautyColeman.ordMatch_weierstrass,
  ChabautyColeman.ordMatch_inf,
  ChabautyColeman.isUnit_coeff_ordAffine_omegaExpansionGeneric,
  ChabautyColeman.isUnit_coeff_ordAffine_omegaExpansionWeierstrass,
  ChabautyColeman.isUnit_coeff_ordInf_omegaExpansionInf,
  ChabautyColeman.isUnit_coeff_of_order_map_eq}
\leanok
\uses{lem:omega-expansion, def:disk}
Let $\bar\omega := \bar h\,dx/y$ on $\bar C_{\Fpbar}$ with $\bar h \ne 0$. For every
$\bar P$:
$\ord_{\bar P}(\bar\omega) = \ord_z(\bar w_{D_{\bar P}}) =: m_{\bar P} < \infty$.
\end{lemma}
\begin{proof}
\leanok
The reduction of the expansion series is the expansion of the reduction, because the
parametrizations of Lemma~\ref{lem:disk-param} reduce to local parameters at $\bar P$ by
construction. The intrinsic order of the differential is computed on the Dedekind
coordinate ring $\A{\bar\kappa}$ at finite points --- the multiplicity of the point
ideal in $(\bar h)$, since $dx/y$ is regular and nonvanishing there --- and at $\infty$
as $\ord_\infty(\bar h) + (2g-2)$, with $\ord_\infty(\bar h) = -2\deg\bar h$ the order
at the infinite place (via the norm to $\bar\kappa(x)$) and $2g-2$ the contribution of
the canonical divisor. \emph{Formalization.} The form consumed by the per-disk count of
Theorem~\ref{thm:disk-bound} is the same statement read as a unit coefficient: in each
of the three charts, the integral expansion has a unit coefficient in degree
$m_{\bar P}$, because the reduced expansion has exactly that order.
\end{proof}

\begin{lemma}[Canonical degree over the residue field]
\label{lem:canon-deg}
\lean{ChabautyColeman.deg_div_omega_bar}
\leanok
\uses{def:hypring, lem:ord-match}
For $0 \ne \bar h$ of degree $\le g-1$ over $\Fpbar$:
$\deg \divisor(\bar h\,dx/y) = 2g-2$ on $\bar C_{\Fpbar}$.
\end{lemma}
\begin{proof}
\leanok
$\divisor(dx/y) = (2g-2)\,\bar\infty$, and $\divisor(\bar h)$ has degree $0$. The degree
is computed as the sum $\sum_{x_0}\sum_{y_0^2 = f(x_0)} \ord_{(x_0,y_0)}(\bar h)$ over
the affine points plus the order at $\infty$: each $x$-fibre contributes twice the root
multiplicity of $\bar h$ at $x_0$ (two points $(x_0,\pm b)$ generically, one doubled
Weierstrass point), so the affine part is $2\deg\bar h$, and the order at $\infty$ is
$2g-2-2\deg\bar h$.
\end{proof}

\begin{remark}[Formalization]
Lemma~\ref{lem:canon-deg} is formalized in
\texttt{HyperellipticCurves.Archive.CanonDegExact}; the main theorems do not depend on it,
since the count only needs the one-sided budget of Corollary~\ref{cor:sum-ord}.
\end{remark}

\begin{corollary}[Total vanishing budget]
\label{cor:sum-ord}
\lean{ChabautyColeman.sum_ord_le, ChabautyColeman.sum_ordPoint_le,
  ChabautyColeman.ordPoint, ChabautyColeman.finsum_ordPoint,
  ChabautyColeman.CurvePoint.equivOption, ChabautyColeman.ordAffine_eq_zero,
  ChabautyColeman.finite_support_ordPoint}
\leanok
\uses{lem:ord-match}
$\sum_{\bar P \in \bar C(\Fp)} m_{\bar P} \le 2g - 2$, with each $m_{\bar P} \ge 0$.
\end{corollary}
\begin{proof}
\leanok
The orders are nonnegative because $\bar\omega$ is regular
(Lemma~\ref{lem:omega-expansion} reduced), and the bound follows from a direct fibre
count over $\Fp$: each $x$-fibre has $0$, $1$ or $2$ rational points, hence contributes
at most twice the root multiplicity of $\bar h$ at $x_0$; summing the multiplicities
gives at most $\deg \bar h$; adding the order $2g - 2 - 2\deg\bar h$ at $\infty$ gives
$\le 2g-2$. The exact canonical degree (Lemma~\ref{lem:canon-deg}) is not needed.

\emph{Formalization.} The form summed in Theorem~\ref{thm:main-count} is: for \emph{any}
finite set $T$ of residue points, $\sum_{\bar P \in T} m_{\bar P} \le 2g-2$
(\texttt{sum\_ordPoint\_le}). The residue field is not assumed finite anywhere in this
chain: for $\bar h \neq 0$ the order vanishes at every point above a non-root of
$\bar h$, and each $x$-fibre is finite, so $\bar\omega$ has finitely many zeros and the
point sums are genuine finite sums over any field. This is what lets the counting
argument compose with Chapter~4, whose base ring has algebraically closed residue
field.
\end{proof}

\begin{corollary}[Vanishing budget of a space of differentials]
\label{cor:sum-ord-space}
\lean{ChabautyColeman.ordSpace, ChabautyColeman.sum_ordSpace_le,
  ChabautyColeman.exists_ordPoint_eq_ordSpace,
  ChabautyColeman.sum_ordPoint_le_two_mul_natDegree,
  ChabautyColeman.ordAffine_generic, ChabautyColeman.ordAffine_weierstrass,
  ChabautyColeman.ordAffine_le_ordAffine,
  Polynomial.finrank_add_natDegree_le_of_forall_dvd,
  Polynomial.exists_forall_natDegree_le_of_ne_bot,
  Polynomial.exists_dvd_forall_mem_and_rootMultiplicity_le}
\leanok
\uses{cor:sum-ord, lem:ord-match}
Let $V \ne 0$ be a subspace of the numerators of the regular differentials --- the
polynomials of degree $< g$ over $\Fpbar$ --- of codimension at most $r$, and put
\[
  n(V, \bar P) \;:=\; \min\{\ord_{\bar P}(\bar h\,dx/y) \;:\; 0 \ne \bar h \in V\} .
\]
Then $\sum_{\bar P \in T} n(V, \bar P) \le 2r$ for every finite set $T$ of points of
$\bar C$. For $\dim V = 1$ this is Corollary~\ref{cor:sum-ord}.
\end{corollary}
\begin{proof}
\leanok
For $x_0 \in \Fpbar$ let $\mu(x_0)$ be the smallest multiplicity of $x_0$ as a root of a
nonzero element of $V$, and let $q := \prod (x - x_0)^{\mu(x_0)}$, the product over the
$x$-coordinates of the affine points of $T$. Every nonzero $\bar h \in V$ is divisible by
$q$, since the factors are coprime and each divides $\bar h$. Let $D$ be the largest
degree of an element of $V$; then every element of $V$ lies in $q \cdot \{\deg \le D -
\deg q\}$, so $\dim V + \deg q \le D + 1$, and with $\dim V \ge g - r$,
\[ \deg q + \bigl(g - 1 - D\bigr) \;\le\; r . \]
At an affine point $\bar P$ over $x_0$ the order of $\bar h\,dx/y$ is the multiplicity of
$x_0$ in $\bar h$, doubled at a Weierstrass point, so it is monotone in that multiplicity
and $n(V, \bar P) = \ord_{\bar P}(q\,dx/y)$; the affine part of the sum is therefore at
most $2\deg q$ by the fibre count of Corollary~\ref{cor:sum-ord}, applied to $q$ with the
point at infinity left out. At infinity, the element of $V$ of maximal degree $D$ gives
$n(V, \infty) \le 2(g - 1 - D)$. Adding the two bounds gives $2r$.
\end{proof}

\begin{remark}[Formalization]
The minimum $n(V, \bar P)$ is \texttt{ordSpace}, an infimum over the nonzero elements of
$V$; it is attained (\texttt{exists\_ordPoint\_eq\_ordSpace}), which is how the per-disk
bound of Theorem~\ref{thm:disk-bound} is fed with an actual differential. The
displayed dimension inequality is \texttt{finrank\_add\_natDegree\_le\_of\_forall\_dvd},
stated for any subspace of polynomials of bounded degree with a common divisor: this is
the elementary substitute, available for a hyperelliptic curve, for the Riemann--Roch
argument of Theorem~6.4 of M.~Stoll, \emph{Independence of rational points on twists of
a given curve}, Compositio Math.\ 142 (2006), 1201--1214.
\end{remark}

\begin{lemma}[Content normalization of functions]
\label{lem:content}
\lean{ChabautyColeman.exists_primitive_smul, ChabautyColeman.exists_primitive_smul_pair,
  ChabautyColeman.exists_primitive_generator}
\leanok
\uses{lem:hypring-units}
Every $h \in \F{K'}^\times$ has a $K'^\times$-multiple $h_0 = a + by$ with
$a, b \in K'(x)$ admitting a \emph{primitive} presentation: numerator data in $O[x]$ of
unit content. The multiple is unique up to $O^\times$, since generators of a principal
ideal differ by $(\A{K'})^\times = K'^\times$.
\end{lemma}
\begin{proof}
\leanok
Scale by the content: primitivity is measured by the Gauss content
$\mathrm{cont}_v(p) = \sup_i v(p_i)$ on $K'[x]$; a common denominator $D$ clears
$h = a + by$ to $(P + Qy)/D$ with $P, Q, D \in K'[x]$, and scaling $D$ and the pair
$(P,Q)$ each to unit content --- attain the finite supremum and divide by the attaining
coefficient --- gives the primitive presentation. The scaling factor is a constant of
$K'^\times$. The statement holds for any valuation $v$; uniqueness up to $O^\times$ is
not needed downstream and is not formalized. \emph{Formalization.} Generators arrive in
the proof as explicit polynomial pairs $(P,Q)$, so only the pair-level normalizations
are used; the function-level statement is formalized in
\texttt{HyperellipticCurves.Archive.ContentGenerator}, and the main theorems do not depend on
it.
\end{proof}

\begin{theorem}[Specialization of principal divisors]
\label{thm:spec}
\leanok
\lean{ChabautyColeman.exists_center_finsum_emultiplicity_generic,
ChabautyColeman.exists_center_finsum_emultiplicity_weierstrass}
\uses{lem:content, lem:disk-param, thm:wprep, lem:two-charts}
Let $O_K$ be the ring of integers of a finite extension $K/K'$ and let
$h \in \F{K}^\times$ be primitive as in Lemma~\ref{lem:content}. Then $\bar h \ne 0$ in
$\F{\bar\kappa}$, and for every \emph{affine} $\bar P \in \bar C(\bar\kappa)$,
\[ \ord_{\bar P}(\bar h) \;=\; \sum_{P \in \divisor(h),\ \red(P) = \bar P} \ord_P(h), \]
the zeros of $h$ on the residue disk being enumerated as actual points over a splitting
discrete valuation ring. Affine points suffice: the reduction map of
Definition~\ref{def:redJ} takes values in the class group of the \emph{affine}
coordinate ring, where points escaping to $\bar\infty$ contribute the unit ideal, so the
statement at $\bar\infty$ is never needed.
\end{theorem}
\begin{proof}
\leanok
Write $h = P + Q\,y$ and expand it on the disk $D_{\bar P}$ (Lemma~\ref{lem:disk-param},
with charts over $O_K$). The number of zeros of the expansion in the open disk --- the
roots of the Weierstrass factor $W$ --- equals the vanishing order of its reduction
(Theorem~\ref{thm:wprep}); the per-disk point sums enumerate these zeros via Hensel
lifting of disk centres. At a Weierstrass point the hyperelliptic involution pairs the
disk parameters ($z \mapsto -z$), so the summed orders factor through the norm and no
individual conjugate order needs to be isolated. \emph{Formalization.} A chart-level
$\Cp$-version at every point including $\bar\infty$ is formalized in
\texttt{HyperellipticCurves.Archive.SpecializationGeneric} and
\texttt{HyperellipticCurves.Archive.SpecializationInf}; the main theorems do not depend on
it.
\end{proof}

\begin{definition}[Reduction of classes]
\label{def:redJ}
\lean{ChabautyColeman.redJ, ChabautyColeman.IsRedJPair, ChabautyColeman.exists_isRedJPair,
  ChabautyColeman.redIdealBar, ChabautyColeman.SplitSupported}
\leanok
\uses{lem:mumford-exists, def:disk, thm:spec}
Define $\redJ \colon \J{K'} \to \Cl(\A{\bar\kappa})$ as follows: represent
$c = \sum_i n_i\,\aj(P_i)$ with $P_i \in C(K_i)$ and $K_i/K'$ finite (possible by
Lemma~\ref{lem:mumford-exists}), and set $\redJ(c) := \sum_i n_i\,\aj(\red P_i)$.
\emph{Formalization.} The map is constructed over a complete discrete valuation ring
$O$ with algebraically closed residue field (the arithmetic version is the composite
with the base change to $\mathrm{Frac}\,O$): a class is represented over a finite
intermediate field $L$ of the algebraic closure by a \emph{split-supported} ideal ---
all prime divisors are maximal ideals of rational points; Mumford representatives become
split-supported over the splitting field of $u$ --- which is reduced pointwise to the
fixed residue curve, running the disk constructions over the integral closure $O_L$,
again a complete discrete valuation ring with the \emph{same} residue field.
\end{definition}

\begin{theorem}[Reduction is a well-defined homomorphism]
\label{thm:redJ-hom}
\lean{ChabautyColeman.redJHom, ChabautyColeman.isRedJPair_unique, ChabautyColeman.redJ_mul,
  ChabautyColeman.redJ_aj, ChabautyColeman.isPrincipal_redIdeal_span,
  ChabautyColeman.redIdealBar_map_of_splitSupported}
\leanok
\uses{def:redJ, thm:spec, lem:content}
$\redJ$ is well defined, is a group homomorphism, and satisfies
$\redJ(\aj(P)) = \aj(\red P)$.
\end{theorem}
\begin{proof}
\leanok
Independence of the representative: two representations differ by a principal divisor
$\divisor(h)$; normalize $h$ primitively (Lemma~\ref{lem:content}); by
Theorem~\ref{thm:spec} its pointwise reduction is $\divisor(\bar h)$, a principal
divisor, hence trivial in $\Cl(\A{\bar\kappa})$. Additivity is then formal.

\emph{Formalization.} Two split-supported representatives are compared over the
compositum and a further finite extension splitting the norms of the connecting
principal generators. There the reductions of the principal parts are principal ---
split the generator into a coprime primitive part, killed by Theorem~\ref{thm:spec},
and a split common factor, handled root by root through
$\mathfrak m_P\,\mathfrak m_{\iota P} = (x - \bar a)$ --- while the reductions of the
representatives are unchanged under base change: a split-supported ideal factors into
rational point ideals, which stay maximal and rational. Compatibility with the
Abel--Jacobi map: the reduction of the maximal ideal of an integral point is the maximal
ideal of the reduced point, and points reducing to $\bar\infty$ contribute the unit
ideal.
\end{proof}

\begin{lemma}[Frobenius-fixed classes are few]
\label{lem:frobfix-finite}
\lean{ChabautyColeman.frobFixed_finite, ChabautyColeman.HypJac.congr}
\leanok
\uses{lem:mumford-exists}
Let $\sigma$ be an automorphism of the base field with finitely many fixed elements that
fixes $\bar f$ (e.g.\ the $q$-power Frobenius on $\Fpbar$). Then the $\sigma$-fixed
subgroup of $\Cl(\A{\Fpbar})$ is finite: a fixed class has a $\sigma$-fixed
\emph{unique} reduced representative (Lemma~\ref{lem:mumford-exists} over $\Fpbar$),
hence Mumford data $(u,v)$ with coefficients in the fixed field, of which there are
finitely many.
\end{lemma}
\begin{proof}
\leanok
The automorphism induces an isomorphism of the coordinate ring and hence of the class
group (formalized in general: a ring isomorphism of Dedekind domains induces an
isomorphism of class groups, mapping the class of an ideal to the class of its image).
Uniqueness of the reduced Mumford representative transports fixedness of a class to
fixedness of its data; count monic $u$ of degree $\le g$ and $v$ of degree $< g$ with
coefficients in the fixed field.
\end{proof}

\begin{theorem}[The kernel of reduction has finite index]
\label{thm:J1-index}
\lean{ChabautyColeman.J1, ChabautyColeman.J1_finiteIndex, ChabautyColeman.redJ_fixed,
  ChabautyColeman.redIdealBar_congr, ChabautyColeman.adjoinRootSetEquiv,
  ChabautyColeman.congr_mapBase}
\leanok
\uses{thm:redJ-hom, lem:frobfix-finite}
$J^1(K') := \ker(\redJ)$ has finite index in $\J{K'}$, bounded by the number of Mumford
data $(u,v)$ over $\Fq$.
\end{theorem}
\begin{proof}
\leanok
Galois-equivariance of $\redJ$ on fixed classes: a class fixed by the arithmetic
automorphism has a fixed unique reduced Mumford representative $(u,v)$; its splitting
field $L = K(\text{roots of }u)$ is stable, and the reduction of the Mumford ideal over
$L$ transports along the induced automorphisms (a split-supported ideal factors into
rational point ideals, whose reductions are compatible with the semilinear
residue-field identification). Classes coming from $\J{K'}$ are fixed, since their
Mumford data live over $K'$, so $\J{K'}/J^1$ embeds into the group of Frobenius-fixed
classes, which is finite by Lemma~\ref{lem:frobfix-finite}.
\end{proof}

\begin{remark}[Formalization]
Since the pairing $\lambda$ of Definition~\ref{def:lambda} is defined on $J^1$
directly, the main theorems depend neither on Lemma~\ref{lem:frobfix-finite} nor on
Theorem~\ref{thm:J1-index}. Both are formalized, together with the equivariance apparatus
(a compatible family of automorphisms of $O$, its fraction field, the algebraic closure
and the residue field, all fixing $f$), in
\texttt{HyperellipticCurves.Archive.J1FiniteIndex},
\texttt{HyperellipticCurves.Archive.Equivariance} and
\texttt{HyperellipticCurves.Archive.StableSplitting}. The kernel $J^1$ itself, the base
change into it, and the isomorphism of class groups induced by a ring isomorphism are part
of the main development.
\end{remark}

\chapter{The abelian logarithm}
This chapter constructs the analytic core of the method: the auxiliary class $\tau$, the
translated formal group law $F_\tau$ on the chart at $\tau$, the logarithm $\Lam$ on the
kernel of reduction, the differentials $\omega'_i$ with their regularity, and the pairing
$\lambda$ with its integral representation on residue disks. Fix $q = p^N$ as produced by
Lemma~\ref{lem:tau-exists} and put $\Oq = W(\Fq)$.

\begin{definition}[Regular differentials]
\label{def:regdiff}
\lean{ChabautyColeman.RegularDiff}
\leanok
$\RegDiff := \bigoplus_{i=0}^{g-1} \Qp\cdot x^i\,dx/y$, a $g$-dimensional $\Qp$-vector
space with its $\Zp$-lattice $\RegDiff_{\ZZ} := \bigoplus \Zp\, x^i\,dx/y$. (In the
formalization, differentials are represented by their numerators $h$ with
$\omega = h\,dx/y$, $\deg h < g$; the packaged submodule is formalized in
\texttt{HyperellipticCurves.Archive.RegularDiff}, and the main theorems do not depend on it.)
\end{definition}

\begin{definition}[Good position]
\label{def:goodpos}
\lean{ChabautyColeman.GoodPosition, ChabautyColeman.GoodPosition.isCoprime_right,
  ChabautyColeman.GoodPosition.isCoprime_two_mul_right}
\leanok
\uses{lem:mumford-exists}
A reduced Mumford datum $(u, v)$ over $\Oq$ is in \emph{good position} if $\deg u = g$,
$\bar u$ is squarefree, and $\gcd(\bar u, \bar f) = 1$; the auxiliary data derived from
$\tau$ below are required to satisfy the analogous conditions where they are used. Each
condition is the non-vanishing modulo $p$ of an explicit resultant or discriminant of the
data. The coprimality conditions involving $\bar v$ --- in particular
$\gcd(\bar u, 2\bar v) = 1$, the invertibility used by the chart --- are consequences.
\end{definition}

\begin{lemma}[Abundance of points]
\label{lem:many-points}
\lean{ChabautyColeman.residue_curve_points_infinite,
  ChabautyColeman.infinite_setOf_eval_ne_zero}
\leanok
$\bar C(\Fpbar)$ is infinite; more precisely, the set of non-Weierstrass affine points is
infinite. Hence for any finite set of conditions as in Definition~\ref{def:goodpos} there
are $g$ distinct non-Weierstrass points avoiding them. (Rationality of the chosen points
over a finite field, needed for the counting argument, is handled at the arithmetic
instantiation in Chapter~5.)
\end{lemma}
\begin{proof}
\leanok
For all but finitely many $\bar a \in \Fpbar$, $\bar f(\bar a)$ is a nonzero square in
$\Fpbar$, and each such $\bar a$ lifts to a non-Weierstrass point
$(\bar a, \sqrt{\bar f(\bar a)})$; the conditions exclude only finitely many points.
\end{proof}

\begin{lemma}[Genericity of the centre]
\label{lem:tau-generic}
\lean{ChabautyColeman.HasGoodCenter, ChabautyColeman.hasGoodCenter_of_pos_genus,
  ChabautyColeman.hasGoodCenter_genus_two,
  ChabautyColeman.IsFoldCerts, ChabautyColeman.IsFoldAvoid,
  ChabautyColeman.exists_centerFoldWitness, ChabautyColeman.exists_centerFoldPack,
  ChabautyColeman.exists_isFoldCerts_isFoldAvoid,
  ChabautyColeman.cantorCompose_point_eq_disjoint,
  ChabautyColeman.cantorCompose_point_eq_double,
  ChabautyColeman.isCoprime_foldChain_of_resultant_ne_zero,
  ChabautyColeman.curveAvoidNum, ChabautyColeman.eval_curveAvoidNum,
  ChabautyColeman.not_curveVanishing_curveAvoidNum,
  ChabautyColeman.not_curveVanishing_curveCondNum_of_squarefree,
  ChabautyColeman.not_curveVanishing_curveAvoidNum_of_squarefree,
  ChabautyColeman.not_curveVanishing_X_snd, ChabautyColeman.not_curveVanishing_X_sub_C,
  ChabautyColeman.exists_curveTuple_ne_zero, ChabautyColeman.curveVanishing_mul,
  ChabautyColeman.not_curveVanishing_of_aeval_ne_zero}
\leanok
\uses{def:goodpos, lem:many-points, def:cantor-map, def:centre-fold, thm:fold-witness,
  lem:fold-witness-genus-one}
There exists a tuple of $g$ affine points of the reduced curve over $\Fpbar$ ---
pairwise distinct $x$-coordinates avoiding a prescribed finite set $T$ of residue
$x$-values, nonzero $y$-coordinates --- whose interpolated Mumford datum
$(\bar u, \bar v)$, with $\bar u = \prod_i (X - \bar x_i)$ and $\bar v$ the interpolant
through the points, is such that the point-by-point fold of
Definition~\ref{def:centre-fold} computing $\tau + \tau - \tau$ is defined throughout:
every step condition of the $2g$-step fold is nonzero, and every intermediate state
avoids the Weierstrass points, $\gcd(U_j, 2V_j) = 1$, the starting datum included. Good
position of $(\bar u, \bar v)$ is automatic: $\bar u$ is squarefree because the
$\bar x_i$ are distinct, and coprime to $\bar f$ because the $\bar y_i$ are nonzero. The
statement holds for every genus $g \ge 1$. For $g \ge 2$ no squarefreeness of $\bar f$
is needed; for $g = 1$ it is needed: for $\bar f = X^3$ in characteristic $3$ the fold
conditions vanish identically on the curve
(cf.\ Lemma~\ref{lem:fold-witness-genus-one}). The avoidance of the prescribed set $T$
is what the per-disk configuration of Lemma~\ref{lem:global-generic} requires: $\tau$
must miss the centres of the residue disks being counted.
\end{lemma}
\begin{proof}
\leanok
\uses{thm:fold-witness, lem:fold-witness-genus-one}
The fold witness --- Theorem~\ref{thm:fold-witness} for $g \ge 2$ and
Lemma~\ref{lem:fold-witness-genus-one} for $g = 1$ --- provides a tuple at which the
$2g$ step conditions, the nonvanishing of the $y$-coordinates, the $T$-avoidance and the
per-stage Weierstrass avoidance all hold simultaneously: each condition is a nonzero
polynomial condition on tuples of curve points, and finitely many such conditions have a
common witness (Lemma~\ref{lem:many-points} and the avoidance principle below; the
stage coprimalities enter as resultants of the isobaric numerator states of
Lemma~\ref{lem:fold-universal}, refuted on the cuspidal degeneration). The B\'ezout
certificates required by the certificate-passing formulation of Cantor's algorithm are
then constructed over the field: at a disjoint step, coprimality of $U_j$ and
$X - \bar x_j$ is the step condition $U_j(\bar x_j) \neq 0$; at the doubling step,
coprimality of $X - \bar x_0$ and $V_0 + \bar y_0$ follows from $2\bar y_0 \neq 0$. With
these certificates the certified computation coincides with the fold step by step: at a
disjoint step the composed square root is determined by Chinese-remainder uniqueness,
and at the doubling step by uniqueness of the square root modulo the composed support.
\end{proof}

\bigskip
\noindent\textbf{The witness for all genera: the point-by-point fold and its moment
certificates.} The following sequence of statements proves the existence of the witness
tuple for every genus; the entire sequence is formalized, and its main steps were
additionally verified symbolically at small genus by the scripts in \texttt{prototype/}: the
fold on the cuspidal curve (\texttt{symfold.m}), the moment characterization of its states
and the lemmas behind it (\texttt{apic.m}, \texttt{p1lemmas.m}), and the mirror
(\texttt{revfold.m}); every condition of every genus $\le 6$ has been certified numerically
from the moment determinants alone (\texttt{detcond.m}).

\begin{definition}[The point-by-point fold]
\label{def:centre-fold}
\uses{def:cantor-map, def:goodpos}
\lean{ChabautyColeman.foldChain, ChabautyColeman.IsFoldState,
  ChabautyColeman.IsFoldState.step, ChabautyColeman.cuspChain_eq_foldChain,
  ChabautyColeman.foldGenH, ChabautyColeman.foldH_eq_of_eval_eq,
  ChabautyColeman.foldH_eq_of_dvd_double, ChabautyColeman.foldXs,
  ChabautyColeman.foldYs}
\leanok
Let $k$ be a field of characteristic $\neq 2$, let $f \in k[X]$ be monic squarefree of
degree $N = 2g+1$, and let $P_1, \dots, P_g$ be affine non-Weierstrass points of
$y^2 = f(x)$ with pairwise distinct $x$-coordinates; let $\tau = (u, v)$ be the reduced
Mumford datum interpolating them. The \emph{fold} computes $\tau + \tau - \tau$ by $2g$
single-point additions: $D_0 := \tau$ and $D_j := \operatorname{red}(D_{j-1} + P_j)$ for
$1 \le j \le g$ (\emph{phase 1}), then
$D_{g+m} := \operatorname{red}(D_{g+m-1} + \iota P_{g+1-m})$ for $1 \le m \le g$
(\emph{phase 2}: the involuted points are added in \emph{reverse} order, so that
$[D_{g+m}] = [2\tau] - \sum_{i > g-m}[P_i - \infty] = [D_{g-m}]$ and the fold retraces
its own phase-1 states backwards; see Lemma~\ref{lem:fold-mirror}). Each addition
composes to a divisor of degree $g+1$ whose $b$-polynomial has degree $\le g$, and a
single Cantor reduction returns to degree $g$; since $2\bigl((g+1)-1\bigr) < \deg f$,
every reduction lies in the unconditional degree regime, so no degree-genericity
condition arises at any step. The \emph{conditions} of the fold are: at each step, the
disjointness value $U(x_P)$ of the current $u$-polynomial (at the single doubling step
$j = 1$ instead $U_0'(x_{P_1})$ and $V_0(x_{P_1}) + y_{P_1}$). Further avoidance
quantities --- e.g.\ the Weierstrass avoidance $\operatorname{res}(U_j, f) \neq 0$ of an
intermediate state --- are polynomial in the state coefficients and are treated by the
same universal method.
\end{definition}

\begin{lemma}[Mirror symmetry]
\label{lem:fold-mirror}
\uses{def:centre-fold, def:cusp-fold, thm:cusp-moment}
\lean{ChabautyColeman.congruence_of_even_mul, ChabautyColeman.foldNodes,
  ChabautyColeman.prod_foldNodes_pairing,
  ChabautyColeman.genNode_zero_atom_ne_zero,
  ChabautyColeman.genNode_trio_atom_ne_zero,
  ChabautyColeman.genNode_tripled_atom_ne_zero,
  ChabautyColeman.genCuspFold, ChabautyColeman.genFoldH,
  ChabautyColeman.genCuspFold_eval_sq_ne_zero,
  ChabautyColeman.exists_genCuspFold_witness}
\leanok
On the cuspidal model the hyperelliptic involution is $t \mapsto -t$, so phase 2 of the
fold is the phase-1 recursion continued along the extended node sequence
$(t_1, \dots, t_g, -t_g, \dots, -t_1)$. Assume the fold conditions hold up to phase-2
step $m \le g$. Then the state $A_{g+m}$ of the extended chain equals the phase-1 state
$A_{g-m}$ of the \emph{same} chain: the fold retraces its own states backwards.
Consequently every phase-2 condition value is the evaluation of a phase-1 state at a
single node --- the disjointness value of phase-2 step $m+1$ is $U_{g-m}(x_{g-m})$, the
stage-$(g-m)$ support evaluated at its own last-added node --- and phase 2 needs no
separate analysis: on the cusp, the full $2g$-step fold is controlled by the phase-1
theory alone.
\end{lemma}
\begin{proof}
\leanok
The stage-$(g+m)$ alphabet of the extended chain factors as
$\prod_{i > g-m}(z^2 - t_i^2)$ times the stage-$(g-m)$ alphabet of the same chain ---
an \emph{even} polynomial factor with nonzero value at $0$ at generic nodes. The master
congruence of Theorem~\ref{thm:cusp-moment} descends along such even factors, so
$A_{g+m}$ satisfies the stage-$(g-m)$ moment congruence, and the value-extraction clause
of Theorem~\ref{thm:cusp-moment} reads off every phase-2 condition value against the
phase-1 determinants at the generic nodes (the avoidance value and the two mirror
values, extending the stage-$(g-m)$ alphabet at $\mp t_{g-m}$); no second chain and no
state identification is needed. Only at the final stage $2g$ is $A_{2g} = A_0$ read off
by the uniqueness clause, to get the last avoidance value from $A_0(0) \neq 0$. The
formal argument is an interleaved induction: conditions up to step $m$ give the
descended congruence at $m$, which gives the conditions at $m+1$. At the level of the
curve the identity $D_{g+m} = D_{g-m}$ also follows from the class computation
$[D_{g+m}] = [2\tau] - \sum_{i > g-m}[P_i - \infty] = [\tau] + \sum_{i \le g-m}
[P_i - \infty] = [D_{g-m}]$ together with uniqueness of reduced representatives, but the
formalized route needs only the cuspidal case, where the congruence argument replaces
the Jacobian.
\end{proof}

\begin{lemma}[Universal isobaric numerators]
\label{lem:fold-universal}
\uses{def:centre-fold}
\lean{ChabautyColeman.MumfordNum, ChabautyColeman.MumfordNum.step,
  ChabautyColeman.foldNum, ChabautyColeman.foldHNum, ChabautyColeman.foldCondNum,
  ChabautyColeman.foldBase, ChabautyColeman.foldNum_succ, ChabautyColeman.den_foldNum,
  ChabautyColeman.map_foldNum, ChabautyColeman.map_foldHNum,
  ChabautyColeman.map_foldCondNum, ChabautyColeman.map_foldBase,
  ChabautyColeman.IsScaledState, ChabautyColeman.IsScaledState.step,
  ChabautyColeman.isScaledState_foldNum, ChabautyColeman.foldCond,
  ChabautyColeman.foldCondNum_eq_den_mul, ChabautyColeman.foldStep_dvd_disjoint,
  ChabautyColeman.foldStep_dvd_double, ChabautyColeman.IsIsobaric,
  ChabautyColeman.IsIsobaric.divByMonic, ChabautyColeman.IsIsobaricState,
  ChabautyColeman.IsIsobaricState.step, ChabautyColeman.foldWt,
  ChabautyColeman.foldCondWt, ChabautyColeman.isIsobaricState_foldNum,
  ChabautyColeman.isWeightedHomogeneous_foldCondNum,
  ChabautyColeman.isWeightedHomogeneous_foldHNum,
  ChabautyColeman.isIsobaricState_foldBase, ChabautyColeman.univWt,
  ChabautyColeman.univF, ChabautyColeman.univXs, ChabautyColeman.univYs,
  ChabautyColeman.univBase, ChabautyColeman.univCondNum,
  ChabautyColeman.isWeightedHomogeneous_univCondNum, ChabautyColeman.cuspSpec,
  ChabautyColeman.isWeightedHomogeneous_aeval_cuspSpec,
  ChabautyColeman.cuspNum, ChabautyColeman.isWeightedHomogeneous_cuspNum,
  ChabautyColeman.cuspNum_eq_foldCondNum}
\leanok
Give $X$ and $x_i$ weight $2$, $y_i$ weight $N$, and the coefficient $a_r$ of $f$ weight
$2(N-r)$. There are \emph{universal} polynomials
$\mathcal{N}_q \in \mathbb{Z}[x_1, \dots, x_g,\, y_1, \dots, y_g,\, a_0, \dots, a_{N-1}]$,
isobaric for these weights and indexed by the quantities $q$ of the fold (the
coefficients of every intermediate state and every condition value), such that over any
field of characteristic $\neq 2$ in which the fold is defined,
$q = \mathcal{N}_q(\text{inputs})\big/\bigl(2^{b_q} \prod \mathcal{N}_{q'}^{\,k_{q'}}\bigr)$
with the product over \emph{earlier} condition quantities $q'$. Under the cuspidal
specialization $a = 0$, $x_i = t_i^2$, $y_i = t_i^N$, each $\mathcal{N}_q$ becomes a
homogeneous polynomial $P_q \in \mathbb{Z}[t_1, \dots, t_g]$ of degree
$\operatorname{wt}(\mathcal{N}_q)$.
\end{lemma}
\begin{proof}
\leanok
Induction along the fold in cofactor form. Composition divides only by the current
disjointness value (at the doubling step, by $2 y_{P} U'(x_P)$); the reduction step
divides by nothing: the cofactor division is by a \emph{monic} polynomial, and the
normalization is by the leading coefficient $-1$, as $f$ dominates. Clearing the
denominators inductively produces the $\mathcal{N}_q$, with denominators as stated.
Every operation is equivariant under the scaling
$(x, y, a_r) \mapsto (\lambda^2 x, \lambda^N y, \lambda^{2(N-r)} a_r)$ --- the equation
$y^2 - f(x)$ is isobaric of weight $2N$ --- so the numerators may be chosen isobaric.
The cuspidal specialization sends an isobaric monomial with exponents $\beta_i$ (on
$x_i$) and $\gamma_i$ (on $y_i$) to $\pm t^{c}$ with
$\sum_i c_i = \sum_i (2\beta_i + N\gamma_i)$, which equals the weight when no $a_r$ is
present; monomials containing some $a_r$ vanish.
\end{proof}

\begin{definition}[The cuspidal model and the $A$-recursion]
\label{def:cusp-fold}
\uses{def:centre-fold}
\lean{ChabautyColeman.cuspChain, ChabautyColeman.IsCuspState}
\leanok
The \emph{cuspidal model} is $y^2 = X^N$ with rational parametrization $x = t^2$,
$y = t^N$; Cantor's algorithm operates on Mumford pairs verbatim. Let $t_1, \dots, t_g$
be indeterminates over the prime field ($\operatorname{char} \neq 2$),
$P_i := (t_i^2,\, t_i^N)$, and run the fold over $\mathbb{Q}(t_1, \dots, t_g)$ resp.\
$\mathbb{F}_p(t_1, \dots, t_g)$. Along phase 1 define $A_j \in \mathbb{Q}(t)[z]$ by
$A_0 := \prod_i (z - t_i)$ and, at the step adding $P_{j+1}$ with composed pair
$(U_c, V_c)$: $A_c := A_j \cdot (z - t_{j+1})$ and $A_{j+1} := (-1)^g\, C_c(-z)$, where
$C_c := \bigl(z^N - V_c(z^2)\bigr)/A_c(z)$, a polynomial by
Lemma~\ref{lem:cusp-divisibility}. On the cusp the avoidance condition is
$U_j(0) \neq 0$, since $\operatorname{res}(U, X^N) = \pm U(0)^N$.
\end{definition}

\begin{lemma}[$z$-lift of the Mumford condition]
\label{lem:cusp-divisibility}
\uses{def:cusp-fold}
\lean{ChabautyColeman.cuspCompose_dvd, ChabautyColeman.cuspNext_dvd,
  ChabautyColeman.cuspStep_mumford, ChabautyColeman.IsCuspState.step,
  ChabautyColeman.isCuspState_cuspChain}
\leanok
At every stage of phase 1, $A \mid z^N - V(z^2)$ --- for the reduced states
$(A_j, V_j)$ and for the composed pairs $(A_c, V_c)$ alike.
\end{lemma}
\begin{proof}
\leanok
Induction along the fold. Base: $v$ interpolates, so $v(t_i^2) = t_i^N$ and each
$z - t_i$ divides $z^N - v(z^2)$; the factors are coprime (distinct $t_i$). Composition:
$V_c = V_j + U_j H$ with $H$ the composition scalar, so
$z^N - V_c(z^2) = \bigl(z^N - V_j(z^2)\bigr) - U_j(z^2)H$ is divisible by $A_j$
(inductively, and $A_j(z) \mid U_j(z^2)$ by Lemma~\ref{lem:cusp-reduction}); the new
factor divides because $V_c(t_{j+1}^2) = t_{j+1}^N$ by the interpolation property of the
composition, and it is coprime to $A_j$ by the disjointness condition. Reduction:
$V_{j+1} \equiv -V_c \bmod U_{j+1}$ and $A_{j+1}(z) \mid U_{j+1}(z^2)$, so it suffices
that $A_{j+1} \mid z^N + V_c(z^2) = -\bigl((-z)^N - V_c((-z)^2)\bigr) =
-A_c(-z)C_c(-z)$, which holds since $A_{j+1} = \pm C_c(-z)$.
\end{proof}

\begin{lemma}[Reduction in the $A$-picture]
\label{lem:cusp-reduction}
\uses{def:cusp-fold, lem:cusp-divisibility}
\lean{ChabautyColeman.cuspCofactor_monic, ChabautyColeman.cuspCofactor_natDegree,
  ChabautyColeman.cuspNext_monic, ChabautyColeman.cuspNext_natDegree,
  ChabautyColeman.cuspNext_U, ChabautyColeman.cuspCompose_U_comp}
\leanok
Every $A_j$ is monic of degree $g$, and $U_j(z^2) = (-1)^g A_j(z) A_j(-z)$.
\end{lemma}
\begin{proof}
\leanok
Base: $u(z^2) = \prod(z^2 - t_i^2) = A_0(z) \cdot \prod(z + t_i) =
(-1)^g A_0(z) A_0(-z)$. Step: from $V_c(z^2)^2 - z^{2N} = U_c(z^2)\, w_c(z^2)$ and the
factorizations $V_c(z^2) - z^N = -A_c(z)C_c(z)$ and $V_c(z^2) + z^N = -A_c(-z)C_c(-z)$,
together with $U_c(z^2) = (-1)^{g+1} A_c(z) A_c(-z)$ (composition multiplies $U_j$ by
$X - t_{j+1}^2$), one gets $w_c(z^2) = (-1)^{g+1} C_c(z) C_c(-z)$. The Cantor reduction
returns $U_{j+1} = w_c/\operatorname{lc}(w_c) = -w_c$ (here $\deg w_c = g$ and
$\operatorname{lc}(w_c) = -1$ because $f = X^N$ dominates $V_c^2$), so
$U_{j+1}(z^2) = (-1)^g C_c(z)C_c(-z) = (-1)^g A_{j+1}(z) A_{j+1}(-z)$; and $C_c$ is
monic of degree $g$ ($z^N - V_c(z^2)$ is monic of degree $N$, $A_c$ monic of degree
$g+1$), so $A_{j+1} = (-1)^g C_c(-z)$ is monic of degree $g$.
\end{proof}

\begin{lemma}[Exactness of the reduction step]
\label{lem:cusp-exact}
\uses{def:cusp-fold, lem:cusp-reduction}
\lean{ChabautyColeman.cuspNext_exact}
\leanok
$A_{j+1}(z)\,A_c(-z) - A_{j+1}(-z)\,A_c(z) = (-1)^{g+1}\, 2 z^N$, exactly.
\end{lemma}
\begin{proof}
\leanok
$(-1)^g\bigl(A_{j+1}(z)A_c(-z) - A_{j+1}(-z)A_c(z)\bigr) = C_c(-z)A_c(-z) - C_c(z)A_c(z)
= \bigl(-z^N - V_c(z^2)\bigr) - \bigl(z^N - V_c(z^2)\bigr) = -2z^N$.
\end{proof}

\begin{theorem}[The moment characterization]
\label{thm:cusp-moment}
\uses{lem:cusp-divisibility, lem:cusp-reduction, lem:cusp-exact}
\lean{ChabautyColeman.cuspMoment_step, ChabautyColeman.cuspMoment_congruence,
  ChabautyColeman.cuspMoment_base, ChabautyColeman.coeff_mul_comp_neg_X_eq_zero,
  ChabautyColeman.momentMatrix, ChabautyColeman.momentRhs,
  ChabautyColeman.momentMatrix_mulVec, ChabautyColeman.det_smul_eq_cramer_of_mulVec,
  ChabautyColeman.momentDet, ChabautyColeman.momentDet_linear,
  ChabautyColeman.moment_system_of_congruence, ChabautyColeman.momentDet_value_atom,
  ChabautyColeman.eq_of_det_momentMatrix_ne_zero}
\leanok
Let $T_j$ be the multiset $\{t_1, \dots, t_g\} \uplus \{t_1, \dots, t_j\}$, $G := g + j$,
$e_r := e_r(T_j)$, and $d_j(z) := \prod_i (t_i - z)^{\epsilon_i}$ with $\epsilon_i = 2$
for $i \le j$ and $1$ otherwise. If the avoidance values $A_i(0)$ for $i < j$ are
nonzero, then:
\begin{enumerate}
\item \textup{(master congruence)}\ \
  $A_j(z)\, d_j(-z) \equiv A_j(-z)\, d_j(z) \pmod{z^{2g+1}}$;
\item equivalently, writing $A_j = z^g + \sum_{m < g} \alpha_m z^m$: \
  $\sum_{m=0}^{g} \alpha_m\, e_{G-k+m} = 0$ for $k = 1, 3, \dots, 2g-1$
  \textup{(the \emph{moment system}; $\alpha_g := 1$)};
\item if moreover $\det M_j \neq 0$ for the system matrix
  $M_j := \bigl(e_{G-2i+1+m}\bigr)_{1 \le i \le g,\ 0 \le m \le g-1}$, then $A_j$ is the
  unique monic degree-$g$ solution, and each condition value of the fold appears as a
  pattern determinant of a once-extended alphabet. The phase-2 condition values are the
  mirror values $A_j(\mp t_j)$ of the retraced phase-1 stages
  (Lemma~\ref{lem:fold-mirror}).
\end{enumerate}
\end{theorem}
\begin{proof}
\leanok
(1) Induction. Base ($j = 0$): $d_0 = (-1)^g A_0$, so
$A_0(z)d_0(-z) = A_0(-z)d_0(z)$ holds \emph{exactly}. Composition multiplies the defect
by $(z^2 - t_{j+1}^2)$:
$A_c(z)d_c(-z) - A_c(-z)d_c(z) =
(z^2 - t_{j+1}^2)\bigl(A_j(z)d_j(-z) - A_j(-z)d_j(z)\bigr)$, where
$d_c := d_j \cdot (t_{j+1} - z)$. Reduction: by Lemma~\ref{lem:cusp-exact},
$A_{j+1}(z) \equiv A_{j+1}(-z)\,A_c(z)\,A_c(-z)^{-1} \pmod{z^{2g+1}}$, since the exact
defect $(-1)^{g+1}2z^N$ vanishes modulo $z^{2g+1}$ ($N = 2g+1$) and $A_c(-z)$ is
invertible modulo $z^{2g+1}$ precisely when $A_c(0) = -t_{j+1}A_j(0) \neq 0$ --- the
avoidance hypothesis. Hence the congruence for $A_c$ transfers to $A_{j+1}$.
(2) $\Gamma(z) := A(z)d(-z) - A(-z)d(z)$ is an odd polynomial, and for odd $k$ its
$z^k$-coefficient equals $2\sum_m \alpha_m e_{G-k+m}$ (write $d(-z) = \sum_r e_{G-r} z^r$
and collect the parity signs, which cancel); in characteristic $\neq 2$ the congruence
modulo $z^{2g+1}$ is thus exactly the system.
(3) Uniqueness is Cramer's rule; the value formulas follow from
$U_j(z^2) = (-1)^g A_j(z)A_j(-z)$ (Lemma~\ref{lem:cusp-reduction}) evaluated at
$z = t_{j+1}$, $t_j$, $0$.
\end{proof}

\begin{remark}[The Schur-polynomial interpretation of the atoms]
\label{rem:schur-atoms}
The determinant values certified in Theorem~\ref{thm:cusp-witness} by the
leading-coefficient induction admit an explicit Schur-polynomial theory. It is
\emph{deliberately left informal} --- the formalized route certifies the determinants
directly and does not use the Schur form --- and is included here as an independent
verification, checked symbolically in Magma, and as the structural explanation of the
$\pm 2^a$ integer contents observed in all numerical data.

\emph{(1) Cramer data are Schur polynomials.}
Let $\Lambda(j) := (g^{\,j+1}, g-1, g-2, \dots, 1)$. Then, for $0 \le m \le g$, the
$m$-th Cramer determinant satisfies
$\mathrm{NA}_m = (-1)^{g-m}\, s_{\Lambda(j)/(m)}(T_j)$ (a skew Schur polynomial, $(m)$
the one-row partition); in particular
$\det M_j = \pm\, s_{(g^j,\, g-1, \dots, 1)}(T_j)$ --- at $j = 0$ the classical
$s_{\delta}(t_1, \dots, t_g) = \prod_{i < l}(t_i + t_l)$ --- and by the one-variable
branching rule $\mathrm{NA}_j(-w) = \pm\, s_{\Lambda(j)}\bigl(T_j \cup \{w\}\bigr)$.
Consequently every fold condition is, up to sign and the factor $\det M_j^{2}$, a
product of two values of the single Schur polynomial $s_{\Lambda(j)}$ at multiset
alphabets: $s_{\Lambda(j)}(T_j \cup \{t\})\, \cdot\, s_{\Lambda(j)}(T_j \cup \{-t\})$
with $t = t_{j+1}$ (step) resp.\ $t = t_j$ (mirror), and $s_{\Lambda(j)}(T_j)^2$
(avoidance). The values at the all-positive alphabets are monomial-positive.
\emph{Why:} the matrix $M_j$ and its column-replaced Cramer variants are minors of the
Toeplitz array $(e_{r-c}(T_j))_{r,c}$; by the dual Jacobi--Trudi identity for skew
shapes, any such minor is $\pm$ a skew Schur polynomial, and matching the index pattern
$e_{G-2i+1+m}$ against $e_{\lambda'_i - \mu'_s - i + s}$ gives
$\lambda' = (G, G-1, \dots, G-g+1)$, i.e.\ $\lambda = \Lambda(j)$, and $\mu' = (1^m)$,
i.e.\ $\mu = (m)$; the sign $(-1)^{g-m}$ is the column transposition moving the
replaced column to its place. For $m = g$ the skew shape drops the full first row,
leaving the straight $(g^j, g-1, \dots, 1)$. The branching rule
$s_{\Lambda}(X \cup \{w\}) = \sum_{r \ge 0} s_{\Lambda/(r)}(X)\, w^r$ (one extra
variable picks up a horizontal strip, necessarily a single row) gives the stated form of
$\mathrm{NA}_j(-w)$; monomial-positivity of (skew) Schur polynomials is
Lindstr\"om--Gessel--Viennot.

\emph{(2) Nested-extreme factorization.}
Let $\Lambda$ be a partition, $\ell(\Lambda) \le n$. Partition the $n$ slots of
$s_\Lambda(x_1, \dots, x_n)$ into $k$ groups of sizes $m_1, \dots, m_k$ carrying
distinct values $v_1 > \dots > v_k$ (a priority order), and let $S(v_1, \dots, v_k)$ be
the resulting polynomial. Define $D_r := \deg_{v_r}$ and the iterated leading
coefficients along the priority order. Then, with $M_r := m_1 + \dots + m_r$ and the
\emph{slices} $\sigma_r := (\Lambda_{M_{r-1}+1}, \dots, \Lambda_{M_r})$:
$D_r = |\sigma_r|$, and the final iterated leading coefficient equals
$\prod_r N_r$ with $N_r := s_{\sigma_r}(1^{m_r}) \ge 1$, the number of semistandard
Young tableaux of the straight shape $\sigma_r$ with $m_r$ letters. Equivalently: the
monomial $v_1^{D_1} \cdots v_k^{D_k}$ is the lexicographically largest monomial of $S$
and has coefficient $\prod_r N_r$.
\emph{Why:} by symmetry assign group $r$ the entry range $\{M_{r-1}+1, \dots, M_r\}$;
the $v_1$-degree of a tableau's monomial is its number of entries $\le m_1$. In a
semistandard tableau the entry in row $i$ is $\ge i$, so entries $\le m_1$ lie in rows
$\le m_1$ and form a Young subdiagram $\mu \subseteq \Lambda$ (left-closed by
row-weakness, up-closed by column-strictness); thus
$|\mu| \le \Lambda_1 + \dots + \Lambda_{m_1}$ with equality iff $\mu$ is the full top
$m_1$ rows, and the bound is attained. A tableau attaining it is exactly a pair: an
arbitrary tableau of the straight shape $\sigma_1$ with $m_1$ letters ($N_1$ choices),
and an arbitrary filling of the lower rows with the larger entries; the lower part, with
entries shifted down by $m_1$, is an arbitrary tableau of the straight shape
$(\Lambda_{m_1+1}, \dots)$ on the remaining letters. Induction on $k$ finishes;
$N_r \ge 1$ since $\ell(\sigma_r) \le m_r$.

\emph{(3) Slice values.}
By Weyl's dimension formula, $s_{(a)}(1) = 1$, $s_{(a,b)}(1,1) = a - b + 1$ and
\[ s_{(a,b,c)}(1,1,1) = \tfrac12\,(a-b+1)(b-c+1)(a-c+2)\,. \]
Since consecutive rows of $\Lambda(j)$ differ by $0$ or $1$: every $2$-slice contributes $1$ or $2$; a
$3$-slice contributes $1$ on $(a,a,a)$, $8$ on $(a, a{-}1, a{-}2)$, and $3$ on the
boundary patterns $(a,a,a{-}1)$, $(a,a{-}1,a{-}1)$ --- the source of the
characteristic-$3$ fingerprint observed in the numerical data.

\emph{(4) Certificates for the positive atoms.}
Each of the values $s_{\Lambda(j)}(T_j)$, $s_{\Lambda(j)}(T_j \cup \{t_{j+1}\})$,
$s_{\Lambda(j)}(T_j \cup \{t_j\})$ ($j \ge 1$, $g \ge 2$) contains an explicit monomial
with coefficient a power of $2$; consequently each is nonzero over every field of
characteristic $\neq 2$, and has $2$-power integer content. Specifically: for the first
two alphabets (doubles and singles only) \emph{every} priority order works, with
$2$-exponent the number of difference-$1$ two-row slices; for the tripled alphabet,
place the triple first if $j \ge 2$ (slice $(g,g,g) \mapsto 1$) and after one single if
$j = 1$ (slice $(g, g{-}1, g{-}2) \mapsto 8$). This is (2) combined with (3).

\emph{(5) Signed group values.}
(a) $s_\mu(t, -t) = 0$ unless $\ell(\mu) \le 2$ and $\mu_1 \equiv \mu_2 \bmod 2$, in
which case it equals $(-1)^{\mu_2} t^{|\mu|}$; in particular
$s_{(g,g)}(t,-t) = (-1)^g t^{2g}$.
(b) $s_\mu(t, t, -t) = t^{|\mu|} s_\mu(1,1,-1)$, and
$s_{(a,a,a)}(1,1,-1) = s_{(a,a,a-1)}(1,1,-1) = (-1)^a$.
(For (a), $s_{(\mu_1,\mu_2)}(x,y) = (xy)^{\mu_2} h_{\mu_1-\mu_2}(x,y)$ and
$h_d(t,-t) = t^d$ for even $d$, $0$ for odd $d$. For (b),
$s_{(a,a,a)} = (xyz)^a \mapsto (-1)^a$;
$s_{(a,a,a-1)} = (xyz)^{a-1} e_2 \mapsto (-1)^{a-1}\cdot(1 - 1 - 1) = (-1)^a$. Note the
contrast with the \emph{unsigned} $s_{(a,a,a-1)}(1,1,1) = 3$: the signed specialization
rescues exactly the boundary pattern.)

\emph{(6) Certificates for the negated atoms.}
For $j \ge 1$, $g \ge 2$ and $\Lambda = \Lambda(j)$:
in $s_{\Lambda}(T_j \cup \{-t_{j+1}\})$ the $t_{j+1}$-degree is exactly $2g$, with
leading coefficient $(-1)^g\, s_{\Lambda''}\bigl(T_j \smallsetminus \{t_{j+1}\}\bigr)$,
$\Lambda''$ the shape with the top two rows removed; in
$s_{\Lambda}(T_j \cup \{-t_j\})$ the $t_j$-degree is exactly $3g$ ($j \ge 2$) resp.\
$3g - 1$ ($j = 1$), with leading coefficient $(-1)^g\, s_{\Lambda'''}(\text{rest})$,
$\Lambda'''$ the shape with the top three rows removed. Continuing with (2) on the
remaining alphabet, each negated atom contains an explicit monomial with coefficient
$\pm 2^a$, hence is nonzero over every field of characteristic $\neq 2$ with $2$-power
content. \emph{Why:} give the signed group --- the pair $\{t, -t\}$ (step case) resp.\
the trio $\{t, t, -t\}$ (mirror case) --- top priority and expand by the coproduct
$s_\Lambda(Y \uplus X) = \sum_\mu s_\mu(Y)\, s_{\Lambda/\mu}(X)$, so the $t$-degree
comes only from $s_\mu(Y)$. In the pair case
$|\mu| \le \Lambda_1 + \Lambda_2 = 2g$ with equality only for $\mu = (g,g)$, which
passes the parity condition of (5)(a) with value $(-1)^g$; in the trio case
$|\mu| \le \Lambda_1 + \Lambda_2 + \Lambda_3$ with equality only for the full top three
rows, with value $(-1)^g$ by (5)(b). In each case $\mu$ consists of full equal-length
top rows, so $\Lambda/\mu$ is the straight shape with those rows removed, its
consecutive-difference pattern is again $0/1$, and no larger $|\mu|$ contributes.

The statements of (1)--(6) were verified symbolically for small $g$ (scripts
\texttt{schurid.m}, \texttt{p2shapes.m}, \texttt{branch.m}, \texttt{p34atoms.m},
\texttt{c3check.m} in \texttt{prototype/}). The formalized counterparts of (4) and (6)
are the leading-coefficient cascades cited with Theorem~\ref{thm:cusp-witness}; the
degree and leading-coefficient statements of (6) are exactly the coefficient identities
of the peeled node there.
\end{remark}

\begin{theorem}[The cuspidal witness, all $g$]
\label{thm:cusp-witness}
\lean{ChabautyColeman.cusp_alphabet_comp_neg,
  ChabautyColeman.momentDet_mul, ChabautyColeman.peelPattern,
  ChabautyColeman.peelStage_single, ChabautyColeman.peelStage_double,
  ChabautyColeman.peelStage_id, ChabautyColeman.peelStage_triple_id,
  ChabautyColeman.peelStage_lead,
  ChabautyColeman.coeff_momentDet_of_stage, ChabautyColeman.degree_prod_coeff_lt,
  ChabautyColeman.coeff_momentDet_single, ChabautyColeman.coeff_momentDet_double,
  ChabautyColeman.coeff_momentDet_id, ChabautyColeman.coeff_momentDet_triple_id,
  ChabautyColeman.coeff_momentDet_lead, ChabautyColeman.peelHom,
  ChabautyColeman.momentDet_ne_zero_of_map,
  ChabautyColeman.momentDet_singles_ne_zero, ChabautyColeman.momentDet_doubles_ne_zero,
  ChabautyColeman.momentDet_X_peelPattern_one,
  ChabautyColeman.eval_ne_zero_of_momentDet_ne_zero,
  ChabautyColeman.momentDet_generic_doubles_ne_zero,
  ChabautyColeman.momentDet_generic_doubles_mul_X_ne_zero,
  ChabautyColeman.momentDet_generic_pair_ne_zero,
  ChabautyColeman.momentDet_generic_trio_ne_zero,
  ChabautyColeman.momentDet_generic_tripled_ne_zero,
  ChabautyColeman.momentDet_generic_trio_id_ne_zero,
  ChabautyColeman.momentDet_generic_tripled_one_ne_zero,
  ChabautyColeman.cuspGenH, ChabautyColeman.cuspChain_cond_of_atoms,
  ChabautyColeman.isCuspState_cuspChain_of_atoms,
  ChabautyColeman.cuspChain_congruence_of_atoms,
  ChabautyColeman.cuspChain_eval_ne_zero_of_atoms,
  ChabautyColeman.cuspStep_cond_base, ChabautyColeman.exists_cuspStep_H,
  ChabautyColeman.genNode, ChabautyColeman.genCuspChain,
  ChabautyColeman.genNode_trio_atom_ne_zero,
  ChabautyColeman.genNode_tripled_atom_ne_zero,
  ChabautyColeman.exists_genCuspChain_witness, ChabautyColeman.genFoldH,
  ChabautyColeman.genCuspFold, ChabautyColeman.exists_genCuspFold_witness,
  ChabautyColeman.genCuspFold_eval_sq_ne_zero,
  ChabautyColeman.isScaledState_foldBase_gen, ChabautyColeman.genFoldH_eq_foldGenH,
  ChabautyColeman.foldCondNum_cusp_ne_zero, ChabautyColeman.cuspNum_ne_zero}
\leanok
\uses{thm:cusp-moment, lem:fold-mirror}
For every $g \ge 2$ and every field $k$ of characteristic $\neq 2$: at the generic
nodes $t_1, \dots, t_g$ the full $2g$-step cuspidal fold is defined, with every
composition, disjointness and avoidance condition of both phases nonzero in
$k(t_1, \dots, t_g)$, and the moment determinants $\det M_j$ are nonzero for all
$j \le g$. Consequently the cleared condition numerators $P_q$ ($q < 2g$) of the
numerator fold at the cuspidal inputs are nonzero polynomials over $k$ --- the
nonvanishing hypothesis that Lemma~\ref{lem:cusp-descent} consumes.
\end{theorem}
\begin{proof}
\leanok
The values to be certified --- the determinants $\det M_j$, the avoidance values
$\mathrm{NA}_j(0)$, the step values $\mathrm{NA}_j(\mp t)$ and the mirror values
$\mathrm{NA}_j(\mp t_j)$ at the last-added node --- are nonzero at the generic nodes by
an exact leading-coefficient induction, one node at a time. Read the pattern determinant
in the polynomial ring in the peeled node $u$ over the remaining generic nodes. For each
of the five factor shapes occurring in the alphabets (a linear factor, its square, a
signed pair, a trio, and a tripled factor), the identity of the corresponding peel stage
states that the top $u$-coefficient of the determinant is a fixed scalar --- $1$, $2$,
$\pm 1$ or $8$ --- times the determinant of the remaining alphabet at the next pattern,
and no other term of the expansion reaches that degree. Nonvanishing therefore follows
by induction on the number of nodes over the integral domain of generic nodes; the
extension atoms peel their distinguished node first. No valuation-theoretic input
enters. The chain induction built on Theorem~\ref{thm:cusp-moment} then discharges
every phase-1 condition from these determinant values, and the mirror
(Lemma~\ref{lem:fold-mirror}) extends the induction through phase 2 along the retraced
phase-1 alphabets. Identifying the resulting scalars with the generic scalars of the
numerator fold and clearing denominators converts each condition into its numerator
$P_q$.

The Schur-side certificates of Remark~\ref{rem:schur-atoms} give an independent
verification; the numerical computation of every condition for $g \le 6$ from the moment determinants
(script \texttt{detcond.m}) agrees with the exact fold wherever both are available.
\end{proof}

\begin{lemma}[Descent from the cusp]
\label{lem:cusp-descent}
\uses{lem:fold-universal, thm:cusp-witness}
\lean{ChabautyColeman.descRate, ChabautyColeman.digits_eq_of_sum_eq,
  MvPolynomial.IsWeightedHomogeneous.aeval,
  ChabautyColeman.orderTop_aeval_descRate_le, ChabautyColeman.le_orderTop_mul_sub_one,
  ChabautyColeman.le_orderTop_prod_sub_one, ChabautyColeman.aeval_family_ne_zero,
  ChabautyColeman.exists_infinityPoint}
\leanok
Fix a field $K$ of characteristic $\neq 2$. Let $\mathcal{D} \ge 1$ bound the
$t$-degrees of the finitely many cuspidal numerators $P_q$ of
Lemma~\ref{lem:fold-universal}, choose a base $B > \mathcal{D}$, set
$M := \mathcal{D}B^{\,g+1}$, and use the \emph{common rates}
$\tilde n_i := M + B^{\,i+1}$ for every fold quantity $q$. If $P_q \neq 0$, then the
value of the universal numerator $\mathcal{N}_q$ at the family of points escaping to
infinity with rates $\tilde n_i$ is a \emph{nonzero} element of $K(\!(s)\!)$: its
valuation is that of the valuation-minimal monomial of
$P_q(t)\rvert_{t_i = s^{-\tilde n_i}}$, which lies in the window
$[\,-dM - \mathcal{D}B^{\,g},\, -dM\,]$ for $d = \operatorname{wt}(\mathcal{N}_q)$.
\end{lemma}
\begin{proof}
\leanok
\emph{(a) Separation and the window.} $P_q$ is homogeneous of some degree
$d \le \mathcal{D}$, so a monomial $t^c$ of $P_q$ specializes to valuation
$-\sum_i c_i \tilde n_i = -dM - \sum_i c_i B^{\,i+1}$: the $M$-part is common to
\emph{all} monomials (homogeneity), and since $0 \le c_i \le d \le \mathcal{D} < B$ the
digits $c_i$ make $c \mapsto \sum_i c_i B^{\,i+1}$ injective. Hence distinct monomials
land on distinct valuations, all inside the window
$[\,-dM - \mathcal{D}B^{\,g},\, -dM\,]$, and
$v^\ast := v\bigl(P_q(s^{-\tilde n})\bigr)$ is attained by a unique monomial; in
particular $P_q(s^{-\tilde n}) \neq 0$.

\emph{(b) The correction budget.} At the family, $y_i = t_i^N (1 + \eta_i)$ with
$v(\eta_i) \ge 2\tilde n_i \ge 2M$ (a Hensel square root), and $x_i = t_i^2$ exactly.
Evaluate $\mathcal{N}_q$ monomial by monomial. The deviation of an $a$-free monomial
$\mu$ from its cuspidal contribution has valuation
$\ge v(\mu(\text{cusp})) + 2M \ge (-dM - \mathcal{D}B^{\,g}) + 2M > -dM \ge v^\ast$,
using that \emph{every} $a$-free monomial of the isobaric $\mathcal{N}_q$ --- whether
or not it survives the collection into $P_q$ --- has $t$-degree exactly $d$ and hence
lies in the window of (a). A monomial containing some $a_r$ has $(x,y)$-weight at most
$d - 2$, so its value has valuation
$\ge -(d-2)(M + B^g) \ge -dM + 2M - \mathcal{D}B^{\,g} > -dM \ge v^\ast$. Hence
$v\bigl(\mathcal{N}_q(\text{family}) - P_q(s^{-\tilde n})\bigr) > v^\ast$, so
$\mathcal{N}_q(\text{family})$ has valuation exactly $v^\ast$; in particular it is
nonzero.
\end{proof}

\begin{theorem}[The fold witness for genus $\ge 2$]
\label{thm:fold-witness}
\lean{ChabautyColeman.exists_foldWitness, ChabautyColeman.curveCondNum,
  ChabautyColeman.curveSpec, ChabautyColeman.not_curveVanishing_curveCondNum,
  ChabautyColeman.eval_curveCondNum, ChabautyColeman.foldXs, ChabautyColeman.foldYs,
  ChabautyColeman.isScaledState_foldBase, ChabautyColeman.foldCond_ne_zero}
\leanok
\uses{lem:fold-universal, thm:cusp-witness, lem:cusp-descent}
Let $g \ge 2$ and let $K$ be algebraically closed of characteristic $\neq 2$, $f$ monic
of degree $N = 2g+1$. Every universal fold numerator $\mathcal{N}_q$ is a nonzero
polynomial function on $g$-tuples of affine curve points; consequently there is a
\emph{single} tuple $P_1, \dots, P_g$ over $K$, with pairwise distinct $x$-coordinates,
at which all $\mathcal{N}_q$ are simultaneously nonzero --- so the point-by-point fold
$\tau + \tau - \tau$ at this tuple is defined throughout and every disjointness and
doubling condition of both phases is nonzero, for every genus $\ge 2$. (Genus $1$,
where the cuspidal specialization degenerates in characteristic $3$, is
Lemma~\ref{lem:fold-witness-genus-one}.)
\end{theorem}
\begin{proof}
\leanok
Each $\mathcal{N}_q$ is nonzero on the variety of point tuples: the family of points
escaping to infinity with the common rates of Lemma~\ref{lem:cusp-descent} is a
$K(\!(s)\!)$-point at which $\mathcal{N}_q$ takes a nonzero value --- a pure polynomial
evaluation, requiring no other condition to hold, hence no circularity. A common
witness tuple then exists because finitely many nonzero polynomial conditions on $C^g$
have a common witness (a product of the conditions is refuted by one witness per
factor); the pairwise distinctness of the $x$-coordinates joins as $g(g-1)$ further
conditions $x_i - x_j$, each refuted at a two-point tuple. At the common tuple the fold
quantities are the ratios of Lemma~\ref{lem:fold-universal} with nonvanishing
numerators and denominators, inductively along the fold. Phase 2 needs no separate
argument: its conditions belong to the same uniform family of numerators, their
cuspidal nonvanishing having been established through the mirror inside
Theorem~\ref{thm:cusp-witness}. Further polynomial avoidance conditions --- the
Weierstrass avoidance $\operatorname{res}(U_j, f) \neq 0$ of the intermediate states,
or $y_i \neq 0$ --- join the same argument by enlarging the finite set of conditions,
their cuspidal nonvanishing being recorded there as well (on the cusp
$\operatorname{res}(U_j, X^N) = \pm U_j(0)^N$ and $U_j(0) = \pm A_j(0)^2 \neq 0$).
\end{proof}

\begin{lemma}[The genus-one fold witness]
\label{lem:fold-witness-genus-one}
\lean{ChabautyColeman.exists_genusOnePoint,
  ChabautyColeman.GenusOne.tangent_factorization,
  ChabautyColeman.GenusOne.foldChain_one, ChabautyColeman.GenusOne.foldChain_two,
  ChabautyColeman.GenusOne.foldGenH_zero_eq, ChabautyColeman.GenusOne.foldGenH_one_eq,
  ChabautyColeman.GenusOne.psi_eval_eq, ChabautyColeman.GenusOne.avoid_eval_eq,
  ChabautyColeman.GenusOne.psi_ne_zero, ChabautyColeman.GenusOne.avoid_ne_zero}
\leanok
\uses{lem:fold-universal}
Let $g = 1$ and let $K$ be algebraically closed of characteristic $\neq 2$, $f$ monic
cubic and \emph{squarefree}. Some affine curve point $(x_0, y_0)$ has $y_0 \neq 0$,
nonzero step-condition numerators at both fold steps, and nonzero stage resultants at
all three fold stages --- the genus-one instance of Theorem~\ref{thm:fold-witness}.

The cuspidal argument does not cover this case: on the cusp the interior condition
value at $g = 1$ is $3t^2/4$, identically zero in characteristic $3$ (the additive
group of the cusp is all $3$-torsion there, so $x(2P) = x(P)$ throughout), and for
$f = X^3$ the conditions genuinely vanish at every curve point --- squarefreeness is
sharp.
\end{lemma}
\begin{proof}
\leanok
A direct tangent--chord computation of the fold
(Lemma~\ref{lem:fold-universal} supplies the transfer to the numerator conditions).
With $\lambda = f'(x_0)/(2y_0)$, $x_2 = \lambda^2 - a - 2x_0$ and
$y_2 = -y_0 - \lambda(x_2 - x_0)$, the single cubic factorization
\[ (y_0 + \lambda(X - x_0))^2 - f \;=\; -(X - x_0)^2 (X - x_2) \]
governs both steps: the chord through $2P$ and $-P$ is the \emph{negated} tangent line,
so the second step needs no new identity, and the chain visits $(X - x_2, y_2)$ and
returns to $(X - x_0, y_0)$ --- $\tau + \tau - \tau = \tau$ on the nose. The condition
values are, up to powers of $2y_0$, the values at $x_0$ of $\psi_3 = 2ff'' - f'^2$ (the
$3$-division polynomial: $\psi_3(x_0) = 4y_0^2\,(x_0 - x_2)$) and of
$f'\psi_3 - 8f^2$ (the numerator of $y(2P)$: $8y_0^3\,y_2$); the stage resultants are
the $y$-values $y_0, y_2, y_0$ of the states, again up to denominators. Both
polynomials are nonzero for squarefree $f$ in \emph{every} characteristic, by
separability: if either vanished, $f \mid f'^2$, hence $f \mid f'$ by coprimality,
hence $f' = 0$ by degrees --- contradicting the coprimality of $f$ and $f' = 0$ for the
nonunit $f$. The witness point takes $x_0$ outside the finitely many roots of
$f \cdot \psi_3 \cdot (f'\psi_3 - 8f^2)$ and $y_0 = \sqrt{f(x_0)}$.
\end{proof}

\begin{lemma}[The auxiliary class $\tau$]
\label{lem:tau-exists}
\lean{ChabautyColeman.TauData, ChabautyColeman.TauData.exists_of_hasGoodCenter,
  ChabautyColeman.exists_interpolating_datum,
  ChabautyColeman.exists_ftauUY_root, ChabautyColeman.exists_isFoldCerts_lift,
  ChabautyColeman.exists_isFoldCerts_ftau, ChabautyColeman.isFtauFoldAvoid_of_center,
  ChabautyColeman.isFoldAvoid_of_map, ChabautyColeman.exists_bezout_lift_map,
  ChabautyColeman.TauData.exists_genus_two}
\leanok
\uses{lem:tau-generic, lem:many-points, lem:disk-param, lem:mumford-exists}
There exist $q$ and a class $\tau \in \J{\Qq}$ whose reduced datum $(u_\tau, v_\tau)$
lifts the witness datum of Lemma~\ref{lem:tau-generic} and satisfies
Definition~\ref{def:goodpos}, together with all the data the constructions below
require: an integral Mumford datum with its cofactor identity and a B\'ezout identity
$c_1 u_\tau + c_2 (2 v_\tau) = 1$ over $\Oq$; the $g$ points splitting it,
$u_\tau = \prod_i (X - a_i)$ with pairwise unit differences $a_i - a_j$ and
$v_\tau(a_i) = b_i$; the \emph{moving points} of the fold --- the Hensel roots of the
perturbed $u$-polynomial over $\Oq[[X;Y]]$, with constant coefficients the $a_i$ ---
and valid B\'ezout data for the fold computation of $F_\tau$ over $\Oq[[X;Y]]$,
avoiding the Weierstrass points at every stage. The residue-level genericity of
Lemma~\ref{lem:tau-generic} travels with the datum: the lifts have prescribed residues,
and the avoidance conditions descend along the residue map. The avoidance of the
prescribed finite set $T$ of residue $x$-values --- $\bar u_\tau(x) \neq 0$ for
$x \in T$ --- also comes with the datum; the per-disk construction of
Lemma~\ref{lem:global-generic} consumes it.
\end{lemma}
\begin{proof}
\leanok
The witness tuple of Lemma~\ref{lem:tau-generic} lifts point by point, by Hensel's
lemma, to disk centres with prescribed residues (Lemma~\ref{lem:disk-param}); $u_\tau$
is the product of the lifted linear factors, and $v_\tau$ is interpolated through the
lifted points by the Chinese remainder theorem --- the lifted $x$-coordinates differ by
units, so the linear factors are pairwise comaximal --- with $\bar v_\tau$ recovered
exactly by uniqueness of interpolation in degree $< g$. The B\'ezout identity over
$\Oq$ lifts by Nakayama's lemma from the stage-$0$ avoidance of the witness. The moving
points exist by Hensel's lemma along the ideal generated by the variables: the
$\tau$-roots are simple roots of $u_\tau$, their derivative values being unit products
of the pairwise differences. The B\'ezout data for the fold lift step by step, again by
Nakayama with prescribed residues against the already-lifted states, and the
Weierstrass avoidance descends from the residue configuration, the residue map
reflecting coprimality of monic polynomials.
\end{proof}

\begin{definition}[Cantor addition as a rational map]
\label{def:cantor-map}
\lean{ChabautyColeman.cantorV, ChabautyColeman.cantorV_sq_sub,
  ChabautyColeman.y_sub_algebraMap_cantorV, ChabautyColeman.cantorCompose,
  ChabautyColeman.cantorReduce}
\leanok
\uses{lem:mumford-exists}
The pair of tuples of rational functions
$\mathrm{Add} = (\mathrm{Add}_u, \mathrm{Add}_v)$ in the coordinates of two Mumford
data of degree $g$, obtained by transcribing Cantor's composition and reduction
symbolically; its domain of definition contains the locus where the finitely many
leading coefficients and resultants appearing in the algorithm are invertible. The
inverse is $\ominus(u, v) := (u, -v)$ (Lemma~\ref{lem:aj-inverse}). The composition is
the coprime branch of Cantor's algorithm with the three-term B\'ezout relation
$s_1 u_1 + s_2 u_2 + s_3(v_1 + v_2) = 1$: unlike the extended Euclidean algorithm,
which is not a rational map near the diagonal, this branch uniformly covers
composition, the near-diagonal case, and doubling.
\end{definition}

\begin{lemma}[Correctness of the rational map]
\label{lem:cantor-correct}
\lean{ChabautyColeman.mumfordIdeal_mul_of_bezout,
  ChabautyColeman.mk0_mumfordIdeal_mul_of_bezout,
  ChabautyColeman.mk0_mumfordIdeal_eq_mk0_cofactor_datum,
  ChabautyColeman.exists_bezout_lift}
\leanok
\uses{def:cantor-map, lem:mumford-exists}
At every point of its domain over any field, $\mathrm{Add}$ evaluates to the reduced
datum of the product class. More precisely, in the certificate-passing form used
throughout: given a B\'ezout identity $s_1 u_1 + s_2 u_2 + s_3(v_1+v_2) = 1$, the
product of the two Mumford ideals \emph{equals} the Mumford ideal of the composed datum
$(u_1 u_2,\ V)$ with $V = s_1 u_1 v_2 + s_2 u_2 v_1 + s_3(v_1 v_2 + f)$, so the
composition computes the product class; and given an inverse of the leading coefficient
of the cofactor $w$ of $v^2 - f = uw$, the normalized reduced datum
$(c\,w,\ (-v) \bmod c\,w)$ is monic, is again a Mumford datum, and represents the same
class (Lemma~\ref{lem:mumford-reduce}). Over a local ring, a B\'ezout identity for the
residues lifts by Nakayama's lemma, since the Mumford $u$-polynomials are monic
(monicity is essential: $1 - \pi X$ over a discrete valuation ring has invertible
residue but is not invertible).
\end{lemma}
\begin{proof}
\leanok
The composed ideal identity is proved directly from the certificate:
$V \equiv v_i \bmod u_i$ and $V^2 \equiv f \bmod u_1 u_2$ with explicit cofactors, and
$y - V$ is an explicit combination of the four product generators; the reduction step
reduces to Lemma~\ref{lem:mumford-reduce} after normalizing by the leading-coefficient
inverse.
\end{proof}

\begin{definition}[The chart at $\tau$]
\label{def:chart}
\lean{ChabautyColeman.chartU, ChabautyColeman.chartV, ChabautyColeman.exists_chartV}
\leanok
\uses{lem:tau-exists, thm:hensel}
Coordinates near $\tau$: the $g$ coefficients $X = (X_1,\dots,X_g)$ of $u - u_\tau$.
There is $V(X) \in \Oq[[X]][x]$ of $x$-degree $< g$ with $V(0) = v_\tau$ and
$V(X)^2 \equiv f \bmod (u_\tau + X)$; the chart point of $X$ is the Mumford datum
$(u_\tau + X,\ V(X))$, i.e.\ the ideal $(u_\tau + X,\ y - V(X))$. \emph{Formalization.}
$V$ is constructed as a single quadratic Hensel root (Theorem~\ref{thm:hensel}) of
$T^2 = f$ in the quotient by $(u_\tau + X)$, a finite free algebra over the power
series ring, which is Henselian for the $(X)$-adic ideal; no topology on $\Oq$ is
needed. The inputs are the cofactor identity $v_\tau^2 - f = u_\tau w$ and the B\'ezout
identity $c_1 u_\tau + c_2 (2 v_\tau) = 1$ carried by the datum of
Lemma~\ref{lem:tau-exists}.
\end{definition}

\begin{definition}[The translated group law $F_\tau$]
\label{def:Ftau}
\lean{ChabautyColeman.FtauFold, ChabautyColeman.ftauFoldPipe,
  ChabautyColeman.FoldCertData, ChabautyColeman.IsFoldCerts,
  ChabautyColeman.IsFtauFoldCerts, ChabautyColeman.ftauFoldPipe_fst_eq,
  ChabautyColeman.ftauFoldXs, ChabautyColeman.ftauFoldYs,
  ChabautyColeman.ftauUY_eq_prod}
\leanok
\uses{def:chart, def:cantor-map, lem:tau-exists}
$F_\tau(X,Y)$ is defined as the chart coordinates of
$\mathrm{ch}(X) + \mathrm{ch}(Y) - \tau$, expanded as a formal power series at $(0,0)$:
in the $\tau$-translated identification $c(X) := \mathrm{ch}(X) - \tau$ this is the
group law, satisfying $\mathrm{ch}(F_\tau(X,Y)) = \mathrm{ch}(X) + \mathrm{ch}(Y) -
\tau$. \emph{Formalization.} The computation is realized as the point-by-point fold of
Definition~\ref{def:centre-fold} over $\Oq[[X;Y]]$: starting from the $X$-side chart
datum, fold in, one point per step, first the $g$ moving points of the $Y$-side chart
--- the Hensel roots of the $Y$-side chart polynomial, which factor it completely ---
and then the $g$ involuted $\tau$-points $(a_i, -b_i)$; each step is one Cantor
composition at a B\'ezout triple followed by the unconditional reduction with leading
coefficient $-1$. The B\'ezout data are a parameter of the definition, one triple per
step; given their validity, the output is a Mumford datum whose $u$-part is exactly the
chart polynomial at $F_\tau(X,Y)$. Existence of valid data over $\Oq[[X;Y]]$ is part of
Lemma~\ref{lem:tau-exists}.
\end{definition}

\begin{lemma}[Integrality, unit laws, commutativity]
\label{lem:Ftau-basic}
\lean{ChabautyColeman.setYZero_FtauFold, ChabautyColeman.setXZero_FtauFold,
  ChabautyColeman.rename_swap_FtauFold, ChabautyColeman.constantCoeff_FtauFold,
  ChabautyColeman.fold_pipeline_eq_self, ChabautyColeman.mk0_foldPipe,
  ChabautyColeman.prod_mk0_foldXs_foldYs, ChabautyColeman.mk0_mumfordIdeal_mul_neg}
\leanok
\uses{def:Ftau, lem:cantor-correct}
$F_\tau \in \Oq[[X;Y]]^g$ (the denominators of Definition~\ref{def:cantor-map} are
units at the centre by Definition~\ref{def:goodpos}); $F_\tau(X,0) = X$,
$F_\tau(0,Y) = Y$, $F_\tau(X,Y) = F_\tau(Y,X)$; and for all
$\mathbf x, \mathbf y \in (\mathfrak m')^g$, $\mathfrak m'$ the maximal ideal of a
complete discrete valuation ring $\Oq'$ over $\Oq$,
$\mathrm{ch}(F_\tau(\mathbf x,\mathbf y))$ represents the class
$\mathrm{ch}(\mathbf x) + \mathrm{ch}(\mathbf y) - \tau$: $F_\tau$ computes the sum in
the $\tau$-translated identification.
\end{lemma}
\begin{proof}
\leanok
Integrality holds by construction, and the vanishing constant term expresses that the
value of the fold at the centre is the $\tau$-datum itself. For the unit law
$F_\tau(X,0) = X$: at $Y = 0$ the moving points degenerate onto the $\tau$-points, so
the fold computes $[\mathrm{ch}(X)] + [\tau] - [\tau] = [\mathrm{ch}(X)]$ --- the
mirrored point block cancels and reduced representatives are unique
(Lemma~\ref{lem:cantor-correct}, Lemma~\ref{lem:mumford-exists}); the identity of
power series follows by evaluating at all points and interpolating
(Theorem~\ref{thm:interpolation}; interpolation over $\Oq$-points suffices).
Commutativity is proved the same way: the swapped computations have the same class,
hence the same reduced output. The last clause is the pointwise class computation of
the fold, transported to evaluation points.
\end{proof}

\begin{theorem}[Associativity]
\label{thm:Ftau-assoc}
\lean{ChabautyColeman.FtauFold_assoc, ChabautyColeman.eval_FtauFold_eq_of_mk0_eq,
  ChabautyColeman.fold_pipeline_eq_of_mk0_eq, ChabautyColeman.mk0_chart_eval_fold,
  ChabautyColeman.chartClass, ChabautyColeman.eval_ftauFoldPipe_eq,
  ChabautyColeman.ftauFoldPipe_constantCoeff_eq, ChabautyColeman.constantCoeff_FtauFold,
  ChabautyColeman.eq_of_dvd_sq_sub_of_bezout}
\leanok
\uses{lem:Ftau-basic, thm:interpolation, lem:cantor-correct}
$F_\tau(F_\tau(X,Y),Z) = F_\tau(X,F_\tau(Y,Z))$; hence $F_\tau$ is a $g$-dimensional
commutative formal group law over $\Oq$ in the sense of Definition~\ref{def:fgl}.
\end{theorem}
\begin{proof}
\leanok
Both sides are power series over $\Oq$ agreeing at every point of the polydisk, because
at points both compute the associative class-group sum
$\mathrm{ch}(x)+\mathrm{ch}(y)+\mathrm{ch}(z)-2\tau$; conclude by
Theorem~\ref{thm:interpolation}. The pointwise statement rests on: the chart at the
evaluated coordinates $F_\tau(\mathbf x, \mathbf y)$ represents
$\mathrm{ch}(\mathbf x) + \mathrm{ch}(\mathbf y) - \tau$ (the evaluated fold output is
the chart datum at the evaluated coordinates --- the $u$-parts agree by the shape of
the result, the $v$-parts by uniqueness of square roots with a B\'ezout certificate);
and computations with class-equal inputs have equal reduced outputs, by uniqueness of
reduced representatives.
\end{proof}

\begin{lemma}[Reduction commutes with reduced representatives at good position]
\label{lem:reduce-commutes}
\lean{ChabautyColeman.red_reducedRep_eq, ChabautyColeman.red_reducedRep_points,
  ChabautyColeman.prod_pointIdeal_eq_reduced,
  ChabautyColeman.eq_of_mumfordIdeal_eq_prod_pointIdeal, ChabautyColeman.redres_sq_curve}
\leanok
\uses{thm:redJ-hom, lem:mumford-exists, def:goodpos}
Let $c \in \J{\Qq}$ be such that $\redJ(c)$ has a reduced datum
$(\bar u_0, \bar v_0)$ over $\Fq$ with $\bar u_0$ squarefree of full degree $g$. Then
the reduced datum $(u, v)$ of $c$ has coefficients in $\Oq$, $\deg u = g$, and
$(\bar u, \bar v) = (\bar u_0, \bar v_0)$.
\end{lemma}
\begin{proof}
\leanok
Compute $\redJ(c)$ from the points of the reduced representative of $c$
(Definition~\ref{def:redJ}): the split Mumford representative over the splitting field
$L$ of $u$ factors into point ideals, and the factorization determines the datum by the
norm identity $\mathfrak a\,\iota(\mathfrak a) = (u)$. Reduce it to the residue curve
and match against $(\bar u_0, \bar v_0)$ by semi-reduction and uniqueness of reduced
representatives (Lemma~\ref{lem:mumford-exists} over $\Fq$): since $\bar u_0$ has full
degree $g$ and is squarefree, no cancellation and no contribution from the disk at
infinity can occur, forcing every point of the representative to be integral with
multiplicity one and distinct reductions. Integrality of $(u, v)$ follows: $u$ has
integral roots, and $v$ is the interpolant through the integral points, the
$x$-differences being units.
\end{proof}

\begin{theorem}[Chart parametrization of the kernel of reduction]
\label{thm:chart-mem}
\lean{ChabautyColeman.J1_equiv_polydisk, ChabautyColeman.polydiskToKer,
  ChabautyColeman.polydiskToKer_bijective, ChabautyColeman.polydiskToKer_eval_FtauFold,
  ChabautyColeman.exists_chartClass_eq, ChabautyColeman.chartClass_injective,
  ChabautyColeman.redJ_chartClass, ChabautyColeman.redJ_tauClass,
  ChabautyColeman.redJ_mk0_mumford, ChabautyColeman.mumfordIdeal_eq_prod_pointIdeal}
\leanok
\uses{thm:redJ-hom, lem:reduce-commutes, thm:hensel, def:chart, lem:Ftau-basic}
The map $c \mapsto{}$chart coordinates of $c + \tau$ is a group isomorphism
$J^1(\Qq) \xrightarrow{\ \sim\ } \big((p\Oq)^g,\ +_{F_\tau}\big)$.
\end{theorem}
\begin{proof}
\leanok
For $c \in J^1$: $\redJ(c+\tau) = \redJ(\tau)$, and the reduction of the class of an
integral datum with squarefree reduced $u$-polynomial is the class of the reduced datum
(split $u$ over $L$ with integral roots --- squarefree $\bar u$ lifts to separability
--- factor the Mumford ideal into point ideals on both levels, and match the
reductions). So $\redJ(c+\tau) = [(\bar u_\tau, \bar v_\tau)]$, and
Lemma~\ref{lem:reduce-commutes} shows the reduced representative of $c + \tau$ is
integral with reduction $(\bar u_\tau, \bar v_\tau)$. Its $u$-part is the chart
polynomial at explicit coordinates $z \in (p\Oq)^g$ (the displacement coefficients),
and its $v$-part is the chart square root, by uniqueness of square roots with a
B\'ezout certificate at the common residue datum. Conversely every chart class reduces
to $\redJ(\tau)$, so chart classes translated by $-\tau$ lie in $J^1$; injectivity is
uniqueness of reduced Mumford representatives (Lemma~\ref{lem:mumford-exists}); the
group compatibility is the pointwise class computation of
Theorem~\ref{thm:Ftau-assoc} together with $[(u,-v)] = [(u,v)]^{-1}$.
\end{proof}

\begin{definition}[The logarithm on the kernel of reduction]
\label{def:Lam}
\lean{ChabautyColeman.Lam, ChabautyColeman.Lam_injective, ChabautyColeman.Lam_add,
  ChabautyColeman.ftauFoldFGL, ChabautyColeman.kerScaledCoords,
  ChabautyColeman.FormalGroupLaw.scaledLog}
\leanok
\uses{thm:Ftau-assoc, thm:log-iso, thm:chart-mem}
$\Lam := \log_{F_\tau} \circ\ (\text{chart coordinates of } \cdot + \tau)\colon
J^1(\Qq) \to \Qq^{\,g}$, a group homomorphism by Theorems~\ref{thm:chart-mem}
and~\ref{thm:log-iso}. ($\Lam$ is moreover injective; for $e = 1$ this is formalized in
\texttt{HyperellipticCurves.Archive.LogInjective}, and the main theorems do not depend on it.)
\emph{Formalization.} $\Lam$ is defined in the $\pi$-scaled coordinates of
Theorem~\ref{thm:log-iso}: its value is the total evaluation of the integral $\pi$-scaled
logarithm at $\mathbf y$, where $\pi\mathbf y$ is the chart coordinate vector of
$c + \tau$; additivity combines the group compatibility of
Theorem~\ref{thm:chart-mem} with Theorem~\ref{thm:log-iso}.
\end{definition}

\begin{lemma}[Rationality of the invariant differential]
\label{lem:delta-rational}
\lean{ChabautyColeman.invariantDiffFold_rational, ChabautyColeman.DatumFamRing,
  ChabautyColeman.ftauFamSpec,
  ChabautyColeman.traceDiff_prod, ChabautyColeman.traceDiff_foldPipe,
  ChabautyColeman.sum_traceMat_ftauFoldPipe,
  ChabautyColeman.sum_traceMat_mul_setYZero_pderiv_FtauFold,
  ChabautyColeman.traceMat_map_setYZero_eq_C,
  ChabautyColeman.det_traceMatrix_ne_zero,
  Polynomial.exists_bezout_of_isUnit_resultant,
  ChabautyColeman.ftauFoldPipe_eq_of_certs, ChabautyColeman.FtauFold_eq_of_certs}
\leanok
\uses{def:Ftau, def:fgl-log}
$M(X) = (\partial F_\tau/\partial Y)(X,0)$ is the power-series expansion of a matrix of
\emph{rational} functions of the full Mumford datum, invertible at $0$; hence
$\delta = M^{-1}dX$ is a $g$-tuple of rational $1$-forms on the chart. (Rationality
holds in the coordinates of the full datum; the chart square root is transcendental in
the chart coordinates alone.)
\end{lemma}
\begin{proof}
\leanok
$F_\tau$ is the expansion of a composite of rational maps
(Definitions~\ref{def:cantor-map} and~\ref{def:Ftau}); differentiation preserves
rationality, and $M(0) = I$. \emph{Formalization.} Each entry of $M$ is exhibited as
the specialization, at the chart data with $Y = 0$, of a universal element of a
three-block coefficient ring localized away from an explicit list of denominators, so
that the entries re-specialize along any family of Mumford data. The route is through
the trace form of Theorem~\ref{thm:omega-prime-global}: the invariance identity of the
fold specializes at $Y = 0$ to $T(X)\cdot M(X) = T_2$ with $T$ the trace matrix of the
$X$-side chart and $T_2$ constant, so $M = T^{-1}T_2$ is determined by the starting
data alone, and all four denominators (three block resultants and the determinant of
the trace matrix) specialize to units unconditionally.
\end{proof}

\begin{definition}[Curve-to-chart maps]
\label{def:phiD}
\leanok
\lean{ChabautyColeman.pairV, ChabautyColeman.pair_dvd, ChabautyColeman.diskY_sq,
  ChabautyColeman.movingPair_dvd, ChabautyColeman.exists_isPipelineCerts_of_center,
  ChabautyColeman.exists_isPipelineCerts_powerSeries, ChabautyColeman.chartCoords,
  ChabautyColeman.phiPipeline, ChabautyColeman.phiZeta,
  ChabautyColeman.phiPipeline_map_constantCoeff, ChabautyColeman.constantCoeff_phiZeta}
\uses{def:cantor-map, def:chart, lem:tau-exists, lem:aj-inverse}
For an auxiliary point $Q$ in good position relative to $\tau$, the rational map
$\phi_Q \colon C \dashrightarrow \text{chart}$ sends $P$ to the chart coordinates of
$[P + \iota Q - 2\infty] + \tau$; it is regular at all $P$ avoiding an explicit proper
closed locus. On a residue disk with parameter $z$ and centre $P_0$, the map is
realized as a three-term Cantor computation over $\Oq[[z]]$ on the data of the moving
pair $[P_z + \iota Q]$, the $\tau$-datum, and the constant pair $[Q + \iota P_0]$; its
chart coordinates form the coordinate series $\zeta(z)$. At $z = 0$ the computation
returns the $\tau$-datum exactly --- the two pair classes cancel and reduced
representatives are unique --- so $\zeta(0) = 0$. As everywhere, the B\'ezout data of
the steps are parameters; their existence reduces to residue-level genericity
(Lemma~\ref{lem:global-generic}) by Nakayama lifting. The computation is ordered so
that all support overlaps are of doubling type: cancelling overlaps (a point meeting
its involute) are genuinely singular for the certificate-based composition, which is
what shapes the good-position locus.
\end{definition}

\begin{lemma}[Expansion injectivity]
\label{lem:exp-inj}
\lean{ChabautyColeman.funcExpansionGeneric_eq_zero_iff,
  ChabautyColeman.funcExpansionGeneric_injective}
\leanok
\uses{lem:disk-param}
A rational differential on $C$ over $\Qq$ whose full expansion at one point of
$C(\Qq)$ vanishes is zero; two rational differentials with equal expansions at one
point are equal.
\end{lemma}
\begin{proof}
\leanok
Stated for functions $P + Q\,y$ ($P, Q \in \Qq[x]$) at an affine non-Weierstrass point,
to which rational functions and differentials reduce by clearing denominators: the
expansion times its conjugate expansion is the Taylor series of the norm polynomial
$P^2 - fQ^2$, which is nonzero for $(P,Q) \neq (0,0)$ by the degree-parity count.
\emph{Formalization.} The main development compares series through uniqueness of
reduced representatives instead; this lemma is formalized in
\texttt{HyperellipticCurves.Archive.ExpInj}, and the main theorems do not depend on it.
\end{proof}

\begin{theorem}[The differentials $\omega'_i$]
\label{thm:omega-prime-global}
\lean{ChabautyColeman.omegaPrime, ChabautyColeman.pderiv_subst_log,
  ChabautyColeman.cantorReduce_fst, ChabautyColeman.map_reduceCert,
  ChabautyColeman.exists_reduceCert_map, ChabautyColeman.natDegree_cantorReduce_of_fullDeg,
  ChabautyColeman.isCoprime_two_mul_of_map, ChabautyColeman.isCoprime_cantorStage₂_of_map,
  ChabautyColeman.isCoprime_cantorPipeline_of_map, ChabautyColeman.isCoprime_phiStage₂_of_residue,
  ChabautyColeman.isCoprime_phiPipeline_of_residue, ChabautyColeman.isCoprime_wStage₂_of_residue,
  ChabautyColeman.isCoprime_wPipeline_of_residue,
  ChabautyColeman.subst_log_eq_antideriv, ChabautyColeman.mk0_phiPipeline_eval,
  ChabautyColeman.phiPipeline_eq_of_aux,
  ChabautyColeman.cantorPipeline_eq_cantorPipeline_of_mk0_eq,
  ChabautyColeman.mk0_pair_mul_mk0_pair_chain,
  ChabautyColeman.FractionTopology.valuedFractionRing,
  ChabautyColeman.FractionTopology.continuous_algebraMap,
  ChabautyColeman.FractionTopology.isTopologicallyNilpotent_algebraMap,
  ChabautyColeman.hasSum_algebraMap_evalT_scaledLog, ChabautyColeman.hasSum_algebraMap_Lam,
  ChabautyColeman.redJ_diskPointDiff_eq_one, ChabautyColeman.aj_eq_mk0_mumfordIdeal,
  ChabautyColeman.FractionTopology.isUniformInducing_algebraMap,
  ChabautyColeman.FractionTopology.instCompleteSpace,
  ChabautyColeman.phiPipeline_monic_dvd_degree_lt,
  ChabautyColeman.map_evalRingHom_eq_add_chartCoords,
  ChabautyColeman.map_evalRingHom_eq_chartV_map,
  ChabautyColeman.chartClass_eq_diskPointDiff_mul_tauClass,
  ChabautyColeman.J1_equiv_polydisk_diskPointDiff, ChabautyColeman.summable_logSubstFamily,
  ChabautyColeman.hasSum_coeff_subst_log, ChabautyColeman.hasSum_Lam_antiderivOmegaPrime,
  ChabautyColeman.HypField.deriv, ChabautyColeman.globPipeline, ChabautyColeman.globZeta,
  ChabautyColeman.globOmega, ChabautyColeman.globExpand,
  ChabautyColeman.map_globExpand_globPipeline, ChabautyColeman.globExpand_globZeta,
  ChabautyColeman.globExpand_deriv, ChabautyColeman.globExpand_globOmega,
  ChabautyColeman.regularDiffBasis, ChabautyColeman.map_chartU_eq_of_chartCoords,
  ChabautyColeman.map_chartV_eq_of_chartCoords, ChabautyColeman.map_chartU_eq_phiPipeline_fst,
  ChabautyColeman.map_chartV_eq_phiPipeline_snd, ChabautyColeman.globExpand_comp_datumFamSpec,
  ChabautyColeman.isUnit_datumFamSpec_glob, ChabautyColeman.exists_globN,
  ChabautyColeman.HypRing.exists_toHypField_mul_eq,
  ChabautyColeman.HypRing.exists_toHypField_mul_eq_finite,
  ChabautyColeman.toHypField_eq_of_bezout, ChabautyColeman.HypField.deriv_toHypField_C,
  Polynomial.topCoeffMod, Polynomial.topCoeffMod_eq_zero_of_natDegree_lt,
  Polynomial.topCoeffMod_mul_monic, Polynomial.topCoeffMod_neg,
  Polynomial.topCoeffMod_congr, Derivation.dotDer, Derivation.coeff_dotDer,
  Derivation.dotDer_eq_zero_of_coeff, ChabautyColeman.traceDiff,
  ChabautyColeman.traceDiff_eq_zero_of_dotDer_eq_zero, ChabautyColeman.dvd_sub_of_bezout,
  ChabautyColeman.traceDiff_mul, ChabautyColeman.traceDiff_reduce,
  ChabautyColeman.sub_mul_topCoeffMod_pair, ChabautyColeman.traceDiff_pair,
  Derivation.natDegree_dotDer_le, Polynomial.natDegree_sq_sub_of_odd,
  ChabautyColeman.traceDiff_congr, ChabautyColeman.exists_bezout_of_dvd_sub,
  ChabautyColeman.isCoprime_two_mul_of_dvd_add,
  ChabautyColeman.traceDiff_reduce_of_natDegree_lt, Derivation.dvd_dotDer_of_dvd_sub,
  ChabautyColeman.traceDiff_reduce_of_dvd_dotDer_of_natDegree_lt,
  ChabautyColeman.traceDiff_reduce_of_dvd_dotDer,
  ChabautyColeman.traceDiff_cantorReduce_of_dvd_dotDer,
  ChabautyColeman.traceDiff_cantorPipeline_of_dvd_dotDer, ChabautyColeman.map_residue_eq,
  ChabautyColeman.traceMat_wPipeline, ChabautyColeman.wU₁_map_seriesResidue,
  ChabautyColeman.wV₁_map_seriesResidue, ChabautyColeman.isPipelineAvoid_w_of_residue,
  ChabautyColeman.cantorCompose_eq_of_isCoprime, ChabautyColeman.isCoprime_wU₁_map,
  ChabautyColeman.two_mul_X_mul_traceDiff_wPipeline_map,
  ChabautyColeman.eq_aeval_mul_of_two_mul_eq,
  ChabautyColeman.X_mul_omegaPrime_eq_aeval_tracePoly_mul_derivative,
  ChabautyColeman.map_omegaExpansionWeierstrass_eq_sum_smul_omegaPrime_tracePoly,
  ChabautyColeman.aeval_ne_zero_of_coeff_zero_ne_zero, ChabautyColeman.aeval_inv_eq_zero_iff,
  ChabautyColeman.invAevalFracHom, ChabautyColeman.sq_inv_pow_mul_X,
  ChabautyColeman.inftyExpand, ChabautyColeman.traceDiff_single, ChabautyColeman.IsAddCerts,
  ChabautyColeman.cantorStage₂_monic_dvd_degree_lt, ChabautyColeman.traceDiff_cantorStage₂,
  ChabautyColeman.IsAddCerts.map, ChabautyColeman.mk0_cantorStage₂,
  ChabautyColeman.map_infSeries_eq, ChabautyColeman.traceMat_infU,
  ChabautyColeman.kzDeriv_aeval, ChabautyColeman.kzDeriv_inv,
  ChabautyColeman.two_mul_y_mul_traceDiff_infU_map, ChabautyColeman.infSeries_mul_omegaPrime_kz,
  ChabautyColeman.map_omegaExpansionInf_eq_sum_smul_omegaPrime_tracePoly,
  ChabautyColeman.HypField.derivation, ChabautyColeman.HypField.deriv_smul,
  ChabautyColeman.HypField.deriv_algebraMap_baseField, MvPowerSeries.pderiv,
  RatFunc.derivative_C, RatFunc.derivative_smul, ChabautyColeman.dvd_cantorCompose_sub_left,
  ChabautyColeman.dvd_cantorCompose_sub_right, ChabautyColeman.isCoprime_cantorCompose,
  ChabautyColeman.traceDiff_cantorCompose, ChabautyColeman.dvd_cantorReduce_add,
  ChabautyColeman.traceDiff_cantorReduce, ChabautyColeman.traceDiff_cantorPipeline,
  ChabautyColeman.HypField.deriv_x, Polynomial.topCoeffMod_sum,
  Polynomial.topCoeffMod_C_mul, Derivation.dotDer_C_mul_X_pow, ChabautyColeman.traceMat,
  ChabautyColeman.traceDiff_eq_sum, ChabautyColeman.chartU_map_rename_eq,
  ChabautyColeman.pderiv_coeff_map_C_eq_zero, ChabautyColeman.traceDiff_chartU_map_rename,
  ChabautyColeman.traceDiff_ftauUX_eq_zero, ChabautyColeman.traceDiff_ftauUY,
  ChabautyColeman.IsFoldCerts.isFoldState, ChabautyColeman.traceDiff_prod,
  ChabautyColeman.traceDiff_foldPipe, ChabautyColeman.sum_traceMat_ftauFoldPipe,
  ChabautyColeman.ftauUY_eq_prod, Polynomial.topCoeffMod_map, ChabautyColeman.traceMat_map,
  ChabautyColeman.ftauFoldPipe_map_setYZero,
  ChabautyColeman.sum_traceMat_mul_setYZero_pderiv_FtauFold,
  Polynomial.exists_topCoeffMod_ne_zero, ChabautyColeman.traceMatrix,
  ChabautyColeman.traceMatrix_comm, ChabautyColeman.sum_mul_traceMatrix,
  ChabautyColeman.eq_zero_of_sum_mul_traceMatrix, ChabautyColeman.det_traceMatrix_ne_zero,
  ChabautyColeman.dvd_sub_of_bezout_of_bezout, ChabautyColeman.traceMat_congr,
  ChabautyColeman.traceMat_map_setYZero_eq_C,
  ChabautyColeman.sum_traceMat_mul_setYZero_pderiv_FtauFold_eq_C,
  ChabautyColeman.traceMat_phiPipeline, ChabautyColeman.globExpand_traceMat,
  ChabautyColeman.eq_mul_of_mul_eq,
  ChabautyColeman.constantCoeff_setYZero, ChabautyColeman.residue_eq_of_constantCoeff_sub_mem',
  ChabautyColeman.map_residue_eq', ChabautyColeman.bezout_ftauFoldPipe_map_setYZero,
  MvPowerSeries.substRingHom, MvPowerSeries.substRingHom_apply, MvPowerSeries.substRingHom_X,
  MvPowerSeries.substRingHom_comp_C, MvPowerSeries.constantCoeff_substRingHom,
  MvPowerSeries.map_substRingHom, MvPowerSeries.hasSubst_map_of_constantCoeff_zero,
  ChabautyColeman.map_substRingHom_diffMatrix, ChabautyColeman.ftauFoldFGL_diffMatrix_eq,
  ChabautyColeman.isCoprime_X_sub_C_of_isUnit_eval, ChabautyColeman.isCoprime_pairV,
  ChabautyColeman.ker_constantCoeff_le_jacobson, ChabautyColeman.isCoprime_chartU,
  ChabautyColeman.isCoprime_chartU_map_rename, ChabautyColeman.isCoprime_cantorPipeline,
  ChabautyColeman.HypField.algebraMap_polynomial_injective, ChabautyColeman.HypField.y_ne_zero,
  ChabautyColeman.isCoprime_two_mul_map, ChabautyColeman.isUnit_two_mul_map,
  ChabautyColeman.isCoprime_pair_const, ChabautyColeman.isCoprime_phiU₁,
  ChabautyColeman.isCoprime_phiU₃, ChabautyColeman.isCoprime_neg_map_C,
  ChabautyColeman.isCoprime_ftauFoldPipe, ChabautyColeman.isCoprime_phiPipeline,
  ChabautyColeman.HypRing.toHypField_apply,
  ChabautyColeman.omegaPrime_taylor, ChabautyColeman.isCoprime_of_map,
  ChabautyColeman.ReflectsCoprime, ChabautyColeman.reflectsCoprime_of_ker_le_jacobson,
  ChabautyColeman.reflectsCoprime_of_field, ChabautyColeman.cantorStage₁_map,
  ChabautyColeman.cantorStage₂_map, ChabautyColeman.cantorStage₃_map,
  ChabautyColeman.mvSeriesResidue, ChabautyColeman.mvSeriesResidue_surjective,
  ChabautyColeman.isUnit_of_mvSeriesResidue_ne_zero,
  ChabautyColeman.ker_mvSeriesResidue_le_jacobson, ChabautyColeman.ftauFoldXs_map,
  ChabautyColeman.ftauFoldYs_map, ChabautyColeman.isFoldAvoid_of_map,
  ChabautyColeman.seriesResidue_apply, ChabautyColeman.seriesResidue_C,
  ChabautyColeman.seriesResidue_X, ChabautyColeman.seriesResidue_C_add_X,
  ChabautyColeman.map_C_map_seriesResidue, ChabautyColeman.phiU₁_map_seriesResidue,
  ChabautyColeman.phiV₁_map_seriesResidue, ChabautyColeman.cantorCompose_eq_of_bezout,
  ChabautyColeman.globExpand_algebraMap_base_comp,
  ChabautyColeman.map_globExpand_map_algebraMap, ChabautyColeman.globExpand_x_eq_diskX,
  ChabautyColeman.globExpand_y_eq_diskY, ChabautyColeman.globExpand_globSlope_eq,
  ChabautyColeman.map_globExpand_globU₁, ChabautyColeman.map_globExpand_globV₁,
  ChabautyColeman.map_globExpand_globU₃, ChabautyColeman.map_globExpand_globV₃,
  ChabautyColeman.IsPipelineAvoid, ChabautyColeman.IsFtauFoldAvoid,
  ChabautyColeman.isFtauFoldAvoid_of_center, ChabautyColeman.isPipelineAvoid_phi_of_residue,
  ChabautyColeman.tracePoly, ChabautyColeman.degree_tracePoly_lt,
  ChabautyColeman.two_smul_tracePoly, ChabautyColeman.sum_smul_tracePoly,
  ChabautyColeman.sum_smul_tracePoly_ne_zero, Derivation.leibniz_mul, IsFractionRing.deriv,
  IsFractionRing.derivation, ChabautyColeman.kzDerivation, ChabautyColeman.kzHom,
  ChabautyColeman.IsPipelineAvoid.map, ChabautyColeman.two_mul_y_mul_traceDiff_phiPipeline_map,
  ChabautyColeman.kzOmega, ChabautyColeman.two_mul_kzOmega_eq_sum,
  ChabautyColeman.two_mul_omegaPrime_kz_eq_sum,
  ChabautyColeman.yseries_mul_omegaPrime_eq_taylor_tracePoly_kz,
  ChabautyColeman.map_omegaExpansionGeneric_eq_sum_smul_omegaPrime,
  ChabautyColeman.aeval_eq_zero_iff_of_aeval_eq_X_sq, ChabautyColeman.aevalFracHom,
  ChabautyColeman.weierstrassExpand, ChabautyColeman.weierstrassExpand_toHypField,
  ChabautyColeman.aeval_map_weierstrassSeries,
  ChabautyColeman.map_weierstrassExpand_globPipeline,
  ChabautyColeman.weierstrassExpand_globZeta, ChabautyColeman.derivative_aeval,
  ChabautyColeman.kzDeriv_weierstrassExpand, ChabautyColeman.weierstrassExpand_globOmega,
  ChabautyColeman.X_mul_omegaPrime_eq_aeval_mul_derivative,
  PowerSeries.map_derivative,
  ChabautyColeman.map_omegaExpansionWeierstrass_eq_sum_smul_omegaPrime}
\uses{lem:delta-rational, def:phiD, thm:chart-mem, def:Lam, lem:omega-expansion}
\leanok
Define $\omega'_i := \phi_Q^*\,\delta_i$, rational differentials on $C_{\Qq}$. Then:
\begin{enumerate}
\item[(a)] $\omega'_i$ is independent of the auxiliary point $Q$;
\item[(b)] for every residue disk $D$ with parameter $z$ and all $P_0, P \in D(\Qq)$,
the value $\Lam_i([P]-[P_0])$ --- defined, since same-disk differences lie in $J^1$ ---
is the integral of the expansion of $\omega'_i$ on $D$ from $z(P_0)$ to $z(P)$;
\item[(c)] each $\omega'_i$ is regular, and $(\omega'_1,\dots,\omega'_g)$ is a basis of
$\RegDiff \otimes \Qq$. Explicitly,
\[
  2\,\omega'_l \;=\; \Bigl(\sum_i (A^{-1})_{li}\,x^i\Bigr)\,\frac{dx}{y},
\]
where $A$ is the invertible \emph{trace matrix} of the $\tau$-datum defined in the
proof --- a Hankel matrix depending on $(u_\tau, c_2)$ alone.
\end{enumerate}
The counting argument consumes clause (b), the regularity and explicit form in (c), and
the linear independence of the numerators; clause (a) and the statement that the
$\omega'_i$ span all of $\RegDiff$ are formalized in
\texttt{HyperellipticCurves.Archive.AuxIndep} and
\texttt{HyperellipticCurves.Archive.RegularDiff}, and the main theorems do not depend on them.
\end{theorem}
\begin{proof}
\leanok
\emph{(a)} Two computations of $\phi_Q$ through different auxiliary points produce the
same Mumford datum over $\Oq[[z]]$: the two auxiliary pair classes cancel, and reduced
representatives of a class over $\operatorname{Frac}(\Oq[[z]])$ are unique. Hence
$\zeta$, and with it $\omega'$, does not depend on $Q$.

\emph{(b)} There are three components. \emph{Formal:} by the chain rule and
$d\Lam = \delta$ (Lemma~\ref{lem:fgl-log-diff}), the composite $\Lam_i(\zeta(z))$ is
the antiderivative of $\omega'_i = \sum_j (M^{-1})_{ij}(\zeta(z))\,\zeta_j'(z)$
vanishing at $0$. \emph{Geometric:} the chart coordinates of the same-disk difference
$[P_{z_0}]-[P_0]$, translated by $\tau$, are exactly $\zeta(z_0)$ --- the evaluated
disk computation is the chart datum at $\zeta(z_0)$, its class is
$([P_{z_0}]+[\iota P_0]) + \tau$, and the difference lies in $J^1$ since both points
reduce to the same residue point. \emph{Analytic:} the value of $\Lam$ is a convergent
sum over $\Qq$ (equipped with its valuation topology, for which it is complete), and
the double series obtained by substituting $\zeta(z)$ into the scaled logarithm is
summable, so it may be summed in either order: summing one way gives the value
$\Lam_i(c)$ at the endpoint, summing the other way gives the coefficients of the
antiderivative of $\omega'_i$. Together these prove (b) for a difference against the
disk centre; the general two-point form follows by additivity of $\Lam$.

\emph{(c)} This is proved by a closed computation of the pullback --- the \emph{trace
form} --- with no Riemann--Roch and no case analysis at infinity. For a Mumford datum
$(U,V)$ over a ring with a derivation $D$ acting on coefficients (write $\dot U$ for
the coefficientwise derivative, so $\deg \dot U < \deg U$ as $U$ is monic), set
\[
  \Theta_k(U,V) \;:=\; -\lambda_U\bigl(X^{k-1}\,\dot U\,s\bigr), \qquad
  s := (2V)^{-1} \bmod U, \qquad
  \lambda_u(h) := \mathrm{coeff}_{\deg u - 1}(h \bmod u).
\]
Classically $\Theta_k = \sum_i x_i^{k-1}D(x_i)/(2y_i)$, summed over the points of the
divisor --- the pullback of $x^{k-1}dx/(2y)$ --- but the coefficient description
requires neither the trace form nor separability, and the only inverse it uses,
$(2V)^{-1} \bmod U$, is the B\'ezout cofactor $c_2$ that every step of the computation
carries. Two elementary properties of $\lambda$ --- $\lambda_u(h) = 0$ for
$\deg h \le \deg u - 2$ (the elementary form of Euler's identity
$\mathrm{Tr}(h/u') = 0$), and $\lambda_{uw}(qw) = \lambda_u(q)$ for monic $u, w$, with
no coprimality required (in particular $w = u$ is allowed) --- yield four facts:
\begin{itemize}
\item[(T1)] $\Theta$ is \emph{additive} under Cantor composition, including the
doubling case, where the Chinese remainder theorem fails but the localization property
of $\lambda$ does not;
\item[(T2)] $\Theta$ is \emph{unchanged} by Cantor reduction: with
$R = V^2 - f = U\widetilde U$ one has $\dot R = 2V\dot V$ ($f$ has constant
coefficients), and the two contributions combine into
$\lambda_R(X^{k-1}\dot V) = 0$ by the vanishing property. The degree count is tight and
uses the odd degree of $f$: $2\deg V = 2g+1$ is impossible by parity, and $k \le g$ is
exactly the right range;
\item[(T3)] on a moving pair $\{P, \iota Q\}$, $\Theta_k = x^{k-1}D(x)/(2y)$ directly;
\item[(T4)] data with constant coefficients have $\dot U = 0$, hence $\Theta_k = 0$.
\end{itemize}
Applying (T1)--(T4) along the disk computation of $\phi_Q$ --- one moving block against
two constant ones --- gives, at the generic point,
\[
  y \cdot \Theta_k(U_\zeta, V_\zeta) \;=\; x^{k-1}/2, \qquad k = 1,\dots,g,
\]
a polynomial of degree $k-1 < g$: regularity, outright. Writing
$C_{kj} := -\lambda_{U}(X^{k+j}s)$, a Hankel matrix with only $2g-1$ distinct entries,
one has $\Theta_{k+1} = \sum_j C_{kj}D(\zeta_j)$. Applying (T1), (T2), (T4) in the
formal setting with the derivations $\partial/\partial Y_k$ turns additivity into the
invariance identity $\sum_j C_{ij}(F(X,Y))\,\partial F_j/\partial Y_k = C_{ik}(Y)$, and
setting $Y = 0$ gives $C(X)\,M(X) = C(0)$, i.e.\ $C = A\,M^{-1}$ with $A := C(0)$
\emph{constant}. The matrix $A$ is invertible: it is the Gram matrix, in the monomial
basis of $\Oq[X]/(u_\tau)$, of the residue pairing $(h,q) \mapsto -\lambda_{u_\tau}(hqs)$,
which is nondegenerate whenever $s$ is invertible modulo $u_\tau$ --- guaranteed by the
B\'ezout identity of the $\tau$-datum. Combining, $2\,\varphi_l = \sum_i
(A^{-1})_{li}\,x^i$ for the numerators $\varphi_l$ of the $\omega'_l$: each is a
polynomial of degree $< g$, and, $A$ being invertible, the $\omega'_l$ form a basis.

The identity of (c) is established on each residue disk directly over
$\operatorname{Frac}(K[[z]])$ with the derivation $d/dz$; on a generic disk it reads
$b\,\eta\,\omega'_i = \operatorname{taylor}_a\bigl(\tfrac12\sum_j (A^{-1})_{ij}
x^j\bigr)$, which is the identification consumed by the per-disk count
(Theorem~\ref{thm:disk-bound}). The analogous identities hold on the Weierstrass and
infinity disks; at a Weierstrass centre the second reduction step passes through a
divisor containing a Weierstrass point, and is crossed by a refinement of (T2) --- if
the constant part $B$ of $U = uB$ has $B \mid \dot V$, which holds automatically when
the second block is constant, then $\Theta$ of the reduced datum equals $\Theta$ of the
moving part alone; the orders of $V^2 - f$ and $U$ at the Weierstrass point cancel, so
the reduced data stay Weierstrass-free. At the infinity disk the configuration
$[P_z - \infty] + \tau$ has only two blocks and needs no auxiliary point at all. The
genericity used in this argument --- $\Theta$ must be defined along the computation,
i.e.\ every intermediate divisor must avoid the Weierstrass points --- is exactly the
condition catalogue of Lemma~\ref{lem:global-generic}.
\end{proof}

\begin{remark}[Why a direct holomorphy proof fails]
One might try to prove clause (c) by verifying the holomorphy of $\omega'_i$ place by
place, using denominator-clearing data for two auxiliary configurations together with a
B\'ezout relation between the two denominators. Such data do not exist: by clause (a),
varying the auxiliary point does not move what is being cleared --- the two configurations
compute the same chart coordinates, so the two denominators are denominators of one and the
same function; a B\'ezout relation between them would force the chart coordinates to be
integral over $\Qq[x]$, hence, being also regular at infinity, \emph{constant}, which
happens only in the degenerate case $\zeta \equiv 0$. (Exact computation over the
function field confirms this: for $y^2 = x^5-x+1$, $\tau = [(0,1)]+[(1,1)]-2\infty$,
centre $(-1,1)$, the chart coordinates have denominator $(x^2-x+3)^2$ --- the same for
every auxiliary point.) The trace-form computation in the proof above replaces the
holomorphy verification: it \emph{exhibits} the regular numerator instead of deducing
its existence.
\end{remark}

\begin{remark}[Formalization]
The trace-form differential $\Theta$ is formalized over an arbitrary commutative ring
with a derivation acting on polynomial coefficients (\texttt{HyperellipticCurves.Curve.TraceDiff});
(T2) is proved for domains, where $\deg(V^2 - f) = \max(2\deg V, \deg f)$ --- over a
general ring a nilpotent leading coefficient of $V$ can collapse the degree, and the
identity fails. $\Theta$ does not depend on the choice of B\'ezout cofactor, so each
step of a computation may carry its own certificate and only existence propagates.
Every reduction in the development is a \emph{single} Cantor step: the fold has no
reduction chains, and in the three-block disk computation the outer blocks are pairs,
so each composition has degree $g+2$ and one reduction returns to degree $g$ for every
genus; the B\'ezout data of a computation therefore consist of two composition triples
and two leading-coefficient inverses. The per-disk identities of clause (c) are proved
over $\operatorname{Frac}(K[[z]])$ with the derivation $d/dz$ on all three chart types
(\texttt{HyperellipticCurves.Log.Trace.TraceKz}, \texttt{TraceKzWeierstrass},
\texttt{TraceKzInf}). What separates a chart's trace identity from $\omega'$ --- transport
of the differential matrix along the substitution, inversion of the constant trace matrix,
and the solve --- does not depend on the chart, and is carried out once
(\texttt{two\_mul\_omegaPrime\_kz\_eq\_sum}), the three capstones supplying only their
chart datum and their numerators. The required Weierstrass-avoidance conditions descend to
residue-level conditions on the reduced configuration, since the residue map reflects
coprimality of monic polynomials (Nakayama). Clause (b) rests on the convergence
infrastructure of \texttt{HyperellipticCurves.Series.FractionTopology} (the valuation
topology on $\Qq$ and its completeness, proved via the complete open subring $\Oq$ and
uniformly separated cosets) and the summable double family of
\texttt{HyperellipticCurves.Log.Analytic.ColemanIntegral}.
\end{remark}

\begin{lemma}[Genericity of the global configuration]
\label{lem:global-generic}
\lean{ChabautyColeman.IsPipelineAvoid, ChabautyColeman.isCoprime_two_mul_iff,
  ChabautyColeman.reduceCanon, ChabautyColeman.ReduceGeneric, ChabautyColeman.reduceGeneric_iff,
  ChabautyColeman.ReduceGeneric.reduceCert, ChabautyColeman.ReduceGeneric.fullDeg,
  ChabautyColeman.ReduceGeneric.reduceCanon_natDegree, ChabautyColeman.ReduceCoprimeCanon,
  ChabautyColeman.reduceCoprimeCanon_iff, ChabautyColeman.genericDeg_pair_compose,
  ChabautyColeman.pipelineCanonData, ChabautyColeman.isPipelineCenterCerts_pipelineCanonData,
  ChabautyColeman.isPipelineAvoid_pipelineCanonData,
  ChabautyColeman.isPipelineCenterCerts_of_pair_blocks,
  ChabautyColeman.cantorCompose_snd_eq_add_mul, ChabautyColeman.cantorPipeline_eq_middle,
  ChabautyColeman.isPipelineAvoid_of_pair_blocks, ChabautyColeman.curveVanishing_mul,
  ChabautyColeman.exists_curveTuple_ne_zero, ChabautyColeman.pencilW_spec,
  ChabautyColeman.pencilTangB_ne_zero, ChabautyColeman.finite_setOf_pencil_multiple_root,
  ChabautyColeman.exists_pencil_param, ChabautyColeman.reduceGeneric_pencil_three,
  ChabautyColeman.exists_isPipelineCerts_disk, ChabautyColeman.exists_lift_curve_point,
  ChabautyColeman.pairV_neg, ChabautyColeman.exists_auxiliary_point,
  ChabautyColeman.exists_isPipelineCerts_phi, ChabautyColeman.exists_auxiliary_point_phi,
  ChabautyColeman.eval_pencilW₀_ne_zero_of_deriv,
  ChabautyColeman.pencilPar_ne_zero_of_nonsingular,
  ChabautyColeman.isCoprime_cantorCofactor_pencil,
  ChabautyColeman.movingPair_dvd_weierstrass,
  ChabautyColeman.exists_auxiliary_point_phi_weierstrass, ChabautyColeman.wPipeline,
  ChabautyColeman.wZeta, ChabautyColeman.wPipeline_map_constantCoeff,
  ChabautyColeman.mk0_wPipeline_eval, ChabautyColeman.redJ_diskPointDiff_eq_one_weierstrass,
  ChabautyColeman.chartClass_eq_diskPointDiff_mul_tauClass_weierstrass,
  ChabautyColeman.J1_equiv_polydisk_diskPointDiff_weierstrass,
  ChabautyColeman.hasSum_Lam_antiderivOmegaPrime_weierstrass, ChabautyColeman.diskDiffClassW,
  ChabautyColeman.lambda_glue_weierstrass,
  ChabautyColeman.weierstrassRoot_eq_eval_weierstrassSeries,
  ChabautyColeman.encard_le_of_lambda_eq_zero_weierstrass,
  ChabautyColeman.encard_inter_residueDisk_le_weierstrass,
  ChabautyColeman.encard_inter_residueDisk_le_ordPoint_weierstrass, ChabautyColeman.infC,
  ChabautyColeman.infV, ChabautyColeman.infU, ChabautyColeman.infZeta,
  ChabautyColeman.aeval_reflect_infV, ChabautyColeman.aeval_reflect_infV_sq_sub,
  ChabautyColeman.infV_sq_sub_factor, ChabautyColeman.infU_monic,
  ChabautyColeman.map_constantCoeff_infU, ChabautyColeman.constantCoeff_infZeta,
  ChabautyColeman.exists_infVred, ChabautyColeman.infRoot_eq_eval_infSeries,
  ChabautyColeman.mk0_mumfordIdeal_of_scaled_factor, ChabautyColeman.mk0_infU_eval,
  ChabautyColeman.redJ_diskPointDiff_eq_one_inf, ChabautyColeman.map_constantCoeff_infVred,
  ChabautyColeman.chartClass_eq_diskPointDiff_mul_tauClass_inf,
  ChabautyColeman.J1_equiv_polydisk_diskPointDiff_inf, ChabautyColeman.diskDiffClassInf,
  ChabautyColeman.hasSum_Lam_antiderivOmegaPrime_inf,
  ChabautyColeman.hasSum_lambdaOne_diskDiffClass_inf,
  ChabautyColeman.mapBase_aj_eq_diskDiffClassInf, ChabautyColeman.lambda_glue_inf,
  ChabautyColeman.encard_le_of_lambda_eq_zero_inf,
  ChabautyColeman.encard_inter_residueDisk_le_inf,
  ChabautyColeman.encard_inter_residueDisk_le_ordPoint_inf}
\uses{thm:omega-prime-global, lem:tau-exists, lem:tau-generic, lem:many-points}
\leanok
The genericity hypotheses of Theorem~\ref{thm:omega-prime-global}(c) are satisfiable at
every residue disk: for every nonsingular disk centre there are an auxiliary curve point
$Q = (x_Q, y_Q)$ with its slope datum, and the class $\tau$ of
Lemma~\ref{lem:tau-exists}, such that
\begin{enumerate}
\item $y_Q$ is a unit --- the auxiliary point does not reduce to a Weierstrass point;
\item every intermediate state of the fold computation of $F_\tau$ avoids the
Weierstrass points (this is carried by the datum of Lemma~\ref{lem:tau-exists}, so no
separate condition remains);
\item the reduced per-disk configuration
$[\bar P_0 + \iota\bar Q] + \bar\tau + [\bar Q + \iota \bar P_0]$ avoids the
Weierstrass points at both of its reduced intermediate data.
\end{enumerate}
Like Lemma~\ref{lem:tau-generic}, this is a purely residue-level genericity statement
followed by a Hensel lift.
\end{lemma}
\begin{proof}
\leanok
Condition (2) is Lemma~\ref{lem:tau-generic}, lifted with the datum of
Lemma~\ref{lem:tau-exists}; a single witness configuration discharges both. For a
Mumford datum, $u$ is coprime to $2v$ exactly when it is coprime to $\bar f$, so every
avoidance condition is a condition against the \emph{fixed} curve --- the intermediate
divisor misses the Weierstrass points --- not against the evolving square root.

For conditions (1) and (3), the key observation is that the candidate composed square
roots of the per-disk configuration form an affine \emph{pencil}: with
$N = u_\tau\,(X - a)$ and $V_0$ the interpolant with $V_0 \equiv v_\tau \bmod u_\tau$
and $V_0(a) = b$, they are $V_t = V_0 + t\,N$, and $V_t^2 - f = N\,W_t$ with $W_t$
quadratic in $t$. The auxiliary point is read off the parameter: for any root $x_Q$ of
$W_t$, the pair $(x_Q, -V_t(x_Q))$ is a curve point. Each genericity condition then
becomes ``avoid finitely many $t$'':
\begin{itemize}
\item the degree condition of the first reduction step is $t \neq 0$;
\item transversality at the centre, the condition $u_\tau(x_Q) \neq 0$, and the
coprimality of the reduced intermediate datum with $\bar f$ are each \emph{affine} in
$t$ with nonvanishing coefficient --- respectively $2b$ (at a Weierstrass centre
$b = 0$, and instead the constant term $W_0(a) = -f'(a)/u_\tau(a)$ is nonzero, by
differentiating the pencil identity), $2v_\tau(r)$ at the roots $r$ of $u_\tau$, and
$N(r)$ at the roots $r$ of $\bar f$, all nonzero because $\bar\tau$ is in good
position;
\item the simple-root condition on $W_t$ is finite by elimination: a multiple root $x$
of $W_t$ off the support of $N$ puts the curve point $(x, V_t(x))$ on the line
$A\,y + B = 0$ with $A = 2(V_0'N - V_0N')$ and $B = 2fN' - f'N$, so $x$ is a root of
$A^2f - B^2$; that polynomial is nonzero by degree parity (odd against even), and
$B \neq 0$ because its top coefficient is the integer $2(g+1) - (2g+1) = 1$, nonzero in
\emph{every} characteristic --- no separability of $u_\tau$ and no discriminant is
needed.
\end{itemize}
The third composition produces $-V_t$ on the nose, so its degree condition coincides
with the first; and the full configuration returns its middle datum --- the two pair
blocks are inverse classes --- so the final avoidance condition reduces to the good
position of $\bar\tau$. A good parameter therefore exists over the algebraically closed
residue field, uniformly at every nonsingular centre ($b \neq 0$, or $a$ a simple root
of $\bar f$); at a Weierstrass centre the reduced intermediate data remain
Weierstrass-free because the orders of $V^2 - f$ and of the $u$-part at the centre
cancel. Finally, the residue-level witness lifts by Hensel's lemma: the centre and the
auxiliary point lift as smooth points, the B\'ezout data lift by Nakayama, and all
conditions are open.
\end{proof}

\begin{remark}[Formalization]
The per-disk statement is
\texttt{exists\_isPipelineCerts\_disk} (residue level, all clauses at a nonsingular
centre) and \texttt{exists\_auxiliary\_point} /
\texttt{exists\_auxiliary\_point\_phi} (Hensel-lifted, with the certified computation
over $\Oq[[z]]$), with the Weierstrass-centre forms
\texttt{exists\_auxiliary\_point\_weierstrass} /
\texttt{exists\_auxiliary\_point\_phi\_weierstrass}. The infinity disk needs no
genericity at all: its centre is $\infty$ itself, so the class to expand is
$[P_z - \infty] + \tau$ --- no auxiliary point and no pair blocks; the computation is
one composition and one reduction in closed form, carried in scaled integral
coordinates because the moving point has non-integral coordinates, and the disk
identities (the chart class, the kernel membership, the value of the pairing, and the
per-disk bound of Theorem~\ref{thm:disk-bound}) are established on it directly. On the
Weierstrass and infinity disks the identification of the two disk parametrizations in
use is part of the statement chain. All conditions are produced for every genus
$\ge 1$.
\end{remark}

\begin{definition}[The pairing $\lambda$ on $J^1(\Qp)$, normalized]
\label{def:lambda1}
\leanok
\lean{ChabautyColeman.lambdaOne, ChabautyColeman.lambdaOne_mul,
  ChabautyColeman.lambdaOne_add_left, ChabautyColeman.lambdaOne_smul_left}
\uses{thm:omega-prime-global, lem:global-generic, def:Lam, def:regdiff}
Write $\omega = \sum_i c_i(\omega)\,\omega'_i$ in the basis of
Theorem~\ref{thm:omega-prime-global}, and set
$\lambda_\omega := \sum_i c_i(\omega)\,\Lam_i$ on $J^1(\Qq)$, restricted to $J^1(\Qp)$.
The pairing is linear in $\omega$ and additive in the class. \emph{Formalization:}
$\lambda$ is implemented in $\omega'$-coordinates --- for a coefficient vector $c$ the
functional $\sum_i c_i\,\Lam_i$ on the kernel of reduction; the identification
$c = c(\omega)$ is the explicit trace-form coordinate identity of
Theorem~\ref{thm:omega-prime-global}(c).
\end{definition}

\begin{lemma}[Rationality of values]
\label{lem:values-rational}
\leanok
\lean{ChabautyColeman.lambdaOne_rational_values,
  ChabautyColeman.FractionTopology.continuous_ringEquiv,
  ChabautyColeman.FractionTopology.valuation_ringEquiv}
\uses{def:lambda1, thm:omega-prime-global}
For $\omega \in \RegDiff$ (over $\Qp$) and $c \in J^1(\Qp)$:
$\lambda_\omega(c) \in \Qp$.
\end{lemma}
\begin{proof}
\leanok
$\lambda_\omega$ is characterized Galois-equivariantly by
Theorem~\ref{thm:omega-prime-global}(b) at $\Qp$-rational disks together with
additivity ($\tau$ drops out), so Frobenius-fixed inputs give fixed outputs. The
analytic content is that a continuous field automorphism transports the convergent
characterization of the disk values: any field automorphism of $\Qq$ over $\Qp$ is an
isometry of the valued field, hence continuous, and uniqueness of limits transports the
values along the induced permutation of the disks. \emph{Formalization.} The
annihilator argument of Chapter~5 runs with values in $\Qq$ directly, so the main theorems
do not depend on this lemma; it is formalized in
\texttt{HyperellipticCurves.Archive.RationalValues}.
\end{proof}

\begin{definition}[The pairing $\lambda$ on $J^1(K_0)$]
\label{def:lambda}
\lean{ChabautyColeman.lambda, ChabautyColeman.J1,
  ChabautyColeman.lambda_mul, ChabautyColeman.lambda_add_left,
  ChabautyColeman.lambda_smul_left, ChabautyColeman.lambdaOne_inv,
  ChabautyColeman.mem_J1_iff, ChabautyColeman.mem_J1_of_mapBase_eq,
  ChabautyColeman.mem_J1_of_mapBase_eq_div}
\leanok
\uses{def:lambda1, thm:redJ-hom, lem:values-rational}
Let $J^1(K_0) := \ker\big(\J{K_0} \to \J{K} \to \J{\Fpbar}\big)$ be the group of
classes over the arithmetic base field whose base change reduces trivially, and let
$\lambda_\omega$ on $J^1(K_0)$ be the value of Definition~\ref{def:lambda1} at the base
change of the class. It is additive in the class and linear in $\omega$.
\end{definition}

\begin{remark}[Formalization]
The pairing $\lambda$ is defined on $J^1$ only, and this suffices: every class the
counting argument pairs is a same-disk difference of rational points, which lies in
$J^1$, and the rank of $J^1$ is bounded by the rank of $\J{K_0}$ by flatness of $\QQ$
over $\ZZ$ (Theorem~\ref{thm:annihilator}). So there is no need to extend $\lambda$ to
all of $\J{\Qp}$ by $\lambda(c) := \tfrac1N \lambda(Nc)$, with $N$ the index of $J^1$
(Theorem~\ref{thm:J1-index}), and the finite-index and Frobenius apparatus of
Theorem~\ref{thm:J1-index} does not enter the proof.
\end{remark}

\begin{theorem}[Interface theorem: linearity and the glue identity]
\label{thm:glue}
\lean{ChabautyColeman.lambda_mul, ChabautyColeman.lambda_add_left,
  ChabautyColeman.lambda_smul_left, ChabautyColeman.diskDiffClass,
  ChabautyColeman.hasSum_lambdaOne_diskDiffClass, ChabautyColeman.lambda_glue,
  ChabautyColeman.mapBase_aj_div_eq_diskDiffClass, ChabautyColeman.mapBase_aj,
  ChabautyColeman.HypJac.congr_aj}
\leanok
\uses{def:lambda, thm:omega-prime-global, def:disk-antider}
$\lambda \colon \RegDiff \to \operatorname{Hom}(J^1(\Qp), \Qp)$ is $\Qp$-linear, and
for every residue disk $D$ with parameter $z$ and all $P, P_0 \in C(\Qp) \cap D$ (the
same-disk difference lies in $J^1(\Qp)$, since $\redJ$ kills it):
\[ \lambda_\omega\big(\aj(P) - \aj(P_0)\big) \;=\; G^{\omega}_D(z(P)) - G^{\omega}_D(z(P_0)). \]
\end{theorem}
\begin{proof}
\leanok
Linearity comes with Definitions~\ref{def:lambda1} and~\ref{def:lambda}. For the glue
identity: same-disk differences lie in $J^1(\Qp)$ (Theorem~\ref{thm:redJ-hom}); apply
Theorem~\ref{thm:omega-prime-global}(b) with the basis expansion of $\omega$; the two
antiderivatives agree up to a constant that cancels in the difference. That the base
change of the difference of the two Abel--Jacobi classes is the same-disk difference
used in Theorem~\ref{thm:omega-prime-global}(b) rests on the naturality of the
Abel--Jacobi map in the base field, composed with base change along the extension.
\end{proof}

\begin{remark}[Genus-two shortcut]
\label{rem:g2-cert}
\uses{def:lambda}
For $g = 2$ the constructions of Definitions~\ref{def:Ftau}--\ref{def:lambda} could be
replaced by the explicit Grant--Flynn formal group and local parameters, imported as
polynomial-identity certificates; this would instantiate Theorem~\ref{thm:glue}
independently of Theorems~\ref{thm:Ftau-assoc} and~\ref{thm:omega-prime-global}. The
route is not needed. The script \texttt{prototype/g2\_formal\_group.py} carries out its
first steps in exact arithmetic --- Cantor composition, the chart series, and $F_\tau$ for
$y^2 = x^5 - x + 1$ at $p = 11$ --- and verifies integrality, the unit laws,
commutativity, associativity and $M(0) = I$ up to total degree~$5$.
\end{remark}

\chapter{The annihilator and the count}

\begin{definition}[Rank hypothesis]
\label{def:rank-hyp}
\lean{ChabautyColeman.HypJac.rank}
\leanok
\uses{def:hypjac}
$\operatorname{rk} \J{\QQ} := \dim_{\QQ} \big(\J{\QQ} \otimes_{\ZZ} \QQ\big)$; the
standing hypothesis of the main theorem is $\operatorname{rk} \J{\QQ} \le r$ for a natural
number $r < g$. No finite generation is assumed: the dimension is a cardinal. The
qualitative form $\operatorname{rk} \J{\QQ} < g$ is equivalent to the existence of such an
$r$, a cardinal below a natural number being itself a natural number.
\emph{Formalization:} the Jacobian is the class group, written multiplicatively, so the
rationalization is $\QQ \otimes_{\ZZ}$ of the additivization; the arithmetic base field
is an arbitrary subfield $K_0$ of the local field, not just $\QQ$.
\end{definition}

\begin{lemma}[The evaluation map factors through $\otimes\,\QQ$]
\label{lem:eval-factors}
\lean{ChabautyColeman.lambdaEval, ChabautyColeman.lambdaEvalRat}
\leanok
\uses{thm:glue, def:basechange, def:rank-hyp}
The bilinear map $(\omega, c) \mapsto \lambda_\omega(c_{\Qp})$, for $c \in \J{\QQ}$
mapped to $\J{\Qp}$ by the homomorphism of Definition~\ref{def:basechange}, induces a
$\QQ$-linear map
$\J{\QQ} \otimes \QQ \to \operatorname{Hom}_{\Qp}(\RegDiff, \Qp)$.
\end{lemma}
\begin{proof}
\leanok
Homomorphisms into the torsion-free uniquely divisible group $\Qp$ factor through
$\otimes\,\QQ$. \emph{Formalization.} The base change is built into $\lambda$
(Definition~\ref{def:lambda}), and the differential is represented by its coordinate
vector in the $\omega'$-basis; the statement is proved for any pairing of an abstract
commutative group that is additive in the class and linear in the coordinates, with
these three properties as hypotheses, and instantiated at $\lambda$.
\end{proof}

\begin{theorem}[Annihilating differential]
\label{thm:annihilator}
\lean{ChabautyColeman.exists_annihilator_submodule,
  ChabautyColeman.HypJac.rank_subgroup_le,
  ChabautyColeman.sum_smul_tracePoly_ne_zero,
  Module.finrank_le_finrank_iInf_ker_add_finrank}
\leanok
\uses{lem:eval-factors, def:rank-hyp, def:regdiff}
If $\operatorname{rk}\J{\QQ} \le r$, the differentials $\omega \in \RegDiff$ with
$\lambda_\omega(c) = 0$ for every $c \in J^1(\QQ)$ form a subspace of codimension at most
$r$; for any such $\omega$, in particular,
$\lambda_\omega\big(\aj(P) - \aj(P_0)\big) = 0$ for all $P, P_0 \in C(\QQ)$ in one
residue disk. For $r < g$ the subspace is nonzero.
\end{theorem}
\begin{proof}
\leanok
$\operatorname{rk} J^1(\QQ) \le \operatorname{rk}\J{\QQ}$, since $\otimes\,\QQ$
preserves the inclusion $J^1 \subseteq \J{\QQ}$ by flatness of $\QQ$ over $\ZZ$. So the
image of Lemma~\ref{lem:eval-factors} spans a $\Qp$-subspace of dimension
$\le \operatorname{rk}\J{\QQ} \le r$, and the annihilator of that span in $\RegDiff$ has
codimension at most $r$. The dimension count is isolated as an abstract statement: for an
$R$-linear map $Q \to \operatorname{Hom}_L(V, L)$, with $R$ a subfield of $L$, the
subspace of $V$ annihilated by the whole image has codimension at most $\dim_R Q$,
because the $L$-span of the image is generated by the images of an $R$-basis of the
source, so the annihilator is the common kernel of $\dim_R Q$ functionals. The
differentials are produced as their coordinate vectors in the $\omega'$-basis; that a
nonzero vector represents a nonzero element of $\RegDiff$ is the linear independence of
the trace-form numerators (Theorem~\ref{thm:omega-prime-global}(c)).
\end{proof}

\begin{lemma}[Integral normalization]
\label{lem:normalize}
\lean{ChabautyColeman.exists_annihilator_residue,
  Polynomial.exists_submodule_map_residue,
  IsLocalRing.linearIndependent_map_of_saturated,
  ChabautyColeman.exists_annihilator_tracePoly}
\leanok
\uses{thm:annihilator, def:regdiff, def:hypcurve}
The reductions $\bar h$ of the integral numerators of the annihilating differentials form
a subspace $V$ of the polynomials of degree $< g$ over $\Fpbar$, again of codimension at
most $r$, and every element of $V$ is the reduction of such a numerator. In particular a
nonzero element of $V$ comes from an $\omega \in \RegDiff_{\ZZ}$ with $\bar\omega \ne 0$
on $\bar C_{\Fpbar}$.
\end{lemma}
\begin{proof}
\leanok
The annihilating differentials with integral numerator form a lattice in the annihilating
subspace of Theorem~\ref{thm:annihilator}: a $\Zp$-submodule of the numerators of degree
$< g$, free of rank equal to the dimension of that subspace, and saturated, since
$\pi^{-1}h$ is again integral and annihilating whenever $\pi \mid h$. Reduction is
therefore injective on a basis of the lattice, and the span of the reduced basis is the
required $V$; an element of $V$ lifts by lifting the coefficients of that basis, scaling
does not affect the vanishing of $\lambda_\omega$, which is linear.
\emph{Formalization.} The differentials are handled through their numerators $h$ (the
polynomials with $\omega = h\,dx/y$, $\deg h < g$), which is what the order computations
of Lemma~\ref{lem:ord-match} and Corollary~\ref{cor:sum-ord-space} consume; the lattice
statement \texttt{exists\_submodule\_map\_residue} is proved for an arbitrary subspace of
the polynomials of bounded degree over the fraction field of a discrete valuation ring.
The composed form used by the main theorem is \texttt{exists\_annihilator\_tracePoly}:
Theorem~\ref{thm:annihilator} instantiated in the trace-form basis, giving for each
$\bar h \in V$ coordinates $c$ \emph{and} an integral numerator $h_0$ with reduction
$\bar h$ and degree $< g$ whose image is
$\sum_i c_i\,\bigl(\text{trace-form numerator}\bigr)_i$ --- a statement about $\tau$
alone, independent of any residue disk, so every residue disk may pick its own
differential out of $V$.
\end{proof}

\begin{lemma}[Disk decomposition of the rational points]
\label{lem:disk-decomp}
\lean{ChabautyColeman.ratPoints_disk_decomp,
  ChabautyColeman.encard_le_of_fiber_encard_le,
  ChabautyColeman.exists_finset_redPoint_rational}
\uses{lem:two-charts, def:disk}
\leanok
$C(\QQ) \subseteq C(\Qp) = \bigsqcup_{\bar P \in \bar C(\Fp)} \big(C(\Qp) \cap
D_{\bar P}\big)$, and rational points reduce to $\Fp$-rational residue points.
\end{lemma}
\begin{proof}
\leanok
Lemma~\ref{lem:two-charts} over $O = \Zp$; coordinates in $\Zp$ (or $t \in p\Zp$)
reduce into $\Fp$. \emph{Formalization.} The decomposition is definitional --- the
residue disk \emph{is} the fibre of the reduction map --- so what the count uses is the
fibrewise bound: a set of points reducing into a finite set $T$ of residue points, with
at most $1 + m_{\bar P}$ points on the disk of each $\bar P \in T$, has at most
$\#T + \sum_{\bar P \in T} m_{\bar P}$ elements. Cardinalities are extended naturals,
so finiteness is not assumed but obtained. The arithmetic half --- that reductions of
$K_0$-rational points land in a set of size at most the point count of the reduced
curve over the residue field of an arithmetic valuation ring $O_0 \subseteq K_0$
compatible with $O$ --- is \texttt{exists\_finset\_redPoint\_rational}.
\end{proof}

\begin{theorem}[Per-disk bound]
\label{thm:disk-bound}
\lean{ChabautyColeman.encard_le_of_lambda_eq_zero, ChabautyColeman.disk_count_le,
  ChabautyColeman.disk_bound_generic, ChabautyColeman.disk_bound_weierstrass,
  ChabautyColeman.disk_bound_inf,
  ChabautyColeman.encard_inter_residueDisk_le,
  ChabautyColeman.encard_inter_residueDisk_le_ordPoint,
  ChabautyColeman.encard_inter_residueDisk_le_ordPoint_tracePoly,
  ChabautyColeman.encard_inter_residueDisk_le_ordPoint_tracePoly_weierstrass,
  ChabautyColeman.encard_inter_residueDisk_le_ordPoint_tracePoly_inf,
  ChabautyColeman.map_omegaExpansionWeierstrass_eq_sum_smul_omegaPrime_tracePoly,
  ChabautyColeman.map_omegaExpansionInf_eq_sum_smul_omegaPrime_tracePoly,
  ChabautyColeman.PSeries.encard_le_of_evalT_const,
  ChabautyColeman.PSeries.hasSum_algebraMap_evalT,
  ChabautyColeman.PSeries.evalT_sub_C,
  ChabautyColeman.PSeries.exists_substituted_antideriv,
  IsLocalRing.natCast_dvd_pow_self,
  ChabautyColeman.PSeries.algebraMap_coeff_eq_coeff_antideriv}
\uses{thm:glue, thm:annihilator, lem:normalize, lem:ord-match, thm:antideriv-zeros,
  def:disk-antider, cor:sum-ord-space, def:hypcurve}
\leanok
Let $\bar P \in \bar C(\Fp)$ and $m := m_{\bar P} = \ord_{\bar P}(\bar\omega)$. Then
$\#\big(C(\QQ) \cap D_{\bar P}\big) \le 1 + m$.
\end{theorem}
\begin{proof}
\leanok
If the intersection is empty, we are done; else fix $P_0 \in C(\QQ) \cap D_{\bar P}$.
For any other rational $P$ in the disk, Theorems~\ref{thm:annihilator}
and~\ref{thm:glue} give
$0 = \lambda_\omega(\aj P) = \lambda_\omega(\aj P_0) + G^{\omega}_D(z(P)) -
G^{\omega}_D(z(P_0))$, so $z(P)$ is a zero in $p\Zp$ of the single function
$G^{\omega}_D(z) + c_0$. The derivative series of $G$ has reduction of order exactly
$m$ (Lemmas~\ref{lem:normalize} and~\ref{lem:ord-match}); choosing for $\omega$ the
differential of smallest order at $\bar P$ gives $m \le 2r$
(Corollary~\ref{cor:sum-ord-space}), and $p > 2r+2$ gives $p > m + 2$; apply
Theorem~\ref{thm:antideriv-zeros}.

\emph{Formalization.} The analytic core is: for points of one disk with distinct
parameters in $\mathfrak m$ at which the integral takes one and the same value,
$\#S \le 1 + m$; a constant value is as good as a zero (subtracting it changes only the
constant coefficient), and the value delivered by Theorem~\ref{thm:glue} --- a
convergent sum over $\mathrm{Frac}(O)$ --- is identified coefficientwise with the value
of the substituted series $H(u) = G(\pi u)$ over $O$ on which
Theorem~\ref{thm:antideriv-zeros} runs; integrality of $H$ is the divisibility of
$\pi^{k+1}$ by $k+1$ in $O$, which is where the bound on the ramification is spent. Rationality of the points enters only through the coordinate
identities matching each $P \in S$ with the parametrized disk point at its parameter:
these give distinctness of the parameters and, through Theorem~\ref{thm:glue} and
Theorem~\ref{thm:annihilator}, the constancy of the integral. The identification of the
disk expansion of $\omega = \sum_i c_i\,\omega'_i$ with the integrand --- the remaining
hypothesis --- is supplied on each of the three chart types by the per-disk identities
of Theorem~\ref{thm:omega-prime-global}(c), fed with the annihilating differential
normalized in the trace-form basis (Lemma~\ref{lem:normalize}); the statements
\texttt{disk\_bound\_generic}, \texttt{disk\_bound\_weierstrass} and
\texttt{disk\_bound\_inf} establish the bound at each chart type from the data of
Lemma~\ref{lem:tau-exists} alone, with the centre lift, the auxiliary point and all
B\'ezout data constructed inside. All of these statements, like Lemma~\ref{lem:normalize}
and the counting step of Theorem~\ref{thm:main-count}, are formulated for an odd-degree
hyperelliptic curve \texttt{C : HyperellipticCurve O} (Definition~\ref{def:hypcurve}) and a
curve \texttt{C₀} over $K_0$: the curve supplies monicity, the degree, oddness and
squarefreeness over $\operatorname{Frac} O$, over $K_0$ and over the residue field through
its instances, and the invertibility of $2$ in $O$ is the instance \texttt{[Invertible (2 : O)]}.
\end{proof}

\begin{theorem}[Chabauty--Coleman, local count]
\label{thm:main-count}
\lean{ChabautyColeman.chabauty_coleman_odd_hyperelliptic,
  ChabautyColeman.ratPoints_encard_le,
  ChabautyColeman.main_affine, ChabautyColeman.ratPoints_finite_of_rank_le,
  ChabautyColeman.CurvePoint.encard_univ,
  Polynomial.separable_map_of_separable_residue}
\uses{thm:disk-bound, lem:disk-decomp, cor:sum-ord-space, def:rank-hyp, def:hypcurve}
\leanok
Under the standing hypotheses, let $K_0$ be a subfield of $\operatorname{Frac} O$ of
characteristic $0$, and let $S$ be a set of points of $C$ with $K_0$-rational coordinates
whose reductions lie in a finite set $T$ of points of $\bar C$. If
$\operatorname{rk}\J{K_0} \le r$ with $r < g$ and $p > 2r + e + 1$, then
\[ \#S\ \le\ \#T\ +\ 2r. \]
The statement over $\QQ$ (Theorem~\ref{thm:main}) follows via the arithmetic form
(Theorem~\ref{thm:main-arith}) and its number-field instantiation
(Theorem~\ref{thm:main-nf}).
\end{theorem}
\begin{proof}
\leanok
Sum Theorem~\ref{thm:disk-bound} over the disks of Lemma~\ref{lem:disk-decomp}, each disk
using the annihilating differential of smallest order at its centre:
$\#S \le \sum_{\bar P \in T} \bigl(1 + n(V, \bar P)\bigr)
\le \#T + 2r$ by Corollary~\ref{cor:sum-ord-space}.

\emph{Formalization.} The full assembly is
\texttt{chabauty\_coleman\_odd\_hyperelliptic}, stated for an odd-degree hyperelliptic curve
\texttt{C : HyperellipticCurve O} (its unit discriminant is the good reduction) and a curve
\texttt{C₀} over $K_0$ agreeing with it over $\operatorname{Frac} O$; the polynomial-level
hypotheses of the inner layers are supplied by the curves. The proof constructs the class $\tau$ avoiding the
$x$-coordinates of $T$ (Lemma~\ref{lem:tau-exists} via Lemma~\ref{lem:tau-generic},
unconditionally for every genus $\ge 1$), produces the space $V$ of reduced numerators of
the annihilating differentials in the trace-form basis (Theorem~\ref{thm:annihilator},
Lemma~\ref{lem:normalize}), picks in each residue disk an element of $V$ of minimal
order there, classifies the residue point as generic affine, Weierstrass or $\infty$,
invokes the matching per-disk bound of Theorem~\ref{thm:disk-bound}, and sums. No genus
lower bound is assumed: the rank hypothesis forces $g \ge 1$, and genus
$1$ passes through Lemma~\ref{lem:tau-generic} since $\bar f$ is squarefree. The
hypotheses of the formal statement are exactly: $O$ a complete discrete valuation ring
of characteristic $0$ with algebraically closed residue field of characteristic
$\neq 2$ --- carrying a uniformizer $\pi$ and a bound $e$ on the absolute ramification
index, with no unramifiedness assumed --- $f$ monic with squarefree reduction of odd
degree, and the rank hypothesis
--- all other instances (oddness and squarefreeness over the fraction field and the
base field, adic completeness) are derived inside; for the squarefreeness this is the
statement that separability of the reduction of a monic polynomial over a local ring
lifts through any ring homomorphism, by the resultant
(\texttt{separable\_map\_of\_separable\_residue}). The affine restatement --- both sides
counting the projective sets with the point at infinity made explicit, of which there is
exactly one on each side by the odd degree --- is \texttt{main\_affine} (via
\texttt{CurvePoint.encard\_univ}), and the qualitative finiteness of any set of
base-rational points is \texttt{ratPoints\_finite\_of\_rank\_le}: the counts are extended
naturals, so the bound caps the cardinality by a natural number.
\end{proof}

\begin{theorem}[Chabauty--Coleman over an arithmetic base]
\label{thm:main-arith}
\lean{ChabautyColeman.chabauty_coleman_odd_hyperelliptic_arith,
  ChabautyColeman.main_affine_arith, ChabautyColeman.ratPoints_finite_arith,
  IsLocalRing.IsUniformizer}
\uses{thm:main-count, lem:disk-decomp, thm:antideriv-zeros}
\leanok
Let $O$ be a complete discrete valuation ring of characteristic $0$ with algebraically
closed residue field of characteristic $\ne 2$, uniformizer $\pi$, residue characteristic
$p$ and absolute ramification index at most $e$ (that is, $v(p) \le e$). Let $K$ be a
subfield of $\operatorname{Frac} O$ of characteristic $0$ and $O_0$ a local ring mapping
compatibly into $O$ and into $K$, with finite residue field $\Fq$, such that every element
of $K$ of valuation $\le 1$ comes from $O_0$. Let $C \colon y^2 = f(x)$ be a hyperelliptic
curve of genus $g$ over $O$ with $f$ monic of odd degree $2g+1$ and coefficients in $O_0$
(so its reduction over the residue field of $O$ is squarefree). If
$\operatorname{rk} J(K) \le r$ with $r < g$ and
\[ p \;>\; 2r + e + 1, \]
then
\[ \#C(K)\ \le\ \#\bar C(\Fq)\ +\ 2r. \]

For a number field $K$ and a prime $\mathfrak p$ of good reduction the
hypotheses hold with $O$ the ring of integers of the completion of the maximal unramified
extension of $K_{\mathfrak p}$, $O_0$ the localization of $O_K$ at $\mathfrak p$,
$q = \#(O_K/\mathfrak p)$ and $e = e(\mathfrak p \mid p)$, so ramified primes are admitted:
this is Corollary~6.7 of Stoll, \emph{loc.\ cit.}; that instantiation is
Theorem~\ref{thm:main-nf}, and its case $K = \QQ$ is Theorem~\ref{thm:main}.
\end{theorem}
\begin{proof}
\leanok
Specialize Theorem~\ref{thm:main-count}: the base field enters only
through the rationality of the coordinates of the points being counted, and the
ramification only through the substitution $z = \pi u$ of
Theorem~\ref{thm:antideriv-zeros}, which asks $p > m + e + 1$ for the order $m$ counted on
the disk; since $m \le 2r$ at every disk (Corollary~\ref{cor:sum-ord-space}), the condition
$p > 2r + e + 1$ suffices. The passage from the residue points over the algebraically closed
residue field of $O$ to the $\Fq$-points of $\bar C$ is the arithmetic half of
Lemma~\ref{lem:disk-decomp}.

\emph{Formalization.} The arithmetic base is carried as a pack: a local ring $O_0$ with
finite residue field together with ring maps $O_0 \to O$ (local) and $O_0 \to K$ that agree
in $\operatorname{Frac} O$, the hypothesis that $O_0$ covers $\{x \in K : v(x) \le 1\}$, and
a polynomial over $O_0$ mapping to $f$; the ramification data is
\texttt{IsLocalRing.IsUniformizer $\pi$ $p$ $e$} ($\mathfrak m = (\pi)$, $p \in \mathfrak m$,
$p \mid \pi^e$). The affine forms \texttt{main\_affine\_arith} and
\texttt{ratPoints\_finite\_arith} state
$\#\{(x,y) \in K^2 : y^2 = f(x)\} + 1 \le \#\{(\bar x,\bar y) \in \mathbb F_q^2 :
\bar y^2 = \bar f(\bar x)\} + 1 + 2r$ and the finiteness of the affine $K$-solution set.
The pack is instantiated in Theorem~\ref{thm:main-nf}, at the local factor of
Chapter~\ref{ch:numberfield} with $O_0$ the localization at $\mathfrak p$; the statements
over $\QQ$ follow from that instantiation (Theorem~\ref{thm:main},
Corollary~\ref{cor:finiteness}).
\end{proof}

\begin{theorem}[Chabauty--Coleman over $\QQ$]
\label{thm:main}
\lean{ChabautyColeman.main_affine_rat_of_rank_le, ChabautyColeman.main_affine_int_of_rank_le,
  ChabautyColeman.main_affine_int, ChabautyColeman.separable_of_squarefree_twist}
\uses{thm:main-nf}
\leanok
Let $C \colon y^2 = f(x)$ be a hyperelliptic curve of genus $g$ over $\QQ$ with $f$ monic of
odd degree $2g+1$, and $p$ a prime such that $f$ has $p$-integral coefficients and
$\bar f \in \Fp[x]$ is squarefree. If $\operatorname{rk} \J{\QQ} \le r$ with $r < g$ and
$p > 2r + 2$, then
\[ \#C(\QQ)\ \le\ \#\bar C(\Fp)\ +\ 2r. \]
\end{theorem}
\begin{proof}
\leanok
This is the case $A = \ZZ$, $K = \QQ$, $\mathfrak p = (p)$ of
Theorem~\ref{thm:main-nf}: the residue field of $(p)$ is $\Fp$, and the ramification
index is $1$ (it is at most $[\QQ:\QQ]$), so the prime condition $p > 2r + e + 1$ reads
$p > 2r + 2$.

\emph{Formalization.} The Lean statement is the affine form --- both sides count the
projective sets with the point at infinity made explicit, of which there is exactly one
on each side by the odd degree --- with the reduction count over $\mathrm{ZMod}\ p$;
$p$-integrality is certified by $d f = F \in \ZZ[x]$ with $p \nmid d$, and the reduced
count is taken on the model $\bar d\, y^2 = \bar F(x)$, to which the specialization
transports squarefreeness and the count along $\ZZ_{(p)}/(p) \cong \mathrm{ZMod}\ p$
(\texttt{main\_affine\_rat\_of\_rank\_le}, stated for \texttt{C : HyperellipticCurve ℚ}; the
comparator builds the curve from $f$ by \texttt{HyperellipticCurve.ofMonicOdd}, the squarefree
reduction making $f$ separable, \texttt{separable\_of\_squarefree\_twist}). The version for
$f \in \ZZ[x]$ is its case $d = 1$ (\texttt{main\_affine\_int\_of\_rank\_le}), whose case $r = g - 1$ is Coleman's
bound \texttt{main\_affine\_int}, where the prime condition reads $p \ge \deg f$.
\end{proof}

\begin{corollary}[Finiteness]
\label{cor:finiteness}
\lean{ChabautyColeman.NumberFieldBase.ratPoints_finite_numberField}
\uses{cor:finiteness-nf}
\leanok
If $f \in \QQ[x]$ is monic, squarefree and of odd degree $2g+1$, and
$\operatorname{rk} \J{\QQ} < g$, then $C(\QQ)$ is finite.
\end{corollary}
\begin{proof}
\leanok
This is the case $A = \ZZ$, $K = \QQ$ of Corollary~\ref{cor:finiteness-nf}; it is not
restated separately in Lean.
\end{proof}

\begin{remark}[Genus one, rank zero]
\label{rem:genus-one}
\uses{thm:main}
Theorem~\ref{thm:main} carries no genus lower bound: the rank hypothesis forces
$g \ge 1$, and the genus-one case is covered unconditionally
(Lemma~\ref{lem:fold-witness-genus-one}). At $g = 1$ the rank hypothesis reads
$\operatorname{rk} = 0$, so $r = 0$, and the bound is $\#E(\QQ) \le \#\bar E(\Fp)$ for a
rank-zero elliptic curve $y^2 = \text{cubic}$ at any odd prime of good reduction --- the
prime condition $p > 2r + 2$ is exactly oddness.
\end{remark}

\chapter{The theorem over number fields}
\label{ch:numberfield}

Theorem~\ref{thm:main-arith} takes as input a complete discrete valuation ring $O$ of
characteristic $0$ with algebraically closed residue field, together with arithmetic
data: a subfield $K$ of $\operatorname{Frac} O$ and a local ring $O_0$ with finite
residue field, mapping into $O$ and into $K$ and covering the elements of $K$ of
valuation $\le 1$. Over $\QQ$ these are $O = W(\Fpbar)$ and $O_0 = \ZZ_{(p)}$; the
theorems over $\QQ$ (Theorem~\ref{thm:main}, Corollary~\ref{cor:finiteness}) are derived
from the number-field statement as the case $A = \ZZ$, $K = \QQ$, $\mathfrak p = (p)$. This chapter supplies
them for an arbitrary number field $K$ and a prime $\mathfrak p$ of its ring of integers; the resulting statement,
Theorem~\ref{thm:main-nf}, is Corollary~6.7 of Stoll, \emph{Independence of rational
points on twists of a given curve}, Compositio Math.\ 142 (2006).

The ring $O$ has to be, up to the choices involved, the ring of integers of the completed
maximal unramified extension of the completion $K_{\mathfrak p}$. Rather than developing
the theory of complete discretely valued fields --- of which Mathlib (as of
version~v4.35.0) has very little: there is no Henselization, no maximal unramified extension, and no structure
theory of complete discrete valuation rings --- we construct $O$ directly, as a direct
factor of the $p$-adically complete ring $W(\Fpbar) \otimes_{\ZZ} O_K$. The merit of this
presentation is that the one genuinely local-arithmetic input --- a uniformizer of
$(O_K)_{\mathfrak p}$ remains a uniformizer of $O$, which is where ``$W(\Fpbar)$ is
unramified over $\Zp$'' enters --- reduces to the separability of the residue extension
$\Fq/\Fp$, in the form of an idempotent splitting of $\Fpbar \otimes_{\Fp} \Fq$.
Everything else is elementary commutative algebra around $p$-adic completeness.

Throughout the chapter, $K$ is a number field of degree $n$ with ring of integers $O_K$;
$p$ is a prime number and $\mathfrak p$ a maximal ideal of $O_K$ containing $p$, with
residue field $\Fq = O_K/\mathfrak p$ and ramification index $e = e(\mathfrak p \mid p)$,
so that $p\,(O_K)_{\mathfrak p} = \mathfrak p^e (O_K)_{\mathfrak p}$; the ideal
$J \le O_K$ denotes the complementary factor, so that $p\,O_K = \mathfrak p^e J$ with
$\mathfrak p + J = O_K$; $\sigma \colon \Fq \to \Fpbar$ is a fixed embedding; and
$W := W(\Fpbar)$ is the ring of Witt vectors, a complete discrete valuation ring with
maximal ideal $pW$ and residue field $\Fpbar$.

\begin{definition}[The base-change ring]
\label{def:witt-tensor}
\lean{ChabautyColeman.NumberFieldBase.rho, ChabautyColeman.NumberFieldBase.extIdeal,
  ChabautyColeman.NumberFieldBase.rho_surjective,
  ChabautyColeman.NumberFieldBase.rho_eq_zero_of_mem_extIdeal}
\leanok
Let $R := W \otimes_{\ZZ} O_K$, a free $W$-module of rank $n$, and let
$\rho \colon R \to \Fpbar$ be the ring homomorphism determined by
$w \otimes a \mapsto \bar w \cdot \sigma(\bar a)$, where $\bar w$ is the residue of $w$
and $\bar a$ the class of $a$ modulo $\mathfrak p$. It is surjective and kills $pR$ and
$\mathfrak p R$. Being module-finite over $W$, the ring $R$ is $p$-adically complete and
Hausdorff, and so is every quotient ring of $R$.
\end{definition}

\begin{lemma}[Separable corners split off]
\label{lem:residue-split}
\lean{exists_isIdempotentElem_ker_eq_span}
\leanok
Let $B$ be a commutative algebra over a field $F$, generated by a single element $t$, and
let $\rho \colon B \to F$ be an $F$-algebra homomorphism. If $t$ satisfies a polynomial
$m \in F[x]$ with $m'(\rho(t)) \neq 0$, then the kernel of $\rho$ is generated by an
idempotent: there is an idempotent $\gamma \in B$ with $\ker \rho = (1 - \gamma)$.

In particular, for a finite separable field extension $k'/k$, a field extension $\Omega$
of $k$ and a $k$-embedding $\sigma \colon k' \to \Omega$, the kernel of the multiplication
map $\mu_\sigma \colon \Omega \otimes_k k' \to \Omega$, $x \otimes a \mapsto x\,\sigma(a)$,
is generated by an idempotent: $\Omega \otimes_k k'$ is generated over $\Omega$ by
$1 \otimes \xi$ for a primitive element $\xi$ of $k'/k$ (for finite $k'$ unconditional:
the unit group is cyclic), which satisfies its separable minimal polynomial, and
separability says exactly that every root is simple.
\end{lemma}
\begin{proof}
\leanok
Write $\alpha := \rho(t)$, so $m(\alpha) = \rho(m(t)) = 0$ and $m = (x - \alpha)\,q$ in
$F[x]$. Since $q(\alpha) = m'(\alpha) \neq 0$, the factors are coprime: there is a
B\'ezout identity $u \cdot (x - \alpha) + v \cdot q = 1$. Set $\gamma := (vq)(t)$. Then
$\gamma^2 - \gamma = (vq \cdot (vq - 1))(t) = -(uv \cdot m)(t) = 0$, so $\gamma$ is an
idempotent, and $1 - \gamma = (u \cdot (x-\alpha))(t)$ lies in the kernel of $\rho$
(evaluate at $\alpha$). Conversely an element of $\ker \rho$ is $g(t)$ with
$g(\alpha) = 0$ (here generation by $t$ enters), so $g = (x - \alpha)h$; and
$(x-\alpha)(t) = \bigl((1 - vq)(x - \alpha)\bigr)(t)$ because the difference is
$v \cdot m$ --- so $g(t) \in (1 - \gamma)$.
\end{proof}

\begin{lemma}[The idempotent]
\label{lem:idempotent}
\lean{ChabautyColeman.NumberFieldBase.exists_isIdempotentElem_ker_rho,
  ChabautyColeman.NumberFieldBase.exists_isIdempotentElem_ker_lift_rho}
\uses{def:witt-tensor, lem:residue-split}
\leanok
There is an idempotent $\varepsilon \in R$ with $\varepsilon \in JR$ whose class modulo
$\mathfrak p R$ is the idempotent $\gamma$ of Lemma~\ref{lem:residue-split}, applied to
the $\Fpbar$-algebra $R/\mathfrak p R$ --- generated by the class $t$ of $1 \otimes \xi$
for a lift $\xi \in O_K$ of a primitive element of $\Fq/\Fp$, which satisfies the
separable minimal polynomial of that element --- and to the homomorphism induced by
$\rho$.
\end{lemma}
\begin{proof}
\leanok
Since $\mathfrak p R + JR = R$, choose $c \in R$ with $c \equiv \gamma \bmod \mathfrak p R$
and $c \in JR$. Its defect $c^2 - c$ lies in
$\mathfrak p R \cap JR = \mathfrak p J R$ (the ideals are coprime). The Newton step
$c \mapsto 3c^2 - 2c^3$ squares the defect ideal --- if $c^2 - c \in I$ then
$(3c^2-2c^3)^2 - (3c^2-2c^3) \in I^2$ --- and moves $c$ only by an element of $I$.
Since $(\mathfrak p J)^{me} \subseteq (\mathfrak p^e J)^m = p^m O_K$ for every $m$, the
iterates form a $p$-adic Cauchy sequence whose defects tend to $0$; by completeness of
$R$ they converge to an idempotent $\varepsilon$ with
$\varepsilon \equiv c \bmod \mathfrak p J R$, which therefore still satisfies
$\varepsilon \equiv \gamma \bmod \mathfrak p R$ and $\varepsilon \in JR$.

\emph{Formalization.} The convergence argument is the existing idempotent-lifting lemma
for adically complete rings (\texttt{IsAdicComplete.exists\_isIdempotentElem\_of\_sq\_sub\_mem}),
applied with the ideal $\mathfrak p J R$, whose adic topology agrees with the $p$-adic
one by the displayed inclusions (the transfer is
\texttt{isAdicComplete\_of\_le\_of\_pow\_le}). The congruence
$\varepsilon \equiv \gamma \bmod \mathfrak p R$ is packaged in the form the next steps
consume: $\ker \rho = \mathfrak p R + (1 - \varepsilon)$, pulled back from
$\ker(\bar\rho) = (1 - \gamma)$ on $R/\mathfrak p R$
(\texttt{exists\_isIdempotentElem\_ker\_lift\_rho}, which applies
Lemma~\ref{lem:residue-split} in an unbundled form to avoid carrying an
$\overline{\mathbb F}_p$-algebra structure on $R/\mathfrak p R$).
\end{proof}

\begin{definition}[The local factor]
\label{def:local-factor}
\lean{ChabautyColeman.NumberFieldBase.Factor, ChabautyColeman.NumberFieldBase.resO,
  ChabautyColeman.NumberFieldBase.rhoO, ChabautyColeman.NumberFieldBase.FactorData,
  ChabautyColeman.NumberFieldBase.map_eq_top_of_mem,
  ChabautyColeman.NumberFieldBase.span_natCast_eq_pow,
  ChabautyColeman.NumberFieldBase.isAdicComplete_map}
\uses{lem:idempotent}
\leanok
Let $O := R/(1 - \varepsilon)$. Since $\rho(\varepsilon) = \mu_\sigma(\gamma) = 1$, the
homomorphism $\rho$ descends to a surjection $\rho \colon O \to \Fpbar$. As a
$W$-module, $O$ is a direct summand of $R$ (the image of $\varepsilon$), hence
torsion-free and module-finite over $W$, and $p$-adically complete and Hausdorff.
Because $\varepsilon \in JR$ maps to $1$, we have $JO = O$, and therefore
\[
  \mathfrak p^e O \;=\; \mathfrak p^e J\, O \;=\; p\,O.
\]
\end{definition}

\begin{lemma}[The factor is a complete discrete valuation ring]
\label{lem:factor-dvr}
\lean{ChabautyColeman.NumberFieldBase.FactorData.isDiscreteValuationRing,
  ChabautyColeman.NumberFieldBase.FactorData.isUniformizer,
  ChabautyColeman.NumberFieldBase.FactorData.exists_unit_mul_pow,
  ChabautyColeman.NumberFieldBase.isLocalRing_factor,
  ChabautyColeman.NumberFieldBase.maximalIdeal_factor,
  ChabautyColeman.NumberFieldBase.map_eq_span,
  ChabautyColeman.NumberFieldBase.hasUniqueDiv_factor,
  ChabautyColeman.NumberFieldBase.charZero_factor,
  ChabautyColeman.NumberFieldBase.isAlgClosed_residueField,
  ChabautyColeman.NumberFieldBase.isAdic_map,
  ChabautyColeman.NumberFieldBase.t2Space_factor,
  ChabautyColeman.NumberFieldBase.completeSpace_factor}
\uses{def:local-factor, lem:residue-split}
\leanok
Choose $\varpi \in \mathfrak p \setminus \mathfrak p^2$. The ring $O$ is a complete
discrete valuation ring of characteristic $0$: it is local with maximal ideal
\[
  \mathfrak m \;=\; \ker \rho \;=\; \mathfrak p\, O \;=\; \varpi\, O,
\]
every element of $O_K \setminus \mathfrak p$ is a unit of $O$, every nonzero $x \in O$ is
$\varpi^a u$ with $u \in O^\times$ and $a = v(x)$ for the normalized valuation $v$, the
residue field is $\Fpbar$ (via $\rho$), and $v(p) = e$; in particular $p \mid \varpi^e$,
and the $\mathfrak m$-adic topology is the $p$-adic one, for which $O$ is complete and
Hausdorff.
\end{lemma}
\begin{proof}
\leanok
Modulo $\mathfrak p$, the ring $O/\mathfrak p O \cong
(\Fpbar \otimes_{\Fp} \Fq)/(1-\gamma)$ is the corner cut out by $\gamma$, which
Lemma~\ref{lem:residue-split} identifies with $\Fpbar$ via $\mu_\sigma$; so
$\mathfrak p O = \ker \rho$ is a maximal ideal. For $x \notin \mathfrak p O$ pick $y$
with $xy \equiv 1 \bmod \mathfrak p O$; then $xy = 1 - t$ with $t \in \mathfrak p O$, and
$\sum_j t^j$ converges $p$-adically because $t^{me} \in \mathfrak p^{me} O = p^m O$; so
$x$ is a unit, and $O$ is local with $\mathfrak m = \mathfrak p O$ and residue field
$\Fpbar$. An element $s \in O_K \setminus \mathfrak p$ has $\rho(s) = \sigma(\bar s)
\neq 0$, so $s \notin \mathfrak m$ is a unit; hence $O$ receives $(O_K)_{\mathfrak p}$
and $\mathfrak p O = \mathfrak p (O_K)_{\mathfrak p} O = \varpi O$.

The two filtrations $(\mathfrak m^a)_a$ and $(p^m O)_m$ are cofinal by
$\mathfrak m^e = pO$, so $O$ is $\mathfrak m$-adically complete and Hausdorff, i.e.\
$\bigcap_a \varpi^a O = 0$. Since $O$ is a nonzero direct summand of a free $W$-module,
$p$ is regular on $O$ and no $p^m$ vanishes; from $pO = \varpi^e O$ we get
$p = \varpi^e w$ and $\varpi^e = p w'$ with $w w' = 1$ (cancel the regular $p$), so
$\varpi^{em} = p^m \cdot \mathrm{unit} \neq 0$ and no power of $\varpi$ vanishes. Now
every nonzero $x \in O$ lies in some largest $\varpi^a O$ (Hausdorffness), so
$x = \varpi^a u$ with $u \notin \varpi O = \mathfrak m$ a unit. Products of nonzero
elements are then of the form $\varpi^{a+b}\,\mathrm{unit} \neq 0$, so $O$ is a domain;
a nonzero ideal is generated by $\varpi^a$ for the minimal occurring $a$, so $O$ is a
principal ideal domain, and a discrete valuation ring with $v(x) = a$, $v(p) = e$ and
$p \mid \varpi^e$. Finally $O$ is torsion-free over $\ZZ$, so of characteristic $0$.

\emph{Formalization.} The factor is a definition (\texttt{Factor}) rather than the raw
quotient, so that it carries one canonical ring structure --- the raw quotient of the
tensor ring has several definitionally equal instance paths, which block rewriting. The
uniformizer is certified by an ideal identity $(\varpi) = \mathfrak p\,C$ with
$C \not\subseteq \mathfrak p$ (equivalent to $\varpi \in \mathfrak p \setminus
\mathfrak p^2$, but directly usable: for $x \in \mathfrak p$ and $s \in C \setminus
\mathfrak p$, the product $xs$ lies in $(\varpi)$ and $s$ is a unit of $O$); all the
certificates are bundled as \texttt{FactorData}. Torsion-freeness is proved on
coordinates for a base-changed basis of $R$, and descends to $O$ by multiplying with
$\varepsilon$; the discrete-valuation structure enters Mathlib through
\texttt{IsDiscreteValuationRing.ofHasUnitMulPowIrreducibleFactorization}.
\end{proof}

\begin{lemma}[The arithmetic data]
\label{lem:factor-arithmetic}
\uses{lem:factor-dvr}
\lean{ChabautyColeman.NumberFieldBase.locFactor,
  ChabautyColeman.NumberFieldBase.isLocalHom_locFactor,
  ChabautyColeman.NumberFieldBase.comp_algebraMap_eq_comp_locFactor,
  ChabautyColeman.NumberFieldBase.resO_injective,
  ChabautyColeman.NumberFieldBase.exists_algebraMap_eq,
  ChabautyColeman.NumberFieldBase.residueFieldEquiv,
  ChabautyColeman.NumberFieldBase.finite_residueField,
  ChabautyColeman.NumberFieldBase.finite_quotient}
\leanok
Set $O_0 := (O_K)_{\mathfrak p}$, a discrete valuation ring with finite residue field
$\Fq$.
\begin{enumerate}
\item The map $O_K \to O$ extends to a local homomorphism $\iota \colon O_0 \to O$.
\item $O_K \to O$ is injective, so $K = \operatorname{Frac} O_K$ embeds into
  $\operatorname{Frac} O$, compatibly with $\iota$ and with the inclusion
  $\iota_0 \colon O_0 \to K$.
\item The valuation of $\operatorname{Frac} O$ restricts on $K$ to the $\mathfrak p$-adic
  valuation: $v(x) = v_{\mathfrak p}(x)$ for $x \in K^\times$. Consequently every element
  of $K$ that is integral for the valuation of $O$ --- the covering condition ``of
  valuation $\le 1$'' of Theorem~\ref{thm:main-arith} --- lies in $O_0$.
\end{enumerate}
\end{lemma}
\begin{proof}
\leanok
(1) is Lemma~\ref{lem:factor-dvr}: the elements of $O_K \setminus \mathfrak p$ become
units of $O$, and $\iota$ is local because $\mathfrak p O_0$ maps into
$\mathfrak p O = \mathfrak m$. For (2), the kernel of $O_K \to O$ is a prime ideal; it is
not maximal, since a maximal kernel $\mathfrak q$ would put the positive characteristic
$\ell$ of $O_K/\mathfrak q \hookrightarrow O$ into the kernel of $\ZZ \to O$,
contradicting characteristic $0$; and $O_K$ has Krull dimension $1$, so the kernel is
$0$. For (3), write $x \in K^\times$ as $\varpi^m u$ with $m = v_{\mathfrak p}(x)$ and
$u \in O_0^\times$; then $\iota(u)$ is a unit of $O$ and $v(x) = m$. The consequence
also follows directly from $O_0$ being a valuation ring of $K$: if $x \notin O_0$, then
$x^{-1}$ is a non-unit of $O_0$, so $\iota(x^{-1})$ is a non-unit of $O$; but if $x$ is
integral for the valuation of $O$, then $x = z$ for some $z \in O$, and
$z\,\iota(x^{-1}) = 1$ exhibits $\iota(x^{-1})$ as a unit.

\emph{Formalization.} $O_0$ is Mathlib's valuation subring of $K$ at $\mathfrak p$
(\texttt{valuationSubringAtPrime}), so $\iota_0$ is the inclusion, and the second
argument for (3) is the one formalized; the residue field of $O_0$ is identified with
$O_K/\mathfrak p$ by \texttt{residueFieldEquiv}.
\end{proof}

\begin{theorem}[Chabauty--Coleman over number fields]
\label{thm:main-nf}
\uses{thm:main-arith, lem:factor-dvr, lem:factor-arithmetic}
\lean{ChabautyColeman.NumberFieldBase.main_affine_numberField,
  ChabautyColeman.NumberFieldBase.main_affine_numberField_int,
  ChabautyColeman.NumberFieldBase.separable_of_squarefree_map_residue,
  ChabautyColeman.NumberFieldBase.ringChar_quotient_prime,
  ChabautyColeman.NumberFieldBase.under_eq_span_ringChar,
  ChabautyColeman.NumberFieldBase.ringChar_residueField,
  ChabautyColeman.NumberFieldBase.nonempty_ringHom_algebraicClosure,
  ChabautyColeman.NumberFieldBase.exists_pow_mul,
  Ideal.exists_span_singleton_eq_mul_of_ne_bot,
  ChabautyColeman.NumberFieldBase.squarefree_map_locFactor,
  ChabautyColeman.NumberFieldBase.encard_residuePairs_nf}
\leanok
Let $K$ be a number field and $C \colon y^2 = f(x)$ a hyperelliptic curve of genus $g$ over
$K$ with $f$ monic of odd degree $2g+1$, and let $\mathfrak p$ be a prime of $O_K$ with residue field $\Fq$, residue characteristic $p$
and ramification index $e$, such that $f$ has $\mathfrak p$-integral coefficients and the
reduction $\bar f \in \Fq[x]$ is squarefree.
If $\operatorname{rk} J(K) \le r$ with $r < g$ and
\[
  p \;>\; 2r + e + 1,
\]
then
\[
  \#C(K) \;\le\; \#\bar C(\Fq) + 2r.
\]
This is Corollary~6.7 of Stoll, \emph{loc.\ cit.}
\end{theorem}
\begin{proof}
\leanok
Instantiate Theorem~\ref{thm:main-arith} at the ring $O$ of
Definition~\ref{def:local-factor}, a complete discrete valuation ring of characteristic
$0$ with algebraically closed residue field $\Fpbar$ (Lemma~\ref{lem:factor-dvr}), with
uniformizer $\varpi$ and $v(p) = e$, and at the arithmetic data $K$,
$O_0 = (O_K)_{\mathfrak p}$ of Lemma~\ref{lem:factor-arithmetic}, whose part (3) is the
required covering of $\{x \in K : v(x) \le 1\}$ by $O_0$. The polynomial $f$ has
coefficients in $O_0$, so it is the image of a monic polynomial over $O_0$, whose reduction
modulo the maximal ideal is $\bar f$; the theorem is applied to that polynomial, and the
residue field of $O_0$ is $\Fq$, so the point count on the right-hand side is
$\#\bar C(\Fq)$. The residue characteristic $p$ is odd, indeed $p > 2r + e + 1 \ge 2$
because $e \ge 1$. The reduction of $f$ over the residue field of $O$ is the image of
$\bar f$ under $\sigma \colon \Fq \to \Fpbar$, and squarefreeness persists: over the
perfect field $\Fq$ a squarefree polynomial is separable, and separability is preserved
by any field extension.

\emph{Formalization.} The Lean statements of this theorem and of
Corollary~\ref{cor:finiteness-nf} are phrased over an arbitrary Dedekind domain $A$,
finite free over $\ZZ$, with fraction field $K$ --- so that $A = \ZZ$ itself qualifies at
$K = \QQ$, and the statements over $\QQ$ become direct instantiations. The prime is a
point $v$ of the height-one spectrum of $A$ (a finite place of $K$), $O_0$ is Mathlib's
ring of $v$-integers of $K$ (\texttt{valuationSubringAtPrime}), and the
$\mathfrak p$-integrality of $f$ is certified by a polynomial over $O_0$ mapping to $f$,
whose reduction modulo the maximal ideal is $\bar f$ and whose points over the residue
field are counted. No auxiliary data appears: the residue characteristic is the
characteristic of that residue field (\texttt{ringChar}), and the ramification index is
taken at the ideal of $\ZZ$ below $\mathfrak p$ (\texttt{Ideal.ramificationIdx'
($\mathfrak p$.under $\ZZ$) $\mathfrak p$}). The version for $f \in A[x]$, with the
reduction taken in $A/\mathfrak p$ (\texttt{main\_affine\_numberField\_int}), is a
corollary.
\end{proof}

\begin{corollary}[Finiteness over number fields]
\label{cor:finiteness-nf}
\uses{thm:main-nf}
\lean{ChabautyColeman.NumberFieldBase.ratPoints_finite_numberField,
  ChabautyColeman.NumberFieldBase.ratPoints_finite_numberField_int,
  ChabautyColeman.NumberFieldBase.exists_prime_squarefree_reduction,
  ChabautyColeman.NumberFieldBase.integralModel,
  ChabautyColeman.NumberFieldBase.dvd_absNorm,
  ChabautyColeman.NumberFieldBase.ramificationIdx'_le_finrank,
  Polynomial.separable_of_isUnit_resultant}
\leanok
Let $K$ be a number field and $C \colon y^2 = f(x)$ a hyperelliptic curve of genus $g$ over
$K$ with $f$ monic of odd degree $2g+1$ (squarefreeness of $f$ is part of being a curve). If
$\operatorname{rk} J(K) < g$, then $C(K)$ is finite. No prime appears in the
statement: a suitable prime exists for every such $f$.
\end{corollary}
\begin{proof}
\leanok
Choose $b \in O_K \setminus \{0\}$ with $g := b f \in O_K[x]$, and let $R \in K$ be the
resultant of $f$ and $f'$; it is nonzero, since $f$ is squarefree, hence separable, over
the perfect field $K$. Write $R = a/s$ with $a, s \in O_K$ nonzero. Let $r < g$ bound the
rank and choose a rational prime $p$ beyond both $2r + [K:\QQ] + 1$ and the absolute norm
$N\bigl((asb)\bigr)$, and a prime $\mathfrak p$ of $O_K$ above $p$ (one exists: $p$
divides $N\bigl((p)\bigr) = p^{[K:\QQ]}$). Then $asb \notin \mathfrak p$ --- otherwise
$p \mid N(\mathfrak p) \mid N\bigl((asb)\bigr)$, where $p \mid N(\mathfrak p)$ because
$O_K/\mathfrak p$ is a finite field of characteristic $p$ --- so none of $a$, $s$, $b$
lies in $\mathfrak p$. Hence $f = b^{-1} g$ has $\mathfrak p$-integral coefficients, and
the resultant of this integral model with its derivative is $a/s$, a unit of
$(O_K)_{\mathfrak p}$; its residue is nonzero, so the resultant of $\bar f$ and
$\bar f'$ over $\Fq$ is a unit and $\bar f$ is squarefree. Moreover the ramification
index satisfies $e \le [K:\QQ]$: taking norms in $p\,O_K = \mathfrak p^e J$ gives
$p^{[K:\QQ]} = N(\mathfrak p)^e\,N(J)$ with $p \mid N(\mathfrak p)$. Hence
$p > 2r + e + 1$, and Theorem~\ref{thm:main-nf} bounds $\#C(K)$ by the finite quantity
$\#\bar C(\Fq) + 2r$.

\emph{Formalization.} The version for $f \in O_K[x]$
(\texttt{ratPoints\_finite\_numberField\_int}) is a corollary: $f$ stays squarefree over
$K$ by Gauss's lemma.
\end{proof}

\chapter{The \texorpdfstring{$x - T$}{x − T} map and \texorpdfstring{$2$}{2}-descent}
\label{ch:descent}

Let $K$ be a field in which $2 \ne 0$, let $f \in K[x]$ be squarefree of odd degree $2g + 1$ with
leading coefficient~$c$, let $L = K[x]/(f)$, and let~$\theta$ be the image of~$x$ in~$L$. The
\emph{descent map}, or $x - T$ map, of Section~5 of M.~Stoll, \emph{Chabauty without the
Mordell--Weil group}, in: Algorithmic and Experimental Methods in Algebra, Geometry, and Number
Theory, Springer, 2017, pp.~623--663 [Sto17], is the homomorphism
\[ \mu \colon J(K) \to L^\times/(L^\times)^2\,, \qquad
   [a, b] \mapsto (-c)^{\deg a}\bigl(a(\theta) - a_1(\theta) f_1(\theta)\bigr)(L^\times)^2\,, \]
where~$a$ is monic and~$a_1$ and~$f_1$ are~$a$ and~$f$ divided by their monic gcd; its kernel
is~$2J(K)$. Over a number field, $\mu$ is the basis of $2$-descent: its image lies in the
$2$-Selmer group, a finite group computable from local data, whose order bounds the rank
of~$J(K)$ when~$J(K)$ is finitely generated. This chapter develops the map, its kernel and the
rank bound, and works out two examples over~$\QQ$.

Section~\ref{sec:hv-classes} introduces the classes $[u, v] \in J(K)$ of Mumford ideals for
arbitrary pairs of polynomials and the identities between them. Section~\ref{sec:ds-map}
defines~$\mu$ without valuation theory, from Mumford representatives prime to~$f$ and an explicit
moving lemma, and proves the formula above, including the correction term~$a_1 f_1$.
Sections~\ref{sec:hv-certificates} and~\ref{sec:hv-cantor} put Proposition~5.1 of [Sto17], which
computes a half of a point of~$2J(K)$, into certificate form: a short list of polynomial
identities and divisibilities between the polynomials that the computation produces implies that
the result is a half, and can be checked by polynomial arithmetic. Section~\ref{sec:ds-kernel}
proves $\ker\mu = 2J(K)$ by the existence half of Proposition~5.1 of [Sto17], which ends in such a
certificate. Section~\ref{sec:ds-basic} shows that the image of~$\mu$ lies in the kernel of the
norm, describes~$J(K)[2]$, and proves the compatibility with base change and the rank formula.
Section~\ref{sec:ds-dedekind}, over the fraction field of a Dedekind domain, bounds the image by
the group $A(\Sigma, 2)$ of classes unramified outside a set~$\Sigma$ of primes, and
Section~\ref{sec:ds-number} defines the $2$-Selmer group and the relaxed Selmer groups over a
number field and proves the rank bound, with the Mordell--Weil theorem and Mattuck's description
of~$J(K_v)$ as cited inputs. Section~\ref{sec:ds-examples} treats two curves of genus~$2$
over~$\QQ$: for the first, $J(\QQ) = J(\QQ)[2]$ and~$C(\QQ)$ is determined without the
Mordell--Weil theorem; for the second, $J(\QQ)$ has rank~$1$.

\emph{Standing hypotheses and conventions.} Throughout the chapter, $K$ is a field in which
$2 \ne 0$, and $f \in K[x]$ is squarefree of odd degree $2g + 1$ with leading coefficient~$c$; as in
Section~5 of [Sto17], $f$ need not be monic. As in Chapter~1, $\A{K} = K[x][y]/(y^2 - f)$ is the
coordinate ring, with the hyperelliptic involution~$\iota$, and $J(K) = \Cl(\A{K})$
(Definitions~\ref{def:hypring} and~\ref{def:hypjac}); the results of Chapters~1 and~4 on Mumford
ideals that are used here need only that~$f$ be squarefree of odd degree. Polynomials are elements
of~$K[x]$; $a \mid b$ means divisibility in~$K[x]$, and $b \equiv b' \bmod a$ means $a \mid b - b'$.
For a nonzero polynomial~$a$, $\operatorname{lc}(a)$ is its leading coefficient. The group~$J(K)$ is
written additively; the groups of square classes are written multiplicatively, so that
$\mu(P + Q) = \mu(P)\mu(Q)$. Sets of primes are denoted by~$\Sigma$, as in Remark~4.2 of [Sto17].

\begin{remark}[Formalization]
The standing hypotheses are a \texttt{variable} block: a field \texttt{K} with
\texttt{[NeZero (2 : K)]} and \texttt{f : K[X]} with \texttt{[Fact (Odd f.natDegree)]} and
\texttt{[Fact (Squarefree f)]}, the instances under which \texttt{HypRing K f} is a Dedekind
domain; the genus is not a variable. The ring~$L$ is \texttt{AdjoinRoot f}, $\theta$ is
\texttt{AdjoinRoot.root f} and~$c$ is \texttt{f.leadingCoeff}. As in Chapter~1, the Jacobian
\texttt{HypJac K f} is Mathlib's multiplicative \texttt{ClassGroup}, so sums, doubles, negatives
and~$0$ in~$J(K)$ become products, squares, inverses and~$1$; $2J(K)$ is the range of
\texttt{powMonoidHom 2}, and $J(K)[2]$ its kernel. The group $L^\times/(L^\times)^2$ is
\texttt{Units.ModPow (AdjoinRoot f) 2}. Lean and Mathlib call a set of primes \texttt{S}. Lean
names without a namespace are in \texttt{ChabautyColeman}. The declarations of this chapter are in
the modules \texttt{HyperellipticCurves.Descent.*} and in \texttt{MumfordValue},
\texttt{MumfordClass} and \texttt{JacobianStructure} of \texttt{HyperellipticCurves.Curve}; the
main theorems of Chapters~5 and~6 do not depend on them.
\end{remark}

\section{Classes of Mumford ideals}
\label{sec:hv-classes}

The Mumford ideal $\mathfrak a_{u,v} = (u, y - v)$ of Definition~\ref{def:mumford-ideal} is
defined for all $u, v \in K[x]$. This section names its class in~$J(K)$ and collects the
identities between such classes that the descent map and the certificates of this chapter
consume. None of them needs a coprimality, degree or monicity condition.

\begin{definition}[Mumford classes]
\label{def:hv-class}
\lean{ChabautyColeman.HypRing.y_sub_algebraMap_ne_zero, ChabautyColeman.mumfordIdeal_ne_bot',
  ChabautyColeman.mumfordIdeal_mem_nonZeroDivisors', ChabautyColeman.mumfordClass,
  ChabautyColeman.mumfordClass_eq_mk0, ChabautyColeman.mumfordClass_eq_one_iff}
\leanok
\uses{def:mumford-ideal, def:hypjac}
For $u, v \in K[x]$, the ideal~$\mathfrak a_{u,v}$ is nonzero, since it contains $y - v \ne 0$.
Its class in~$J(K)$ is denoted by~$[u, v]$; so $[u, v] = 0$ if and only if~$\mathfrak a_{u,v}$ is
principal.
\end{definition}

Every point of~$J(K)$ is~$[u, v]$ for exactly one reduced pair $(u, v)$, that is, one with~$u$
monic, $\deg v < \deg u \le g$ and $u \mid v^2 - f$ (Lemmas~\ref{lem:mumford-reduce}
and~\ref{lem:mumford-exists}). For every pair with~$u$ monic and $u \mid v^2 - f$, $[u, v]$ is the
point denoted in the same way in [Sto17], the class of $(u, v) - (\deg u)\infty$ for the effective
divisor~$(u, v)$ with this Mumford representation. The class~$[u, v]$ is defined for every pair,
not only for pairs with $u \mid v^2 - f$; for instance, $[0, v] = 0$, since
$\mathfrak a_{0,v} = (y - v)$ is principal.

\begin{remark}[Formalization]
\texttt{mumfordClass f u v} is \texttt{ClassGroup.mk0} applied to \texttt{mumfordIdeal f u v} with
the proof \texttt{mumfordIdeal\_mem\_nonZeroDivisors' f u v}, which needs no hypothesis $u \ne 0$.
By proof irrelevance it is the class \texttt{ClassGroup.mk0} of Chapter~1 for any other membership
proof (\texttt{mumfordClass\_eq\_mk0}); the class-level lemmas of Chapter~1 apply after a rewrite
with this lemma, and proofs never unfold the definition. \texttt{mumfordClass\_eq\_one\_iff} is
Mathlib's \texttt{ClassGroup.mk0\_eq\_one\_iff} for \texttt{mumfordClass}.
\texttt{HypRing.y\_sub\_algebraMap\_ne\_zero} holds over every nontrivial commutative ring and needs
nothing about~$f$.
\end{remark}

\begin{lemma}[Basic rules]
\label{lem:hv-class-basic}
\lean{ChabautyColeman.mumfordClass_eq_of_dvd_sub, ChabautyColeman.mumfordClass_one,
  ChabautyColeman.mumfordClass_C_mul, ChabautyColeman.mumfordClass_neg,
  ChabautyColeman.mumfordClass_eq_inv_of_dvd_add, ChabautyColeman.mumfordClass_zero_sq}
\leanok
\uses{def:hv-class}
Let $u, v, v' \in K[x]$.
\begin{enumerate}
  \item If $v \equiv v' \bmod u$, then $[u, v] = [u, v']$. Moreover, $[1, v] = 0$, and
    $[\lambda u, v] = [u, v]$ for $\lambda \in K^\times$.
  \item If $u \mid v^2 - f$, then $[u, -v] = -[u, v]$, and $[u, v'] = -[u, v]$ whenever
    $v' \equiv -v \bmod u$. In particular, $2[u, 0] = 0$ if $u \mid f$.
\end{enumerate}
\end{lemma}

\begin{proof}
\leanok
\uses{lem:mumford-norm, lem:iota-norm}
(1) Already the ideals agree: $y - v = (y - v') - (v - v')$ with $v - v' \in (u)$ gives
$\mathfrak a_{u,v} = \mathfrak a_{u,v'}$; $1 \in \mathfrak a_{1,v}$; and $\lambda u$ and~$u$
generate the same ideal. (2) Here $u \ne 0$, as $v^2 - f \ne 0$: $v^2$ is zero or of even degree,
and~$f$ has odd degree. By Lemma~\ref{lem:mumford-norm},
$\mathfrak a_{u,v}\,\mathfrak a_{u,-v} = \mathfrak a_{u,v}\,\iota(\mathfrak a_{u,v}) = (u)$ is
principal, so $[u, -v] = -[u, v]$; the second claim follows with~(1), and the last one is the case
$v = 0$, where $[u, 0] = [u, -0] = -[u, 0]$.

\emph{Formalization.} The ideal-level facts in~(1) and $v^2 - f \ne 0$ are those of Chapter~1
(\texttt{mumfordIdeal\_eq\_of\_dvd\_sub}, \texttt{mumfordIdeal\_one},
\texttt{mumfordIdeal\_C\_mul}, \texttt{sq\_sub\_ne\_zero}), and the inversion in~(2) is
\texttt{classGroup\_mk0\_mumfordIdeal\_map\_star}. In multiplicative notation, (2) consists of
\texttt{mumfordClass\_neg} (the class of $(u, -v)$ is the inverse of that of $(u, v)$),
\texttt{mumfordClass\_eq\_inv\_of\_dvd\_add} (the same for $(u, v')$ under $u \mid v + v'$) and
\texttt{mumfordClass\_zero\_sq} ($[u, 0]^2 = 1$).
\end{proof}

\begin{lemma}[A B\'ezout identity]
\label{lem:hv-bezout}
\lean{Squarefree.span_eq_top_of_mul_dvd_sq_sub}
\leanok
Let $u_1, u_2, v \in K[x]$ with $u_1 u_2 \mid v^2 - f$. Then~$u_1$, $u_2$ and~$v$ generate the
unit ideal of~$K[x]$: there are $s_1, s_2, s_3 \in K[x]$ with $s_1 u_1 + s_2 u_2 + s_3 v = 1$.
\end{lemma}

\begin{proof}
\leanok
Let~$t$ generate the ideal $(u_1, u_2, v)$ of the principal ideal domain~$K[x]$. Then~$t^2$
divides~$v^2$ and~$u_1 u_2$, hence also $f = v^2 - (v^2 - f)$. Since~$f$ is squarefree, $t$~is a
unit.

\emph{Formalization.} The statement is the equation of ideals $(u_1, u_2, v) = (1)$, for elements
$f, u_1, u_2, v$ of a commutative ring in which every finitely generated ideal is principal
(\texttt{IsBezout}); the proof applies the definition of \texttt{Squarefree f} to a generator,
with no factorization into primes. The triple $(u, 2v, w)$ with $v^2 - f = uw$ in the proof of
Lemma~\ref{lem:mumford-norm} is a special case, as~$2$ is a unit.
\end{proof}

\begin{lemma}[Factorization]
\label{lem:hv-factor}
\lean{ChabautyColeman.mumfordIdeal_mul_mumfordIdeal, ChabautyColeman.mumfordClass_mul,
  ChabautyColeman.mumfordClass_mul_of_dvd_sub, ChabautyColeman.mumfordClass_pow,
  ChabautyColeman.mumfordClass_sq_mul}
\leanok
\uses{def:hv-class}
Let $v \in K[x]$.
\begin{enumerate}
  \item If $u_1, u_2 \in K[x]$ with $u_1 u_2 \mid v^2 - f$, then
    $\mathfrak a_{u_1,v}\,\mathfrak a_{u_2,v} = \mathfrak a_{u_1 u_2, v}$; hence
    $[u_1 u_2, v] = [u_1, v] + [u_2, v]$.
  \item In~(1), $[u_1 u_2, v] = [u_1, v_1] + [u_2, v_2]$ for all $v_1, v_2 \in K[x]$ with
    $v \equiv v_1 \bmod u_1$ and $v \equiv v_2 \bmod u_2$.
  \item If $u \in K[x]$ and $n \ge 0$ with $u^n \mid v^2 - f$, then $[u^n, v] = n[u, v]$. In
    particular, if $A = a_0^2 a$ and $A \mid B^2 - f$, then $[A, B] = 2[a_0, B] + [a, B]$.
\end{enumerate}
\end{lemma}

\begin{proof}
\leanok
\uses{lem:hv-bezout, def:cantor-map, lem:cantor-correct, lem:hv-class-basic}
(1) By Lemma~\ref{lem:hv-bezout}, and since~$2$ is a unit, there are $s_1, s_2, s_3 \in K[x]$
with $s_1 u_1 + s_2 u_2 + s_3 (v + v) = 1$. Cantor's composition in certificate form
(Lemma~\ref{lem:cantor-correct}), with $v_1 = v_2 = v$, gives
$\mathfrak a_{u_1,v}\,\mathfrak a_{u_2,v} = \mathfrak a_{u_1 u_2, V}$, where
\[ V = (s_1 u_1 + s_2 u_2) v + s_3 (v^2 + f) = (1 - 2 s_3 v) v + s_3 (v^2 + f)
     = v - s_3 (v^2 - f)\,. \]
So $V \equiv v \bmod u_1 u_2$, and $\mathfrak a_{u_1 u_2, V} = \mathfrak a_{u_1 u_2, v}$, as in
the proof of Lemma~\ref{lem:hv-class-basic}\,(1). (2) follows from~(1) and
Lemma~\ref{lem:hv-class-basic}\,(1). (3) By induction on~$n$: the case $n = 0$ is
Lemma~\ref{lem:hv-class-basic}\,(1), and for $n \ge 1$ apply~(1) to $u_1 = u$ and
$u_2 = u^{n-1}$. The last claim is~(1) for $u_1 = a_0^2$ and $u_2 = a$, followed by the case
$n = 2$.

\emph{Formalization.} No coprimality of~$u_1$ and~$u_2$ is assumed: a common prime factor~$\pi$
satisfies $\pi^2 \mid v^2 - f$, so $\pi \nmid f$, and the case $u_1 = u_2$ gives
$[u^2, v] = 2[u, v]$. The ideal-level statement takes \texttt{Squarefree f} and $2 \ne 0$ as
explicit hypotheses, the class-level ones the instances. In \texttt{mumfordClass\_pow} the
exponent~$n$ is implicit, as the hypothesis $u^n \mid v^2 - f$ determines it. The last claim
of~(3), \texttt{mumfordClass\_sq\_mul}, is the reduction to squarefree~$a$ in Section~5 of [Sto17]:
$[A, B] \in 2J(K)$ if and only if $[a, B] \in 2J(K)$, and the halves of~$[A, B]$ are the points
$[a_0, B] + Q$ with $2Q = [a, B]$.
\end{proof}

\begin{lemma}[Finite products]
\label{lem:hv-split}
\lean{ChabautyColeman.mumfordClass_prod_of_dvd_sub}
\leanok
\uses{def:hv-class}
Let $n \ge 0$ and $u_1, \dots, u_n, v, v_1, \dots, v_n \in K[x]$ with
$u_1 \cdots u_n \mid v^2 - f$ and $v \equiv v_i \bmod u_i$ for $i = 1, \dots, n$. Then
$[u_1 \cdots u_n, v] = [u_1, v_1] + \dots + [u_n, v_n]$. In particular, if
$b, d_f, d_a, w_1, R, w, r \in K[x]$ satisfy $w = d_f d_a w_1$, $w \mid r^2 - f$,
$r \equiv 0 \bmod d_f$, $r \equiv b \bmod d_a$ and $r \equiv R \bmod w_1$, then
$[w, r] = [d_f, 0] + [d_a, b] + [w_1, R]$.
\end{lemma}

\begin{proof}
\leanok
\uses{lem:hv-factor, lem:hv-class-basic}
By induction on~$n$: the case $n = 0$ is Lemma~\ref{lem:hv-class-basic}\,(1), and for $n \ge 1$,
Lemma~\ref{lem:hv-factor}\,(2), applied to~$u_1 \cdots u_{n-1}$ and~$u_n$, reduces the claim to
the case~$n - 1$. The last claim is the case $n = 3$.

\emph{Formalization.} The factors are indexed by a finite set, and no coprimality is assumed; the
case of distinct linear factors~$u_i$ is \texttt{mk0\_mumfordIdeal\_prod}. The last
claim has no declaration of its own: the proof of Theorem~\ref{thm:hv-general} splits off~$w_1$
and then~$d_f$ and~$d_a$ by Lemma~\ref{lem:hv-factor} (\texttt{mumfordClass\_mul},
\texttt{mumfordClass\_mul\_of\_dvd\_sub}).
\end{proof}

\begin{lemma}[Changing signs on a factor]
\label{lem:hv-twist}
\lean{ChabautyColeman.dvd_sq_sub_of_dvd_sub_of_dvd_add,
  ChabautyColeman.mumfordClass_mul_eq_mul_sq}
\leanok
\uses{def:hv-class}
Let $a_+, a_-, b, \tilde b \in K[x]$ and $a = a_+ a_-$ with $a \mid b^2 - f$,
$\tilde b \equiv b \bmod a_+$ and $\tilde b \equiv -b \bmod a_-$. Then $a \mid \tilde b^2 - f$
and $[a, b] = [a, \tilde b] + 2[a_-, b]$.
\end{lemma}

\begin{proof}
\leanok
\uses{lem:hv-factor, lem:hv-class-basic}
As $\tilde b^2 - f = (b^2 - f) + (\tilde b - b)(\tilde b + b)$ and $a_+ a_-$ divides the product,
$a \mid \tilde b^2 - f$. By Lemma~\ref{lem:hv-factor}, $[a, b] = [a_+, b] + [a_-, b]$ and
$[a, \tilde b] = [a_+, b] + [a_-, -b]$; and $[a_-, -b] = -[a_-, b]$ by
Lemma~\ref{lem:hv-class-basic}\,(2), since $a_- \mid b^2 - f$.

\emph{Formalization.} The lemma does not assume that~$a_+$ and~$a_-$ are coprime. The
divisibility, \texttt{dvd\_sq\_sub\_of\_dvd\_sub\_of\_dvd\_add}, holds in every commutative ring.
\end{proof}

\begin{lemma}[Common factors]
\label{lem:hv-cancel}
\lean{ChabautyColeman.mumfordClass_mul_mumfordClass_eq_of_dvd_add}
\leanok
\uses{def:hv-class}
Let $t, p_1, p_2, v_1, v_2 \in K[x]$ with $t p_1 \mid v_1^2 - f$, $t p_2 \mid v_2^2 - f$ and
$t \mid v_1 + v_2$. Then $[t p_1, v_1] + [t p_2, v_2] = [p_1, v_1] + [p_2, v_2]$.
\end{lemma}

\begin{proof}
\leanok
\uses{lem:hv-factor, lem:hv-class-basic}
By Lemma~\ref{lem:hv-factor}\,(1), $[t p_i, v_i] = [t, v_i] + [p_i, v_i]$ for $i = 1, 2$. Since
$t \mid v_1^2 - f$ and $v_2 \equiv -v_1 \bmod t$, Lemma~\ref{lem:hv-class-basic}\,(2) gives
$[t, v_2] = -[t, v_1]$.

\emph{Formalization.} No coprimality is assumed; the hypothesis $t \mid v_1 + v_2$ gives the
name its suffix \texttt{\_of\_dvd\_add}.
\end{proof}

At the level of ideals, the same argument (Lemmas~\ref{lem:hv-factor}\,(1)
and~\ref{lem:mumford-norm}) gives
\[ \mathfrak a_{tp_1,v_1}\,\mathfrak a_{tp_2,v_2} =
   (t)\,\mathfrak a_{p_1,v_1}\,\mathfrak a_{p_2,v_2}\,, \]
as $\mathfrak a_{t,v_1}\,\mathfrak a_{t,v_2} = \mathfrak a_{t,v_1}\,\iota(\mathfrak a_{t,v_1})
= (t)$: the common factor~$t$ collects the points of the first effective divisor whose images
under~$\iota$ lie on the second, and Cantor's composition removes them together with their images
as the divisor of the function~$t$.

\begin{lemma}[Cantor's reduction steps, certificate form]
\label{lem:hv-reduce}
\lean{ChabautyColeman.mumfordClass_eq_of_sq_sub_eq,
  ChabautyColeman.mumfordClass_eq_of_forall_sq_sub_eq}
\leanok
\uses{def:hv-class}
Let $u, v, u', v' \in K[x]$ and $\lambda \in K$ with $v^2 - f = \lambda u u'$ and
$v' \equiv -v \bmod u'$. Then $[u, v] = [u', v']$. In particular, if $n \ge 0$,
$u_0, v_0, \dots, u_n, v_n \in K[x]$ and $\lambda_0, \dots, \lambda_{n-1} \in K$ satisfy
$v_i^2 - f = \lambda_i u_i u_{i+1}$ and $v_{i+1} \equiv -v_i \bmod u_{i+1}$ for $0 \le i < n$,
then $[u_0, v_0] = [u_n, v_n]$.
\end{lemma}

\begin{proof}
\leanok
\uses{lem:mumford-divisor-y, lem:hv-class-basic}
Since $v^2 - f \ne 0$ (proof of Lemma~\ref{lem:hv-class-basic}\,(2)), $\lambda \ne 0$.
Lemma~\ref{lem:mumford-divisor-y}, applied with $w = \lambda u'$, gives
$[u, v] = [\lambda u', -v]$; by Lemma~\ref{lem:hv-class-basic}\,(1), this is $[u', -v]$, which is
$[u', v']$. The last claim follows by induction on~$n$.

\emph{Formalization.} The step of Lemma~\ref{lem:mumford-divisor-y} used is
\texttt{mk0\_mumfordIdeal\_eq\_mk0\_cofactor\_neg}, with \texttt{mumfordIdeal\_C\_mul};
$\lambda \ne 0$ is not a hypothesis. Unlike the normalized reduction step of
Lemma~\ref{lem:cantor-correct}, which makes the cofactor monic and takes the remainder of~$-v$
modulo it, the lemma allows the cofactor divided by any constant and replaces the remainder by a
congruence. In the second statement the chain consists of sequences $\mathbb N \to K[x]$ and
$\mathbb N \to K$, of which only the values at $0, \dots, n$ (at $0, \dots, n - 1$ for~$\lambda$)
matter, with one hypothesis for all $i < n$ for each condition.
\end{proof}

In Cantor's algorithm, $u'$ is the cofactor $(v^2 - f)/u$ divided by its leading
coefficient~$\lambda$, and~$v'$ is the remainder of~$-v$ modulo~$u'$; if $\deg v < \deg u$ and
$\deg u > g$, then $\deg u' < \deg u$, so that the algorithm terminates
(Lemma~\ref{lem:mumford-reduce}). The lemma needs neither the degree bound nor the normalization.

\section{The descent map}
\label{sec:ds-map}

\begin{setting}
\label{set:ds-setting}
Let $L = K[x]/(f)$, let $\theta \in L$ be the image of~$x$, and let
$L^\square = L^\times/(L^\times)^2$ and $K^\square = K^\times/(K^\times)^2$. Let
$N \colon L^\square \to K^\square$ be the homomorphism induced by the norm~$N_{L/K}$. For a nonzero
$a \in K[x]$, let $d = \gcd(a, f)$ be the monic gcd, and $a = d a_1$, $f = d f_1$, as in
Section~5 of [Sto17].
\end{setting}

Since~$f$ is squarefree, the Chinese remainder theorem gives an isomorphism
$L \cong \prod_p L_p$, $L_p = K[x]/(p)$, over the monic irreducible factors~$p$ of~$f$, which
maps~$\theta$ to the family of the images~$\theta_p$ of~$x$. An element of~$L$ is a unit if and
only if all its components are nonzero, and a unit is a square if and only if all its components
are squares. For $h \in K[x]$, the component of~$h(\theta)$ at~$p$ is~$h(\theta_p)$, which
vanishes if and only if $p \mid h$.

\begin{definition}[The value of a Mumford polynomial]
\label{def:ds-value}
\lean{ChabautyColeman.mumfordValue}
\leanok
For a nonzero $a \in K[x]$ let
\[ m(a) = (-c)^{\deg a}\operatorname{lc}(a)^{-1}\bigl(a(\theta) - a_1(\theta) f_1(\theta)\bigr)
   \in L\,. \]
If~$a$ is prime to~$f$, then $d = 1$, $f_1 = f$, and
$m(a) = (-c)^{\deg a}\operatorname{lc}(a)^{-1} a(\theta)$.
\end{definition}

For monic~$a$, the class of~$m(a)$ is the class $\tilde\mu(a)$ of Section~5 of [Sto17]. The factor
$\operatorname{lc}(a)^{-1}$ makes $m(\lambda a) = m(a)$ for $\lambda \in K^\times$ (multiplying~$a$
by~$\lambda$ multiplies~$a_1$ by~$\lambda$ and leaves~$d$ and~$f_1$ unchanged), in accordance with
$[\lambda a, b] = [a, b]$ (Lemma~\ref{lem:hv-class-basic}\,(1)). The value depends only on~$a$; a
second polynomial~$b$ enters only through the hypothesis $a \mid b^2 - f$ of the next lemma, which
makes~$m(a)$ a unit.

\begin{lemma}[Units and multiplicativity]
\label{lem:ds-value}
\lean{ChabautyColeman.isUnit_mumfordValue, ChabautyColeman.mumfordValue_mul,
  Squarefree.dvd_of_dvd_of_dvd_sq_sub, Squarefree.isRelPrime_of_dvd_of_mul_dvd_sq_sub}
\leanok
\uses{def:ds-value}
\begin{enumerate}
  \item Let $a, b \in K[x]$ with $a \mid b^2 - f$. Then~$d$ is squarefree, $d$ and~$a_1$ are
    coprime, $d \mid b$, and~$m(a)$ is a unit of~$L$: at a monic irreducible factor~$p$ of~$f$,
    its component is $-(-c)^{\deg a}\operatorname{lc}(a)^{-1}a_1(\theta_p)f_1(\theta_p)$ if
    $p \mid d$, and $(-c)^{\deg a}\operatorname{lc}(a)^{-1}a(\theta_p)$ if $p \mid f_1$; both are
    nonzero.
  \item If $a, a' \in K[x]$ are nonzero and prime to~$f$, then $m(aa') = m(a)m(a')$.
\end{enumerate}
\end{lemma}

\begin{proof}
\leanok
(1) $a \ne 0$, as $b^2 - f \ne 0$ ($f$ has odd degree), and~$d$ divides the squarefree~$f$. Every
prime factor of~$d$ divides~$f$ and $b^2 - f$, hence~$b$; so the squarefree~$d$ divides~$b$. If a
prime~$p$ divided both~$d$ and~$a_1$, then $p^2 \mid a \mid b^2 - f$ and $p \mid b$, so
$p^2 \mid f$, which is impossible. Let $p \mid f$ be monic irreducible. If $p \mid d$, then
$p \mid a$, $p \nmid a_1$ and $p \nmid f_1$, as~$f$ is squarefree. If $p \mid f_1$, then
$p \nmid d$, so $p \nmid a$, as a common factor of~$a$ and~$f$ divides~$d$; and
$f_1(\theta_p) = 0$. (2) $aa'$ is prime to~$f$, and degrees and leading coefficients are
multiplicative.

\emph{Formalization.} \texttt{mumfordValue f a} is the image in \texttt{AdjoinRoot f} of
$(-c)^{\deg a}\operatorname{lc}(a)^{-1}(a - a_1 f_1)$, with the normalized (monic) gcd
of~$K[x]$, taken classically; its value at $a = 0$ is~$0$, as $\operatorname{lc}(0)^{-1} = 0$ in
Lean. Proofs rewrite with \texttt{mumfordValue\_eq\_of\_mul\_eq}, which computes~$m(a)$ from any
factorizations $a = d a_1$, $f = d f_1$ with monic~$d$ and $a_1$ prime to~$f_1$. The proof of~(1)
needs no decomposition of~$L$: it writes $a - a_1 f_1 = a_1(d - f_1)$, where~$a_1$ is prime to
$f = d f_1$ and $d - f_1$ is prime to $d f_1$
(\texttt{Squarefree.isCoprime\_sub\_mul\_of\_dvd\_sq\_sub}), and an element prime to~$f$ has a unit
image in \texttt{AdjoinRoot f} (\texttt{AdjoinRoot.isUnit\_mk\_iff}). The claims that~$d$
divides~$b$ and is prime to~$a_1$ hold in every commutative ring that is a decomposition monoid.
\end{proof}

\begin{lemma}[Well-definedness]
\label{lem:ds-coprime}
\lean{ChabautyColeman.isSquare_mumfordValue_mul_of_mumfordClass_eq}
\leanok
\uses{def:ds-value, def:hv-class}
Let $a, b, a', b' \in K[x]$ with $a \mid b^2 - f$ and $a' \mid b'^2 - f$, with~$a$ and~$a'$ prime
to~$f$, and with $[a, b] = [a', b']$. Then $m(a)m(a')$ is a square in~$L^\times$.
\end{lemma}

\begin{proof}
\leanok
\uses{lem:mumford-norm, lem:iota-norm, lem:normform, lem:hv-class-basic, lem:ds-value}
By Lemma~\ref{lem:hv-class-basic}\,(2), $[a', -b'] = -[a', b'] = -[a, b]$, so the ideal
$\mathfrak a_{a,b}\,\mathfrak a_{a',-b'}$ is principal, generated by $h = p + qy$ with
$p, q \in K[x]$. By Lemma~\ref{lem:mumford-norm},
$\mathfrak a_{a,b}\,\iota(\mathfrak a_{a,b}) = (a)$ and
$\mathfrak a_{a',-b'}\,\iota(\mathfrak a_{a',-b'}) = (a')$, so the principal ideal of~$\A{K}$
generated by the norm $h\,\iota(h) = p^2 - fq^2 \in K[x]$ (Lemma~\ref{lem:normform}) is generated
by~$aa'$. As~$\A{K}$ is a free $K[x]$-module with basis $1, y$, an ideal of~$K[x]$ is the
contraction of the ideal of~$\A{K}$ that it generates; hence $p^2 - fq^2 = \lambda aa'$ with
$\lambda \in K^\times$. Evaluating at~$\theta$, where~$f$ vanishes, gives
$p(\theta)^2 = \lambda a(\theta)a'(\theta)$, so
\[ m(a)m(a') = (-c)^D\bigl(\lambda\operatorname{lc}(a)\operatorname{lc}(a')\bigr)^{-1}
   p(\theta)^2\,, \qquad D = \deg a + \deg a'\,. \]
The degree of~$p^2$ is even, and that of~$fq^2$ is odd if $q \ne 0$, so the leading term of
$p^2 - fq^2$ is that of~$p^2$ if~$D$ is even and that of~$-fq^2$ if~$D$ is odd. In the first case
$\lambda\operatorname{lc}(a)\operatorname{lc}(a') = \operatorname{lc}(p)^2$, and $(-c)^D$ is a
square; in the second $\lambda\operatorname{lc}(a)\operatorname{lc}(a') =
-c\operatorname{lc}(q)^2$, and
$(-c)^D/(-c\operatorname{lc}(q)^2) = (-c)^{D-1}\operatorname{lc}(q)^{-2}$ is a square. In both
cases $m(a)m(a')$ is a square constant times~$p(\theta)^2$, and~$p(\theta)$ is a unit, as
$a(\theta)$ and~$a'(\theta)$ are (Lemma~\ref{lem:ds-value}\,(1)).

\emph{Formalization.} The norm identity of Lemma~\ref{lem:mumford-norm} carries the
well-definedness; no valuation theory of~$\A{K}$ is needed. The contraction is Mathlib's
\texttt{Ideal.map\_injective\_of\_faithfullyFlat}, as~$\A{K}$ is free over~$K[x]$
(\texttt{associated\_norm\_of\_mumfordIdeal\_mul\_mumfordIdeal\_eq\_span}). Associated polynomials
have the same value (\texttt{mumfordValue\_eq\_of\_associated}), and the case distinction on the
leading term is
\texttt{Polynomial.isSquare\_neg\_leadingCoeff\_pow\_mul\_leadingCoeff\_sq\_sub\_mul\_sq}; the step
that~$p(\theta)$ is a unit is not needed. The conclusion is that $m(a)m(a')$ is a square in
\texttt{AdjoinRoot f}; in~$L^\square$, \texttt{Units.ModPow.mk\_eq\_one\_iff\_isSquare} translates
it.
\end{proof}

\begin{lemma}[Moving, {[Sto17, Section~5]}]
\label{lem:ds-move}
\lean{ChabautyColeman.mumfordClass_eq_of_mul_eq_sub_sq,
  ChabautyColeman.exists_isCoprime_mumfordClass_eq}
\leanok
\uses{def:hv-class}
Let $a, b \in K[x]$ with $a \mid b^2 - f$.
\begin{enumerate}
  \item Let $h, a'' \in K[x]$ with $aa'' = f - (b + ha)^2$. Then $a'' \mid (b + ha)^2 - f$ and
    $[a, b] = [a'', -(b + ha)]$.
  \item For every nonzero $g \in K[x]$ there is $h \in K[x]$ such that the polynomial~$a''$
    of~(1) is prime to~$fg$.
\end{enumerate}
\end{lemma}

\begin{proof}
\leanok
\uses{lem:hv-class-basic, lem:hv-reduce}
(1) The divisibility is clear. By Lemma~\ref{lem:hv-class-basic}\,(1), $[a, b] = [a, b + ha]$;
and since $(b + ha)^2 - f = (-1) \cdot a \cdot a''$, Lemma~\ref{lem:hv-reduce} with
$\lambda = -1$ and $v' = -(b + ha)$ gives $[a, b + ha] = [a'', -(b + ha)]$.

(2) Write $f - b^2 = a a_0''$; then $a'' = a_0'' - 2hb - h^2 a$. For a monic irreducible
factor~$p$ of~$fg$, let~$\theta_p$ be the image of~$x$ in the field $k_p = K[x]/(p)$ and
\[ Q_p(t) = a_0''(\theta_p) - 2b(\theta_p)t - a(\theta_p)t^2 \in k_p[t]\,, \]
so that $a''(\theta_p) = Q_p(h(\theta_p))$. The polynomial~$Q_p$ is not zero: otherwise~$p$
divides~$a$, $b$ and~$a_0''$, hence $p^2 \mid aa_0'' + b^2 = f$, which is impossible. So~$Q_p$ has
at most two roots in~$k_p$, while~$k_p$ has at least three elements, as $2 \ne 0$ in~$K$. Choose
$t_p \in k_p$ with $Q_p(t_p) \ne 0$ for each of the finitely many~$p$, and~$h$ with
$h(\theta_p) = t_p$ for all~$p$, by the Chinese remainder theorem. Then no~$p$ divides~$a''$,
so~$a''$ is prime to~$fg$.

\emph{Formalization.} In~(1), \texttt{k} stands for~$h$. The proof of~(2) avoids the residue
fields: \texttt{Squarefree.exists\_mul\_eq\_sub\_sq\_and\_isCoprime} holds in every principal
ideal domain in which~$2$ is a unit, with any nonzero~$N$ in place of~$fg$. No prime~$p$ divides
the values of $a_0'' - 2bt - at^2$ at all three of $t = 0, 1, -1$, as it would then divide~$a$,
$b$ and~$a_0''$; an induction over the prime factors of~$N$, with the Chinese remainder theorem for
two coprime elements at each step (\texttt{IsCoprime.exists\_dvd\_sub\_and\_dvd\_sub}), gives~$h$.
\end{proof}

In the notation of [Sto17], after Proposition~5.1, $f - b^2 = ac'$ and
$[c' - 2hb - h^2 a, -b - ha] = [a, b]$; so $a'' = c' - 2hb - h^2 a$, and [Sto17] takes~$h$
constant to make~$a''$ prime to~$f$. If~$K$ is infinite, all but finitely many constants
$h = t \in K$ make~$a''$ prime to~$fg$, as in the proof of Lemma~\ref{lem:ds-move}\,(2):
$a''(\theta_p) = Q_p(t)$, and each~$Q_p$ has at most two roots in $K \subseteq k_p$. This remark
is not formalized.

\begin{definition}[The $x - T$ map]
\label{def:ds-mu}
\lean{ChabautyColeman.HypJac.descentMap,
  ChabautyColeman.HypJac.descentMap_mumfordClass_of_isCoprime}
\leanok
\uses{def:ds-value, def:hv-class}
There is a unique homomorphism $\mu \colon J(K) \to L^\square$ with
$\mu([a, b]) = m(a)(L^\times)^2$ for all $a, b \in K[x]$ with $a \mid b^2 - f$ and~$a$ prime
to~$f$.
\end{definition}

\begin{proof}
\leanok
\uses{lem:ds-move, lem:ds-coprime, lem:ds-value, lem:mumford-reduce, lem:hv-factor}
Every point of~$J(K)$ is $[a, b]$ with $a \mid b^2 - f$ (Lemma~\ref{lem:mumford-reduce}), and by
Lemma~\ref{lem:ds-move}\,(2) with $g = 1$ also with~$a$ prime to~$f$. So the formula determines a
map, which does not depend on the representative by Lemma~\ref{lem:ds-coprime}, as~$L^\square$
has exponent~$2$. Let $P = [a, b]$ and $Q = [a', b']$ with~$a$ and~$a'$ prime to~$f$; by
Lemma~\ref{lem:ds-move}\,(2) with $g = a$, the polynomial~$a'$ can be chosen prime to~$a$ as well.
By the Chinese remainder theorem there is~$B$ with $B \equiv b \bmod a$ and
$B \equiv b' \bmod a'$; then $aa' \mid B^2 - f$, as~$a$ and~$a'$ are coprime and each divides
$B^2 - f$, and $[aa', B] = P + Q$ by Lemma~\ref{lem:hv-factor}\,(2). As~$aa'$ is prime to~$f$,
$\mu(P + Q)$ is the class of $m(aa') = m(a)m(a')$ (Lemma~\ref{lem:ds-value}\,(2)).

\emph{Formalization.} \texttt{HypJac.descentMap f} chooses, for each point, a representative prime
to~$f$ (\texttt{HypJac.exists\_isCoprime\_mumfordClass\_eq}, from
\texttt{exists\_mk0\_mumfordIdeal\_reduced} and the moving lemma), and
\texttt{HypJac.descentMap\_mumfordClass\_of\_isCoprime} characterizes it: the value is the class of
any unit~$u$ with $u = m(a)$. Later proofs use only this lemma and
Theorem~\ref{thm:ds-formula}, never the choice. Two units have the same class if and only if their
product is a square (\texttt{Units.ModPow.mk\_eq\_mk\_iff\_isSquare\_mul}), which turns
Lemma~\ref{lem:ds-coprime} into the independence of the choice.
\end{proof}

\begin{theorem}[The formula for~$\mu$, {[Sto17, Section~5]}]
\label{thm:ds-formula}
\lean{ChabautyColeman.HypJac.descentMap_mumfordClass,
  ChabautyColeman.HypJac.descentMap_mumfordClass_X_sub_C}
\leanok
\uses{def:ds-mu, def:ds-value}
For all $a, b \in K[x]$ with $a \mid b^2 - f$, $\mu([a, b])$ is the class of~$m(a)$.
\end{theorem}

\begin{proof}
\leanok
\uses{lem:ds-move, lem:ds-value}
By Lemma~\ref{lem:ds-move}, there are~$h$ and~$a''$ with $aa'' = f - v^2$, $v = b + ha$, such
that~$a''$ is prime to~$f$ and $[a, b] = [a'', -v]$; so~$\mu([a, b])$ is the class of~$m(a'')$.
By Lemma~\ref{lem:ds-value}\,(1), $d \mid b$, so $v = dv_1$ with $v_1 = b/d + ha_1$, and
$a_1a'' = f_1 - dv_1^2$ after cancelling~$d$. Hence the polynomial identity
\[ (a - a_1f_1)\,a'' = f\,(1 + v_1)^2 - (v + f_1)^2\,: \]
indeed $(a - a_1f_1)a'' = aa'' - f_1 \cdot a_1a'' = f - v^2 - f_1^2 + fv_1^2$ and
$(v + f_1)^2 = v^2 + 2fv_1 + f_1^2$. Evaluating at~$\theta$ gives
$\bigl(a(\theta) - a_1(\theta)f_1(\theta)\bigr)a''(\theta) = -s^2$ with $s = (v + f_1)(\theta)$,
so
\[ m(a)m(a'') = -(-c)^D\bigl(\operatorname{lc}(a)\operatorname{lc}(a'')\bigr)^{-1}s^2\,, \qquad
   D = \deg a + \deg a''\,. \]
If $2\deg v > 2g + 1$, then~$D$ is even and
$\operatorname{lc}(a)\operatorname{lc}(a'') = -\operatorname{lc}(v)^2$, so
$m(a)m(a'') = c^D\operatorname{lc}(v)^{-2}s^2$; otherwise $D = 2g + 1$ and
$\operatorname{lc}(a)\operatorname{lc}(a'') = c$, so $m(a)m(a'') = c^{2g}s^2$. In both cases
$m(a)m(a'')$ is a square in~$L^\times$ (both factors are units by
Lemma~\ref{lem:ds-value}\,(1), so~$s$ is one), and the classes of~$m(a)$ and~$m(a'')$ agree.

\emph{Formalization.} The comparison of~$m(a)$ and $m(a'')$ is
\texttt{isSquare\_mumfordValue\_mul\_of\_mul\_eq\_sub\_sq}; the identity is checked by
\texttt{linear\_combination} after substituting $a = da_1$, $f = df_1$ and $v = dv_1$. Only one
move is needed, the moving lemma with~$f$ as the polynomial to avoid.
\texttt{HypJac.descentMap\_mumfordClass\_X\_sub\_C} is the case $a = x - x_0$ for a point
$(x_0, y_0)$ with $f(x_0) \ne 0$, where $m(a) = c(x_0 - \theta)$.
\end{proof}

Section~5 of [Sto17] states this formula and refers to Schaefer for it; the identity above proves
it. For~$a$ prime to~$f$, Theorem~\ref{thm:ds-formula} is the definition.

\begin{corollary}[Values at the points of order~$2$]
\label{cor:ds-torsion-values}
\lean{ChabautyColeman.HypJac.descentMap_mumfordClass_zero_of_mul_eq,
  ChabautyColeman.HypJac.map_liftAlgHom_descentMap_mumfordClass_X_sub_C_zero_of_ne,
  ChabautyColeman.HypJac.map_liftAlgHom_descentMap_mumfordClass_X_sub_C_zero_self}
\leanok
\uses{def:ds-mu}
Let $h \in K[x]$ be a nonzero divisor of~$f$. Then $\mu([h, 0])$ is the class of
\[ (-c)^{\deg h}\operatorname{lc}(h)^{-1}\bigl(h(\theta) - (f/h)(\theta)\bigr)\,, \]
whose component at a monic irreducible factor~$p$ of~$f$ is
$-(-c)^{\deg h}\operatorname{lc}(h)^{-1}(f/h)(\theta_p)$ if $p \mid h$ and
$(-c)^{\deg h}\operatorname{lc}(h)^{-1}h(\theta_p)$ if $p \nmid h$.
\end{corollary}

\begin{proof}
\leanok
\uses{thm:ds-formula}
$h \mid 0^2 - f$. Here $d = \operatorname{lc}(h)^{-1}h$, $a_1 = \operatorname{lc}(h)$ and
$f_1 = \operatorname{lc}(h)f/h$, so $a_1f_1 = \operatorname{lc}(h)^2f/h$. Hence~$m(h)$ differs from
the displayed element by the square factor~$\operatorname{lc}(h)^2$ in the components at the
factors of~$h$ and agrees with it in the others; apply Theorem~\ref{thm:ds-formula}.

\emph{Formalization.} The Lean statement takes the cofactor~$e$ with $he = f$ in place of~$f/h$,
so that no division appears. For a root~$x_0$ of~$f$ in~$K$ and $h = x - x_0$, the components of
the value at the roots~$x_1$ of~$f$ are $\operatorname{lc}(f)(x_0 - x_1)$ for $x_1 \ne x_0$ and
$\operatorname{lc}(f)f'(x_0)$ for $x_1 = x_0$ (the two declarations with
\texttt{map\_liftAlgHom}, for the evaluation $L \to K$ at~$x_1$).
\end{proof}

For monic~$h$ the two elements are equal. The values $\mu([h, 0])$ describe~$\mu$ on~$J(K)[2]$
(Theorem~\ref{thm:ds-two-torsion}).

\begin{remark}[The elliptic case]
\label{rem:ds-ec}
For $g = 1$ and monic~$f$, $\mu$ is the classical $x - T$ map of the elliptic curve
$y^2 = f(x)$. A point $P = (x_0, y_0)$ corresponds to $[x - x_0, y_0]$ (the Abel--Jacobi map of
Definition~\ref{def:hypjac}), and $m(x - x_0) = x_0 - \theta$ if $f(x_0) \ne 0$, while for
$f(x_0) = 0$, where $d = x - x_0$, $a_1 = 1$ and $f_1 = f/(x - x_0)$,
$m(x - x_0) = x_0 - \theta + f_1(\theta)$, whose component at $x - x_0$ is $f'(x_0)$.
\end{remark}

\section{Halving certificates}
\label{sec:hv-certificates}

For a point $P \in 2J(K)$, Proposition~5.1 of [Sto17] computes a point~$Q$ with $2Q = P$ from the
solution of a system of linear congruences. This section and the next turn that computation into
a certificate: Theorems~\ref{thm:hv-generic}, \ref{thm:hv-general} and~\ref{thm:hv-add} and
Corollary~\ref{cor:hv-record} say that a short list of polynomial identities and divisibilities
between the polynomials that the computation produces implies $2Q = P$. These conditions can be
checked by polynomial arithmetic, so neither the linear system, nor the square root~$s$ of the
proposition, nor the minimality of its solution, nor Cantor's algorithm has to be trusted. The
one new idea is Lemma~\ref{lem:hv-principal}: the function~$uy - v$ in the proof of
Proposition~5.1 of [Sto17] generates a Mumford ideal. This turns the computation of the divisor
of~$uy - v$ there into an identity of ideals, so that the halving itself needs no reduction step;
Cantor's composition and reduction enter only through the sum that forms the half when the first
polynomial of~$P$ is not squarefree (Section~\ref{sec:hv-cantor}). The proof of
$\ker\mu = 2J(K)$ in Section~\ref{sec:ds-kernel} ends in a certificate of this section.

\begin{lemma}[The function $uy - v$]
\label{lem:hv-principal}
\lean{ChabautyColeman.dvd_mul_sq_sub_of_dvd_sq_sub_sq_mul,
  ChabautyColeman.span_singleton_mul_y_sub_eq_mumfordIdeal,
  ChabautyColeman.mumfordClass_eq_one_of_sq_sub_eq}
\leanok
\uses{def:mumford-ideal, def:hv-class}
Let $u, v, e \in K[x]$, and let $N = v^2 - u^2 f$ be the norm of~$uy - v$.
\begin{enumerate}
  \item If $eu \equiv 1 \bmod N$, then the principal ideal $(uy - v)$ of~$\A{K}$ is the Mumford
    ideal~$\mathfrak a_{N, ev}$.
  \item Let $A \in K[x]$ and $\lambda \in K$ with $N = \lambda A$ and $eu \equiv 1 \bmod A$. Then
    $A \mid (ev)^2 - f$. If moreover $\lambda \ne 0$, then $[A, ev] = 0$.
\end{enumerate}
\end{lemma}

\begin{proof}
\leanok
\uses{lem:hv-class-basic, lem:mumford-exists}
(1) Write $eu = 1 + tN$ with $t \in K[x]$. Then $uy - v = u (y - ev) + tvN$ lies
in~$\mathfrak a_{N, ev}$. Conversely, $(uy - v)(uy + v) = u^2 f - v^2 = -N$, so~$N$ lies
in~$(uy - v)$, and so does $y - ev = (uy - v)\bigl(e + ty(uy + v)\bigr)$. Hence
$(uy - v) = \mathfrak a_{N, ev}$.
(2) Write $eu = 1 + tA$ with $t \in K[x]$. Then
\[ (ev)^2 - f = e^2 (v^2 - u^2 f) + (e^2 u^2 - 1) f
               = A\bigl(\lambda e^2 + t (eu + 1) f\bigr)\,. \]
If $\lambda \ne 0$, then $eu \equiv 1 \bmod N$, so $[N, ev] = 0$ by~(1), and
$[A, ev] = [\lambda A, ev] = [N, ev]$ by Lemma~\ref{lem:hv-class-basic}\,(1).

\emph{Formalization.} The ideal statement~(1) is stated for the norm~$N$ itself,
without~$\lambda$, and holds over every field, with no hypothesis on~$f$; its proof checks
membership in the Mumford ideal with \texttt{mem\_mumfordIdeal\_iff}. The statements of~(2), which
the certificates use, keep~$\lambda$, called \texttt{lam} as~$\lambda$ is a keyword of Lean; the
class-level one also assumes $\lambda \ne 0$ and passes from~$\lambda A$ to~$A$ by
\texttt{mumfordClass\_C\_mul}. For $u = e = 1$, (1) says that~$(y - v)$ is the Mumford ideal
of~$v^2 - f$; with Lemma~\ref{lem:hv-factor} this gives the ideal identity of
Lemma~\ref{lem:mumford-divisor-y}.
\end{proof}

The second polynomial of the Mumford ideal is $ev \equiv v/u \bmod N$, as it must be: at a
root~$\xi$ of~$N$, the function~$uy - v$ vanishes at the point $(\xi, v(\xi)/u(\xi))$.

\begin{definition}[Halving certificates]
\label{def:hv-cert}
\lean{ChabautyColeman.HalvingCertificate}
\leanok
Let $a, b, w, r \in K[x]$. A \emph{halving certificate} for $(a, b)$ with half $(w, r)$ consists
of polynomials $u, v, e, e' \in K[x]$ and a constant $\lambda \in K^\times$ such that
\begin{description}
  \item[\textmd{(H1)}] $u^2 f = v^2 - \lambda a w^2$,
  \item[\textmd{(H2)}] $v \equiv ub \bmod a$,
  \item[\textmd{(H3)}] $ur \equiv -v \bmod w$, and
  \item[\textmd{(H4)}] $eu + e' a w^2 = 1$.
\end{description}
\end{definition}

In Proposition~5.1 of [Sto17], let $d = \gcd(a, f)$, so that $a = d a_1$ and $f = d f_1$, let
$(u, v, w)$ be the solution of the congruences~(6) there with the degree bounds~(7) and~$w$ monic
of minimal degree, and let $d_1 = \gcd(u, w)$. In the generic case $d_1 = 1$, the
polynomials~$u$, $dv$, $w$, the polynomial~$r$ of the proposition and $\lambda = (-c)^{\deg a}$
satisfy (H1) to~(H3), with~$dv$ in the role of~$v$: (H1) is relation~(8) multiplied by~$d$, (H2)
is the second congruence in~(6) modulo~$a_1$ and holds modulo~$d$ since $d \mid b$, and~(H3)
defines~$r$. Also, $u$ is coprime to~$a$: a common prime divisor~$\pi$ of~$u$ and~$a$
divides~$dv$ by~(H2), hence~$\pi^2$ divides $\lambda a w^2$ by~(H1), so $\pi \mid w$, since~$a$ is
squarefree; this contradicts $d_1 = 1$. So cofactors as in~(H4) exist.

\begin{remark}[Formalization]
\texttt{HalvingCertificate f a b w r} is a structure with the data fields \texttt{u}, \texttt{v},
\texttt{e}, \texttt{e'} and~\texttt{lam} and one field for each condition, including
$\lambda \ne 0$; it needs only a field~$K$. It is a structure because its data come from a
computation and each condition is checked separately. The field for~(H1) is the norm identity
$v^2 - u^2 f = \lambda (a w^2)$, in the form that Lemma~\ref{lem:hv-principal} consumes; (H2)
and~(H3) are the divisibilities $a \mid v - ub$ and $w \mid ur + v$, which a concrete check proves
by supplying the quotients; (H4) is the coprimality of~$u$ and~$aw^2$ with its witnesses as data
fields. There are no monicity, degree or squarefreeness conditions: Theorem~\ref{thm:hv-generic}
needs none of them. They matter for the existence of a certificate (Proposition~5.1 of [Sto17])
and for reading $[w, r]$ as the point with the Mumford representation $(w, r)$
(Lemma~\ref{lem:mumford-exists}).
\end{remark}

\begin{lemma}[The pairs of a certificate]
\label{lem:hv-cert-mumford}
\lean{ChabautyColeman.HalvingCertificate.dvd_mul_sub_one,
  ChabautyColeman.HalvingCertificate.dvd_mul_sq_sub,
  ChabautyColeman.HalvingCertificate.dvd_mul_sub,
  ChabautyColeman.HalvingCertificate.dvd_mul_add_half,
  ChabautyColeman.HalvingCertificate.dvd_sq_sub,
  ChabautyColeman.HalvingCertificate.dvd_sq_sub_half}
\leanok
\uses{def:hv-cert}
Let $(u, v, e, e', \lambda)$ be a halving certificate for $(a, b)$ with half $(w, r)$. Then
$eu \equiv 1 \bmod a w^2$ and
\[ a w^2 \mid (ev)^2 - f\,, \qquad ev \equiv b \bmod a\,, \qquad ev \equiv -r \bmod w\,; \]
in particular $a \mid b^2 - f$ and $w \mid r^2 - f$.
\end{lemma}

\begin{proof}
\leanok
\uses{lem:hv-principal}
The first claim is~(H4). By~(H1), $v^2 - u^2 f = \lambda (a w^2)$, so
$a w^2 \mid (ev)^2 - f$ by Lemma~\ref{lem:hv-principal}\,(2). Since~$a$ and~$w$ divide~$eu - 1$,
the identities $ev - b = e(v - ub) + (eu - 1)b$ and $ev + r = e(ur + v) - (eu - 1)r$ give the two
congruences, by~(H2) and~(H3). Finally,
$b^2 - f = \bigl((ev)^2 - f\bigr) - (ev - b)(ev + b)$ and
$r^2 - f = \bigl((ev)^2 - f\bigr) - (ev + r)(ev - r)$.

\emph{Formalization.} In the namespace \texttt{HalvingCertificate}, the six statements are
separate declarations, as the proof of Theorem~\ref{thm:hv-general} uses one of them as well; the
suffix \texttt{\_half} marks the statements on the half~$(w, r)$ that have a counterpart
for~$(a, b)$.
\end{proof}

\begin{theorem}[Halving, generic case, certificate form]
\label{thm:hv-generic}
\lean{ChabautyColeman.HalvingCertificate.mumfordClass_eq_sq}
\leanok
\uses{def:hv-cert, def:hv-class}
If $(a, b)$ has a halving certificate with half $(w, r)$, then $[a, b] = 2[w, r]$ in~$J(K)$. This
is the generic case of Proposition~5.1 of [Sto17].
\end{theorem}

\begin{proof}
\leanok
\uses{lem:hv-principal, lem:hv-cert-mumford, lem:hv-factor, lem:hv-class-basic}
Let $(u, v, e, e', \lambda)$ be the certificate. By Lemma~\ref{lem:hv-principal}\,(2), applied as
in the proof of Lemma~\ref{lem:hv-cert-mumford}, $[a w^2, ev] = 0$. By
Lemma~\ref{lem:hv-cert-mumford}, $a w^2 \mid (ev)^2 - f$, $ev \equiv b \bmod a$ and
$ev \equiv -r \bmod w$; so Lemma~\ref{lem:hv-factor}, applied to~$a$ and~$w^2$ and then to~$w$
and~$w$, gives
\[ 0 = [a w^2, ev] = [a, ev] + 2[w, ev] = [a, b] + 2[w, -r]\,. \]
Since $w \mid r^2 - f$ (Lemma~\ref{lem:hv-cert-mumford}), $[w, -r] = -[w, r]$ by
Lemma~\ref{lem:hv-class-basic}\,(2). Hence $[a, b] = 2[w, r]$.

\emph{Formalization.} The conclusion reads $[a, b] = [w, r]^2$ in the multiplicative group
\texttt{HypJac K f}. Passed as terms, the field for~(H1) and \texttt{dvd\_mul\_sub\_one} determine
all implicit arguments of \texttt{mumfordClass\_eq\_one\_of\_sq\_sub\_eq}; with that lemma and those
of Lemma~\ref{lem:hv-cert-mumford}, the proof is a chain of rewrites.
\end{proof}

At the level of ideals, the same argument (Lemmas~\ref{lem:hv-principal}\,(1),
\ref{lem:hv-cert-mumford} and~\ref{lem:hv-factor}\,(1), and the ideal identities in the proof of
Lemma~\ref{lem:hv-class-basic}\,(1)) gives
$(uy - v) = \mathfrak a_{a,b}\,\mathfrak a_{w,-r}^2$; this is the identity
$\divisor(uy - v) = (a, b) + 2(w, -r) - (\deg a + 2\deg w)\infty$ in the proof of
Proposition~5.1 of [Sto17], where $(a, b)$ and~$(w, -r)$ denote divisors, as there. Neither the
degree bounds~(7) there nor the minimality of~$w$ is used.

\begin{remark}[Halves of sums]
\label{rem:hv-sums}
Definition~\ref{def:hv-cert} puts no condition on the degree of~$a$, and Proposition~5.1
of [Sto17] admits $\deg a \le 2g + 1$. So Theorem~\ref{thm:hv-generic} also certifies the halves
of a sum that is never formed. Let $P = [a_P, b_P]$ and $S = [a_S, b_S]$, put $A = a_P a_S$, and
let~$b_A$ satisfy $A \mid b_A^2 - f$, $b_A \equiv b_P \bmod a_P$ and $b_A \equiv b_S \bmod a_S$.
Then $[A, b_A] = P + S$ by Lemma~\ref{lem:hv-factor}\,(2), and a halving certificate for
$(A, b_A)$ with half $(w, r)$ shows $P + S = 2[w, r]$. This is the second variant in Remark~5.2
of [Sto17]; the first one, with $0 \ne S \in J(K)[2]$, has $\gcd(A, f) \ne 1$ and is covered by
Theorem~\ref{thm:hv-generic} if $d_1 = 1$ and by Theorem~\ref{thm:hv-general} in general.
\end{remark}

In general, $d_1 = \gcd(u, w)$ need not be~$1$ in Proposition~5.1 of [Sto17]. There
$d_1 = d_f d_a$ with $d_f = \gcd(d_1, f)$ and $d_a = \gcd(d_1, a_1)$, and
\[ u_1 = u/d_1\,, \qquad v_1 = v/d_1\,, \qquad w_1 = w/d_1\,; \]
the proposition chooses~$r$ with $r u_1 \equiv -v_1 d \bmod w_1 d_a$ and $r \equiv 0 \bmod d_f$.
Let~$R$ be the remainder of~$r$ modulo~$w_1$, and let~$\tilde b$ satisfy
$\tilde b \equiv b \bmod a/d_a$ and $\tilde b \equiv -b \bmod d_a$. Then~$u_1$, $d v_1$
and~$\lambda$, with B\'ezout cofactors for~$u_1$ and~$a w_1^2$, form a halving certificate for
$(a, \tilde b)$ with half $(w_1, R)$: (H1) is relation~(8) divided by~$d_1^2$ and multiplied
by~$d$, and~(H2) and~(H3) follow from the congruences in the proof of the proposition.
Lemmas~\ref{lem:hv-twist} and~\ref{lem:hv-split}, the content of the divisor computation at the
end of that proof, transfer $[a, \tilde b] = 2[w_1, R]$ to $[a, b] = 2[w, r]$.

\begin{theorem}[Halving in general, certificate form]
\label{thm:hv-general}
\lean{ChabautyColeman.mumfordClass_eq_sq_of_halvingCertificate}
\leanok
\uses{def:hv-cert, def:hv-class}
Let $a, b, w, r \in K[x]$ with $a \mid b^2 - f$. Assume that there are polynomials~$a'$, $d_a$,
$d_f$, $\tilde b$, $w_1$ and~$R$ in~$K[x]$ such that
\begin{enumerate}
  \item $a = a' d_a$, $\tilde b \equiv b \bmod a'$ and $\tilde b \equiv -b \bmod d_a$;
  \item $w = d_f d_a w_1$, $w \mid r^2 - f$, $r \equiv 0 \bmod d_f$, $r \equiv b \bmod d_a$
    and $r \equiv R \bmod w_1$;
  \item $(a, \tilde b)$ has a halving certificate with half $(w_1, R)$.
\end{enumerate}
Then $[a, b] = 2[w, r]$. This is Proposition~5.1 of [Sto17] in certificate form.
\end{theorem}

\begin{proof}
\leanok
\uses{thm:hv-generic, lem:hv-twist, lem:hv-factor, lem:hv-class-basic}
Lemma~\ref{lem:hv-twist} gives $[a, b] = [a, \tilde b] + 2[d_a, b]$. By
Theorem~\ref{thm:hv-generic}, $[a, \tilde b] = 2[w_1, R]$. By Lemma~\ref{lem:hv-factor}, applied
to $w = (d_f d_a) w_1$ and then to~$d_f d_a$, $2[w, r] = 2[d_f, 0] + 2[d_a, b] + 2[w_1, R]$, and
$2[d_f, 0] = 0$ by Lemma~\ref{lem:hv-class-basic}\,(2), since~$d_f$ divides~$r$ and~$r^2 - f$,
hence~$f$. So $2[w, r] = 2[d_a, b] + 2[w_1, R] = [a, b]$.

\emph{Formalization.} The certificate in~(3) is a hypothesis of type
\texttt{HalvingCertificate f a Bt W R}, with \texttt{Bt} for~$\tilde b$ and~\texttt{W}
for~$w_1$, and each condition in~(1) and~(2) is a hypothesis of its own; $a'$, $d_a$, $d_f$,
$\tilde b$, $w_1$ and~$R$ are implicit, since the hypotheses determine them. The only coprimality
used anywhere is~(H4): neither~$a'$ and~$d_a$ nor the three factors of~$w$ need to be coprime.
\end{proof}

If $d_1 = 1$, then $d_a = d_f = 1$, and the theorem with $a' = a$, $\tilde b = b$, $w_1 = w$ and
$R = r$ is Theorem~\ref{thm:hv-generic}.

\section{Cantor's algorithm and halving records}
\label{sec:hv-cantor}

When the first polynomial~$A$ of a point $[A, B]$ with $A \mid B^2 - f$ is not squarefree, write
$A = a_0^2 a$ with~$a$ squarefree; then $[A, B] = 2[a_0, B] + [a, B]$
(Lemma~\ref{lem:hv-factor}\,(3)), and a half of~$[A, B]$ is the sum of~$[a_0, B]$ and a
half~$[w, r]$ of~$[a, B]$. To obtain it as a reduced pair, the sum is computed by Cantor's
algorithm: a composition followed by reduction steps until the degree is at most~$g$. This section
turns a run of that algorithm into a certificate, as Section~\ref{sec:hv-certificates} does for
the halving. The composition needs no B\'ezout cofactors, since by Lemma~\ref{lem:hv-factor}\,(2)
three divisibilities determine the class of the composed pair. Its one special case, a common
factor of the first polynomials of the two pairs on which the pairs are opposite, cancels
(Lemma~\ref{lem:hv-cancel}). Each reduction step is an instance of Lemma~\ref{lem:hv-reduce},
which determines the new first polynomial only up to a constant factor and the new second
polynomial only by a congruence, so that a certificate needs neither the normalization of
Cantor's algorithm nor its degree conditions.

\begin{definition}[Addition certificates]
\label{def:hv-add-cert}
\lean{ChabautyColeman.AdditionCertificate}
\leanok
Let $u_1, v_1, u_2, v_2, u, v \in K[x]$. An \emph{addition certificate} for $(u_1, v_1)$
and~$(u_2, v_2)$ with sum $(u, v)$ consists of polynomials $t, p_1, p_2 \in K[x]$, an integer
$n \ge 0$, polynomials $u'_0, v'_0, \dots, u'_n, v'_n \in K[x]$ and constants
$\lambda_0, \dots, \lambda_{n-1} \in K$ such that
\begin{description}
  \item[\textmd{(A1)}] $u_1 = t p_1$, $u_2 = t p_2$ and $t \mid v_1 + v_2$;
  \item[\textmd{(A2)}] $u'_0 = p_1 p_2$, $u'_0 \mid {v'_0}^2 - f$, $v'_0 \equiv v_1 \bmod p_1$
    and $v'_0 \equiv v_2 \bmod p_2$;
  \item[\textmd{(A3)}] ${v'_i}^2 - f = \lambda_i u'_i u'_{i+1}$ and
    $v'_{i+1} \equiv -v'_i \bmod u'_{i+1}$ for $0 \le i < n$; and
  \item[\textmd{(A4)}] $(u'_n, v'_n) = (u, v)$.
\end{description}
\end{definition}

Cantor's algorithm produces such a certificate for any two pairs $(u_1, v_1)$ and~$(u_2, v_2)$
with $u_1 \mid v_1^2 - f$ and $u_2 \mid v_2^2 - f$. It takes $t = \gcd(u_1, u_2, v_1 + v_2)$;
then~$p_1$, $p_2$ and $v_1 + v_2$ generate the unit ideal of~$K[x]$. Indeed, let~$\pi$ be a common
prime divisor of the three. Then~$\pi$ divides~$t$; let~$\pi^k$, with $k \ge 1$, be the exact
power of~$\pi$ that divides~$t$. Then~$\pi^{k+1}$ divides $u_1 = t p_1$ and~$u_2$, hence
$v_1^2 - f$ and $v_2^2 - f$ and so their difference $(v_1 - v_2)(v_1 + v_2)$, whereas
$v_1 + v_2$ is divisible by~$\pi^k$, but not by~$\pi^{k+1}$, as~$t$ is not. So~$\pi$ divides
$v_1 - v_2$ as well as $v_1 + v_2$, hence~$2v_1$ and, as $2 \ne 0$ in~$K$, also~$v_1$.
Then~$\pi^2$ divides~$v_1^2$ and, as $k + 1 \ge 2$, also $v_1^2 - f$, hence~$f$, which is
impossible. So Cantor's composition (Lemma~\ref{lem:cantor-correct}), applied to $(p_1, v_1)$
and~$(p_2, v_2)$, gives~$V$ with $V \equiv v_i \bmod p_i$ for $i = 1, 2$ and
$p_1 p_2 \mid V^2 - f$, and~(A2) holds for $u'_0 = p_1 p_2$ and the remainder~$v'_0$ of~$V$
modulo~$u'_0$. The reduction steps follow the paragraph after Lemma~\ref{lem:hv-reduce}, until
$\deg u'_n \le g$.

\begin{theorem}[Cantor's algorithm, certificate form]
\label{thm:hv-add}
\lean{ChabautyColeman.AdditionCertificate.mumfordClass_mul_mumfordClass}
\leanok
\uses{def:hv-add-cert, def:hv-class}
Let $u_1, v_1, u_2, v_2, u, v \in K[x]$ with $u_1 \mid v_1^2 - f$ and $u_2 \mid v_2^2 - f$. If
$(u_1, v_1)$ and~$(u_2, v_2)$ have an addition certificate with sum $(u, v)$, then
$[u_1, v_1] + [u_2, v_2] = [u, v]$.
\end{theorem}

\begin{proof}
\leanok
\uses{lem:hv-cancel, lem:hv-factor, lem:hv-reduce}
By~(A1) and Lemma~\ref{lem:hv-cancel}, the sum $[u_1, v_1] + [u_2, v_2]$ is
$[p_1, v_1] + [p_2, v_2]$. By~(A2) and Lemma~\ref{lem:hv-factor}\,(2), this is
$[p_1 p_2, v'_0] = [u'_0, v'_0]$; by~(A3) and Lemma~\ref{lem:hv-reduce}, it is $[u'_n, v'_n]$,
which is $[u, v]$ by~(A4).

\emph{Formalization.} \texttt{AdditionCertificate}, with the six polynomials as parameters, is a
structure with the data $t$, $p_1$, $p_2$ and~$n$, the sequences $\mathbb N \to K[x]$ and
$\mathbb N \to K$ of the chain, as in Lemma~\ref{lem:hv-reduce}, and one field for each condition,
those of~(A3) for all $i < n$; it needs only a field~$K$. The fields for $t \mid v_1 + v_2$,
$v'_0 \equiv v_1 \bmod p_1$ and $v'_0 \equiv v_2 \bmod p_2$ are primed (\texttt{dvd\_add'},
\texttt{dvd\_sub\_left'}, \texttt{dvd\_sub\_right'}), as Mathlib marks such fields
(\texttt{Subgroup.mul\_mem'}), so that they do not hide Mathlib's lemmas of the same names in the
namespace. The divisibilities $u_1 \mid v_1^2 - f$ and $u_2 \mid v_2^2 - f$ are hypotheses of the
theorem, not fields: in a halving record they follow from~(R1) and~(R3). A concrete certificate
writes the sequences as lists. The certificate puts no condition on degrees or leading coefficients,
so it does not show that $(u, v)$ is reduced; that is a separate check, needed only to read~$[u, v]$
as the point with the Mumford representation $(u, v)$ (Lemma~\ref{lem:mumford-exists}). Compared
with the composition of Lemma~\ref{lem:cantor-correct}, the certificate needs no B\'ezout cofactors,
which the proof of Lemma~\ref{lem:hv-factor} obtains from Lemma~\ref{lem:hv-bezout}, and it allows a
common factor~$t$.
\end{proof}

A halving computation following Proposition~5.1 of [Sto17], for a point $[A, B]$ with
$A \mid B^2 - f$, thus produces the factorization $A = a_0^2 a$, a representative~$b$ of~$B$
modulo~$a$, the data of Theorem~\ref{thm:hv-general} for a half~$[w, r]$ of~$[a, b]$, and a run of
Cantor's algorithm for $[a_0, B] + [w, r]$. The record of such a computation is the following.

\begin{definition}[Halving records]
\label{def:hv-record}
\lean{ChabautyColeman.HalvingData}
\leanok
\uses{def:hv-cert, def:hv-add-cert}
Let $A, B, U_H, V_H \in K[x]$. A \emph{halving record} for $(A, B)$ with half $(U_H, V_H)$
consists of polynomials~$a_0$, $a$, $b$, $a'$, $d_a$, $d_f$, $\tilde b$, $w_1$, $R$, $w$ and~$r$
in~$K[x]$, a halving certificate for $(a, \tilde b)$ with half $(w_1, R)$ and an addition
certificate for $(a_0, B)$ and $(w, r)$ with sum $(U_H, V_H)$ such that
\begin{description}
  \item[\textmd{(R1)}] $A = a_0^2 a$, $A \mid B^2 - f$ and $b \equiv B \bmod a$;
  \item[\textmd{(R2)}] $a = a' d_a$, $\tilde b \equiv b \bmod a'$ and
    $\tilde b \equiv -b \bmod d_a$;
  \item[\textmd{(R3)}] $w = d_f d_a w_1$, $w \mid r^2 - f$, $r \equiv 0 \bmod d_f$,
    $r \equiv b \bmod d_a$ and $r \equiv R \bmod w_1$.
\end{description}
\end{definition}

\begin{corollary}[A halving record certifies a half]
\label{cor:hv-record}
\lean{ChabautyColeman.HalvingData.mumfordClass_eq_sq}
\leanok
\uses{def:hv-record, def:hv-class}
If $(A, B)$ has a halving record with half $(U_H, V_H)$, then $[A, B] = 2[U_H, V_H]$.
\end{corollary}

\begin{proof}
\leanok
\uses{thm:hv-general, thm:hv-add, lem:hv-factor, lem:hv-class-basic}
By~(R1), $a$ divides~$A$ and hence~$B^2 - f$; as $b^2 - f = (B^2 - f) + (b - B)(b + B)$, also
$a \mid b^2 - f$. By Lemma~\ref{lem:hv-factor}\,(3) and Lemma~\ref{lem:hv-class-basic}\,(1),
$[A, B] = 2[a_0, B] + [a, B] = 2[a_0, B] + [a, b]$, and $[a, b] = 2[w, r]$ by
Theorem~\ref{thm:hv-general}. Finally, $a_0 \mid A \mid B^2 - f$ by~(R1) and $w \mid r^2 - f$
by~(R3), so $[a_0, B] + [w, r] = [U_H, V_H]$ by Theorem~\ref{thm:hv-add}, and
$[A, B] = 2\bigl([a_0, B] + [w, r]\bigr) = 2[U_H, V_H]$.

\emph{Formalization.} \texttt{HalvingData f A B UH VH} is a structure with the data fields of the
definition (\texttt{Bt} for~$\tilde b$, \texttt{W} for~$w_1$ and \texttt{aq} for~$a'$), the
fields \texttt{cert}, a \texttt{HalvingCertificate}, and \texttt{add}, an
\texttt{AdditionCertificate}, and one field for each condition (the one for
$b \equiv B \bmod a$ is \texttt{dvd\_sub'}, primed for the reason given after
Theorem~\ref{thm:hv-add}). As in a halving certificate, the half is a parameter, so the
conclusion $[A, B] = [U_H, V_H]^2$ mentions no field of the record. If moreover $(U_H, V_H)$ is
reduced, a separate check, then $[U_H, V_H]$ is the point with the Mumford representation
$(U_H, V_H)$ (Lemma~\ref{lem:mumford-exists}).
\end{proof}

If $d_1 = 1$, then $d_a = d_f = 1$, $a' = a$, $\tilde b = b$, $w_1 = w$ and $R = r$. If~$A$ is
squarefree, then $a_0 = 1$ and $a = A$, and the addition certificate is trivial: $t = p_1 = 1$,
$p_2 = w$, $n = 0$ and $(U_H, V_H) = (w, r)$.

\section{The kernel of the descent map}
\label{sec:ds-kernel}

The inclusion $2J(K) \subseteq \ker\mu$ is formal, as~$L^\square$ has exponent~$2$. The converse
is the existence half of Proposition~5.1 of [Sto17]: if the class of~$m(a)$ is trivial, the linear
system of the proposition has a solution, and a descent ending in a halving certificate
(Theorem~\ref{thm:hv-generic} with Lemma~\ref{lem:hv-twist}) shows $[a, b] \in 2J(K)$.

\begin{theorem}[Existence of halves, {[Sto17, Proposition~5.1]}]
\label{thm:ds-halving-exists}
\lean{ChabautyColeman.exists_mumfordClass_eq_sq_of_isSquare_mumfordValue,
  ChabautyColeman.exists_sq_mul_eq_mul_sq_sub_mul_sq,
  ChabautyColeman.exists_mumfordClass_eq_sq_of_sq_sub_eq}
\leanok
\uses{def:ds-value, def:hv-class}
Let $a, b, s \in K[x]$ with~$a$ monic and squarefree, $\deg a \le 2g + 1$, $a \mid b^2 - f$ and
$s^2 \equiv (-c)^{\deg a}(a - a_1f_1) \bmod f$. Then $[a, b] \in 2J(K)$.
\end{theorem}

\begin{proof}
\leanok
\uses{thm:hv-generic, lem:hv-twist, lem:ds-value}
As in [Sto17], the polynomials~$d$, $a_1$ and~$f_1$ are squarefree and pairwise coprime, and~$b$
is prime to~$a_1$ (Lemma~\ref{lem:ds-value}\,(1)). Let $\lambda = (-c)^{\deg a}$.
\begin{enumerate}
  \item \emph{A solution.} The congruences~(6) of [Sto17],
    \[ vd \equiv ws \bmod f_1\,, \qquad vd \equiv ub \bmod a_1\,, \qquad
       uf_1 \equiv ws \bmod d\,, \]
    are linear in the coefficients of~$u$, $v$ and~$w$. Under the degree bounds~(7),
    $\deg u < \deg a/2$, $\deg v \le g + \deg a/2 - \deg d$ and $\deg w \le g$, there are
    $\deg f_1 + \deg a_1 + \deg d + 1$ coefficients and $\deg f_1 + \deg a_1 + \deg d$ linear
    conditions, so there is a nonzero solution.
  \item \emph{$w \ne 0$.} If $w = 0$, then $f_1 \mid v$, so $v = 0$ by the degree bounds (here
    $\deg a \le 2g + 1$ is used); then $a_1 \mid u$ and $d \mid u$, so $a \mid u$, and $u = 0$ by
    the degree bound. So $w \ne 0$, and~$w$ may be taken monic.
  \item \emph{The relation~(8)} $u^2f_1 = dv^2 - \lambda a_1w^2$ holds modulo~$f_1$, $a_1$ and~$d$
    by the three congruences and $s^2 \equiv \lambda(a - a_1f_1) \bmod f$, hence
    modulo~$f_1a_1d$, and all its terms have degree less than $\deg(f_1a_1d)$.
  \item \emph{The function $Uy - V$.} Let $U = u$, $V = vd$ and $W = w$. By~(8),
    $V^2 - U^2f = d(dv^2 - u^2f_1) = \lambda aW^2$, and $V \equiv Ub \bmod a$: modulo~$a_1$ by the
    second congruence, and modulo~$d$ as~$d$ divides~$b$ and~$V$.
  \item \emph{Descent.} We show by induction on $\deg W$ that $[a, b] \in 2J(K)$ whenever
    $a \mid b^2 - f$ and $U, V, W \in K[x]$ satisfy $W \ne 0$, $V^2 - U^2f = \lambda aW^2$ and
    $V \equiv Ub \bmod a$. If~$U$ is prime to~$W$, it is prime to~$aW^2$: a prime~$p$
    dividing~$U$ and~$a$ divides~$V$, so $p^2 \mid \lambda aW^2$, and $p \mid W$ as~$a$ is
    squarefree. Then~$U$, $V$, $\lambda$ and B\'ezout cofactors $e, e'$ with $eU + e'aW^2 = 1$
    form a halving certificate for $(a, b)$ with half $(W, R)$, $R = -eV$, and
    Theorem~\ref{thm:hv-generic} gives $[a, b] = 2[W, R]$. Otherwise let~$p$ be a common prime
    factor of~$U$ and~$W$. Then $p^2 \mid V^2$, so $p \mid V$, and $(U/p, V/p, W/p)$ satisfies the
    relation, with $a \mid p(V/p - (U/p)b)$. If $p \nmid a$ or $p \mid V/p - (U/p)b$, then
    $V/p \equiv (U/p)b \bmod a$, and the induction hypothesis applies. Otherwise $p \mid a$, and
    $p \mid V/p + (U/p)b$, as~$p$ divides $(V/p)^2 - (U/p)^2f$ and $b^2 - f$. Choose~$b'$ with
    $b' \equiv b \bmod a/p$ and $b' \equiv -b \bmod p$; then $a \mid b'^2 - f$,
    $V/p \equiv (U/p)b' \bmod a$, and $[a, b] = [a, b'] + 2[p, b]$ by Lemma~\ref{lem:hv-twist};
    the induction hypothesis applies to $(a, b')$.
\end{enumerate}

\emph{Formalization.} The statement takes the hypothesis that~$m(a)$ is a square in~$L$ in place
of the congruence for~$s$ (\texttt{AdjoinRoot.isSquare\_mk\_iff}), and it does not assume~$a$
monic: for general~$a$, the factor $\operatorname{lc}(a)^{-1}$ of~$m(a)$ enters~$\lambda$, which
becomes $(-c)^{\deg a}\operatorname{lc}(a)^{-1}$. Steps (1) to~(3) are
\texttt{exists\_sq\_mul\_eq\_mul\_sq\_sub\_mul\_sq}: the solution space is a product of spaces of
polynomials of bounded degree (\texttt{Polynomial.degreeLT}), the conditions are a $K$-linear map
to $K[x]/(f_1) \times K[x]/(a_1) \times K[x]/(d)$, whose dimension is
$\deg f_1 + \deg a_1 + \deg d$ (Mathlib's \texttt{AdjoinRoot.powerBasis}), and a linear map to a
space of smaller dimension has a nonzero kernel
(\texttt{Polynomial.exists\_ne\_zero\_forall\_mem\_degreeLT\_and\_dvd\_vecMul}). Step~(2) is a
parity argument there: for $w = 0$, (8) reads $u^2f_1 = dv^2$, and as $\deg(df_1)$ is odd,
$u = v = 0$ (\texttt{Polynomial.eq\_zero\_of\_mul\_sq\_eq\_mul\_sq}). Step~(5) is
\texttt{exists\_mumfordClass\_eq\_sq\_of\_sq\_sub\_eq}, by strong induction on the degree of~$W$;
it replaces the minimal solution and the factors $d_f$, $d_a$ of [Sto17], and needs only~(8) and
the second congruence. The conclusion is that $[a, b]$ is a square in the multiplicative group
\texttt{HypJac K f}.
\end{proof}

\begin{theorem}[The kernel of~$\mu$]
\label{thm:ds-kernel}
\lean{ChabautyColeman.HypJac.ker_descentMap}
\leanok
\uses{def:ds-mu}
The kernel of~$\mu$ is~$2J(K)$.
\end{theorem}

\begin{proof}
\leanok
\uses{thm:ds-halving-exists, thm:ds-formula, lem:mumford-reduce, lem:hv-factor}
$\mu(2Q) = \mu(Q)^2 = 1$. Conversely, let $\mu(P) = 1$, and let $P = [a, b]$ with a reduced pair
(Lemma~\ref{lem:mumford-reduce}): $a$ is monic of degree at most~$g$. Write $a = a_0^2a'$
with~$a'$ monic and squarefree. By Lemma~\ref{lem:hv-factor}\,(3), $P = 2[a_0, b] + [a', b]$, so
$\mu([a', b]) = 1$, and by Theorem~\ref{thm:ds-formula} there is $s \in K[x]$ with
$s(\theta)^2 = m(a') = (-c)^{\deg a'}\bigl(a'(\theta) - a'_1(\theta)f'_1(\theta)\bigr)$, that is,
$s^2 \equiv (-c)^{\deg a'}(a' - a'_1f'_1) \bmod f$, where $a' = d'a'_1$ and $f = d'f'_1$ with
$d' = \gcd(a', f)$. As $\deg a' \le g \le 2g + 1$ and $a' \mid b^2 - f$,
Theorem~\ref{thm:ds-halving-exists} gives $[a', b] \in 2J(K)$, hence $P \in 2J(K)$.

\emph{Formalization.} The kernel of \texttt{HypJac.descentMap f} is the range of
\texttt{powMonoidHom 2}. The decomposition $a = a_0^2a'$ is Mathlib's
\texttt{exists\_sq\_mul\_squarefree}; $a'$ need not be made monic, as the Lean form of
Theorem~\ref{thm:ds-halving-exists} does not assume it. The steps are
\texttt{exists\_mumfordClass\_eq\_sq\_of\_descentMap\_eq\_one} and its squarefree case
\texttt{exists\_mumfordClass\_eq\_sq\_of\_squarefree\_of\_descentMap\_eq\_one}.
\end{proof}

\begin{corollary}[The injection of $J(K)/2J(K)$]
\label{cor:ds-c0}
\lean{ChabautyColeman.HypJac.ker_descentMap_le_sup}
\leanok
\uses{def:ds-mu}
The map~$\mu$ induces an injective homomorphism $J(K)/2J(K) \to L^\square$. In particular,
$\ker\mu \subseteq \Gamma + 2J(K)$ for every subgroup $\Gamma \subseteq J(K)$.
\end{corollary}

\begin{proof}
\leanok
\uses{thm:ds-kernel}
$\ker\mu = 2J(K) \subseteq \Gamma + 2J(K)$, by Theorem~\ref{thm:ds-kernel}.

\emph{Formalization.} The Lean statement is the second claim: the kernel of
\texttt{HypJac.descentMap f} lies in the supremum of~$\Gamma$ and the range of
\texttt{powMonoidHom 2}. Restricting the codomain of~$\mu$ to a subgroup that contains its image
does not change the kernel.
\end{proof}

In the proof of Theorem~2.1 of [Sto17], the map $A(k) \to \Sel_p A$ to the $p$-Selmer group of an
abelian variety~$A$ plays the role of~$\mu$, and the inclusion of Corollary~\ref{cor:ds-c0} for it
follows from the injectivity of $A(k)/pA(k) \to \Sel_p A$. For odd-degree hyperelliptic curves,
Algorithm~4.1 and Remark~4.2 of [Sto17] use~$\mu$ in this role, with $p = 2$.

\section{Norm, \texorpdfstring{$2$}{2}-torsion, base change and rank}
\label{sec:ds-basic}

\begin{lemma}[Norms and resultants]
\label{lem:ds-resultant}
\lean{AdjoinRoot.leadingCoeff_pow_mul_norm_mk_eq_resultant}
\leanok
For every $p \in K[x]$, $c^{\deg p}N_{L/K}\bigl(p(\theta)\bigr) = \operatorname{Res}(f, p)$.
\end{lemma}

\begin{proof}
\leanok
$L = K[x]/(c^{-1}f)$, and $c^{-1}f$ is monic, so the formula for monic moduli gives
$N_{L/K}(p(\theta)) = \operatorname{Res}(c^{-1}f, p) = c^{-\deg p}\operatorname{Res}(f, p)$, by
the homogeneity of the resultant in its first argument.

\emph{Formalization.} The formula for monic moduli is \texttt{AdjoinRoot.norm\_mk\_eq\_resultant}.
As~$f$ and $c^{-1}f$ are associated, the two quotients are isomorphic by Mathlib's
\texttt{AdjoinRoot.algEquivOfAssociated}, but they are different types; the lemma transports the
norm along this isomorphism (\texttt{Algebra.norm\_eq\_of\_algEquiv}) and compares the resultants
by \texttt{Polynomial.resultant\_C\_mul\_left}. The Lean statement holds over every commutative
ring, for~$f$ with invertible leading coefficient.
\end{proof}

\begin{theorem}[The norm condition]
\label{thm:ds-norm}
\lean{ChabautyColeman.HypJac.descentMap_mem_ker_norm}
\leanok
\uses{def:ds-mu}
$N(\mu(P)) = 1$ for every $P \in J(K)$.
\end{theorem}

\begin{proof}
\leanok
\uses{thm:ds-formula, lem:ds-move, lem:ds-resultant}
Let $P = [a, b]$ with~$a$ prime to~$f$ (Lemma~\ref{lem:ds-move}). By
Lemma~\ref{lem:ds-resultant} and the symmetry of the resultant,
\[ N_{L/K}\bigl(m(a)\bigr) = \bigl((-c)^{\deg a}\operatorname{lc}(a)^{-1}\bigr)^{2g+1}
   c^{-\deg a}(-1)^{(2g+1)\deg a}\operatorname{Res}(a, f)\,. \]
Over a splitting field of~$a$, $\operatorname{Res}(a, f) =
\operatorname{lc}(a)^{2g+1}\prod_\alpha f(\alpha) = \operatorname{lc}(a)^{2g+1}\bigl(\prod_\alpha
b(\alpha)\bigr)^2$, the products running over the roots~$\alpha$ of~$a$ with multiplicity, since
$f(\alpha) = b(\alpha)^2$; and $\prod_\alpha b(\alpha) = \operatorname{lc}(a)^{-\deg
b}\operatorname{Res}(a, b) \in K$. Collecting the factors, $N_{L/K}(m(a)) = \bigl(c^{g\deg
a}\operatorname{lc}(a)^{-\deg b}\operatorname{Res}(a, b)\bigr)^2$.

\emph{Formalization.} The statement is that~$\mu(P)$ lies in the kernel of the map
\texttt{Units.ModPow.map (Algebra.norm K) 2} induced by the norm. The proof needs no splitting
field: for $a \mid b^2 - f$, the resultant $\operatorname{Res}(a, f)$ is computed in an even
degree from $f \equiv b^2 \bmod a$
(\texttt{Polynomial.isSquare\_leadingCoeff\_pow\_mul\_resultant\_of\_dvd\_sq\_sub}), and
\texttt{isSquare\_norm\_mumfordValue} collects the constants.
\end{proof}

\begin{theorem}[$2$-torsion]
\label{thm:ds-two-torsion}
\lean{ChabautyColeman.HypJac.sq_eq_one_iff, ChabautyColeman.mumfordClass_zero_eq_of_mul_eq,
  ChabautyColeman.mumfordClass_zero_inj, ChabautyColeman.HypJac.card_ker_powMonoidHom_two}
\leanok
\uses{def:hv-class}
\begin{enumerate}
  \item A point $P \in J(K)$ satisfies $2P = 0$ if and only if $P = [h, 0]$ for a divisor~$h$
    of~$f$.
  \item $[h, 0] = [f/h, 0]$ for every nonzero divisor~$h$ of~$f$, and the points $[h, 0]$ with
    $h \mid f$ monic and $\deg h \le g$ are pairwise distinct.
  \item $J(K)[2]$ has~$2^{m-1}$ elements, where~$m$ is the number of irreducible factors of~$f$.
\end{enumerate}
\end{theorem}

\begin{proof}
\leanok
\uses{lem:hv-class-basic, lem:mumford-reduce, lem:mumford-exists, lem:hv-reduce}
(1) If $h \mid f$, then $2[h, 0] = 0$ by Lemma~\ref{lem:hv-class-basic}\,(2). Conversely, let
$P = [u, v]$ with a reduced pair (Lemma~\ref{lem:mumford-reduce}). Then $(u, -v)$ is reduced as
well, and $[u, -v] = -P$; so $2P = 0$ gives $v = -v$ by the uniqueness of reduced representatives
(Lemma~\ref{lem:mumford-exists}), hence $v = 0$ and $u \mid -f$. (2) As
$0^2 - f = (-1) \cdot h \cdot (f/h)$, Lemma~\ref{lem:hv-reduce} with $\lambda = -1$ gives
$[h, 0] = [f/h, 0]$. A pair $(h, 0)$ with~$h$ monic and $\deg h \le g$ is reduced, so these points
are distinct by Lemma~\ref{lem:mumford-exists}. (3) The~$2^m$ monic divisors~$h$ of~$f$ fall into
pairs $\{h, c^{-1}f/h\}$, of which exactly one member has degree at most~$g$, as the degrees add
up to $2g + 1$; by~(1) and~(2), the points of order dividing~$2$ are the $2^{m-1}$ distinct points
$[h, 0]$ with $\deg h \le g$.

\emph{Formalization.} Mathlib has no subgroup of $n$-torsion points of a multiplicative group, so
$J(K)[2]$ is the kernel of the squaring map \texttt{powMonoidHom 2}, and~$m$ is the cardinality of
the finite set \texttt{f.Factors} of monic irreducible factors. In~(3), a monic divisor of~$f$ is a
product over a subset of \texttt{f.Factors} (\texttt{Polynomial.Factors.exists\_prod\_eq}), a subset
and its complement give the pair $\{h, c^{-1}f/h\}$, and
\texttt{Finset.natCard\_subtype\_eq\_two\_pow\_of\_compl\_iff\_not} counts the subsets with
$2\deg h < \deg f$. That $[h, 0]$ has order dividing~$2$ is \texttt{mumfordClass\_zero\_sq}.
In~(2), \texttt{mumfordClass\_zero\_eq\_of\_mul\_eq} takes the cofactor~$e$ with $he = f$ and is
one Cantor step, and \texttt{mumfordClass\_zero\_inj} is the distinctness for monic~$h$ with
$2\deg h < \deg f$, from \texttt{mumfordClass\_inj}, the form for classes of
\texttt{mumfordIdeal\_reduced\_unique}, which also gives~(1).
\end{proof}

\begin{theorem}[Base change]
\label{thm:ds-naturality}
\lean{ChabautyColeman.HypJac.map_descentMap, ChabautyColeman.HypJac.map_comp_descentMap}
\leanok
\uses{def:ds-mu, def:basechange}
Let $K \to K'$ be a field extension over which~$f$ stays squarefree (for instance, $K$ perfect),
let $L' = K'[x]/(f) = L \otimes_K K'$, and let $\mu' \colon J(K') \to L'^\square$ be the map of
Definition~\ref{def:ds-mu} over~$K'$. Then $\mu'(P_{K'}) = \mu(P)_{L'}$ for every $P \in J(K)$,
where $P \mapsto P_{K'}$ is the map $J(K) \to J(K')$ of Definition~\ref{def:basechange} and
$\xi \mapsto \xi_{L'}$ the map $L^\square \to L'^\square$.
\end{theorem}

\begin{proof}
\leanok
\uses{thm:ds-formula}
Let $P = [a, b]$ with~$a$ prime to~$f$. Then $P_{K'} = [a, b]$ over~$K'$, $a$ is prime to~$f$
over~$K'$ by the same B\'ezout identity, and~$m(a)$ computed in~$L'$ is the image of~$m(a)$
computed in~$L$.

\emph{Formalization.} $L'$ is \texttt{AdjoinRoot (f.map (algebraMap K K'))}, the map $L \to L'$
is \texttt{AdjoinRoot.map}, and on square classes \texttt{Units.ModPow.map}; the map on Jacobians
is \texttt{HypJac.map}, and \texttt{HypJac.map\_mumfordClass} computes it on Mumford classes.
\texttt{HypJac.map\_descentMap} is the pointwise statement and \texttt{HypJac.map\_comp\_descentMap}
the square of homomorphisms; as~$m(a)$ commutes with every homomorphism of fields
(\texttt{map\_mumfordValue}), the proof uses Theorem~\ref{thm:ds-formula} on any representative,
not one prime to~$f$. The hypotheses are the instances
\texttt{Fact (Squarefree (f.map (algebraMap K K')))}, which
\texttt{instFactSquarefreeMapAlgebraMap} provides for a perfect field~$K$, and
\texttt{NeZero (2 : K')}, which holds for every extension of~$K$.
\end{proof}

For a number field~$K$ and a place~$v$, with $K' = K_v$, the theorem says that the local descent
map~$\mu_v$ takes the image of~$P$ in~$J(K_v)$ to the restriction of~$\mu(P)$ to $L \otimes_K
K_v$. The map $J(K) \to J(K')$ is injective (Theorem~\ref{thm:basechange-inj}).

\begin{theorem}[Rank]
\label{thm:ds-rank}
\lean{ChabautyColeman.HypJac.card_range_descentMap,
  ChabautyColeman.HypJac.pow_rank_le_card_of_range_descentMap_le}
\leanok
\uses{def:ds-mu}
Assume that~$J(K)$ is finitely generated, of rank~$\operatorname{rk} J(K)$. Then
\[ \#\mu(J(K)) = 2^{\operatorname{rk} J(K)}\,\#J(K)[2]\,. \]
Consequently $2^{\operatorname{rk} J(K)}\,\#J(K)[2] \le \#S$ for every finite subgroup
$S \subseteq L^\square$ containing~$\mu(J(K))$.
\end{theorem}

\begin{proof}
\leanok
\uses{thm:ds-kernel}
$\mu(J(K)) \cong J(K)/2J(K)$ by Theorem~\ref{thm:ds-kernel}, and a finitely generated abelian
group $G \cong \ZZ^n \times T$ with~$T$ finite has $\#G/2G = 2^n\,\#T/2T = 2^n\,\#G[2]$.

\emph{Formalization.} The hypothesis is the instance \texttt{IsMulFG (HypJac K f)}, the rank is
\texttt{Module.finrank} over~$\ZZ$ of the additive group \texttt{Additive (HypJac K f)}, and the
count is \texttt{CommGroup.card\_modPow}. For finitely generated~$J(K)$ this rank is
$\dim_{\QQ}(J(K) \otimes_{\ZZ} \QQ)$, the rank of Definition~\ref{def:rank-hyp} in the hypothesis of
the main theorems; that definition assumes no finite generation, and the comparison of the two
forms is not formalized.
\end{proof}

\section{Over a Dedekind domain}
\label{sec:ds-dedekind}

\begin{setting}
\label{set:ds-dedekind}
In Setting~\ref{set:ds-setting}, let~$K$ be a perfect field (Remark~\ref{rem:ds-perfect}), the
fraction field of a Dedekind domain~$R$, and let~$\Sigma$ be a set of nonzero primes of~$R$. For
a finite extension field~$F$ of~$K$ with integral closure~$B$ of~$R$ and $n \ge 1$, let
$F(\Sigma, n) \subseteq F^\times/(F^\times)^n$ be the group of classes whose valuation is
divisible by~$n$ at every prime of~$B$ not above~$\Sigma$ (the Selmer group of~$B$ relative to the
primes above~$\Sigma$). For each monic irreducible factor~$p$ of~$f$, let~$B_p$ be the integral
closure of~$R$ in $L_p = K[x]/(p)$.
\end{setting}

\begin{definition}[The group $A(\Sigma, 2)$]
\label{def:ds-selmer}
\lean{AdjoinRoot.selmerGroup, AdjoinRoot.finite_selmerGroup}
\leanok
Let $A(\Sigma, 2) \subseteq L^\square$ be the group of classes whose component in~$L_p^\square$
lies in~$L_p(\Sigma, 2)$ for every~$p$, under $L^\square \cong \prod_p L_p^\square$. If~$\Sigma$
is finite and each~$B_p$ has a finite class group and a finitely generated unit group (for
instance if~$K$ is a number field), then $A(\Sigma, 2)$ is finite.
\end{definition}

\begin{proof}
\leanok
Each $L_p(\Sigma, 2)$ is finite, by the exact sequence
$1 \to B_{p,\Sigma}^\times/2 \to L_p(\Sigma, 2) \to \Cl(B_{p,\Sigma})[2] \to 1$ for the
ring~$B_{p,\Sigma}$ of $\Sigma$-integers and the $S$-unit theorem.

\emph{Formalization.} \texttt{AdjoinRoot.selmerGroup} is defined for every exponent~$n$ and used
with $n = 2$. It is the pullback along \texttt{AdjoinRoot.modPowEquivPiFactors} of the product,
over \texttt{f.Factors}, of the Selmer groups \texttt{IsDedekindDomain.selmerGroupAbove} of the
integral closures of~$R$ in the factors \texttt{AdjoinRoot p}, relative to the primes above~$\Sigma$
(\texttt{IsDedekindDomain.HeightOneSpectrum.primesAbove}); \texttt{AdjoinRoot.mem\_selmerGroup\_iff}
and \texttt{AdjoinRoot.mk\_mem\_selmerGroup\_iff} are the componentwise conditions. The
separability of the factor fields is an instance family, which
\texttt{AdjoinRoot.instIsSeparableFactors} provides for a perfect~$K$.
\texttt{AdjoinRoot.finite\_selmerGroup} takes a finite~$\Sigma$ and instances for the finiteness of
the class groups and the finite generation of the unit groups of the integral closures. Mathlib's
\texttt{IsDedekindDomain.selmerGroup} is the case of the fraction field itself.
\end{proof}

\begin{remark}[Separability]
\label{rem:ds-perfect}
The integral closure of~$R$ in~$L_p$ is a Dedekind domain finite over~$R$ by Mathlib's results for
finite \emph{separable} extensions, and Lemma~\ref{lem:ds-norm-selmer} assumes the extension
separable, as Mathlib's relative norm \texttt{Algebra.intNorm} does; in the proof of
Lemma~\ref{lem:ds-point-mumford} it is applied to $F[x]/(a)$ for a finite extension~$F$ of~$K$ and
an irreducible~$a \in F[x]$. Both separability conditions hold if~$K$ is perfect (then so is~$F$),
in particular in characteristic~$0$, which covers number fields and their completions. Over an
imperfect~$K$, $f$ is assumed squarefree, not separable, and an irreducible~$a$ can be
inseparable. So~$K$ is assumed perfect in this section and the next.
\end{remark}

\begin{lemma}[Norms map Selmer groups to Selmer groups]
\label{lem:ds-norm-selmer}
\lean{IsDedekindDomain.norm_mem_selmerGroup}
\leanok
Let $R \subseteq B$ be Dedekind domains with fraction fields $K \subseteq F$, where~$F/K$ is finite
and separable and~$B$ is the integral closure of~$R$ in~$F$, let~$\Sigma$ be a set of nonzero
primes of~$R$, and $n \ge 1$. If $x \in F^\times$ has valuation divisible by~$n$ at every prime
of~$B$ not above~$\Sigma$, then~$N_{F/K}(x)$ has valuation divisible by~$n$ at every prime of~$R$
not in~$\Sigma$. So~$N_{F/K}$ maps $F(\Sigma, n)$ into $K(\Sigma, n)$.
\end{lemma}

\begin{proof}
\leanok
For a prime $v \notin \Sigma$, $v(N_{F/K}x) = \sum_{w \mid v} f(w \mid v)\,w(x)$, the sum over the
primes~$w$ of~$B$ above~$v$ with their residue degrees: for $x \in B$, the ideal norm maps $xB$ to
$N_{F/K}(x)R$ and each~$w$ to $v^{f(w \mid v)}$, and is multiplicative; in general write
$x = y/z$ with $y, z \in B$.

\emph{Formalization.} The map induced by the norm sends \texttt{selmerGroupAbove} into Mathlib's
\texttt{IsDedekindDomain.selmerGroup}. The valuation formula is
\texttt{IsDedekindDomain.toAdd\_valuationOfNeZero\_norm}, from \texttt{Ideal.multiplicity\_relNorm}
(the multiplicity of~$v$ in the relative norm of an ideal is $\sum_{w \mid v} f(w \mid v)$ times
that of~$w$), which rests on Mathlib's \texttt{Ideal.relNorm}, \texttt{Ideal.relNorm\_singleton}
and \texttt{Ideal.relNorm\_eq\_pow\_of\_isMaximal}. The lemma is missing from Mathlib.
\end{proof}

\begin{lemma}[The point lemma]
\label{lem:ds-point}
\lean{AdjoinRoot.of_sub_root_mem_selmerGroup,
  IsDedekindDomain.mem_selmerGroupAbove_of_sq_eq_aeval}
\leanok
\uses{def:ds-selmer}
Assume that~$\Sigma$ contains every prime~$v$ of~$R$ at which a coefficient of~$f$ is not
integral, $c$ is not a unit or $\operatorname{disc} f$ is not a unit. Let $x_0, y_0 \in K$ with
$y_0^2 = f(x_0) \ne 0$. Then $x_0 - \theta$ is a unit of~$L$, and its class lies in
$A(\Sigma, 2)$.
\end{lemma}

\begin{proof}
\leanok
For each~$p$, $p(x_0) \ne 0$, so the component $x_0 - \theta_p$ is not zero. Let~$w$ be a prime
of~$B_p$ above a prime $v \notin \Sigma$. As $c^{-1}f$ is monic with $v$-integral coefficients,
$w(\theta_p) \ge 0$. If $w(x_0) < 0$, then $w(x_0 - \theta_p) = w(x_0)$ and
$2w(y_0) = w(f(x_0)) = (2g + 1)w(x_0)$, as the leading term of~$f(x_0)$ dominates ($c$ is a
$v$-unit); so~$w(x_0)$ is even. If $w(x_0) \ge 0$, write
$f(x_0) = f(x_0) - f(\theta_p) = (x_0 - \theta_p)G$ with~$G$ integral at~$w$ and
$G \equiv f'(\theta_p) \bmod (x_0 - \theta_p)$. As $\operatorname{disc} f$ and~$c$ are $v$-units,
so is $N_{L/K}(f'(\theta))$, and~$f'(\theta_p)$ is a $w$-unit. So either $w(x_0 - \theta_p) = 0$,
or $w(G) = 0$ and $w(x_0 - \theta_p) = w(f(x_0)) = 2w(y_0)$; in both cases $w(x_0 - \theta_p)$ is
even.

\emph{Formalization.} \texttt{AdjoinRoot.of\_sub\_root\_mem\_selmerGroup} is stated for the class
of any unit with value $x_0 - \theta$ (\texttt{AdjoinRoot.isUnit\_of\_sub\_root} supplies one). The
condition on~$\Sigma$ is that~$f$ has separable reduction at the primes not in~$\Sigma$
(\texttt{Polynomial.HasSeparableReduction}: integral coefficients, a unit leading coefficient and a
unit discriminant, Mathlib's \texttt{Polynomial.discr}). The valuation core is
\texttt{Valuation.exists\_map\_sub\_eq\_sq}, for a polynomial of odd degree with a unit leading
coefficient: for non-integral~$x_0$ the leading term dominates; otherwise a Taylor expansion over
the integers (\texttt{Polynomial.exists\_mem\_eval\_eq\_add\_mul}) and the unit~$f'(\theta_p)$
(\texttt{Valuation.map\_derivative\_eval\_eq\_one}, from a B\'ezout identity for the resultant
of~$f$ and~$f'$ over the valuation ring, in place of the norm argument) give
\texttt{Valuation.map\_sub\_eq\_one\_or\_eq\_map\_eval}. The local form
\texttt{Polynomial.HasSeparableReduction.exists\_map\_sub\_eq\_sq} is applied over the extension
field, to the image of~$f$, which keeps separable reduction
(\texttt{Polynomial.HasSeparableReduction.map}); so~$x_0$ and~$y_0$ may lie in the extension field,
as the proof of Lemma~\ref{lem:ds-point-mumford} needs, and
\texttt{IsDedekindDomain.mem\_selmerGroupAbove\_of\_sq\_eq\_aeval} is the point lemma in a single
field.
\end{proof}

\begin{lemma}[The point lemma for Mumford polynomials]
\label{lem:ds-point-mumford}
\lean{IsDedekindDomain.mem_selmerGroupAbove_of_dvd_sq_sub,
  IsDedekindDomain.mem_selmerGroup_of_dvd_sq_sub}
\leanok
Let~$\Sigma$ be as in Lemma~\ref{lem:ds-point}, let~$F$ be a finite extension field of~$K$,
$\theta \in F$ a root of~$f$, and $a, b \in F[x]$ with $a \mid b^2 - f$ and~$a$ prime to~$f$. Then
$\operatorname{lc}(a)^{-1}a(\theta) \ne 0$, and its class lies in~$F(\Sigma, 2)$.
\end{lemma}

\begin{proof}
\leanok
\uses{lem:ds-point, lem:ds-norm-selmer, lem:ds-resultant}
As~$a$ is prime to~$f$ and $f(\theta) = 0$, $a(\theta) \ne 0$. If $a = a'a''$, then~$a'$
and~$a''$ divide $b^2 - f$ and are prime to~$f$, and $\operatorname{lc}(a)^{-1}a(\theta)$ is the
product of the values for~$a'$ and~$a''$; a constant~$a$ gives~$1$. So, by induction on the
irreducible factors, we may assume that~$a$ is monic and irreducible. Let $F' = F[x]/(a)$, a
finite extension field of~$F$, and let~$x_1 \in F'$ be the image of~$x$. Then
$b(x_1)^2 = f(x_1)$, and $f(x_1) \ne 0$ as~$a$ is prime to~$f$. Let~$w$ be a prime of the integral
closure of~$R$ in~$F'$ above a prime $v \notin \Sigma$. The proof of Lemma~\ref{lem:ds-point},
with~$F'$, $\theta$, $x_1$ and~$b(x_1)$ in place of~$L_p$, $\theta_p$, $x_0$ and~$y_0$, shows that
$w(x_1 - \theta)$ is even: it uses only that~$\theta$ is a root of~$f$ and that
$y_0^2 = f(x_0) \ne 0$. So the class of $\theta - x_1$ lies in~$F'(\Sigma, 2)$, and by
Lemma~\ref{lem:ds-norm-selmer}, for $F'/F$ and the primes of the integral closure of~$R$ in~$F$
above~$\Sigma$, its norm lies in~$F(\Sigma, 2)$. As the characteristic polynomial of~$x_1$ over~$F$
is~$a$, this norm is $N_{F'/F}(\theta - x_1) = a(\theta)$; equivalently, by
Lemma~\ref{lem:ds-resultant} for~$a$ over~$F$, it is $\operatorname{Res}(a, \theta - x) =
a(\theta)$.

\emph{Formalization.} The general form is
\texttt{IsDedekindDomain.mem\_selmerGroup\_of\_dvd\_sq\_sub}, over a Dedekind domain~$B$ with
perfect fraction field (in the application the integral closure of~$R$ in~$L_p$), for a polynomial
with separable reduction outside a set of primes of~$B$: the class of a unit~$u$ with
$u = \operatorname{lc}(a)^{-1}a(\theta)$ lies in Mathlib's \texttt{IsDedekindDomain.selmerGroup}.
The induction is \texttt{WfDvdMonoid.induction\_on\_irreducible}. The step for an irreducible~$a$,
monic or not, works in $F' = \texttt{AdjoinRoot a}$ with the integral closure of~$B$ in~$F'$: the
point lemma in a field and Lemma~\ref{lem:ds-norm-selmer} give
\texttt{IsDedekindDomain.norm\_sub\_mem\_selmerGroup\_of\_sq\_eq\_aeval}, and
\texttt{AdjoinRoot.leadingCoeff\_mul\_norm\_of\_sub\_root} computes
$N_{F'/F}(\theta - x_1) = \operatorname{lc}(a)^{-1}a(\theta)$. The lemma as stated, with~$f$
over~$K$ and~$\Sigma$, is \texttt{IsDedekindDomain.mem\_selmerGroupAbove\_of\_dvd\_sq\_sub}:
separable reduction moves up to~$B$ by \texttt{Polynomial.HasSeparableReduction.map}.
\end{proof}

\begin{theorem}[The image of~$\mu$]
\label{thm:ds-image-selmer}
\lean{ChabautyColeman.HypJac.range_descentMap_le_selmerGroup}
\leanok
\uses{def:ds-mu, def:ds-selmer}
If~$\Sigma$ is as in Lemma~\ref{lem:ds-point}, then $\mu(J(K)) \subseteq A(\Sigma, 2)$.
\end{theorem}

\begin{proof}
\leanok
\uses{lem:ds-move, lem:ds-point-mumford}
Let $P \in J(K)$. By Lemma~\ref{lem:ds-move}, $P = [a, b]$ with~$a$ prime to~$f$, and by
Definition~\ref{def:ds-mu}, $\mu(P)$ is the class of
$m(a) = (-c)^{\deg a}\operatorname{lc}(a)^{-1}a(\theta)$. Its component at~$p$ is
$(-c)^{\deg a}\operatorname{lc}(a)^{-1}a(\theta_p)$. By Lemma~\ref{lem:ds-point-mumford} for
$F = L_p$ and $\theta = \theta_p$ ($a$ stays prime to~$f$ in~$L_p[x]$), the class of
$\operatorname{lc}(a)^{-1}a(\theta_p)$ lies in~$L_p(\Sigma, 2)$, and so does the class
of~$(-c)^{\deg a}$, as~$c$ is a unit at every prime not in~$\Sigma$. So every component
of~$\mu(P)$ lies in~$L_p(\Sigma, 2)$, that is, $\mu(P) \in A(\Sigma, 2)$.

\emph{Formalization.} The proof works one component~$L_p$ at a time; only norms of field extensions
$L_p[x]/(a')$ of~$L_p$ occur, in Lemma~\ref{lem:ds-point-mumford}. It combines
\texttt{HypJac.exists\_isCoprime\_mumfordClass\_eq},
\texttt{HypJac.descentMap\_mumfordClass\_of\_isCoprime} and
\texttt{AdjoinRoot.mem\_selmerGroup\_iff} with \texttt{algHom\_mumfordValue\_mem\_selmerGroupAbove}:
for a $K$-algebra homomorphism $\varphi \colon L \to F$ (here the projection to~$L_p$) the class of
$\varphi(m(a))$ lies in~$F(\Sigma, 2)$, by Lemma~\ref{lem:ds-point-mumford} for the
root~$\varphi(\theta)$ of~$f$ and
\texttt{IsDedekindDomain.mem\_selmerGroupAbove\_of\_valuation\_eq\_one} for the constant. Each~$L_p$
is perfect by \texttt{Algebra.IsAlgebraic.perfectField}.
\end{proof}

\section{Number fields}
\label{sec:ds-number}

\begin{setting}
\label{set:ds-number}
In Setting~\ref{set:ds-setting}, let~$K$ be a number field and $R = O_K$. For a place~$v$ of~$K$,
let~$K_v$ be the completion, $L_v = K_v[x]/(f) = L \otimes_K K_v$,
$\mu_v \colon J(K_v) \to L_v^\square$ the map of Definition~\ref{def:ds-mu} over~$K_v$ ($f$ stays
squarefree, as~$K$ has characteristic~$0$), $r_v \colon J(K) \to J(K_v)$ the natural map and
$\operatorname{res}_v \colon L^\square \to L_v^\square$ the restriction. Let~$\Sigma_f$ be the
finite set of primes~$v$ of~$R$ at which a coefficient of~$f$ is not integral, $c$ is not a unit
or $\operatorname{disc} f$ is not a unit.
\end{setting}

\begin{definition}[The $2$-Selmer group]
\label{def:ds-sel2}
\lean{ChabautyColeman.HypJac.selmerGroupTwo, ChabautyColeman.HypJac.localCondition}
\leanok
\uses{def:ds-mu}
The $2$-Selmer group is
\[ \Sel_2 = \{\xi \in L^\square : N(\xi) = 1 \text{ and } \operatorname{res}_v(\xi) \in
   \mu_v(J(K_v)) \text{ for every place~$v$ of~$K$}\}\,. \]
\end{definition}

\begin{theorem}[Finiteness of the $2$-Selmer group]
\label{thm:ds-sel2}
\lean{ChabautyColeman.HypJac.range_descentMap_le_selmerGroupTwo,
  ChabautyColeman.HypJac.selmerGroupTwo_le_selmerGroup,
  ChabautyColeman.HypJac.selmerGroupTwo_le_ker_norm,
  ChabautyColeman.HypJac.finite_selmerGroupTwo}
\leanok
\uses{def:ds-sel2, def:ds-selmer}
\begin{enumerate}
  \item $\mu(J(K)) \subseteq \Sel_2$.
  \item $\Sel_2 \subseteq A(\Sigma_f, 2) \cap \ker N$.
  \item $\Sel_2$ is finite.
\end{enumerate}
\end{theorem}

\begin{proof}
\leanok
\uses{thm:ds-norm, thm:ds-naturality, thm:ds-image-selmer}
(1) Theorem~\ref{thm:ds-norm}, and $\operatorname{res}_v(\mu(P)) = \mu_v(r_v(P))$ by
Theorem~\ref{thm:ds-naturality}. (2) For a prime $v \notin \Sigma_f$,
Theorem~\ref{thm:ds-image-selmer} over~$K_v$, with $R = O_v$ and $\Sigma = \emptyset$, shows
that $\mu_v(J(K_v))$ consists of unramified classes; and a class in~$L^\square$ lies in
$A(\Sigma_f, 2)$ if and only if its restriction is unramified at every prime $v \notin \Sigma_f$.
(3) follows from~(2) and Definition~\ref{def:ds-selmer}.

\emph{Formalization.} \texttt{HypJac.selmerGroupTwo} is defined for a Dedekind domain~$R$ whose
fraction field~$K$ has characteristic~$0$ and a family of field extensions of~$K$ (for a number
field, the completions at the infinite places): it is the intersection of $\ker N$, of the local
conditions at the completions \texttt{v.adicCompletion K} at the primes~$v$ of~$R$, and of the
local conditions at the members of the family. The local condition \texttt{HypJac.localCondition}
at a field extension~$F$ of~$K$ is the preimage of the image of~$\mu_F$ under the restriction
to~$F$. (2) is stated for every set~$\Sigma$ of primes outside which~$f$ has separable reduction
(\texttt{Polynomial.HasSeparableReduction}), and of the equivalence in its proof only the
implication from right to left is used:
\texttt{AdjoinRoot.mem\_selmerGroup\_of\_forall\_adicCompletion}, for every squarefree~$f$ and every
exponent. There~$O_v$ is the valuation ring of the completion, a discrete valuation ring, and~$f$
has separable reduction at its prime because it has at~$v$. In~(3), $\Sigma_f$ is finite by
\texttt{Polynomial.finite\_setOf\_not\_hasSeparableReduction}, as $\operatorname{disc} f \ne 0$, and
\texttt{HypJac.finite\_selmerGroupTwo} takes the instances on class groups and unit groups of
\texttt{AdjoinRoot.finite\_selmerGroup}.
\end{proof}

\begin{definition}[Cited structures]
\label{def:ds-cited}
\lean{ChabautyColeman.HypJac.MordellWeil, ChabautyColeman.HypJac.LocalStructure}
\leanok
\begin{description}
  \item[\textmd{(MW)}] For a number field~$K$, the group~$J(K)$ is finitely generated (the
    Mordell--Weil theorem: A.~Weil, \emph{L'arithm\'etique sur les courbes alg\'ebriques}, Acta
    Math.~52 (1929), 281--315; M.~Hindry and J.~H.~Silverman, \emph{Diophantine geometry. An
    introduction}, Springer, 2000, Theorem~C.0.1).
  \item[\textmd{(Loc)}] For a finite extension~$K_v$ of~$\QQ_\ell$ with ring of
    integers~$O_v$, there is an isomorphism of groups $J(K_v) \cong O_v^g \times T$ with~$T$
    finite (A.~Mattuck, \emph{Abelian varieties over $p$-adic ground fields}, Ann.\ of Math.~(2)
    62 (1955), 92--119; for $K_v = \QQ_2$ also B.~Poonen and M.~Stoll, \emph{Most odd degree
    hyperelliptic curves have only one rational point}, Ann.\ of Math.~(2) 180 (2014),
    1137--1166, Section~4).
\end{description}
\end{definition}

Mattuck's isomorphism is one of topological groups. The statements below use only the group
isomorphism, and no topology on~$J(K_v)$ is involved.

\begin{remark}[Formalization]
The two results are cited, not proved: each is a class or structure whose docstring gives the
references. \texttt{HypJac.MordellWeil K f}, for a number field~$K$, is a class with the single
field that \texttt{HypJac K f} is finitely generated (\texttt{IsMulFG}), which it provides as an
instance. \texttt{HypJac.LocalStructure p F f g}, for a field~$F$ that is an algebra over the
field~$\QQ_p$ of Mathlib, is a structure with values in \texttt{Prop}: $\deg f = 2g + 1$, and there
is a finite group~$T$ with a group isomorphism from \texttt{HypJac F f} to the product of
$(\ZZ_p^d)^g$ and~$T$, where $d = [F : \QQ_p]$. As a group, $O_v \cong \ZZ_p^d$; this form
of~(Loc) avoids the integral closure of~$\ZZ_p$ in~$F$.
\end{remark}

\begin{lemma}[Consequences of (Loc)]
\label{lem:ds-loc}
\lean{ChabautyColeman.HypJac.LocalStructure.card_modPow_two,
  ChabautyColeman.HypJac.LocalStructure.card_modPow_two_of_ne_two,
  ChabautyColeman.HypJac.LocalStructure.mem_iInf_range_powMonoidHom_pow_iff}
\leanok
\uses{def:ds-cited, def:hypjac}
Assume~(Loc) for a finite extension~$K_v$ of~$\QQ_\ell$.
\begin{enumerate}
  \item $\#J(K_v)/2J(K_v) = 2^{g[K_v : \QQ_2]}\,\#J(K_v)[2]$ if $\ell = 2$, and
    $\#J(K_v)/2J(K_v) = \#J(K_v)[2]$ if~$\ell$ is odd.
  \item $\bigcap_{n \ge 0} \ell^nJ(K_v)$ is the subgroup of elements of finite order prime
    to~$\ell$.
\end{enumerate}
\end{lemma}

\begin{proof}
\leanok
Let $J(K_v) \cong O_v^g \times T$. As a group, $O_v \cong \ZZ_\ell^{[K_v : \QQ_\ell]}$ is
torsion-free, so $J(K_v)[2] \cong T[2]$, and $\#T/2T = \#T[2]$ as~$T$ is finite.
(1) $\#O_v/2O_v$ is $2^{[K_v : \QQ_2]}$ if $\ell = 2$ and~$1$ if~$\ell$ is odd, as~$2$ is then a
unit. (2) In every abelian group, an element of finite order~$m$ prime to~$\ell$ is an
$\ell^n$-th multiple for every~$n$: if $k\ell^n \equiv 1 \bmod m$, then $x = \ell^n(kx)$.
Conversely, $\bigcap_n \ell^nO_v = 0$, and the subgroups $\ell^nT$ of the finite group~$T$
decrease and become constant, equal to the subgroup~$T_{\ell'}$ of elements of order prime
to~$\ell$ (on the $\ell$-primary part of~$T$, multiplication by~$\ell^n$ is zero for large~$n$, and
on~$T_{\ell'}$ it is bijective); as $\ell^n(O_v^g \times T) = \ell^nO_v^g \times \ell^nT$, the
intersection is $0 \times T_{\ell'}$.

\emph{Formalization.} With the prime~$\ell$ named~$p$, (1) is
\texttt{HypJac.LocalStructure.card\_modPow\_two} for $p = 2$ and
\texttt{HypJac.LocalStructure.card\_modPow\_two\_of\_ne\_two} for $p \ne 2$, about the cardinality
of the quotient by the range of \texttt{powMonoidHom 2} and of its kernel; both go through the
index of $2\ZZ_p$ in~$\ZZ_p$, which is~$2$ for $p = 2$ (\texttt{PadicInt.index\_range\_nsmul\_self})
and~$1$ otherwise (\texttt{PadicInt.range\_nsmul\_eq\_top\_of\_not\_dvd}), and
\texttt{AddSubgroup.index\_range\_nsmul\_prod\_of\_injective}. (2) says that an element lies in the
intersection of the ranges of the maps \texttt{powMonoidHom (p\textasciicircum{}n)} if and only if
its order is prime to~$p$: the order~$0$ of an element of infinite order is not prime to a prime.
The implication from right to left holds in every commutative group
(\texttt{CommGroup.mem\_range\_powMonoidHom\_of\_coprime}); for the other, Krull's intersection
theorem gives $\bigcap_n p^n M = 0$ for every finitely generated $\ZZ_p$-module~$M$
(\texttt{PadicInt.eq\_zero\_of\_forall\_mem\_range\_nsmul\_pow}), and
\texttt{CommGroup.forall\_mem\_range\_powMonoidHom\_pow\_prod\_iff} treats the product with a finite
group.
\end{proof}

\begin{lemma}[The case $\ell = 2$]
\label{lem:ds-loc-two}
\lean{ChabautyColeman.HypJac.LocalStructure.mem_iInf_range_powMonoidHom_two_pow_iff}
\leanok
\uses{lem:ds-loc, def:hypjac}
Assume~(Loc) for a finite extension~$K_v$ of~$\QQ_2$. Then $\bigcap_{n \ge 0} 2^nJ(K_v)$ is the
subgroup of elements of odd order.
\end{lemma}

\begin{proof}
\leanok
Lemma~\ref{lem:ds-loc}\,(2) with $\ell = 2$.
\end{proof}

For odd~$\ell$, by contrast, $2$ is a unit of~$O_v$, and $\bigcap_n 2^nJ(K_v)$ contains a copy
of~$O_v^g$, so it is not torsion. The proof of Theorem~\ref{thm:ds-e1} uses
Lemma~\ref{lem:ds-loc-two} for $K_v = \QQ_2$.

\begin{definition}[Relaxed Selmer groups]
\label{def:ds-relaxed}
\lean{ChabautyColeman.relaxedSelmerGroup, ChabautyColeman.range_descentMap_le_relaxedSelmerGroup,
  ChabautyColeman.finite_relaxedSelmerGroup}
\leanok
\uses{def:ds-selmer, def:ds-mu, def:ds-sel2}
For a finite set $\Sigma \supseteq \Sigma_f$ of primes and a set~$V$ of places of~$K$, let
\[ S(\Sigma, V) = \{\xi \in A(\Sigma, 2) : N(\xi) = 1 \text{ and } \operatorname{res}_v(\xi)
   \in \mu_v(J(K_v)) \text{ for } v \in V\}\,, \]
with the local conditions of Definition~\ref{def:ds-sel2} at the places in~$V$ only. Then
$S(\Sigma, V)$ is finite, and $\mu(J(K)) \subseteq S(\Sigma, V)$.
\end{definition}

\begin{proof}
\leanok
\uses{thm:ds-image-selmer, thm:ds-norm, thm:ds-naturality}
Definition~\ref{def:ds-selmer} for the finiteness, and Theorems~\ref{thm:ds-image-selmer},
\ref{thm:ds-norm} and~\ref{thm:ds-naturality} for the inclusion.

\emph{Formalization.} \texttt{relaxedSelmerGroup} takes the places of~$V$ as a family of field
extensions of~$K$ (their completions), with the local conditions \texttt{HypJac.localCondition}
of \texttt{HypJac.selmerGroupTwo}; as these are taken at arbitrary field extensions, for $K = \QQ$
a condition at~$2$ can use Mathlib's field~$\QQ_2$ itself. The inclusion takes the hypothesis
that~$f$ has separable reduction outside~$\Sigma$ (\texttt{Polynomial.HasSeparableReduction}), as
Theorem~\ref{thm:ds-image-selmer} does; \texttt{finite\_relaxedSelmerGroup} is the finiteness for
finite~$\Sigma$, with the instances on class groups and unit groups of
\texttt{AdjoinRoot.finite\_selmerGroup}.
\end{proof}

For $K = \QQ$, $2 \in \Sigma$ and $V = (\Sigma \cup \{\infty\}) \setminus \{2\}$, $S(\Sigma, V)$ is
the group~$S$ of Remark~4.2 of [Sto17], which leaves out the $2$-adic condition, with~$\Sigma$ in
place of the set of bad primes for $2$-descent used there; $S(\Sigma, \emptyset) = A(\Sigma, 2) \cap
\ker N$ is coarser, and $\Sel_2 = S(\Sigma, V)$ when~$V$ is the set of all places, by
Theorem~\ref{thm:ds-sel2}\,(2). Each of them contains~$\Sel_2$, and so it can replace the
$2$-Selmer group in the algorithm of Section~4 of [Sto17], by Remark~4.2 there.

\begin{remark}[Local conditions are needed]
\label{rem:ds-dimension}
Let $K = \QQ$ and~$f$ irreducible, and let $(r_1, r_2)$ be the signature of the field~$L$, so that
$r_1 + 2r_2 = 2g + 1$. The unramified classes include $O_L^\times/(O_L^\times)^2$, of dimension
$r_1 + r_2$ (the roots of unity of~$L$ are~$\pm 1$, as $r_1 \ge 1$), on which the norm is onto
$\{\pm 1\}$, as $N(-1) = -1$; so the unramified classes of norm~$1$ have dimension at least
$r_1 + r_2 - 1$, whatever the class group of~$L$. At the real place,
$J(\mathbb R) \cong (\mathbb R/\ZZ)^g \times (\ZZ/2)^k$ and
$\#J(\mathbb R)[2] = 2^{r_1 + r_2 - 1}$ (Theorem~\ref{thm:ds-two-torsion} over~$\mathbb R$), so
$k = r_1 + r_2 - 1 - g = (r_1 - 1)/2$ is the dimension of $\mu_\infty(J(\mathbb R))$
(Theorem~\ref{thm:ds-kernel} over~$\mathbb R$), inside the space of dimension $r_1 - 1$ of sign
vectors of norm~$1$. So the classes that are unramified everywhere, of norm~$1$ and satisfy the
real condition form a space of dimension at least
$r_1 + r_2 - 1 - (r_1 - 1)/2 = (r_1 - 1)/2 + r_2 = g$, and so does every group defined by these
conditions with ramification allowed at some primes, such as $S(\Sigma, \{\infty\})$, which
contains that space. As $J(\QQ)[2] = 0$, the rank bound of Theorem~\ref{thm:ds-rank-bound} from
such a group is never below~$g$, and every example of rank less than~$g$ with irreducible~$f$
needs a local condition at~$2$ or at an odd bad prime, hence the local image size of
Lemma~\ref{lem:ds-loc}\,(1) and~(Loc).
\end{remark}

\begin{theorem}[The rank bound]
\label{thm:ds-rank-bound}
\lean{ChabautyColeman.HypJac.pow_rank_le_card_selmerGroupTwo,
  ChabautyColeman.pow_rank_le_card_relaxedSelmerGroup}
\leanok
\uses{def:ds-cited, def:ds-sel2, def:ds-relaxed}
Assume~(MW). Then $2^{\operatorname{rk} J(K)}\,\#J(K)[2] \le \#S$ for every finite subgroup
$S \subseteq L^\square$ containing~$\mu(J(K))$, in particular for $S = \Sel_2$ and for the
groups~$S(\Sigma, V)$ of Definition~\ref{def:ds-relaxed}.
\end{theorem}

\begin{proof}
\leanok
\uses{thm:ds-rank, thm:ds-sel2, def:ds-relaxed}
Theorem~\ref{thm:ds-rank}, with Theorem~\ref{thm:ds-sel2} and Definition~\ref{def:ds-relaxed} for
the inclusions and the finiteness.

\emph{Formalization.} Both statements take the instance \texttt{IsMulFG (HypJac K f)}, which the
class \texttt{HypJac.MordellWeil} provides over a number field, and the instances of the
finiteness statements; they are \texttt{HypJac.pow\_rank\_le\_card\_of\_range\_descentMap\_le} of
Theorem~\ref{thm:ds-rank} applied to the inclusions.
\end{proof}

\section{Examples}
\label{sec:ds-examples}

\begin{remark}[E1: genus~$2$, $f$ split]
\label{rem:ds-e1}
The first example is the curve
\[ C \colon y^2 = f(x) = x(x^2 - 1)(x^2 - 9) = x^5 - 10x^3 + 9x\,, \]
with the roots $\alpha = -3, -1, 0, 1, 3$ and $\operatorname{disc} f = 2^{16} \cdot 3^6$, so
$\Sigma_f = \{2, 3\}$. No smaller set is possible: among five integers, two are congruent
modulo~$2$ and two modulo~$3$. Here $L = \QQ^5$, and $A(\Sigma_f, 2) \cap \ker N$ has dimension
$4(\#\Sigma_f + 1) = 12$ over~$\FF_2$, whereas $\mu(J(\QQ)[2])$ has dimension~$4$. A local
condition at one place cuts out at most $4$ dimensions (an odd prime) or~$6$ (the prime~$2$), so
no single place suffices; the set $V = \{2, 3\}$ does: $S(\Sigma_f, V) = \mu(J(\QQ)[2])$
(Proposition~\ref{prop:ds-e1-sel}). Theorem~\ref{thm:ds-e1} shows $J(\QQ) = J(\QQ)[2] \cong
(\ZZ/2)^4$ and $C(\QQ) = \{\infty, (0, 0), (\pm 1, 0), (\pm 3, 0)\}$. Its proof needs neither~(MW)
nor any $2$-adic analysis on~$C$: the descent and~(Loc) for~$\QQ_2$ and~$\QQ_3$ give
$J(\QQ) = J(\QQ)[2] + T$ with~$T$ of odd order, and reduction at the good prime~$5$ kills~$T$.
\end{remark}

\begin{setting}
\label{set:ds-e1}
\lean{ChabautyColeman.DescentExample1.f, ChabautyColeman.DescentExample1.F,
  ChabautyColeman.DescentExample1.S}
\leanok
Let $f = x^5 - 10x^3 + 9x \in \QQ[x]$, with the roots
$\alpha_1, \dots, \alpha_5 = -3, -1, 0, 1, 3$, and $L = \QQ[T]/(f) \cong \QQ^5$,
$\xi \mapsto (\xi(\alpha_i))_i$. Let $\Sigma = \{2, 3\}$, let~$V$ consist of~$\QQ_2$ and~$\QQ_3$,
and $S = S(\Sigma, V)$ (Definition~\ref{def:ds-relaxed}). Let $\bar f = x^5 - x \in \FF_5[x]$ be
the reduction of~$f$ modulo~$5$.
\end{setting}

\begin{remark}[Formalization]
In the module \texttt{HyperellipticCurves.Descent.Example1}: \texttt{f} over~$\QQ$ and the
integral model \texttt{F} for the reduction modulo~$5$ (\texttt{DescentExample1.map\_F}), the
roots as a vector of integers, and~$\Sigma$ as~\texttt{S}, the primes $(2)$, $(3)$ of~$\ZZ$
(\texttt{Rat.HeightOneSpectrum.primesEquiv}). The family~$V$ consists of the $p$-adic fields of
Mathlib for $p = 2, 3$. $f$ has separable reduction outside~$\Sigma$
(\texttt{DescentExample1.hasSeparableReduction\_of\_notMem}), as
$\operatorname{disc} F = 2^{16} \cdot 3^6$ (\texttt{DescentExample1.discr\_F}, from Mathlib's
\texttt{Polynomial.resultant\_deriv}: the product of the values $144, -16, 9, -16, 144$ of~$F'$ at
the roots).
\end{remark}

\begin{lemma}[Square classes of local fields]
\label{lem:ds-padic-squares}
\lean{HenselianLocalRing.isSquare_iff_isSquare_residue, IsDiscreteValuationRing.modPowTwoEquiv,
  PadicInt.isSquare_iff_toZModPow_three_eq_one, Padic.modPowTwoEquivZModEight,
  Padic.modPowTwoEquiv_mk_of_eq_pow_mul}
\leanok
\uses{thm:hensel}
\begin{enumerate}
  \item Let~$O$ be a Henselian local ring in which~$2$ is a unit, with residue field~$k$. A unit
    of~$O$ is a square if and only if its image in~$k$ is a square.
  \item Let~$O$ moreover be a discrete valuation ring, with fraction field~$F$ and a
    uniformizer~$\pi$. Then $\pi^n u \mapsto (n \bmod 2, \bar u)$, for $u \in O^\times$, induces
    an isomorphism $F^\times/(F^\times)^2 \cong \ZZ/2 \times k^\times/(k^\times)^2$; for
    finite~$k$, both have order~$4$. This covers the finite extensions of~$\QQ_p$ for odd~$p$,
    among them~$\QQ_3$.
  \item A unit~$u$ of~$\ZZ_2$ is a square if and only if $u \equiv 1 \bmod 8$, and
    $2^n u \mapsto (n \bmod 2, u \bmod 8)$ induces an isomorphism
    $\QQ_2^\times/(\QQ_2^\times)^2 \cong \ZZ/2 \times (\ZZ/8)^\times$, of order~$8$.
\end{enumerate}
\end{lemma}

\begin{proof}
\leanok
(1) The image of a square is a square. Conversely, if $\bar u = \bar a^2$, Hensel's lemma applies
to $X^2 - u$ at~$a$, whose derivative~$2a$ is a unit. (2) $F^\times = \pi^{\ZZ} \times O^\times$,
and $\pi^n u$ is a square if and only if~$n$ is even and~$u$ is a square, which is~(1). (3) The
square of an odd integer is~$\equiv 1 \bmod 8$. Conversely, Hensel's lemma in the form for
$\lvert g(a)\rvert < \lvert g'(a)\rvert^2$ applies to $g = X^2 - u$ at $a = 1$, as
$\lvert 1 - u\rvert_2 \le 2^{-3} < \lvert 2\rvert_2^2$. The isomorphism follows as in~(2).

\emph{Formalization.} (1) is proved from the defining property of \texttt{HenselianLocalRing}; a
local ring that is adically complete at its maximal ideal, such as~$\ZZ_p$, is Henselian
(\texttt{IsAdicComplete.henselianLocalRing}). (2) comes from
$F^\times/(F^\times)^n \cong \ZZ/n \times O^\times/(O^\times)^n$ for every~$n$
(\texttt{IsDiscreteValuationRing.modPowEquiv}, from \texttt{IsDiscreteValuationRing.unitsEquivProd})
and $O^\times/(O^\times)^2 \cong k^\times/(k^\times)^2$ by~(1)
(\texttt{HenselianLocalRing.modPowTwoEquiv}), stated with the classes \texttt{Units.ModPow}; the
order~$4$ for finite~$k$ is not formalized, the examples compute with the isomorphism itself.
For~$\QQ_p$, $p$ odd, the isomorphism is \texttt{Padic.modPowTwoEquiv}, with the uniformizer~$p$
and the residue field $\ZZ/p$ (\texttt{PadicInt.residueField}); its value at~$p^n u$ for a unit~$u$
of~$\ZZ_p$ is \texttt{Padic.modPowTwoEquiv\_mk\_of\_eq\_pow\_mul}. (3) comes from Mathlib's
\texttt{hensels\_lemma}; the values of the isomorphisms on classes $p^n m$ with $m \in \ZZ$ are
\texttt{Padic.modPowTwoEquivZModEight\_mk\_of\_val\_eq} and
\texttt{Padic.modPowTwoEquiv\_mk\_of\_val\_eq}. Mathlib has none of these statements.
\end{proof}

\begin{lemma}[$A(\Sigma, 2)$ for split~$f$ over~$\QQ$]
\label{lem:ds-e1-selmer}
\lean{Rat.selmerGroup_eq_closure, AdjoinRoot.mem_selmerGroup_iff_of_splits}
\leanok
\uses{def:ds-selmer}
\begin{enumerate}
  \item For a finite set~$\Sigma$ of primes, $\QQ(\Sigma, n)$ is generated by the classes of~$-1$
    and of the primes in~$\Sigma$.
  \item If $f \in \QQ[x]$ is the product of distinct factors $x - \alpha_i$, a class
    $\xi \in L^\times/(L^\times)^n$ lies in $A(\Sigma, n)$ if and only if every
    component~$\xi(\alpha_i)$ lies in~$\QQ(\Sigma, n)$.
\end{enumerate}
\end{lemma}

\begin{proof}
\leanok
\uses{lem:ds-norm-selmer}
(1) Unique factorization in~$\ZZ$, with $\ZZ^\times = \{\pm 1\}$. (2) The field factors of~$L$ are
the $\QQ[x]/(x - \alpha_i) \cong \QQ$, with integral closure~$\ZZ$, and a class of~$\QQ$ lies
in~$\QQ(\Sigma, n)$ if and only if its image in such a factor lies in its Selmer group: one
direction is restriction, the other the norm (Lemma~\ref{lem:ds-norm-selmer}), here of degree~$1$.

\emph{Formalization.} (1) is stated for a finite set of primes of~$\ZZ$, with Mathlib's
\texttt{Rat.HeightOneSpectrum.natGenerator} (the prime number of a prime of~$\ZZ$), as an equality
of subgroups; the proof adjoins the primes one at a time
(\texttt{IsDedekindDomain.selmerGroup\_insert}) to $\ZZ^\times/(\ZZ^\times)^n$, which is all
of~$\QQ(\emptyset, n)$ as~$\ZZ$ has trivial class group. (2) is stated for a field~$K$ with a
Dedekind domain~$R$ (the field factors of $K[x]/(f)$ are copies of~$K$, the integral closure
is~$R$, \texttt{IsDedekindDomain.map\_algebraMap\_mem\_selmerGroupAbove\_iff} for degree~$1$), and
$f$ with at least $\deg f$ distinct roots. The components~$\xi(\alpha_i)$ are those of the
isomorphism \texttt{AdjoinRoot.modPowEquivPiOfRoots} from $L^\times/(L^\times)^n$ to the product
of copies of $K^\times/(K^\times)^n$, given by evaluation at the roots
(\texttt{AdjoinRoot.equivPiOfRoots}, Lagrange interpolation); it commutes with base change
(\texttt{AdjoinRoot.modPowEquivPiOfRoots\_map}).
\end{proof}

\begin{lemma}[The local images of E1]
\label{lem:ds-e1-local}
\lean{ChabautyColeman.DescentExample1.range_descentMap_three,
  ChabautyColeman.DescentExample1.range_descentMap_two}
\leanok
\uses{lem:ds-loc, thm:ds-two-torsion, thm:ds-kernel, thm:ds-formula,
  cor:ds-torsion-values, lem:ds-padic-squares}
Assume~(Loc) for~$\QQ_2$ and~$\QQ_3$. Then $\mu_3(J(\QQ_3))$ is the group of order~$16$ generated
by the images of the five points $(\alpha_j, 0)$, and $\mu_2(J(\QQ_2))$ is the group of order~$64$
generated by these and the images of the points of~$C(\QQ_2)$ with $x = 4$ and $x = 1/5$.
\end{lemma}

\begin{proof}
\leanok
The group $\mu_v(J(\QQ_v)) \cong J(\QQ_v)/2J(\QQ_v)$ (Theorem~\ref{thm:ds-kernel} over~$\QQ_v$) has
order $\#J(\QQ_v)[2] = 16$ for $v = 3$ and $2^2 \cdot 16 = 64$ for $v = 2$, by
Lemma~\ref{lem:ds-loc}\,(1) and Theorem~\ref{thm:ds-two-torsion}, as~$f$ has five linear factors
over~$\QQ_v$. The listed classes lie in it, the classes of $4 - \theta$ and $1/5 - \theta$ by
Theorem~\ref{thm:ds-formula}: the two points over~$\QQ_2$ exist by
Lemma~\ref{lem:ds-padic-squares}\,(3), as $f(4) = 2^2 \cdot 105$ and
$f(1/5) = 2^8 \cdot 21/5^5$ with $105 \equiv 21 \cdot 5^5 \equiv 1 \bmod 8$. The isomorphisms of
Lemma~\ref{lem:ds-padic-squares}\,(2) for~$\QQ_3$ and~(3) for~$\QQ_2$, applied to the five
components of the values (Corollary~\ref{cor:ds-torsion-values} and
Theorem~\ref{thm:ds-formula}), separate~$16$, resp.~$64$, products of the listed classes.

\emph{Formalization.} The classes of $4 - \theta$ and $1/5 - \theta$ enter as units of
$\QQ_2[T]/(f)$ with these values; they lie in the image by
\texttt{HypJac.descentMap\_mumfordClass\_X\_sub\_C}, the value at a point $(x_0, y_0)$ with
$f(x_0) \ne 0$. The components of the values $\mu((\alpha_j, 0))$ at the roots~$\alpha_i$ are
$\operatorname{lc}(f)(\alpha_j - \alpha_i)$ for $i \ne j$ and $\operatorname{lc}(f) f'(\alpha_j)$
for $i = j$ (Corollary~\ref{cor:ds-torsion-values}), written as products of~$-1$, $2$, $3$ with a
table of exponents that \texttt{decide} checks against the integers. The order of the image is
\texttt{HypJac.card\_range\_descentMap\_eq\_card\_quotient} with Lemma~\ref{lem:ds-loc}; the lower
bound for the subgroup generated by the listed classes is
\texttt{Subgroup.card\_closure\_range\_eq\_two\_pow} (elements of order dividing~$2$ of which no
nonempty product is trivial generate a group of order~$2^k$), applied to the images in root
coordinates, where \texttt{decide} checks the products. Four of the five points $(\alpha_j, 0)$
suffice: the five values have product~$1$.
\end{proof}

\begin{proposition}[The relaxed Selmer group of~E1]
\label{prop:ds-e1-sel}
\lean{ChabautyColeman.DescentExample1.relaxedSelmerGroup_eq}
\leanok
\uses{def:ds-relaxed, lem:ds-e1-selmer, lem:ds-e1-local}
Assume~(Loc) for~$\QQ_2$ and~$\QQ_3$. Then $S(\Sigma, V) = \mu(J(\QQ)[2])$.
\end{proposition}

\begin{proof}
\leanok
\uses{cor:ds-torsion-values, thm:ds-naturality, lem:ds-padic-squares}
The inclusion~$\supseteq$ is Definition~\ref{def:ds-relaxed}, as~$f$ has separable reduction
outside~$\Sigma$. Let $\xi \in S(\Sigma, V)$. By Lemma~\ref{lem:ds-e1-selmer}, the components of
the elements of~$A(\Sigma, 2)$ lie in the group generated by the classes of~$-1$, $2$ and~$3$,
which stay independent in~$\QQ_2^\times/(\QQ_2^\times)^2$ (Lemma~\ref{lem:ds-padic-squares}\,(3));
so the map from~$A(\Sigma, 2)$ to $L_2^\times/(L_2^\times)^2$, $L_2 = \QQ_2[T]/(f)$, is injective.
The classes~$P$ and~$Q$ of~$A(\Sigma, 2)$ with components $(-1, -3, 1, 3, 1)$ and
$(-3, -2, -3, 3, -6)$ map to the classes of $4 - \theta$ and $1/5 - \theta$, as the components
agree in~$\QQ_2^\times/(\QQ_2^\times)^2$. By Lemma~\ref{lem:ds-e1-local} for~$\QQ_2$, as~$\mu$
commutes with the map to~$\QQ_2$ (Theorem~\ref{thm:ds-naturality}), $\xi$ therefore lies in the
group generated by~$P$, $Q$ and the values $\mu((\alpha_j, 0))$
(Corollary~\ref{cor:ds-torsion-values}), which lie in~$A(\Sigma, 2)$: $\xi = wz$ with
$w \in \mu(J(\QQ)[2])$ and $z \in \{1, P, Q, PQ\}$. Both~$\xi$ and~$w$ satisfy the local condition
at~$\QQ_3$, hence so does~$z$. By Lemma~\ref{lem:ds-e1-local} for~$\QQ_3$, the image of~$z$ is
then a product of the images of the values $\mu((\alpha_j, 0))$, which the isomorphism of
Lemma~\ref{lem:ds-padic-squares}\,(2), applied to the five components, excludes for $z \ne 1$.
So $\xi = w$.

\emph{Formalization.} The classes of~$A(\Sigma, 2)$ and their images over~$\QQ_2$ and~$\QQ_3$ are
compared in root coordinates (\texttt{AdjoinRoot.modPowEquivPiOfRoots} and its naturality
\texttt{AdjoinRoot.modPowEquivPiOfRoots\_map}), every component a product of the classes of~$-1$,
$2$, $3$ given by a table of exponents, and the last step is a check by \texttt{decide}. The
elements of a subgroup generated by elements of order dividing~$2$ are the products over the finite
sets of generators (\texttt{Subgroup.mem\_closure\_range\_iff\_of\_sq\_eq\_one}). The norm
condition is not used.
\end{proof}

\begin{definition}[A cited input: reduction]
\label{def:ds-red}
\lean{ChabautyColeman.HypJac.TorsionReduction}
\leanok
\begin{description}
  \item[\textmd{(Red)}] Let $F \in \ZZ[x]$ be monic of odd degree and squarefree, and~$p \ge 3$
    a prime not dividing $\operatorname{disc} F$. Then reduction modulo~$p$ maps the torsion
    subgroup of~$J(\QQ)$ injectively into~$\tilde J(\FF_p)$, the Jacobian of $y^2 = \bar F(x)$
    (N.~M.~Katz, \emph{Galois properties of torsion points on abelian varieties}, Invent.\
    Math.~62 (1981), 481--502, Appendix: the kernel of reduction is torsion-free, as
    $e = 1 < p - 1$; for torsion of order prime to~$p$ also M.~Hindry and J.~H.~Silverman,
    \emph{loc.\ cit.}, Theorem~C.1.4).
\end{description}
\end{definition}

\begin{remark}[Formalization]
\texttt{HypJac.TorsionReduction p F}, for $F \in \ZZ[x]$, is a structure with values in
\texttt{Prop}, like~(Loc): there is a homomorphism from the Jacobian over~$\QQ$ of the image
of~$F$ to the Jacobian over~$\ZZ/p$ of its reduction that is injective on the elements of finite
order, in the form that an element of finite order in its kernel is trivial. The reduction map
is not defined; the example uses only such a homomorphism, through
\texttt{HypJac.TorsionReduction.pow\_eq\_one\_of\_isOfFinOrder}: an exponent of the group over
$\ZZ/p$ kills every element of finite order of~$J(\QQ)$.
\end{remark}

\begin{lemma}[$\tilde J(\FF_5)$]
\label{lem:ds-e1-five}
\lean{ChabautyColeman.DescentExample1.exponent_hypJac_zmod_five,
  ChabautyColeman.DescentExample1.sq_eq_one_of_eq_X_pow_five_sub_X}
\leanok
\uses{def:hypjac}
The group $\tilde J(\FF_5)$ of $y^2 = \bar f(x) = x^5 - x$ has exponent~$2$; so it is
$\tilde J(\FF_5)[2]$, of order~$16$.
\end{lemma}

\begin{proof}
\leanok
\uses{thm:ds-two-torsion, lem:mumford-reduce}
Every class has a reduced Mumford representative $(u, v)$: $u$ monic, $\deg v < \deg u \le 2$ and
$u \mid v^2 - \bar f$ (Lemma~\ref{lem:mumford-reduce}). For $u = x - a$, $v^2 = \bar f(a) = 0$.
For $u = x^2 + bx + c$ and $v = dx + e$, the remainder of $v^2 - \bar f$ modulo~$u$ has two
coefficients, polynomials in $b, c, d, e$, which vanish over~$\FF_5$ only for $d = e = 0$
(\texttt{decide} over the~$5^4$ tuples). So $v = 0$, and the class of $(u, 0)$ has order
dividing~$2$. The order~$16$ is Theorem~\ref{thm:ds-two-torsion}.

\emph{Formalization.} The statements are for any polynomial over $\ZZ/5$ equal to $x^5 - x$, so
that they apply to the reduction of the integral model (\texttt{DescentExample1.map\_F\_zmod\_five},
by \texttt{reduce\_mod\_char}). The form that Theorem~\ref{thm:ds-e1} consumes is
\texttt{DescentExample1.sq\_eq\_one\_of\_eq\_X\_pow\_five\_sub\_X}; the order~$16$ is not stated.
\end{proof}

\begin{theorem}[E1]
\label{thm:ds-e1}
\lean{ChabautyColeman.DescentExample1.ker_powMonoidHom_two_eq_top,
  ChabautyColeman.DescentExample1.card_hypJac, ChabautyColeman.DescentExample1.ratPoints_eq}
\leanok
\uses{def:ds-cited, def:ds-red, def:hypjac}
Assume~(Loc) for~$\QQ_2$ and~$\QQ_3$ and~(Red) for $p = 5$. Then
$J(\QQ) = J(\QQ)[2] \cong (\ZZ/2)^4$, and $C(\QQ) = \{\infty, (0, 0), (\pm 1, 0), (\pm 3, 0)\}$.
\end{theorem}

\begin{proof}
\leanok
\uses{prop:ds-e1-sel, thm:ds-kernel, def:ds-relaxed, lem:ds-loc, lem:ds-loc-two,
  thm:basechange-inj, lem:ds-e1-five, thm:ds-two-torsion, lem:mumford-exists}
Write $J = J(\QQ)$ and $\Gamma = J[2]$, a group of order~$16$ by Theorem~\ref{thm:ds-two-torsion}.
By Definition~\ref{def:ds-relaxed} and Proposition~\ref{prop:ds-e1-sel},
$\mu(J) \subseteq \mu(\Gamma)$, and as $\ker\mu = 2J$ (Theorem~\ref{thm:ds-kernel}),
$J = \Gamma + 2J$. Multiplying by~$2^n$ gives $2^nJ = 2^n\Gamma + 2^{n+1}J \subseteq \Gamma +
2^{n+1}J$, so $J = \Gamma + 2^nJ$ for every~$n$, by induction. Let $a \in J$, and for each~$n$
choose $\gamma_n \in \Gamma$ with $a - \gamma_n \in 2^nJ$. As~$\Gamma$ is finite, some
$\gamma \in \Gamma$ equals~$\gamma_n$ for infinitely many~$n$; as the subgroups~$2^nJ$ decrease,
$t = a - \gamma$ lies in $\bigcap_n 2^nJ$. Every element of $\bigcap_n 2^nJ$ has odd order: its
image in~$J(\QQ_2)$ does, by Lemma~\ref{lem:ds-loc-two}, and $J \to J(\QQ_2)$ is injective
(Theorem~\ref{thm:basechange-inj}). By~(Red) and Lemma~\ref{lem:ds-e1-five}, the image of~$t$
in~$\tilde J(\FF_5)$ has odd order and order dividing~$2$, so it is~$0$, and $t = 0$. So
$J = \Gamma$. Finally let $P = (x_0, y_0) \in C(\QQ)$. As $[P - \infty]$ has order dividing~$2$,
$[P - \infty] = [(x_0, -y_0) - \infty]$; the reduced Mumford representatives $(x - x_0, y_0)$
and $(x - x_0, -y_0)$ of these classes agree (Lemma~\ref{lem:mumford-exists}), so $y_0 = 0$ and
$f(x_0) = 0$.

\emph{Formalization.} In the module \texttt{HyperellipticCurves.Descent.Example1}. $C(\QQ)$ is
stated for the affine points, $\{(x, y) \in \QQ^2 : y^2 = f(x)\} = \{(\alpha_i, 0)\}$. The
hypotheses are \texttt{HypJac.LocalStructure} for~$\QQ_2$ and~$\QQ_3$ and
\texttt{HypJac.TorsionReduction 5}; \texttt{HypJac.MordellWeil} does not occur.
$J(\QQ) = J(\QQ)[2] \cong (\ZZ/2)^4$ is stated as \texttt{ker\_powMonoidHom\_two\_eq\_top}, the
kernel of the squaring map is everything, with \texttt{card\_hypJac}, $\#J(\QQ) = 16$. The step
from $\Gamma + 2J$ to $\Gamma + 2^nJ$ is \texttt{Subgroup.sup\_range\_powMonoidHom\_pow\_eq\_top},
and the reduced representatives are compared by \texttt{mumfordClass\_inj}. The Lean concludes
through the torsion of~$J(\QQ)$: by the argument with the finite group~$\Gamma$
(\texttt{Subgroup.coprime\_orderOf\_mk\_of\_mem\_iInf\_sup\_range\_powMonoidHom\_pow}),
$J(\QQ)/\Gamma$ is a torsion group, hence so is~$J(\QQ)$ (Mathlib's
\texttt{IsMulTorsion.extension\_closed}), and by~(Red) and Lemma~\ref{lem:ds-e1-five} every element
of~$J(\QQ)$ has order dividing~$2$.
\end{proof}

\begin{remark}[E2: rank~$1$]
\label{rem:ds-e2}
The second example is the curve
\[ C \colon y^2 = f(x) = (x + 4)(x + 3)(x^3 - x + 1) = x^5 + 7x^4 + 11x^3 - 6x^2 - 5x + 12\,, \]
with $\operatorname{disc} f = -23^3 \cdot 59^2$, so $\Sigma_f = \{23, 59\}$. Here
$L = \QQ \times \QQ \times F$, where $F = \QQ[t]/(t^3 - t + 1)$ is the cubic field of
discriminant~$-23$, of class number~$1$ and unit rank~$1$. The places $V = \{23, 59\}$ suffice for
the descent: the prime~$2$ is not needed, and the components in~$\QQ_v$ of the local conditions
suffice, apart from one class at~$23$ (Lemma~\ref{lem:ds-e2-local}).

The group $J(\QQ)[2] \cong (\ZZ/2)^2$ (Theorem~\ref{thm:ds-two-torsion}) is the torsion subgroup
of~$J(\QQ)$ (by a computation with Magma), and the class of $(-11/9, 1020/729) - \infty$ has
infinite order (Lemma~\ref{lem:ds-e2-points}). Under~(MW) and~(Loc) for~$\QQ_{23}$ and~$\QQ_{59}$,
$\operatorname{rk} J(\QQ) = 1$, and $\Sel_2 = S(\Sigma_f, V)$ has dimension~$3$, as it lies between
$\mu(J(\QQ))$ and $S(\Sigma_f, V)$ (Theorem~\ref{thm:ds-e2}). The known rational points are
$\infty$, $(-4, 0)$, $(-3, 0)$ and $(-11/9, \pm 1020/729)$. As $\operatorname{rk} J(\QQ) = 1 < g$,
Theorem~\ref{thm:main} with $r = 1$ gives $\#C(\QQ) \le \#\bar C(\FF_p) + 2$ for the primes
$p \ge 5$ of good reduction; this bound is at least~$8$ at every such prime (it is~$8$ for $p = 5$
and $p = 7$, and for $p > 19$ the Weil bound gives more). So it does not determine~$C(\QQ)$; that
needs further methods, such as the Selmer group Chabauty method of [Sto17], and is not formalized
here.
\end{remark}

\begin{setting}
\label{set:ds-e2}
\lean{ChabautyColeman.DescentExample2.g, ChabautyColeman.DescentExample2.f,
  ChabautyColeman.DescentExample2.S, ChabautyColeman.DescentExample2.P}
\leanok
Let $g = x^3 - x + 1$ and $f = (x + 4)(x + 3)g \in \QQ[x]$, let $F = \QQ[t]/(g)$, with ring of
integers~$O_F$, and $L = \QQ[T]/(f) \cong \QQ \times \QQ \times F$,
$\xi \mapsto (\xi(-4), \xi(-3), \xi(t))$; we write classes in $L^\times/(L^\times)^2$ as triples
of their components. Let $\Sigma = \{23, 59\}$, let~$V$ consist of~$\QQ_{23}$ and~$\QQ_{59}$, and
$S = S(\Sigma, V)$ (Definition~\ref{def:ds-relaxed}). Let $T_1, T_2, P \in J(\QQ)$ be the classes
of $(-4, 0) - \infty$, $(-3, 0) - \infty$ and $(-11/9, 1020/729) - \infty$.
\end{setting}

\begin{remark}[Formalization]
In the modules \texttt{HyperellipticCurves.Descent.Example2.Basic} (Setting~\ref{set:ds-e2} to
Lemma~\ref{lem:ds-e2-padic}\,(1), and Lemma~\ref{lem:ds-e2-points}) and
\texttt{HyperellipticCurves.Descent.Example2.Rank} (Lemma~\ref{lem:ds-e2-local},
Proposition~\ref{prop:ds-e2-sel}, Theorem~\ref{thm:ds-e2}), as for~E1. $f$ has separable reduction
outside~$\Sigma$, also at~$2$, where \texttt{Polynomial.HasSeparableReduction} asks nothing special:
$\operatorname{disc} f = -23^3 \cdot 59^2$, from
$\operatorname{disc}(pq) = \operatorname{disc}(p) \operatorname{disc}(q) \operatorname{Res}(p, q)^2$
with $\operatorname{disc} g = -23$, $g(-4) = -59$ and $g(-3) = -23$. The classes are
\texttt{mumfordClass} values, $P$ that of $(x + 11/9, 1020/729)$. The factors $x + 4$, $x + 3$, $g$
of~$f$ are indexed by $\{0, 1, 2\}$ (\texttt{DescentExample2.factorsEquiv}).
\end{remark}

\begin{lemma}[The cubic field]
\label{lem:ds-e2-field}
\lean{ChabautyColeman.DescentExample2.subsingleton_classGroup,
  ChabautyColeman.DescentExample2.card_modPow_two,
  ChabautyColeman.DescentExample2.range_map_eq_closure,
  ChabautyColeman.DescentExample2.selmerGroupAbove_eq_closure}
\leanok
\uses{def:ds-selmer}
\begin{enumerate}
  \item $O_F$ has trivial class group, and $O_F^\times/(O_F^\times)^2$ has order~$4$ and is
    generated by the classes of~$-1$ and~$t$, a unit with $t^{-1} = 1 - t^2$.
  \item The elements $t + 3$, $2t - 3$, $t + 4$, $5t + 6$ and $t^2 + 4t - 3$ of~$O_F$, of norms
    $23$, $-23$, $59$, $-59$ and~$-59$, generate prime ideals, and these are the primes of~$O_F$
    above~$23$ and~$59$:
    \[ 23\,t^3 = -(t + 3)(2t - 3)^2\,, \qquad 59 = t^6\,(t + 4)(5t + 6)(t^2 + 4t - 3)\,. \]
  \item $F(\Sigma, 2)$ is generated by the classes of~$-1$, $t$ and the five elements of~(2).
\end{enumerate}
\end{lemma}

\begin{proof}
\leanok
\uses{lem:ds-e2-padic, lem:ds-padic-squares, lem:ds-resultant}
(1) The power basis of~$t$ has discriminant $\operatorname{disc} g = -23$, so the field
discriminant divides~$-23$ and is negative; the Minkowski bound gives class number~$1$, and
Dirichlet's unit theorem gives unit rank~$1$, the signature being~$(1, 1)$. As the roots of unity
are~$\pm 1$, $O_F^\times/(O_F^\times)^2$ has order~$2^{1 + 1}$. The classes of~$-1$, $t$ and~$-t$
are nontrivial: under $t \mapsto r_1$, resp.~$r_2$, the roots of~$g$ in~$\QQ_{59}$ with
$r_1 \equiv -4$ and $r_2 \equiv -13 \bmod 59$ (Lemma~\ref{lem:ds-e2-padic}), the images $-1$,
$r_1$ and~$-r_2$ are units that are not squares modulo~$59$
(Lemma~\ref{lem:ds-padic-squares}\,(2)). (2) The norm of $at + b$ is $b^3 - a^2b - a^3$
(Lemma~\ref{lem:ds-resultant}), and an element of norm~$\pm p$, for a prime~$p$, generates a prime
ideal of absolute norm~$p$. The norm of~$t^2 + 4t - 3$ follows from the identity for~$59$ by
multiplicativity, as $N(t) = -1$. A prime of~$O_F$ above~$23$ contains $23t^3$, hence $t + 3$ or
$2t - 3$, and so equals the prime ideal generated by it; likewise for~$59$. (3) As the class group
is trivial, $F(\emptyset, 2)$ is the image of the units, given by~(1), and adding the primes
of~(2) one at a time adds their generators.

\emph{Formalization.} $O_F$ is the integral closure of~$\ZZ$ in the field factor \texttt{AdjoinRoot
g} of~$L$. (1) uses general lemmas on number fields that are missing from Mathlib: for the class
number, the Minkowski bound for cubic fields
(\texttt{NumberField.RingOfIntegers.isPrincipalIdealRing\_of\_finrank\_eq\_three\_of\_abs\_discr\_le})
with \texttt{NumberField.isPrincipalIdealRing\_integralClosure\_iff}, and for the unit rank
\texttt{NumberField.Units.finrank\_eq\_one\_of\_finrank\_eq\_three\_of\_discr\_neg} with
\texttt{NumberField.finrank\_additive\_units\_integralClosure}; the order~$4$ is
\texttt{Units.card\_modPow\_two}. The classes of~$-1$, $t$, $-t$ are nonsquares by
\texttt{Padic.isSquare\_iff\_modPowTwoCoords\_eq\_zero} after the evaluations $t \mapsto r_1, r_2$
(\texttt{AdjoinRoot.lift}). (2) uses \texttt{NumberField.prime\_of\_abs\_norm\_eq} (an element of
norm~$\pm p$ is prime); the norms of linear elements come from
\texttt{AdjoinRoot.leadingCoeff\_mul\_norm\_of\_sub\_root}. (3) comes from
\texttt{IsDedekindDomain.selmerGroup\_image\_eq\_range\_sup\_closure} and
\texttt{IsDedekindDomain.selmerGroup\_range\_eq\_range\_sup\_closure}, the generalization of
Lemma~\ref{lem:ds-e1-selmer}\,(1): for trivial class group and a finite set of primes with given
generators, $K(\Sigma, n)$ is generated by the image of~$R^\times$ and the generators (by
\texttt{IsDedekindDomain.selmerGroup\_insert} and
\texttt{IsDedekindDomain.selmerGroup.fromUnit\_surjective}); \texttt{Rat.selmerGroup\_eq\_closure}
is proved through them. The consumed form of~(3),
\texttt{DescentExample2.selmerGroupAbove\_eq\_closure}, is for any polynomial equal to~$g$, so that
it applies to the field factor of~$f$ at the factor~$g$.
\end{proof}

\begin{lemma}[$A(\Sigma, 2)$ of~E2]
\label{lem:ds-e2-selmer}
\lean{ChabautyColeman.DescentExample2.selmerGroup_le_closure,
  ChabautyColeman.DescentExample2.card_selmerGroup_le}
\leanok
\uses{lem:ds-e1-selmer, lem:ds-e2-field}
A class $\xi \in L^\times/(L^\times)^2$ lies in $A(\Sigma, 2)$ if and only if
$\xi(-4), \xi(-3) \in \QQ(\Sigma, 2)$ and $\xi(t) \in F(\Sigma, 2)$. So $A(\Sigma, 2)$ is
generated by the $13$ classes with two trivial components and the third one equal to~$-1$, $23$
or~$59$ (at~$-4$ or~$-3$) or to one of the seven generators of Lemma~\ref{lem:ds-e2-field}\,(3)
(in~$F$), and $\#A(\Sigma, 2) \le 2^{13}$.
\end{lemma}

\begin{proof}
\leanok
\uses{def:ds-selmer}
The field factors of~$L$ are those of the three irreducible factors of~$f$, and $A(\Sigma, 2)$ is
defined componentwise (Definition~\ref{def:ds-selmer}). For the two linear factors,
Lemma~\ref{lem:ds-e1-selmer}\,(2) identifies the condition with
$\xi(-4), \xi(-3) \in \QQ(\Sigma, 2)$, generated by~$-1$, $23$, $59$
(Lemma~\ref{lem:ds-e1-selmer}\,(1)); for~$g$, Lemma~\ref{lem:ds-e2-field}\,(3).

\emph{Formalization.} \texttt{Polynomial.Factors.equivOfProd} identifies the factors of a product
of distinct monic irreducible polynomials with their index set, as
\texttt{Polynomial.Factors.equivOfRoots} does for linear factors. The $13$ classes are the images,
under the inverse of \texttt{AdjoinRoot.modPowEquivPiFactors}, of families with one nontrivial
component. A subgroup generated by~$n$ elements of order dividing~$2$ has at most~$2^n$ elements
(\texttt{Subgroup.card\_closure\_range\_le\_two\_pow}). Only the inclusion of $A(\Sigma, 2)$ in the
subgroup generated by the $13$ classes is formalized, which is what
Proposition~\ref{prop:ds-e2-sel} uses. The classes are given by data, the factor and the
coefficients of the nontrivial component, so that their coordinates in
Proposition~\ref{prop:ds-e2-sel} are computed by one lemma.
\end{proof}

\begin{lemma}[The local fields of E2]
\label{lem:ds-e2-padic}
\lean{PadicInt.exists_aeval_eq_zero_toZModPow_eq, Padic.not_isSquare_adjoinRoot_of_odd_valuation,
  Padic.modPowTwoEquiv_mk_of_toZModPow_eq, Padic.modPowTwoEquivZModTwo_mk_of_toZModPow_two_eq,
  ChabautyColeman.DescentExample2.exists_isRootData_fiftyNine,
  ChabautyColeman.DescentExample2.exists_isRootData_twentyThree}
\leanok
\uses{lem:ds-padic-squares}
\begin{enumerate}
  \item $g$ has a root $r \in \ZZ_{23}$ with $r \equiv 181 \bmod 23^2$, and roots
    $r_1, r_2, r_3 \in \ZZ_{59}$ with $r_1, r_2, r_3 \equiv 3182, 2406, 1374 \bmod 59^2$; so
    $r \equiv -3 \bmod 23$ and $r_1, r_2, r_3 \equiv -4, -13, 17 \bmod 59$.
  \item Over~$\QQ_{59}$, $f$ has the five distinct roots $-4, -3, r_1, r_2, r_3$.
    Over~$\QQ_{23}$, $f = (x + 4)(x + 3)(x - r)h$ with $h = x^2 + rx + r^2 - 1$ irreducible: its
    discriminant $4 - 3r^2 \equiv 5 \cdot 23 \bmod 23^2$ has odd valuation.
  \item Let~$p$ be odd and $x \in \ZZ_p$ with $x \equiv p^n m \bmod p^{n+1}$ for an integer~$m$
    prime to~$p$. Then~$x$ and $p^n m$ have the same class in $\QQ_p^\times/(\QQ_p^\times)^2$. In
    particular, for $Q \in \ZZ[x]$ and a root~$\rho$ of~(1) with $\rho \equiv a \bmod p^2$, the
    class of~$Q(\rho)$ is determined by $Q(a) \bmod p^2$ if $p^2 \nmid Q(a)$.
  \item Let $E = \QQ_{23}[x]/(h)$, with~$\theta_h$ the class of~$x$, and let
    $u, w \in \QQ_{23}$ with $u - wr/2 \in \ZZ_{23}^\times$ not a square modulo~$23$. Then
    $u + w\theta_h$ is not a square in~$E$. In particular, $-\theta_h$ is not a square, as
    $r/2 \equiv 10 \bmod 23$.
\end{enumerate}
\end{lemma}

\begin{proof}
\leanok
(1) Hensel's lemma, in the form with $\lvert \rho - a\rvert = \lvert g(a)\rvert$ for
$\lvert g'(a)\rvert = 1$, as $g(a) \equiv 0 \bmod p^2$ and $p \nmid g'(a)$ for these~$a$. (2) The
roots are distinct modulo~$59$ except $r_1 \equiv -4$, and $r_1 \ne -4$ as $g(-4) = -59$. A
root~$x_0$ of~$h$ in~$\QQ_{23}$ would make $4 - 3r^2 = (2x_0 + r)^2$ a square. (3) $x/(p^n m)$ is
a unit $\equiv 1 \bmod p$, hence a square (Lemma~\ref{lem:ds-padic-squares}\,(1)). (4) Let
$s = 2\theta_h + r$, so that $s^2 = D = 4 - 3r^2$ and
$u + w\theta_h = (u - wr/2) + (w/2)s$. If $u + w\theta_h = (\alpha + \beta s)^2$ with
$\alpha, \beta \in \QQ_{23}$, then $u - wr/2 = \alpha^2 + D\beta^2$. The two terms have
valuations of different parity, so~$\alpha^2$ is a unit and $D\beta^2 \in 23\ZZ_{23}$, and
$u - wr/2 \equiv \alpha^2 \bmod 23$.

\emph{Formalization.} (1) is Hensel's lemma with a congruence modulo~$p^k$
(\texttt{PadicInt.exists\_aeval\_eq\_zero\_toZModPow\_eq}, from Mathlib's
\texttt{hensels\_lemma}). (3) is the value of \texttt{Padic.modPowTwoEquiv} at $p^n u$ for a
unit~$u$ of~$\ZZ_p$ with a given residue (\texttt{Padic.modPowTwoEquiv\_mk\_of\_eq\_pow\_mul}, which
\texttt{Padic.modPowTwoEquiv\_mk\_of\_val\_eq} uses) and at $x \equiv p^n m \bmod p^k$, $n < k$
(\texttt{Padic.modPowTwoEquiv\_mk\_of\_toZModPow\_eq}). (4) is stated for a monic quadratic
over~$\QQ_p$, $p$ odd, whose discriminant has odd valuation
(\texttt{Padic.not\_isSquare\_adjoinRoot\_of\_odd\_valuation}); the irreducibility in~(2) comes from
\texttt{Polynomial.irreducible\_X\_sq\_add\_C\_mul\_X\_add\_C\_of\_not\_isSquare}. The example
uses~(3) in the form \texttt{Padic.modPowTwoEquivZModTwo\_mk\_of\_toZModPow\_two\_eq}: the class
of $c^2 y$, for $y \in \ZZ_p$ with residue $z \ne 0$ modulo~$p^2$, has the coordinates
\texttt{Padic.modPowTwoCoords p z} in~$(\ZZ/2)^2$ (the valuation modulo~$2$, and whether the unit
part is a nonsquare modulo~$p$), which \texttt{decide} computes. The roots enter only inside
proofs, through \texttt{DescentExample2.exists\_isRootData\_fiftyNine} and
\texttt{DescentExample2.exists\_isRootData\_twentyThree}; the lemmas on coordinates take a vector
of roots with their residues modulo~$p^2$ (\texttt{DescentExample2.IsRootData}), and no main
statement mentions them.
\end{proof}

\begin{lemma}[The local images of E2]
\label{lem:ds-e2-local}
\lean{ChabautyColeman.DescentExample2.range_descentMap_twentyThree,
  ChabautyColeman.DescentExample2.range_descentMap_fiftyNine}
\leanok
\uses{lem:ds-loc, thm:ds-two-torsion, thm:ds-kernel, thm:ds-formula, cor:ds-torsion-values,
  thm:ds-naturality, lem:ds-e2-padic}
Assume~(Loc) for~$\QQ_{23}$ and~$\QQ_{59}$. Then $\mu_{59}(J(\QQ_{59}))$ is the group of
order~$16$ generated by the images of~$T_1$, $T_2$ and of the points of~$C(\QQ_{59})$ with
$x = -1/2$ and $x = 4/3$, and $\mu_{23}(J(\QQ_{23}))$ is the group of order~$8$ generated by the
images of~$T_1$, $T_2$ and of the points of~$C(\QQ_{23})$ with $x = 0$.
\end{lemma}

\begin{proof}
\leanok
\uses{lem:ds-padic-squares}
As for Lemma~\ref{lem:ds-e1-local}: the orders are $\#J(\QQ_v)[2] = 2^4$ and~$2^3$
(Lemma~\ref{lem:ds-loc}\,(1) and Theorem~\ref{thm:ds-two-torsion}), as~$f$ has five, resp.~four,
irreducible factors over~$\QQ_v$ (Lemma~\ref{lem:ds-e2-padic}\,(2)). The points exist:
$f(-1/2) = 385/32$ and $f(4/3) = 11440/243$ are units of~$\ZZ_{59}$ with residues $6^2$ and~$4^2$,
and $f(0) = 12 \equiv 9^2 \bmod 23$. Their classes, and those of~$T_1$ and~$T_2$
(Corollary~\ref{cor:ds-torsion-values}), are compared through their components at the roots
of~$f$ in~$\QQ_v$, each mapped by Lemma~\ref{lem:ds-padic-squares}\,(2) to its valuation
modulo~$2$ and the class of its residue, computed by Lemma~\ref{lem:ds-e2-padic}\,(3). At~$59$
the four classes are independent. At~$23$ the classes of~$T_1$ and~$T_2$ are independent, and
the class~$\delta$ of the point with $x = 0$ has trivial components at the roots $-4$, $-3$
and~$r$, as $4$, $3$ and $-r \equiv 3$ are squares modulo~$23$; but its component $-\theta_h$
in~$E$ is not a square (Lemma~\ref{lem:ds-e2-padic}\,(4)), so the three classes are independent.

\emph{Formalization.} As in~E1, with three general lemmas:
\texttt{HypJac.LocalStructure.range\_descentMap\_eq\_of\_le\_of\_card\_ge} (for odd~$v$, a subgroup
of the local image of order at least $\#J(\QQ_v)[2]$ is all of it),
\texttt{HypJac.mk\_mem\_range\_descentMap\_of\_isSquare} (the class of $x_0 - \theta$ lies in the
local image if $f(x_0)$ is a square) and \texttt{Padic.isSquare\_iff\_modPowTwoCoords\_eq\_zero};
the lower bounds come from \texttt{Subgroup.card\_closure\_range\_eq\_two\_pow\_of\_monoidHom}. The
components at the roots are evaluations $L_v \to \QQ_v$ (\texttt{AdjoinRoot.liftAlgHom}), with the
coordinates \texttt{DescentExample2.rootBits}, and, for~$E$, the map
\texttt{AdjoinRoot.algHomOfDvd} to $\QQ_{23}[x]/(h)$; over~$\QQ_{23}$ no isomorphism of
$L_v^\times/(L_v^\times)^2$ with a product is needed. The number of factors over~$\QQ_{23}$ comes
from \texttt{Polynomial.Factors.card\_eq\_of\_associated\_prod}, over~$\QQ_{59}$ from
\texttt{Polynomial.Factors.card\_eq\_of\_injective} as in~E1. Only components of odd valuation
need the roots modulo~$p^2$: $-4 - r_1$ of~$\mu_{59}(T_1)$ and $-3 - r$ of~$\mu_{23}(T_2)$ here,
and the generators of the primes of Lemma~\ref{lem:ds-e2-field}\,(2) at their own roots in
Proposition~\ref{prop:ds-e2-sel}.
\end{proof}

\begin{proposition}[The relaxed Selmer group of~E2]
\label{prop:ds-e2-sel}
\lean{ChabautyColeman.DescentExample2.card_relaxedSelmerGroup_le}
\leanok
\uses{def:ds-relaxed, lem:ds-e2-selmer, lem:ds-e2-local}
Assume~(Loc) for~$\QQ_{23}$ and~$\QQ_{59}$. Then $\#S(\Sigma, V) \le 8$.
\end{proposition}

\begin{proof}
\leanok
\uses{thm:ds-naturality, lem:ds-e2-padic, lem:ds-padic-squares}
Each root~$\rho$ of~$f$ in~$\QQ_v$, for $v = 23, 59$, gives a homomorphism
$L^\times/(L^\times)^2 \to \QQ_v^\times/(\QQ_v^\times)^2 \to \FF_2^2$: the component at~$\rho$, then
its valuation modulo~$2$ and whether its residue is a nonsquare. Together, the roots $-4$, $-3$,
$r_1$, $r_2$, $r_3$ of~$f$ in~$\QQ_{59}$ and $-4$, $-3$, $r$ in~$\QQ_{23}$ give
\[ \beta = (\beta_{59}, \beta_{23}) \colon L^\times/(L^\times)^2 \to \FF_2^{10} \times
   \FF_2^6\,, \]
with~$\beta_v$ factoring through~$\operatorname{res}_v$ (Theorem~\ref{thm:ds-naturality}). By
Lemma~\ref{lem:ds-e2-local}, the images of the local images $\mu_v(J(\QQ_v))$ under~$\beta_{59}$
and~$\beta_{23}$ have dimensions~$4$ and~$2$, as the class~$\delta$ of the point with $x = 0$ lies
in the kernel of~$\beta_{23}$. So there are ten linearly independent linear forms on
$\FF_2^{10} \times \FF_2^6$ vanishing on both, which give a homomorphism
$\Theta \colon L^\times/(L^\times)^2 \to \FF_2^{10}$ with $S \subseteq \ker\Theta$. The forms can
be chosen so that~$\Theta$ maps the ten generators $(-1, 1, 1)$, $(23, 1, 1)$, $(1, 23, 1)$,
$(1, 59, 1)$ and $(1, 1, \varphi)$, $\varphi = t$, $t + 3$, $2t - 3$, $t + 4$, $5t + 6$,
$t^2 + 4t - 3$, of~$A(\Sigma, 2)$ (Lemma~\ref{lem:ds-e2-selmer}) to the standard basis
of~$\FF_2^{10}$. So $\Theta(A(\Sigma, 2)) = \FF_2^{10}$, and
$\#S \le \#(A(\Sigma, 2) \cap \ker\Theta) = \#A(\Sigma, 2)/2^{10} \le 2^3$.

\emph{Formalization.} The forms are a $16 \times 10$ matrix over~$\ZZ/2$, computed with Magma and
checked in Lean by \texttt{decide}: it kills the values of~$\beta$ on the generators of the local
images and maps the ten generators to the standard basis. The map~$\beta_v$ is
\texttt{DescentExample2.globalBits}, with values in~$(\ZZ/2)^2$ per root
(\texttt{Padic.modPowTwoEquivZModTwo}; the second coordinate is the quadratic character
of~$\ZZ/p$, Mathlib's \texttt{quadraticChar}, through
\texttt{FiniteField.modPowTwoEquivZModTwo}). The count runs in the subgroup~$B$ generated by the
$13$ classes, which contains $A(\Sigma, 2)$ (\texttt{DescentExample2.selmerGroup\_le\_closure}), so
the classes need not be shown to lie in~$A(\Sigma, 2)$: $\#(B \cap \ker\Theta) \cdot
\#\Theta(B) = \#B \le 2^{13}$, by Mathlib's \texttt{Subgroup.relIndex\_mul\_card} and
\texttt{Subgroup.relIndex\_ker}, and \texttt{Subgroup.card\_closure\_range\_le\_two\_pow}. The norm
condition is not used, and the remaining three generators $(59, 1, 1)$, $(1, -1, 1)$,
$(1, 1, -1)$ need no values.
\end{proof}

\begin{lemma}[$P$ has infinite order]
\label{lem:ds-e2-points}
\lean{ChabautyColeman.DescentExample2.not_isOfFinOrder_P,
  ChabautyColeman.DescentExample2.card_closure_descentMap_eq}
\leanok
\uses{thm:ds-formula, cor:ds-torsion-values, thm:ds-naturality, thm:ds-two-torsion,
  lem:ds-e2-padic}
The classes $\mu(T_1)$, $\mu(T_2)$ and $\mu(P)$ generate a subgroup of order~$8$
of~$L^\times/(L^\times)^2$, and~$P$ has infinite order.
\end{lemma}

\begin{proof}
\leanok
\uses{lem:ds-padic-squares}
The values are (Corollary~\ref{cor:ds-torsion-values} and Theorem~\ref{thm:ds-formula})
\[ \mu(T_1) = (59, -1, -4 - t)\,, \qquad \mu(T_2) = (1, -23, -3 - t)\,, \qquad
   \mu(P) = (25/9, 16/9, -11/9 - t)\,. \]
Their images under the map~$\beta$ of the proof of Proposition~\ref{prop:ds-e2-sel} are
independent: only $\mu(T_1)$ has odd valuation at the root~$-4$ in~$\QQ_{59}$; $\mu(P)$ has trivial
components in~$\QQ_{23}$, unlike~$\mu(T_2)$; and the component $-11/9 - r_2$ of~$\mu(P)$ is a unit
of~$\ZZ_{59}$ that is not a square modulo~$59$. If~$P$ had finite order,
$B = J(\QQ)[2] + \langle P\rangle$ would be a finite group with $B[2] = J(\QQ)[2]$, of order~$4$
(Theorem~\ref{thm:ds-two-torsion}, as~$f$ has three irreducible factors). As~$\mu$ kills~$2B$ and
$\#(B/2B) = \#B[2]$ for a finite abelian group, $\#\mu(B) \le 4$, whereas~$\mu(B)$ contains the
group of order~$8$ above.

\emph{Formalization.} Neither~(MW) nor~(Loc) is used. The counting step is
\texttt{Subgroup.card\_map\_le\_card\_inf\_ker\_powMonoidHom}: for a finite subgroup~$B$ of a
commutative group~$G$ and a homomorphism~$\varphi$ to a group of exponent~$n$,
$\#\varphi(B) \le \#(B \cap G[n])$, here with $n = 2$; a subgroup generated by finitely many
elements of finite order is finite (\texttt{Subgroup.finite\_closure\_of\_isOfFinOrder}, from
Mathlib's \texttt{CommGroup.finite\_of\_fg\_isMulTorsion}). The Lean takes for~$B$ the subgroup
generated by~$T_1$, $T_2$ and~$P$ and bounds $\#B[2]$ by $\#J(\QQ)[2] = 4$, so it does not need
$J(\QQ)[2] = \langle T_1, T_2\rangle$.
\end{proof}

\begin{theorem}[E2]
\label{thm:ds-e2}
\lean{ChabautyColeman.DescentExample2.finrank_hypJac,
  ChabautyColeman.DescentExample2.relaxedSelmerGroup_eq}
\leanok
\uses{def:ds-cited, thm:ds-rank-bound, prop:ds-e2-sel, lem:ds-e2-points}
Assume~(MW) and~(Loc) for~$\QQ_{23}$ and~$\QQ_{59}$. Then $\operatorname{rk} J(\QQ) = 1$, and
$S(\Sigma, V) = \mu(J(\QQ)[2] + \langle P\rangle)$, of order~$8$.
\end{theorem}

\begin{proof}
\leanok
\uses{thm:ds-two-torsion, def:ds-relaxed}
By Theorem~\ref{thm:ds-rank-bound} and Proposition~\ref{prop:ds-e2-sel},
$2^{\operatorname{rk} J(\QQ)} \cdot 4 \le 8$, as $\#J(\QQ)[2] = 2^{3 - 1}$
(Theorem~\ref{thm:ds-two-torsion}). By Lemma~\ref{lem:ds-e2-points}, $P$ has infinite order, so
$\operatorname{rk} J(\QQ) \ge 1$, as $J(\QQ)$ is finitely generated. Finally
$\mu(J(\QQ)[2] + \langle P\rangle) \subseteq S$ (Definition~\ref{def:ds-relaxed}) has order~$8$ by
Lemma~\ref{lem:ds-e2-points}.

\emph{Formalization.} The rank is \texttt{Module.finrank} over~$\ZZ$ of \texttt{Additive (HypJac
$\QQ$ f)}; the hypotheses are the instance that \texttt{HypJac $\QQ$ f} is finitely generated (which
\texttt{HypJac.MordellWeil} provides) and \texttt{HypJac.LocalStructure} for~$\QQ_{23}$
and~$\QQ_{59}$; the statement on $S(\Sigma, V)$, \texttt{relaxedSelmerGroup\_eq}, needs~(Loc) only.
The instances of Theorem~\ref{thm:ds-rank-bound} on the class groups and unit groups of the field
factors are found for $K = \QQ$ in general, as the instances
\texttt{NumberField.instFiniteClassGroupIntegralClosure} and
\texttt{NumberField.instIsMulFGUnitsIntegralClosure}; $\operatorname{rk} J(\QQ) \ge 1$ comes through
Mathlib's \texttt{Module.finrank\_eq\_zero\_iff\_isTorsion}. The order~$8$ of $S(\Sigma, V)$ is not
part of the Lean statement; it is \texttt{DescentExample2.card\_closure\_descentMap\_eq} with the
proof.
\end{proof}

